\documentclass[11pt,letterpaper]{article}
\usepackage[margin=0.86in,headheight=16pt]{geometry}
\usepackage{fontspec}
\setmainfont{DejaVu Serif}
\setsansfont{DejaVu Sans}
\setmonofont{DejaVu Sans Mono}[Scale=0.82]
\newfontfamily\chinesefont{Noto Serif CJK SC}
\usepackage{amsmath,amssymb,amsthm,mathtools,mathrsfs}
\usepackage{microtype,enumitem,booktabs,longtable,array}
\usepackage{xcolor,fancyhdr,fvextra,needspace}
\usepackage[unicode,colorlinks=true,linkcolor=black,urlcolor=blue!50!black,citecolor=black]{hyperref}
\usepackage{bookmark}
\DefineVerbatimEnvironment{verbatim}{Verbatim}{breaklines=true,breakanywhere=true,fontsize=\small,baselinestretch=1.03,xleftmargin=0.8em,xrightmargin=0.5em,breaksymbolleft={}}
\providecommand{\tightlist}{\setlength{\itemsep}{0pt}\setlength{\parskip}{0pt}}
\setlist{topsep=4pt,itemsep=2pt,parsep=1pt,leftmargin=*}
\setlength{\parskip}{4pt}
\setlength{\parindent}{0pt}
\setlength{\emergencystretch}{2em}
\hyphenation{yields domains weights}
\makeatletter
\renewcommand{\numberline}[1]{\hb@xt@2.2em{#1\hfil}}
\makeatother
\setcounter{tocdepth}{1}
\setcounter{secnumdepth}{1}
\pagestyle{fancy}\fancyhf{}
\fancyhead[L]{\small Finite rank BFSS exterior estimates}
\fancyhead[R]{\small Technical review draft v2.1}
\fancyfoot[C]{\thepage}
\renewcommand{\headrulewidth}{0pt}
\newtheorem{theorem}{Theorem}[section]
\newtheorem{lemma}[theorem]{Lemma}
\newtheorem{proposition}[theorem]{Proposition}
\newtheorem{corollary}[theorem]{Corollary}
\theoremstyle{definition}\newtheorem{hypothesis}[theorem]{Hypothesis}
\newcommand{\Tr}{\operatorname{Tr}}
\newcommand{\ad}{\operatorname{ad}}
\newcommand{\Rea}{\operatorname{Re}}
\newcommand{\Ran}{\operatorname{Ran}}
\newcommand{\supp}{\operatorname{supp}}
\newcommand{\phys}{\mathrm{phys}}
\newcommand{\exc}{\mathrm{exc}}
\newcommand{\dd}{\,\mathrm{d}}
\newcommand{\hh}{\mathfrak h}
\newcommand{\QQ}{\mathcal Q}
\newcommand{\FF}{\mathcal F}
\newcommand{\norm}[1]{\left\lVert #1\right\rVert}
\newcommand{\ip}[2]{\left\langle #1,#2\right\rangle}
\numberwithin{equation}{section}
\begin{document}
\begin{titlepage}
\vspace*{0.2in}
{\LARGE\bfseries Exterior Hardy bounds for\\[4pt]finite rank BFSS Hamiltonians\par}
\vspace{0.25in}
{\large Global induction and a singlet decay improvement\par}
\vspace{0.35in}
{\large Yicheng Pan\quad {\chinesefont 潘奕成}\par}
{5 October 2026\quad Version 2.1\quad Local review copy\\[2pt]Typesetting revision 6 October 2026\par}
\vspace{0.3in}
\textbf{Technical review draft.} This document presents a proposed analytical proof for specialist checking. The global implication is stated with its exact local inputs, and twelve appendices contain the baseline derivations and the added rotation argument. Those inputs are substantive parts of the proposed BFSS proof; their inclusion is not a substitute for verifying them. No independent human review or proof-assistant verification is represented as having occurred.

\textbf{AI disclosure.} OpenAI AI systems made substantial contributions to research, derivation, code, manuscript preparation, and mathematical checking of this draft. The author name identifies the person under whose direction the work was prepared; this version does not assert that the named author has personally checked every argument or approved a public submission. AI-generated checking is not independent expert peer review.

\vspace{0.2in}
\begin{abstract}
For the trace-normalized sixteen-supercharge $SU(N)$ matrix Hamiltonian, we assemble a candidate exterior quadratic-form proof at every fixed finite rank. The baseline bound has every coefficient $c<49/4$ on arbitrary physical exterior vectors. This revision adds $c<121/4$ on physical $\operatorname{Spin}(9)$ singlets, using a bounded-radius internal nonsinglet uncertainty estimate and rotation-equivariant cluster localization. Internal continua are retained; no internal spectral gap or ground-state projection is required. Combined with the existing Hasler--Hoppe theorem that every physical $L^2$ zero mode is rotation invariant, the proposed full argument gives ordinary radial moments of every order $0\le m<9$ for each actual zero mode. These conclusions depend on the complete baseline local and global analytical inputs, whose candidate derivations remain subject to specialist review. Existence, uniqueness, the endpoint $m=9$, uniformity in $N$, and full soft or scattering theorems are not established. Polchinski already stated the below-nine expectation through a Born--Oppenheimer tail argument; the rotation identities and weighted method also have substantial prior art. No priority or independent human or formal proof verification is claimed.
\end{abstract}
\vfill
Original documentation: CC BY 4.0. Original accompanying code: MIT. Third-party papers are cited and linked, not redistributed. No affiliation or external endorsement is claimed.
\end{titlepage}
\begingroup
\small
\setlength{\parskip}{0pt}
\tableofcontents
\endgroup
\clearpage

\section{Scope and prior work}
The configuration space has nine traceless Hermitian matrices and noncompact commuting valleys. Separated-channel asymptotics are informative, but an exterior form inequality must hold for every physical test vector, including states whose internal cluster coordinates have unrestricted energy and configurations with several simultaneous or nested clusters. A local expansion with all eigenvalue gaps comparable does not by itself give that estimate.

The proposed statement is deliberately finite-rank. For each fixed $N$, an exterior radius may be chosen after all rank-dependent constants: $c<49/4$ for arbitrary physical vectors, and $c<121/4$ for rotation singlets. There is no bound on how that radius grows with $N$. A conclusion valid at each finite $N$ cannot be exchanged with a large-$N$ limit without further estimates.

Sethi and Stern derive weighted $L^2$ estimates from exterior supercharge coercivity and analyze the excited-oscillator correction and inverse-square cancellation in the two-particle problem \cite[manuscript pp.~20--21 and 37--39]{SS}. The present bounded-weight multiplier argument is an instance of that established Hardy--Agmon method \cite{Agmon}. Konechny computes a free leading effective Hamiltonian in an all-rank large-separation regime \cite{K}. Hasler and Hoppe establish an asymptotic factorization in a region where the selected eigenvalue gaps are all of the same large order \cite{HH}. These are essential precursors, with hypotheses different from a collision-uniform exterior form bound on arbitrary states.

The $SU(2)$ wavefunction analysis of Fr\"ohlich, Graf, Hasler, Hoppe and Yau is explicitly formulated for a formal asymptotic series \cite{FGHHY}. Lin and Yin propose a recursive cluster form of the leading asymptotic wavefunction \cite{LY}. These results motivate the below-nine tail benchmark. The new singlet comparison addresses the angular part of an actual-state estimate; the baseline errors and domains must still be controlled in every channel.

Polchinski already explicitly states convergence of the fourth moment and convergence for orders $L<9$, using a Hamiltonian Born--Oppenheimer argument and an $l=2$ two-block tail \cite[printed p.~12, Eqs.~(7.8)--(7.10)]{P}. The $L<9$ expectation, the $l=2$ angular mechanism, and the associated constants are therefore not claimed as new. The distinction under examination is a global all-channel theorem for actual $L^2$ zero modes. The references inspected here do not provide a priority certificate or establish the absence of a stronger theorem elsewhere.

Sethi and Stern proved rotational invariance results for supersymmetric bound states \cite{SSI}. Hasler and Hoppe gave explicit rotation primitives and proved $\operatorname{Spin}(d)$ invariance for $SU(N)$ in dimensions $d=3,5,9$ \cite{HHI}. We use their kernel-invariance theorem and corrected primitive, rather than presenting them as new. The quantitative support-radius estimate and its insertion into the proposed collision-uniform ledger are the candidate addition under review.

\subsection{How to read the proof status}
Theorem~\ref{thm:global} proves the baseline all-vector implication from the realization and local inputs stated below. Theorem~\ref{thm:singlet} gives the added singlet implication, and Corollary~\ref{cor:moments} applies it to actual zero modes. Appendices A--J supply candidate derivations of the baseline inputs; Appendices~\ref{app:rotation}--\ref{app:model} detail the added rotation and model-domain argument. This is a review packet containing the proposed full argument, not a claim that its difficult local estimates have been independently certified. The finite tests check specified algebraic identities and constants; they do not verify an infinite-dimensional operator inequality.

The most important review points are the actual nonzero-internal source inventory, the colored complete-form reserve, differentiated orbit remainders, the singular-stratum incidence localization, and form-domain compatibility. A reviewer can inspect them directly without relying on a private report or an acceptance label. The relevant derivations, including the corrected $\QQ(H)$ interpretation and explicit module identities, are reproduced here.

\section{Operator and domain conventions}\label{sec:operator}
Let $T_A$, $1\le A\le N^2-1$, be Hermitian generators with
\begin{equation}
 \Tr(T_AT_B)=\delta_{AB},\qquad [T_A,T_B]=if_{ABC}T_C,
 \qquad X_i=X_i^AT_A\quad(1\le i\le9).
\end{equation}
The Euclidean metric is $\sum_i\Tr(\dd X_i^2)$, and
\begin{equation}
 \rho_N^2=\sum_i\Tr(X_i^2)=\sum_{i,A}(X_i^A)^2.
\end{equation}
Take Hermitian Majorana generators $\theta^A_\alpha$ with
$\{\theta^A_\alpha,\theta^B_\beta\}=\delta_{AB}\delta_{\alpha\beta}$, on a complex Clifford module $\FF_N$ of dimension $2^{8(N^2-1)}$. The matrices $\Gamma_i$ are real symmetric $16\times16$ matrices satisfying $\{\Gamma_i,\Gamma_j\}=2\delta_{ij}$; set $\Gamma_{ij}=\frac12[\Gamma_i,\Gamma_j]$.

For a fixed positive length $\ell$, define
\begin{align}
 C_{Aij}&=f_{ABC}X_i^BX_j^C,\\
 (\ad X_i)_{AB}&=\Tr(T_A[X_i,T_B])=-if_{CAB}X_i^C,\\
 D_{\ad}(X)&=\sum_i\Gamma_i\otimes\ad X_i.
\end{align}
With the displayed ordering of spin and color indices, the kinetic-one expression is
\begin{align}
 h_N={}&-\Delta_X+\ell^{-6}\sum_{i<j,A}C_{Aij}^2
              +\ell^{-3}\theta^T D_{\ad}(X)\theta.\label{eq:operator}
\end{align}
The Yukawa term is therefore
$-i\ell^{-3}\sum f_{CAB}X_i^C\theta^A_\alpha(\Gamma_i)_{\alpha\beta}\theta^B_\beta$.
This convention fixes the sign by defining $\ad X$ explicitly. The matrix $D_{\ad}$ is Hermitian and purely imaginary antisymmetric in a Hermitian trace-orthonormal color basis.

The corresponding normalized real charges are
\begin{equation}
 q_\alpha=\sqrt2\sum_{A,\beta}
 \left(\sum_i(\Gamma_i)_{\alpha\beta}p_i^A
   -\ell^{-3}\sum_{i<j}(\Gamma_{ij})_{\alpha\beta}C_{Aij}\right)
 \theta^A_\beta,\qquad p_i^A=-i\partial_{X_i^A}.\label{eq:charges}
\end{equation}
Their spinor-averaged square gives \eqref{eq:operator} on the common smooth core. Individual squares have gauge terms off the physical sector. On physical vectors the usual gauge terms vanish; on colored internal modules only the averaged identity is used. The physical Hamiltonian and charges in units with scalar kinetic coefficient $R_{\rm BFSS}/2$ are
\begin{equation}
 H_N=\frac{R_{\rm BFSS}}2 h_N,
 \qquad Q_\alpha=\sqrt{R_{\rm BFSS}/2}\,q_\alpha.
\end{equation}
The rest of the proof uses $\ell=1$, obtainable by a fixed dilation. Restoring $\ell$ changes the exterior radius, not the dimensionless Hardy coefficient in kinetic-one units.

\begin{hypothesis}[Supersymmetric realization]\label{hyp:realization}
Use the gauge-invariant subspace
$\mathscr H_N=L^2(\mathbb R^{9(N^2-1)};\FF_N)^{SU(N)}$.
On $\mathscr C_N=C_c^\infty(\mathbb R^{9(N^2-1)};\FF_N)^{SU(N)}$ take the supersymmetric form
\begin{equation}
 \hh_N[u]=\frac1{16}\sum_{\alpha=1}^{16}\norm{q_\alpha u}^2.
\end{equation}
Assume its compatible closed supercharge-stack realization and the bounded Lipschitz, gauge-invariant scalar multiplier rule. The operator $h_N$ is the nonnegative self-adjoint operator associated to this closed form. On unrestricted internal-color spaces use the corresponding closed averaged form. The differential expression and averaged identity agree on their common core.
\end{hypothesis}
This is the precise realization used throughout. It is not a Dirichlet realization in a fixed valley. The exterior domain $\QQ_{N,R}$ is the closure of physical compactly supported smooth vectors in $\{\rho_N>R\}$ in the norm $(\norm{u}^2+\hh_N[u])^{1/2}$.
The notation $\QQ(h_N)$ means the quadratic-form domain; it is distinct from the operator domain $D(h_N)$. No moving projection is claimed to preserve $D(h_N)$. The standard physical realization carries the natural unitary $\operatorname{Spin}(9)$ action, which preserves its closed form and physical core. The kernel-invariance input below refers to this compatible realization, not an arbitrary self-adjoint extension.

For real scalar multipliers with $\sum_j\chi_j^2=1$, the kinetic-one product rule gives the IMS identity
\begin{equation}
 \hh_N[u]=\sum_j\hh_N[\chi_ju]
                -\int\sum_j|\nabla\chi_j|^2|u|^2.\label{eq:ims}
\end{equation}
If $\hh_N[\Psi]=0$, every closed charge annihilates $\Psi$, and for bounded real Lipschitz $w$,
\begin{equation}
 \hh_N[w\Psi]=\int|\nabla w|^2|\Psi|^2.\label{eq:zeroidentity}
\end{equation}
The Clifford normalization in \eqref{eq:charges} makes the coefficient in both identities exactly one.

\section{The local inputs and their proof dependencies}\label{sec:inputs}
Fix a proper partition $\pi=(n_1,\ldots,n_k)$ of $N$, $k\ge2$. In an intrinsic stationary-flag branch write
\begin{equation}
 X_i=Z_i+A_i+Y_i,\qquad Z_i=\operatorname{diag}(z_{1i}I_{n_1},\ldots,z_{ki}I_{n_k}),
 \qquad\sum_an_az_a=0.
\end{equation}
The $A_i^{(a)}$ are traceless internal matrices, and $Y$ is off-block with the intrinsic normal constraint $B_Z^*Y=0$, where $B_Q\xi=(i[\xi,Q_i])_i$. Put
\begin{equation}
 s^2=\sum_an_a|z_a|^2,\quad
 \alpha_a^2=\sum_i\Tr((A_i^{(a)})^2),\quad
 r_{ab}=|z_a-z_b|,\quad r_* =\min_{a<b}r_{ab},\quad W_p=\sum_{a<b}r_{ab}^{-p}.
\end{equation}
A local tube has $\alpha_a+\alpha_b\le\kappa r_{ab}$ and $|Y|<\sigma r_*$. The first condition is pairwise relative; it does not bound an unrelated large block by the smallest gap elsewhere. All local constants below may depend on fixed $N$, but not on $r_{\max}/r_*$.

\subsection{L1 Exact geometry and physical projection}
The intrinsic slice has internal kinetic coefficient exactly one, exact physical density $J_{\phys}=\det M$, and exact same-branch unitary transitions. Flattening by $J_{\phys}^{1/2}$ includes its differentiated scalar term. A normalized even cutoff vacuum $U_\sigma$ defines
\begin{equation}
 u=U_\sigma f+v,\qquad U_\sigma^*v=0.\label{eq:decomp}
\end{equation}
It preserves the branch form domain with padded slow and fast supports. For every slow scalar multiplier $m\ge0$,
\begin{equation}
 \int m|u|^2=\int m\bigl(|f|^2+\norm{v}_{\rm fiber}^2\bigr).\label{eq:masssplit}
\end{equation}
At $A=0$, eight negative spin modes are occupied for each bifundamental color. The determinant character restricts trivially to each internal $SU(n_a)$; the normal-ordering offset gives zero block-center charge, and the positive Gaussian is invariant. Equivariant internal-radial transport preserves this residual identification. Thus the low coefficient is physical for every internal $SU(n_a)$ factor. The complement need only obey total residual equivariance and can be colored internally. Appendix A derives the metric, density, projector and transition identities; Appendix I treats closed forms and the retained Clifford module. Spatial equivariance, the anchor connection, and the two-block reference model used for the singlet refinement are checked separately in Section~\ref{sec:singletmodel} and Appendix~\ref{app:model}.

\subsection{L2 Source and fast energy control}
Let $H_{\exc}=\sum_e(H_{{\rm fast},e}-E_{0,e})$, with the frozen ground energy subtracted pair by pair. On the cutoff complement,
\begin{equation}
 H_{\exc}[v]\ge g r_*\norm{v}_{\rm fiber}^2\label{eq:gap}
\end{equation}
for a fixed $g>0$ beyond a sufficiently large radius. This is a fast-fiber inequality at each slow coordinate, integrated over slow coordinates when used as a full form. Set $H_s=T_C+H_I$, $H_I=\sum_ah_{n_a}$, omitting singleton internal factors. $T_C$ keeps the normal connection on full fibers. On low coefficients it uses the internal-radial identification with the $A=0$ line.

The actual leading source is $S=H_1U$, and the actual diagonal multiplier $W$ obeys
\begin{equation}
 \norm{W-S^*H_{\exc}^{-1}S}\le C_N\kappa W_2.\label{eq:source}
\end{equation}
Only a shifted fast excitation inverse occurs here. It is not an inverse of $H_I$ or of the full interacting Hamiltonian. The frozen energy is paid separately:
\begin{equation}
 |E_0[u]|\le C_NR_*^{-3}H_I[u]+C_N\kappa W_2[u].\label{eq:e0}
\end{equation}
Appendix B proves a commutator-square bound for the kinetic-aware frozen frequencies. Appendix C gives the centered cancellation, its first-jet condition, graph excitation denominators, and complete slow-form estimates. Appendix D computes the actual adapted $H_1,H_2$ coefficients and their $O(\kappa W_2)$ comparison.

At $A=0$, one pair source has excitation energy $4r_e$, while a triangle source has energy $2(r_e+r_f+r_g)$. The summed triangle denominator is essential when an opposite pair is much closer than a distant block. Freezing $A$ never permits discarding internal derivatives of the moving vacuum.

\subsection{L3 Colored local reserve and mixed ledger}
There is a genuine postcompletion nonnegative square $J[v]$, and
$B[v]=H_s[v]+H_{\exc}[v]+J[v]$ is nonnegative. For any prescribed small local losses, the actual mixed ledger is
\begin{align}
 \hh_N[U_\sigma f+v]\ge{}&\hh_N[U_\sigma f]
 +c_sH_s[v]+c_JJ[v]+(1-d-\eta)H_{\exc}[v]
 +2\Rea\ip{Sf}{v}\nonumber\\
 &-a_\eta H_s[f]-b_\eta W_2[f],\label{eq:mixed}
\end{align}
where $c_s,c_J>0$, and
\begin{align}
 a_\eta&\le C_\eta(\sigma^2+R_*^{-3}),\\
 b_\eta&\le C_\eta(\kappa+\kappa^2+\sigma^2+R_*^{-3}).
\end{align}
The diagonal part is $H_s^{\rm Berry}+W+E_0$ plus errors bounded by
\begin{equation}
 \theta H_s+C_\theta(\kappa+\sigma+\sigma^2+R_*^{-3})W_2+C_NR_*^{-3}H_s.
\end{equation}
The Berry comparison costs
\begin{equation}
 |H_s^{\rm Berry}-H_s|\le\theta H_s+C_N(\theta^{-1}\kappa^4+\kappa)W_2.
\end{equation}
Normalized cutoff errors decay as a polynomial times $\exp(-c\sigma_0^2R_*^3)$ with rootwise weights. Appendices E--I give the connection estimates, nonflat Berry comparison, averaged colored-charge identity, complementary reserve, mixed errors, cutoff bounds and form-domain argument. The square $J$ is the actual retained postcompletion square, not a second copy of the original Gauss square.

\subsection{L4 Finite incidence localization}
Appendix J supplies a finite-sheet, gauge-invariant localization near the commuting cone. With fixed tree endpoint $\kappa_{\rm tree}$, logarithmic width $L$, inner ratio $\delta=\kappa_{\rm tree}e^{-L}$, and a positive resulting gap $\Delta_N$, its leaf supports satisfy
\begin{equation}
 \alpha_a+\alpha_b\le\beta r_{ab},\qquad
 r_*\ge\Delta_Ns,\qquad s\le\rho_N\le D_Ns.\label{eq:shapes}
\end{equation}
The normalized incidence weights $w_j$ obey $\sum_jw_j^2=1$. For any desired $e>0$, a later sufficiently thin near-cone neighborhood has
\begin{equation}
 E:=\sum_j|\nabla_Xw_j|^2\le\sum_jw_j^2E_j+e\rho_N^{-2}.\label{eq:imsdefect}
\end{equation}
This is a summed inequality, not a leafwise normalized identity. The path cost has a slow-only majorant
\begin{align}
 E_j\le M_\pi&=\gamma\left[
 \sum_{n_a>1}\alpha_a^{-2}\mathbf1_{\{\alpha_a\ge T(s)\}}
 +\sum_{I}s_I^{-2}\right],\label{eq:majorant}\\
 \gamma&=C_NL^{-2},\qquad T(s)=d_N\delta\Delta_Ns,\qquad M_\pi\le D_{\rm path}s^{-2}.
\end{align}
The second sum runs over proper unions of at least two final blocks; $s_I$ is their relative center radius. The constant in $\gamma$ is independent of $L,\delta$. The possibly large $D_{\rm path}$ is used only on complementary vectors and is paid at a later exterior radius.

\subsection{L5 Scalar outer bounds}
For $\mathcal B(X)=\sum_{i<j}\norm{[X_i,X_j]}_F^2$ and $t=\mathcal B/\rho_N^4$, an outer cutoff has gradient cost bounded by $C\mathcal B/\rho_N^6$. On every padded proper branch,
\begin{equation}
 \mathcal B\le C(V_I+s^2|Y|^2),\qquad
 V_I\le H_I+C_N\sum_a\alpha_a,
\end{equation}
as weighted forms, including on colored internal modules. Positive oscillator estimates give
\begin{equation}
 \norm{Yg}^2\le C_N\sum_er_e^{-2}H_{\exc,e}[g]
                       +C_N\sum_er_e^{-1}\norm{g}^2.\label{eq:osc}
\end{equation}
These elementary potential and oscillator bounds are used explicitly below. They require the full retained oscillator and colored internal forms, not merely their norm gaps.

\section{Recovering the all vector local reserve}
\begin{lemma}\label{lem:reserve}
Under L1--L3, for every $\varepsilon>0$ one can choose fixed positive tube widths and an exterior threshold such that
\begin{equation}
 \hh_N[U_\sigma f+v]\ge(1-\varepsilon)H_s[f]-\varepsilon W_2[f]+bB[v]\label{eq:reserve}
\end{equation}
for some $b>0$, uniformly in unrelated gap ratios.
\end{lemma}
\begin{proof}
Choose $b<c_s,c_J$ and $d+\eta+b<1/2$. Reserve $bH_{\exc}[v]$ before completing the remaining source square. Since $S$ is orthogonal to the uncut fast kernel,
\begin{equation}
 (1-d-\eta-b)H_{\exc}[v]+2\Rea\ip{Sf}{v}
 \ge-\frac{\ip{f}{S^*H_{\exc}^{-1}Sf}}{1-d-\eta-b}.
\end{equation}
The virtual multiplier is at most $C_NW_2$, so replacing its coefficient by one costs at most $2C_N(d+\eta+b)W_2$. Combine this with \eqref{eq:source}, \eqref{eq:e0}, and the diagonal, Berry and mixed errors in L3. First assign small $\theta,d,\eta,b$, then choose $\kappa$, then fixed $\sigma$ and a smaller cutoff width $\sigma_0$, and finally enlarge $R_*$. Every loss fits $\varepsilon(H_s+W_2)$.

The complement was defined using $U_\sigma$, while source completion uses the uncut kernel. This does not invalidate completion: $S$ lies in the uncut complementary range, and the nonnegative quadratic form can be completed there for any $v$. The cutoff-complement gap is obtained from the exponentially small angle between the cut and uncut vacuum lines. The result retains all of $bB[v]$.
\end{proof}
For fixed $N$ there are finitely many partition types, so take finite maxima of pre-parameter constants and minima of admitted widths and reserves. Their uniformity in unrelated gap ratios is required before any logarithmic hierarchy is chosen. Merely knowing a local Schur-infimum inequality would not retain the complementary mass needed globally.

\section{Shape independent Hardy payments}
Set $h_0=49/4$. In the mass metric $\sum_an_a|\dd z_a|^2$, the center dimension is $D_\pi=9(k-1)$. Radial integration by parts, valid for Hilbert-valued sections with a metric connection, yields
\begin{equation}
 T_C[f]\ge h_\pi\int s^{-2}|f|^2,
 \qquad h_\pi=\frac{(D_\pi-2)^2}{4}\ge h_0.\label{eq:centerhardy}
\end{equation}
For compact Dirichlet sections there is no boundary term; form closure extends the inequality. One can prove the scalar constant by expanding
$\int|\partial_rf+(D_\pi-2)f/(2r)|^2r^{D_\pi-1}\dd r\ge0$.
Metric compatibility gives the same real cross term for a unitary connection.

The linear map $Z\mapsto z_a-z_b$ has squared norm $\nu_{ab}=1/n_a+1/n_b$. Orthogonal relative-coordinate Hardy and $T_{ab}\le T_C$ give
\begin{equation}
 \int r_{ab}^{-2}|f|^2\le (h_0\nu_{ab})^{-1}T_C[f],
 \qquad W_2[f]\le K_\pi T_C[f],\quad
 K_\pi=h_0^{-1}\sum_{a<b}\nu_{ab}^{-1}.\label{eq:rootpayment}
\end{equation}
Thus $K_N=\max_\pi K_\pi$ is finite and contains no inverse shape gap. For any proper union $I$ of at least two blocks, the orthogonal relative-center subspace has dimension $9(|I|-1)$, whence
\begin{equation}
 \int s_I^{-2}|f|^2\le h_0^{-1}T_C[f].\label{eq:unionpayment}
\end{equation}
All other coordinates and fermionic variables remain spectators.

\begin{lemma}[Internal soft Hardy]\label{lem:soft}
Suppose a lower-rank physical form has the exterior coefficient $9$. Then there are $A_n<\infty,T_n$ such that for $T\ge T_n$,
\begin{equation}
 \hh_n[g]\ge a_0\int_{\alpha\ge T}\alpha^{-2}|g|^2
                  -A_nT^{-4}\norm{g}^2,\qquad a_0=5/2.\label{eq:soft}
\end{equation}
\end{lemma}
\begin{proof}
Start with $w_T^0(\alpha)=\alpha^2(T^4+\alpha^4)^{-1/2}$, and multiply by a fixed exterior cutoff $\beta$ that equals one beyond its annulus. Put $w_T=\beta w_T^0$ and $v_T=(1-w_T^2)^{1/2}$. Exact IMS, nonnegativity of $\hh_n[v_Tg]$, and exterior Hardy for $w_Tg$ imply
\begin{equation}
 \hh_n[g]\ge(9-4)\int\alpha^{-2}w_T^2|g|^2-O(T^{-4})\norm{g}^2.
\end{equation}
Away from the fixed annulus,
$|\nabla w_T|^2+|\nabla v_T|^2\le4\alpha^{-2}w_T^2$;
the annular derivatives are $O(T^{-4})$. For large $T$, $w_T^2\ge1/2$ on $\alpha\ge T$, proving the claim. It remains valid through $\alpha=0$ and does not require an internal gap.
\end{proof}
The cost $\gamma\alpha^{-2}\mathbf1_{\alpha\ge T}$ is therefore paid by $(\gamma/a_0)\hh_n+(\gamma A_n/a_0)T^{-4}$. Apply this only to physical low coefficients, never to a colored complement.

\section{Incidence localization and exact global IMS}
For clarity, the geometry behind L4 is recorded here; its full derivation is Appendix J. A commuting configuration has cluster scales $a_B$ and regularized separation scales
\begin{equation}
 d_B^{-2}=\sum_{j\notin B}(|\lambda_j-c_B|^2+a_B^2)^{-1}.
\end{equation}
Use select/reject angles
$\vartheta_B=\Theta(\log(a_B/(\delta d_B))/L)$ and their cosine/sine factors. Thin selected clusters are laminar. A final partition has the symmetric leaf weight given by the product of selected final-block cosines and rejected proper-union sines. Crossing omitted nodes have their sine identically one on the nonzero leaf support. Its derivative cost is $C_NL^{-2}a_B^{-2}$ on transition regions.

The decision tree orders proper subsets by decreasing cardinality and skips a subset if a previously selected block overlaps it; otherwise it splits into select and reject children. The elementary pointwise identity
\begin{equation}
 |\nabla(a\cos\vartheta)|^2+|\nabla(a\sin\vartheta)|^2
 =|\nabla a|^2+a^2|\nabla\vartheta|^2
\end{equation}
shows recursively that the leaf square sum is one and that its total gradient cost is the leaf-weighted sum of visited angle costs. Laminarity identifies each effective visited node as either a final block or a proper union of final blocks; a crossing node has zero transition derivative on the nonzero leaf support. This proves the pathwise accounting used by the symmetric leaf formula.

Off the commuting cone, keep every admitted ordered stationary flag, rather than asserting a globally unique projector. For a sheet $j$ of $k$ blocks let
\begin{equation}
 b_j=(k!)^{-1/2}\chi_Yq_\pi,\qquad
 Q=\sum_jb_j^2,\qquad w_j=b_jQ^{-1/2}.
\end{equation}
The factor $k!$ accounts for ordering; there is no additional $N!$ factor. Padded supports, a nondegenerate flag Hessian, and a finite rank-only multiplicity bound give smooth sums even with monodromy. On the whole commuting cone the identities $Q=1$, $\dd Q=0$, and $E=\sum_jw_j^2E_j$ hold for ambient derivatives, including singular strata. Finite-sheet continuity after fixing the hierarchy yields \eqref{eq:imsdefect} in a sufficiently thin neighborhood.

One quantitative support choice is
\begin{equation}
 \beta=\frac{2\kappa_{\rm tree}}{\sqrt{1-2\kappa_{\rm tree}^2}},\quad
 M_N=\frac{N\sqrt{1+\beta^2}+1}{1-\beta},\quad
 \Delta_N=\frac{(\delta/M_N)^{N-2}}{\sqrt{N(1+\beta^2)}}.
\end{equation}
Padding decreases this positive gap if necessary. Since
$\rho_N^2=s^2+\sum_a\alpha_a^2+|Y|^2$ and $r_{ab}\le\sqrt2s$, one can use $D_N^2=1+2N\beta^2+2\sigma^2$ before the harmless enlargement for padded supports.

Choose $0<t_0<t_1$ so that the near-cone cutoff lies inside that neighborhood and all contributing fast support cutoffs equal one. Write
\begin{equation}
 \chi=\cos\vartheta_{\rm out}(t),\quad
 \chi_O=\sin\vartheta_{\rm out}(t),\quad
 u_j=\chi w_ju,\quad u_O=\chi_Ou,\quad
 E_{\rm out}=|\nabla\vartheta_{\rm out}(t)|^2.
\end{equation}
The exact scalar IMS identity and L4 give
\begin{align}
 \hh_N[u]\ge{}&\sum_j\left\{\hh_N[u_j]
 -\int(E_{\rm out}+E_j+e\rho_N^{-2})|u_j|^2\right\}\nonumber\\
 &+\hh_N[u_O]-\int E_{\rm out}|u_O|^2.\label{eq:globalims}
\end{align}
Every error is assigned on both sides of each split. No agreement between projections on different sheets is required.

After flattening, project each $u_j$ as in \eqref{eq:decomp}. The normalized weights and $\chi$ can depend on $Y$. The projected coefficient has only the leaf's slow support; $U_\sigma f$ may leave the original near-cone set, but stays in its padded branch. This is why all majorants and local estimates are required on the entire padded fast fiber. Slow support has strict margins, so zero extension to canonical center variables is legitimate for the Hardy inequalities. The projection bound is in the form norm only.

\section{Paying the path and outer errors}
Assume the theorem at all lower ranks. Fix a small $a>0$. Equations \eqref{eq:soft} and \eqref{eq:unionpayment} applied to \eqref{eq:majorant} give, after choosing $L$ large enough,
\begin{equation}
 M_\pi[f]\le aH_I[f]+aT_C[f]+C_{\rm soft}\int s^{-4}|f|^2.\label{eq:pathlow}
\end{equation}
For example $\gamma\le aa_0$ and $\gamma2^N/h_0\le a$ suffice. The soft constant contains the now fixed inverse powers of $d_N\delta\Delta_N$ and lower-rank constants; its large size affects only a later radius. Enforce $T(s)\ge\max_{n<N}T_n$ at that radius.

For $v$, use the other part of L4:
\begin{equation}
 M_\pi[v]\le D_{\rm path}\int s^{-2}\norm{v}_{\rm fiber}^2.
\end{equation}
This is absorbed by any assigned fraction of $H_{\exc}[v]$ once $s$ is large, since $r_*\ge\Delta_Ns$. For both low and complementary pieces,
$e\rho_N^{-2}\le es^{-2}$; the latter is a slow multiplier, so \eqref{eq:masssplit} applies exactly.

Differentiating the quartic polynomial gives
\begin{equation}
 |\nabla\mathcal B|\le C_N\rho_N\sqrt{\mathcal B},\qquad
 |\nabla t|\le C_N\sqrt{\mathcal B}/\rho_N^3,\qquad
 E_{\rm out}\le C_{\rm out}\mathcal B/\rho_N^6.
\end{equation}
On each branch, L5 bounds this by
$C_{\rm out}'(s^{-6}V_I+s^{-4}|Y|^2)$ on the full padded support. Because the $s^{-6}$ weight commutes with internal differentiation and $\sum\alpha_a\le C_Ns$, the internal contribution costs
\begin{equation}
 C_{\rm out}'s^{-6}H_I+C_{\rm out}''s^{-5}
\end{equation}
as fiberwise integrated forms, for low or colored complementary inputs. The oscillator estimate \eqref{eq:osc} yields
\begin{align}
 \int s^{-4}|Yv|^2&\le C_D\int s^{-6}H_{\exc}[v\text{ at fixed slow variables}]
                         +C_D\int s^{-5}\norm{v}_{\rm fiber}^2,\\
 \int s^{-4}|YU_\sigma f|^2&\le C_D\int s^{-5}|f|^2.
\end{align}
The $V_I$ multiplier has zero low/complement cross term. The positive $|Y|^2$ cross is bounded by the two diagonal pieces, with a factor two. At sufficiently large radius these estimates imply
\begin{equation}
 \int E_{\rm out}|u_j|^2\le aH_I[f]
     +C_{\rm outer}\int s^{-5}|f|^2+\frac b4B[v].\label{eq:outerpay}
\end{equation}
A norm gap alone would not control the unbounded $|Y|^2$ multiplier; the retained full excitation form is indispensable.

On the outside branch $t\ge t_0$, the bosonic potential is at least $v_Nt_0\rho_N^4$, whereas the fermionic norm is at most $F_N\rho_N$. The outer gradient cost is lower order. Hence, beyond a larger radius,
\begin{equation}
 \hh_N[u_O]-\int E_{\rm out}|u_O|^2
       \ge c\int\rho_N^{-2}|u_O|^2.\label{eq:outside}
\end{equation}
There is no branch-weight cost on $u_O$ in \eqref{eq:globalims}.

\section{Baseline all vector theorem and rank induction}
\begin{theorem}[Baseline all-vector global implication]\label{thm:global}
Assume Hypothesis~\ref{hyp:realization} and L1--L5 for every proper partition at each fixed finite rank, with the parameter uniformities specified above. Then for every integer $N\ge2$ and every $0<c<49/4$, some finite $R_{N,c}$ satisfies
\begin{equation}
 \hh_N[u]\ge c\int\rho_N^{-2}|u|^2,\qquad u\in\QQ_{N,R_{N,c}}.\label{eq:main}
\end{equation}
In physical units the right side is multiplied by $R_{\rm BFSS}/2$.
\end{theorem}
\begin{proof}
Fix $N,c$ and take
\begin{equation}
 0<a<\min\left\{\frac1{12},\frac{h_0-c}{8(h_0+1)}\right\},
 \qquad\varepsilon(1+K_N)\le a.
\end{equation}
Lemma~\ref{lem:reserve} and \eqref{eq:rootpayment} give
\begin{equation}
 \hh_N[u_j]\ge(1-a)(H_I[f]+T_C[f])+bB[v].
\end{equation}
Choose $e\le a$ and $L$ for \eqref{eq:pathlow}. Enlarge the final radius until
$C_{\rm soft}s^{-4}+C_{\rm outer}s^{-5}\le as^{-2}$,
the complementary path and defect costs each use at most $(b/8)H_{\exc}[v]$, and \eqref{eq:outerpay} holds. Subtracting all assigned errors in \eqref{eq:globalims} gives the single ledger
\begin{align}
 \hh_N[u_j]&-\int(E_{\rm out}+E_j+e\rho_N^{-2})|u_j|^2\nonumber\\
 &\ge(1-3a)H_I[f]+(1-2a)T_C[f]
          -2a\int s^{-2}|f|^2+\frac b2B[v].\label{eq:ledger}
\end{align}
The complementary payments are exactly $b/4+b/8+b/8=b/2$. There is no second Gauss reserve. Positivity of $H_I$ and \eqref{eq:centerhardy} give low coefficient
$[(1-2a)h_0-2a]>c$.
Increase the radius so that $(b/2)g\Delta_Ns\ge cs^{-2}$. Then the last term controls $c\int s^{-2}\norm{v}_{\rm fiber}^2$. Thus the whole right side of \eqref{eq:ledger} is at least
\begin{equation}
 c\int s^{-2}(|f|^2+\norm{v}_{\rm fiber}^2)
 =c\int s^{-2}|u_j|^2\ge c\int\rho_N^{-2}|u_j|^2.
\end{equation}
This exact slow-multiplier orthogonality prevents replacing full localized mass by low mass. Sum over leaves and use \eqref{eq:outside}, \eqref{eq:globalims}, and the square-sum identity. Gauge averaging, compact exhaustion and exterior form closure extend the result from the physical core to the stated domain.

At $N=2$ the only partition is $1+1$; there are no internal forms or selected-block costs. The same proof is a base case with center dimension nine and no inductive input. If all ranks below $N$ are proved, each nontrivial block of a proper partition has smaller rank, giving exactly Lemma~\ref{lem:soft}. This is ordinary finite-rank induction.
\end{proof}
The fully split leaf has center Hardy coefficient $[9(N-1)-2]^2/4$, but the common bound is limited by possible two-block channels. No sharpness theorem is inferred from this bottleneck. Simultaneous $2+2$, $2+1+1$, and $2+2+1$ clusters are covered by the symmetric leaf construction; no relative ordering of disjoint internal sizes is imposed.

\subsection{Acyclic parameter order}
The proof requires the following order, rather than mutually dependent smallness choices.
\begin{enumerate}
\item Fix $N,c,a$, lower-rank soft constants, and the rank-only $K_N$.
\item Fix $\varepsilon\le a/(1+K_N)$ and the local Young allocations, completion loss, and retained $b$.
\item Choose the local endpoint $\kappa$, a smaller tree endpoint, then fixed fast width $\sigma$ and vacuum cutoff width $\sigma_0$.
\item Choose $L$ for the $L^{-2}$ payment, then $\delta$, the positive $\Delta_N$, padded supports and their finite continuity constants.
\item Choose the near-cone thickness for the summed defect, then $t_1,t_0$ and the finite outer-cutoff constant.
\item Finally enlarge the physical exterior radius to meet every gap, lower-rank threshold, Gaussian tail, $s^{-4}$, $s^{-5}$ and complementary absorption requirement.
\end{enumerate}
Inverse powers of the small hierarchy gaps occur only in constants already fixed before the final radius. The local $\varepsilon W_2$ cost was paid by center Hardy before the hierarchy was chosen, avoiding a circular $\varepsilon\Delta_N^{-2}$ requirement.

\section{Rotation invariance and internal uncertainty}\label{sec:rotation}
The extra input for the singlet improvement is rotational representation theory on complete physical internal forms. It does not replace them by internal ground states. We first record an older theorem and then a quantitative consequence of its rotation identity.

\begin{theorem}[Kernel invariance of Hasler and Hoppe]\label{thm:invariance}
For the standard physical realization in Hypothesis~\ref{hyp:realization}, every existing $L^2$ zero mode of $h_N$ is invariant under $\operatorname{Spin}(9)$, at every fixed finite $N$.
\end{theorem}
This is \cite[Theorem 1(a), Section 2]{HHI}, building on \cite{SSI}; it is not a new assertion of this draft. Its realization and domain are important. Appendix~\ref{app:rotation} recalls the cutoff argument: linear growth of the rotation primitive cancels the inverse-radius cutoff derivative, and $L^2$ shell tails suffice. No positive moment, uniqueness, or internal bound state is assumed.

For an arbitrary finite collection of physical internal blocks, write
$H_I=\sum_a\hh_{n_a}$ and $\alpha^2=\sum_a\alpha_a^2$.
Let $P_I$ be Haar averaging for their simultaneous internal $\operatorname{Spin}(9)$ action; it fixes all external spectators.
\Needspace{12\baselineskip}
\Needspace{13\baselineskip}
\begin{lemma}[Internal nonsinglet uncertainty on a ball]\label{lem:ball}
For physical form vectors supported in $\alpha\le R$,
\begin{equation}
 H_I[g]\ge\frac1{4R^2}\norm{(1-P_I)g}^2.\label{eq:ball}
\end{equation}
The estimate also holds for Hilbert-valued spectators and in the Dirichlet form closure of the internal ball. It is a radius-of-support estimate, not a positive spectral gap on the full internal space and not the operator inequality $H_I\ge\frac14\alpha^{-2}(1-P_I)$.
\end{lemma}
\begin{proof}
Use charges $\widetilde q_\alpha=q_\alpha/\sqrt2$, where $q_\alpha$ are the coefficient-$16$ charges of Section~\ref{sec:operator}. On physical vectors,
$\sum_\alpha\norm{\widetilde q_\alpha g}^2=8\hh[g]$.
For one fixed generator $j=(u,v)$, the corrected Hasler--Hoppe primitive is
\begin{align}
 \widetilde a_{j,\alpha}&=\frac{5b_{j,\alpha}+2c_{j,\alpha}}7,
 &J_j&=\sum_\alpha\{\widetilde q_\alpha,\widetilde a_{j,\alpha}\},\label{eq:rotationidentity}\\
 b_{j,\alpha}&=\frac1{16}(X_u\Gamma_v-X_v\Gamma_u)_{\alpha\epsilon}\theta_\epsilon,
 &c_{j,\alpha}&=\frac1{16}X_w(\Gamma_w\Gamma_{uv})_{\alpha\epsilon}\theta_\epsilon.
\end{align}
Color contractions and the sum over $w$ are understood. The bare $b$ alone has the wrong orbital-to-spin ratio in nine dimensions. The CAR and Clifford identities give, for this one $j$,
\begin{equation}
 \sum_\alpha\widetilde a_{j,\alpha}^2
 =\left(\frac{|X_u|^2+|X_v|^2}{32}
       +\frac{\sum_{w\notin\{u,v\}}|X_w|^2}{392}\right)I
 \le\frac{\alpha^2}{32}I.\label{eq:fixedprimitive}
\end{equation}
The expression summed over all $36$ generators instead equals $(9/28)\alpha^2I$. These are different summation scopes in the same convention. In the baseline convention the primitive must also be rescaled: $a_{j,\alpha}=\widetilde a_{j,\alpha}/\sqrt2$. Then $J_j=\sum_\alpha\{q_\alpha,a_{j,\alpha}\}$; the fixed-generator bound is $\alpha^2/64$ and the all-generator identity is $9\alpha^2/56$. Appendix~\ref{app:rotation} gives both computations and a stronger Casimir variant.

Sum the blockwise fixed-$j$ identities. On the physical compact core, Cauchy--Schwarz gives
\begin{equation}
 |\ip{g}{J_j^Ig}|\le
 2\sqrt{8H_I[g]}\sqrt{R^2\norm{g}^2/32}
 =R\sqrt{H_I[g]}\norm{g}.\label{eq:uncertainty}
\end{equation}
The coefficient is unchanged in the baseline convention because $2\sqrt{16/64}=1$. Only a fixed-generator norm has been used.

The ball and the closed form are rotation invariant. Decompose into compact-group isotypes. On each irreducible representation factor the invariant form is the identity tensored with a multiplicity-space form. Every nontrivial $\operatorname{Spin}(9)$ irrep has a highest-weight vector with $J_{12}$ eigenvalue $m\ge1/2$. Apply \eqref{eq:uncertainty} to that vector tensored with any multiplicity-space form vector; it gives $H_I\ge m^2/R^2\ge1/(4R^2)$. Invariance then gives the same inequality on the full isotype. Sum the nontrivial isotypes and use nonnegativity on the invariant part. Compact-group smoothing, finite isotypic truncation, and form closure justify the passage from the core; no unbounded rotation generator is applied to an uncontrolled form vector. Singleton blocks add no internal variables.
\end{proof}

\section{The singlet slow model and its domain}\label{sec:singletmodel}
For a two-block partition $n_1+n_2=N$, let $\mathscr H_a$ be its physical internal Hilbert space, with $\mathscr H_a=\mathbb C$ for a singleton. In the mass-normalized center variable $x\in\mathbb R^9$ put
\begin{equation}
 s=|x|,\qquad r_{12}=\nu s,\qquad
 \nu^2=\frac1{n_1}+\frac1{n_2}\le2.
\end{equation}
The reference slow space and closed form are
\begin{align}
 \mathscr K&=L^2(\mathbb R_x^9;
 \FF_{\rm rel}\otimes\mathscr H_1\otimes\mathscr H_2),\\
 q_s&=T_x+H_I,\qquad
 \QQ(q_s)=\QQ(T_x)\cap\QQ(H_I).\label{eq:modeldomain}
\end{align}
Here $T_x$ is the ordinary Hilbert-valued center Dirichlet form. The internal spaces contain their complete bosonic and fermionic variables. Residual gauge and equal-block exchange constraints restrict the admissible subspace and do not add angular states.

\subsection{Why the model has the ordinary diagonal action}
The stationary-flag functional obeys
\begin{equation}
 \Phi_{RX}(P)=\sum_{a,i}n_a^{-1}\bigl(\Tr(P_a(RX)_i)\bigr)^2
 =\Phi_X(P)
\end{equation}
for the same ordered flag $P$. This equality holds as a function on the entire flag manifold; its flag gradient and Hessian are equal as well. Hence spatial rotations preserve stationary flags, Hessian admission and incidence sheets. All scalar weights in L4 are rotation invariant. Color frames can be chosen from the flag alone and are constant along spatial rotation orbits; gauge transitions commute with spatial rotations. A global singlet therefore localizes to invariant sections without an angular partition.

At $A=0$ the two-block fast ground line has stabilizer $\operatorname{Spin}(8)$. Its genuine one-dimensional unitary representation is trivial because that group is connected and semisimple. This includes its central element $-1$; the filled determinant has $8n_1n_2$ occupied modes, an even number. The resulting global equivariant normalized anchor has zero angular Berry one-form, since $S^8=\operatorname{Spin}(9)/\operatorname{Spin}(8)$ has no invariant tangent covector. Its radial Berry form also vanishes. This gives the ordinary $T_x$ in \eqref{eq:modeldomain}.

The negative fast spectral projector along $A\mapsto tA$ is covariant. Its Kato generator $[\dot P_t,P_t]$ is covariant, so uniqueness of the transport equation makes the radial transport equivariant. The positive Gaussian and normalized scalar cutoff have the same property. The transported injection $U_\sigma$ thus intertwines the actual action with
\begin{equation}
 J_{\rm slow}=L_x+S_{\rm rel}+J_I.\label{eq:slowaction}
\end{equation}
Finite-$A$ Berry curvature is not asserted to vanish. Its comparison and the differentiated vacuum sources remain exactly the terms already paid in L2--L3. Appendix~\ref{app:model} supplies the global anchor and transition details.

\subsection{Zero extension precedes internal averaging}
The bounded projection in L1 places the actual coefficient $f$ in the slow Dirichlet form closure of its padded branch. Its zero extension consequently belongs to \eqref{eq:modeldomain}; this is the same extension needed for the baseline Hardy payments. Arbitrary distributional sections are not being extended across a boundary. The extended $f$ is supported in the larger radial tube
\begin{equation}
 \mathcal C_{\beta}=\{(x,A):\alpha_1+\alpha_2\le\beta_{\rm pad}\nu|x|\}.
 \label{eq:radialtube}
\end{equation}
The detailed leaf may have additional direction-dependent tests and need not be invariant under internal-only rotations.

Haar averaging $P_I$ acts only on the simultaneous internal variables and fermions. It commutes with $T_x$, the complete $H_I$, center radial weights, and the diagonal action after averaging. It is a contraction in the model form norm and preserves \eqref{eq:radialtube}. Thus it gives orthogonality separately in $T_x$, $H_I$, and the $s^{-2}$ norm.

The order of application is essential:
\begin{enumerate}
\item Apply the paid local ledger \eqref{eq:ledger} to the actual localized vector and its actual coefficient $f$.
\item Regard that same $f$ as its zero extension in $\mathscr K$.
\item Split $f=P_If+(1-P_I)f$ in this reference space.
\item Estimate the full-space model form $q_s$ on these summands.
\end{enumerate}
Averaging may leave the detailed leaf but cannot leave the larger radial tube. No local Hamiltonian comparison, detailed-leaf multiplier bound, or projection theorem is reapplied to an averaged component outside its original leaf. Both components remain diagonal singlets when $f$ is one. Fubini, compact slow exhaustion, and closed-form approximation justify the fiberwise estimates below.

\begin{proposition}[Two-block singlet comparator]\label{prop:singletslow}
Choose the padded endpoint width initially so that $\beta_{\rm pad}\le1/12$. A diagonal singlet $f$ in \eqref{eq:modeldomain} supported in \eqref{eq:radialtube} satisfies
\begin{equation}
 (T_x+H_I)[f]\ge h_*\norm{s^{-1}f}^2,
 \qquad h_*:=\frac{121}{4}.\label{eq:singletslow}
\end{equation}
\end{proposition}
\begin{proof}
Set $f_0=P_If$ and $f_1=(1-P_I)f$. The relative sixteen-Majorana module is
$\FF_{\rm rel}=\mathbf{44}\oplus\mathbf{84}\oplus\mathbf{128}$, as recalled in \cite[Section 2.4]{LY}. Scalar harmonics on $S^8$ have highest weight $(l,0,0,0)$. Among those three summands only $\mathbf{44}$ matches a harmonic, at $l=2$. Thus the internal-invariant diagonal singlet $f_0$ has angular eigenvalue $2(2+7)=18$. This applies to its arbitrary internal invariant multiplicity space, not only an internal ground state. Radial Hardy gives
\begin{equation}
 T_x[f_0]\ge(49/4+18)\norm{s^{-1}f_0}^2.\label{eq:angularzero}
\end{equation}
For each fixed $x$, $f_1$ has internal support $\alpha\le\beta_{\rm pad}\nu s$. Lemma~\ref{lem:ball}, followed by Fubini and form closure, gives
\begin{equation}
 H_I[f_1]\ge\frac1{4\beta_{\rm pad}^2\nu^2}\norm{s^{-1}f_1}^2
 \ge18\norm{s^{-1}f_1}^2.\label{eq:angularnonzero}
\end{equation}
Ordinary center Hardy supplies $49/4$ on $f_1$ as well. Orthogonal additivity proves \eqref{eq:singletslow}. Half-integral internal representations, gapless continua, and entanglement between blocks are included in Lemma~\ref{lem:ball}. Equal-block exchange constraints only restrict this space. At $N=2$ both internal factors are scalar and $f_1=0$.
\end{proof}
For $k\ge3$, the center dimension is at least $18$, and \eqref{eq:centerhardy} alone gives $h_\pi\ge64>h_*$. No angular classification or internal averaging is needed on those leaves.

\section{Singlet exterior estimate}\label{sec:sharptheorem}
\Needspace{11\baselineskip}
\begin{theorem}[Singlet global implication]\label{thm:singlet}
Assume the same realization and local analytical inputs as in Theorem~\ref{thm:global}, with the equivariant constructions and Dirichlet model identification detailed in Sections~\ref{sec:rotation}--\ref{sec:singletmodel} and Appendices~\ref{app:rotation}--\ref{app:model}. For every fixed $N\ge2$ and every $0<c<121/4$, there is a finite $R^{\rm sing}_{N,c}$ such that
\begin{equation}
 \hh_N[u]\ge c\norm{\rho_N^{-1}u}^2,
 \qquad u\in\QQ_{N,R^{\rm sing}_{N,c}}^{\operatorname{Spin}(9)}.
 \label{eq:singletmain}
\end{equation}
The superscript restricts the physical exterior form domain to rotation singlets. No coefficient above $49/4$ is asserted for arbitrary nonsinglet vectors.
\end{theorem}
\begin{proof}
Keep the lower-rank all-vector Theorem~\ref{thm:global} for all internal soft-Hardy payments. There is no induction using the stronger conclusion. Given $c<h_*$, choose
\begin{equation}
 0<a<\min\left\{\frac1{12},\frac{h_*-c}{2(3h_*+2)}\right\}.
\end{equation}
At the initial local-width selection, also require $\beta_{\rm pad}\le1/12$ for two-block leaves. Choose the tree endpoint small enough to retain this padded restriction. This precedes $L,\delta,\Delta_N$ and introduces no new inverse-gap dependency; follow the acyclic parameter order of the baseline thereafter.

Apply \eqref{eq:ledger} before internal averaging. Its low part has the smaller internal coefficient $1-3a$. Discarding only the nonnegative extra $aT_C$ gives
\begin{align}
 (1-3a)H_I[f]+(1-2a)T_C[f]-2a\norm{s^{-1}f}^2
 &\ge\bigl[(1-3a)h_*-2a\bigr]\norm{s^{-1}f}^2\nonumber\\
 &>c\norm{s^{-1}f}^2.\label{eq:sharpledger}
\end{align}
For two blocks this uses the model comparator after zero extension. For three or more blocks it uses $h_\pi\ge64$ and $H_I\ge0$. Replacing $h_0$ mechanically while keeping the coefficient $1-2a$ would be incorrect, since $H_I$ now contributes to coercivity.

The colored complement receives no rotational improvement. Increase the final radius until its retained $(b/2)H_{\exc}$ controls $c\norm{s^{-1}v}^2$, exactly as before. The slow mass identity, $\rho_N\ge s$, outside-cone estimate, and incidence square sum then prove \eqref{eq:singletmain}. Gauge averaging, rotation averaging, compact exhaustion and exterior form closure preserve the stated domain. Every IMS, Berry and Schur loss remains charged once, and no internal spectral gap or bound-state projection was introduced.
\end{proof}
This theorem inherits all substantive analytical dependencies of Appendices A--J. The rotation argument closes an additional implication if those baseline inputs hold; it does not independently certify them. The number $121/4$ is the threshold supplied by this comparator, not a proved sharpness result for the full Hamiltonian.

\section{Moments of actual zero modes}
\begin{corollary}\label{cor:moments}
Under Theorem~\ref{thm:singlet}, every $\Psi\in\mathscr H_N$ with $\hh_N[\Psi]=0$ satisfies
\begin{equation}
 \int\rho_N^m|\Psi|^2<\infty\qquad(0\le m<9).
\end{equation}
No existence of such a $\Psi$ is asserted.
\end{corollary}
\begin{proof}
By Theorem~\ref{thm:invariance}, $\Psi$ is a rotation singlet. Bounded radial multipliers preserve that property. Choose a fixed smooth radial inner cutoff $\eta$ supported in the exterior of \eqref{eq:singletmain}, equal to one farther out. For $a_p\ge1$, use bounded truncated powers $g_M(\rho)=\min(\rho^{a_p},M^{a_p})$, or smooth Lipschitz approximants, so that $|\nabla g_M|\le a_pg_M/\rho$. Apply \eqref{eq:singletmain} and \eqref{eq:zeroidentity} to $\eta g_M\Psi$. For any $\zeta>0$, Young's inequality gives
\begin{equation}
 \bigl[c-(1+\zeta)a_p^2\bigr]\norm{\rho_N^{-1}\eta g_M\Psi}^2
 \le C_{\eta,\zeta,a_p}\norm{\Psi}^2,
\end{equation}
with the right side independent of large $M$: all derivatives of $\eta$ lie in one fixed annulus. Choose $a_p^2<c<h_*$ and $\zeta$ small, then let $M\to\infty$ by monotone convergence. This proves exterior integrability with exponent $m=2a_p-2$. The bounded interior is harmless since $m\ge0$ and $\Psi\in L^2$. Every $0\le m<9$ is obtained by taking $a_p=(m+2)/2<11/2$ and choosing $c$ between $a_p^2$ and $121/4$.
\end{proof}
The endpoint $m=9$ is not obtained because this coercive margin vanishes there. No pointwise asymptotic expansion or nonzero leading tail is proved. Polchinski's expected free $l=2$ two-cluster behavior $r^{-9}$, with radial density $r^8\dd r$, already predicts the same below-nine threshold \cite{P}. The candidate contribution is the actual-state implication through the stated all-channel form estimates, not that asymptotic prediction.

At fixed rank the corollary includes the eighth ordinary radial moment. Every coordinate polynomial of degree at most four, with bounded fermion-matrix coefficients, times $\Psi$ belongs to $L^2$; polynomial expectations through degree eight are absolutely finite. A gauge-invariant such observable therefore has a well-defined centered physical $L^2$ source $(O-\langle O\rangle)\Psi$. This gives no rank-uniform source norm, reduced-resolvent estimate at zero energy, or inverse spectral moment. Momentum insertions, all Ward identities, all bootstrap levels, numerical radius bounds, large-$N$ tightness, scattering completeness and full soft theorems require additional arguments.

\section{Reproducibility and review boundaries}
The accompanying bundle contains this editable \LaTeX\ source, twelve mathematical source notes, eleven executable regression programs, a runner, dependency versions, finite-check outputs, a separate Chinese summary, and SHA-256 manifests. It excludes private coordination and review transcripts, unrelated research, email, and downloaded third-party papers. The source notes retain the mathematical equations and local scope boundaries; procedural review labels and superseded status commentary are removed. Provenance records the exact input hashes and included sections.

The regression suite covers finite-matrix geometry, adapted Gauss Taylor coefficients, frozen-energy identities, source excitation bookkeeping, square completion, regularized tree/IMS examples, and the scalar global budget. The global budget uses exact rational arithmetic over 2,693 proper integer partition types through $N=20$, six target coefficients including $49/4-10^{-6}$, the retained complementary fraction, and moment arithmetic. Floating-point checks include hierarchical gap stress cases. Two added programs check exact Clifford coefficients for all 36 rotation generators, the reciprocal normalization conversion, the two-block width payment, and the revised singlet ledger and moment arithmetic. Successful execution is evidence for those finite calculations only.

A specialist review should check, in order:
\begin{enumerate}
\item the trace/CAR normalization and realization assumptions in Section~\ref{sec:operator};
\item exact geometry and density, Appendix A; the frozen energy bound, Appendix B;
\item centered first jets and actual adapted source cancellation, Appendices C--D;
\item rootwise differential estimates, full colored-charge comparison, and the non-double-counted reserve, Appendices E--I;
\item normalized incidence localization at singular commuting strata and its padded domains, Appendix J;
\item the common ledger and all-vector induction;
\item the corrected rotation primitive, bounded-radius isotypic argument, equivariant vacuum line and zero-extension order, Appendices~\ref{app:rotation}--\ref{app:model};
\item the smaller retained internal coefficient in the singlet ledger and the bounded-weight passage.
\end{enumerate}
The most recent domain correction is implemented consistently: the Schur infimum and any minimizing complement belong to $\QQ(h_N)$, and projection boundedness is in the form norm, equivalently the graph norm of the closed charge stack. Operator-domain preservation is not inferred. The global lower bound does not depend on minimizer attainment.

This draft is ready to be examined, not certified as correct or submission-ready by the existence of the bundle. A future public version should reflect the named author's own review and a specialist check of the full analytical dependencies. No priority, affiliation, human endorsement, formal verification, or journal outcome is asserted.

\begingroup
\small
\begin{thebibliography}{9}
\bibitem{SS} S. Sethi and M. Stern, \emph{D-Brane Bound States Redux}, Commun. Math. Phys. 194 (1998), 675--705. \href{https://arxiv.org/abs/hep-th/9705046}{arXiv:hep-th/9705046}.
\bibitem{SSI} S. Sethi and M. Stern, \emph{Invariance Theorems for Supersymmetric Yang--Mills Theories}, Adv. Theor. Math. Phys. 4 (2000), 487--501. \href{https://arxiv.org/abs/hep-th/0001189}{arXiv:hep-th/0001189}.
\bibitem{HHI} D. Hasler and J. Hoppe, \emph{Zero Energy States of Reduced Super Yang--Mills Theories in $d+1=4,6$ and $10$ Dimensions Are Necessarily $\operatorname{Spin}(d)$ Invariant}. \href{https://arxiv.org/abs/hep-th/0211226}{arXiv:hep-th/0211226} (2002).
\bibitem{Agmon} S. Agmon, \emph{Lectures on Exponential Decay of Solutions of Second-Order Elliptic Equations}, Princeton University Press, Princeton, 1982.
\bibitem{K} A. Konechny, \emph{On Asymptotic Hamiltonian for $SU(N)$ Matrix Theory}, JHEP 10 (1998), 018. \href{https://arxiv.org/abs/hep-th/9805046}{arXiv:hep-th/9805046}.
\bibitem{HH} D. Hasler and J. Hoppe, \emph{Asymptotic Factorisation of the Ground-State for $SU(N)$-invariant Supersymmetric Matrix-Models}. \href{https://arxiv.org/abs/hep-th/0206043}{arXiv:hep-th/0206043} (2002).
\bibitem{FGHHY} J. Fr\"ohlich, G. M. Graf, D. Hasler, J. Hoppe and S.-T. Yau, \emph{Asymptotic Form of Zero Energy Wave Functions in Supersymmetric Matrix Models}. \href{https://arxiv.org/abs/hep-th/9904182}{arXiv:hep-th/9904182} (1999).
\bibitem{LY} Y.-H. Lin and X. Yin, \emph{On the Ground State Wave Function of Matrix Theory}. \href{https://arxiv.org/abs/1402.0055}{arXiv:1402.0055} (2014).
\bibitem{P} J. Polchinski, \emph{M-Theory and the Light Cone}, Prog. Theor. Phys. Suppl. 134 (1999), 158--170. \href{https://arxiv.org/abs/hep-th/9903165}{arXiv:hep-th/9903165}.
\end{thebibliography}
\endgroup

\clearpage
\appendix
\renewcommand{\theequation}{\thesection.\arabic{equation}}
\fancyhead[L]{\small Local technical derivations}
\begin{center}\Large\bfseries Technical appendices\end{center}
The first ten appendices reproduce the substantive mathematical derivations on which the BFSS specialization depends. Their subsection and displayed equation numbers are local to each source note. Historical source labels are retained for traceability, including gaps and nonsequential labels; these do not indicate omitted PDF pages. These are included for direct review, not as external certified lemmas. Their local scope boundaries should be read together with the dependency map in Section~\ref{sec:inputs}; no single appendix purports to prove the global theorem by itself. The last two appendices add the corrected rotation algebra and equivariant reference-domain construction for the singlet refinement.

The main text and appendices use typeset mathematics. The longer derivations in Appendices A--J have been transcribed from the source-level formulas, with their coefficient inventory, operator order, and local labels retained. The original editable source notes and an expression-by-expression transcription record accompany this edition for comparison. This production revision does not provide an additional mathematical certification.
\clearpage
\section{Intrinsic multiblock
geometry}\label{a-intrinsic-multiblock-geometry}

Technical derivation for review. Equation and subsection numbers in this
appendix are local to the source derivation. Historical local labels are
retained, including numbering gaps or nonsequential labels. The
accompanying source notes retain the original expressions for
comparison.

\subsection{1. Setting and precise
scope}\label{a-setting-and-precise-scope}

Fix \(N\) and a partition \(N=n_{1}+\ldots +n_{k}, k\ge 2\). Hermitian
matrix tuples have the real trace inner product. Retain the residual
group \(K=S(U(n_{1}) \times  \ldots  \times  U(n_{k}))\), including its
action on the fermions; do not divide by any internal nonabelian orbit.
Write

% DISPLAY a.1
% SOURCE X_i=Z_i+A_i+Y_i,
% SOURCE Z_i=diag(z_1i I_(n_1),...,z_ki I_(n_k)),
% SOURCE sum_a n_a z_a=0,   Tr A_i^(a)=0,
% SOURCE b_ab=z_a-z_b,   r_ab=|b_ab|,   r_*=min_(a<b) r_ab>0.
\begin{gather*}
\begin{aligned}
& X_{i}=Z_{i}+A_{i}+Y_{i}, \\
& Z_{i}=\operatorname{diag} (z_{1i} I_{n_{1}},\ldots ,z_{ki} I_{n_{k}}), \\
& \sum _{a} n_{a} z_{a}=0, \quad  \operatorname{Tr}  A_{i}^{(a)}=0, \\
& b_{ab}=z_{a}-z_{b}, \quad  r_{ab}=\vert b_{ab}\vert , \quad  r_{*}=\min _{a<b} r_{ab}>0.
\end{aligned}\notag
\end{gather*}

A is block diagonal and \(Y\) is off block. Put
\(\alpha _{a}^{2}=\sum _{i} \Vert A_{i}^{(a)}\Vert _{F}^{2}\). Assume
only

% DISPLAY a.2
% SOURCE alpha_a+alpha_b <= kappa r_ab  for every a<b.             (1)
\begin{gather*}
\begin{aligned}
& \alpha _{a}+\alpha _{b} \le  \kappa  r_{ab} \quad  \text{ for every } a<b.
\end{aligned}\tag{1}
\end{gather*}

There is no hypothesis \(\alpha _{a} \ll  r_{*}\) for a block unrelated
to a shortest pair. Set \(\kappa \le 1/100\). For the fast tube use

% DISPLAY a.3
% SOURCE |Y| <= sigma r_*,   sigma<=1/100.                         (2)
\begin{gather*}
\begin{aligned}
& \vert Y\vert  \le  \sigma  r_{*}, \quad  \sigma \le 1/100.
\end{aligned}\tag{2}
\end{gather*}

Condition (2) is deliberately a fast-coordinate condition. It is
compatible with the Gaussian widths at sufficiently large \(r_{*}\) and
does not impose the forbidden common scale on A. A smooth smaller cutoff
may use the harmonic minimum of the \(r_{ab}\) instead of the nonsmooth
literal minimum.

Let \(C\) be the Euclidean block-center tuple space, I the space of
blockwise traceless internal tuples, and \(m\) the off-block Hermitian
gauge-parameter space. Its real dimension is

% DISPLAY a.4
% SOURCE q=2 sum_(a<b) n_a n_b.
\begin{gather*}
\begin{aligned}
& q=2 \sum _{a<b} n_{a} n_{b}.
\end{aligned}\notag
\end{gather*}

For a Hermitian tuple \(Q\) let \(B_{Q} \xi =(i[\xi ,Q_{i}])_{i}\).
Define

% DISPLAY a.5
% SOURCE R=(B_Z^* B_Z)^(1/2)=direct_sum_(a<b) r_ab I_(2n_an_b),
% SOURCE E=B_Z R^(-1),   E^*E=I_m,
% SOURCE N_Z=ker B_Z^* in the off-block tuple space.
\begin{gather*}
\begin{aligned}
& R=(B_{Z}^{*} B_{Z})^{1/2}=\bigoplus _{a<b} r_{ab} I_{2n_{a} n_{b}}, \\
& E=B_{Z} R^{-1}, \quad  E^{*}E=I_{m}, \\
& N_{Z}=\ker  B_{Z}^{*} \text{ in the off-block tuple space }.
\end{aligned}\notag
\end{gather*}

The intrinsic normal constraints are \(B_{Z}^{*}Y=0\), equivalently
\(b_{ab} \cdot  Y_{ab}=0\) for every pair. The slice \(S\) consists of
the triples \((Z,A,Y)\) satisfying these constraints. Statements below
are tensor identities on this slice and its local gauge-saturated image.
A simultaneous smooth projector branch exists locally by the nonsingular
Hessian proved below. This does not assert a globally unique branch
across incompatible cluster partitions.

\subsection{2. Horizontal differential and the exact internal
coefficient}\label{a-horizontal-differential-and-the-exact-internal-coefficient}

Use the normal connection obtained by orthogonal projection onto
\(N_{Z}\). For \(h=(\delta  Z,\delta  A,\delta  Y_{N})\) in
\(H=C \oplus  I \oplus  N_{Z}\), horizontal differentiation of
\(B_{Z}^{*}Y=0\) gives

% DISPLAY a.6
% SOURCE delta Y=delta Y_N+E K h,
% SOURCE K h= -R^(-1) B_(delta Z)^*Y.                             (3)
\begin{gather*}
\begin{aligned}
& \delta  Y=\delta  Y_{N}+E K h, \\
& K h= -R^{-1} B_{\delta  Z}^{*}Y.
\end{aligned}\tag{3}
\end{gather*}

\(K\) vanishes identically on I and \(N_{Z}\). The slice embedding
differential, in the orthogonal ambient splitting
\(H \oplus  \operatorname{range}  E\), is

% DISPLAY a.7
% SOURCE S h=(h,K h),
% SOURCE g_S=S^*S=I_H+K^*K
% SOURCE    =g_C direct_sum I_I direct_sum I_(N_Z).                (4)
\begin{gather*}
\begin{aligned}
& S h=(h,K h), \\
& g_{S}=S^{*}S=I_{H}+K^{*}K \\
& =g_{C} \oplus  I_{I} \oplus  I_{N_{Z}}.
\end{aligned}\tag{4}
\end{gather*}

Thus the induced slice has internal kinetic coefficients \textbf{exactly
one}, and there are no internal/center or internal/fast induced-metric
cross terms. The full reduced co-metric additionally contains the
positive orbital contribution from the Gauss square; it is not claimed
that its total AA block equals I. The coefficient-one \(t_{A}\) summand
is retained before that square is estimated. This is an exact identity
for every A and \(Y\) in the coordinate domain, not an oscillator
expectation.

Skew-adjointness of the gauge action gives a useful equivalent identity.
For \(O=B_{X}, O_{C}=P_{C} O\) and \(O_{H}=P_{H} O\),

% DISPLAY a.8
% SOURCE O_C=P_C B_Y,   K=R^(-1) O_C^*.                            (5)
\begin{gather*}
\begin{aligned}
& O_{C}=P_{C} B_{Y}, \quad  K=R^{-1} O_{C}^{*}.
\end{aligned}\tag{5}
\end{gather*}

The elementary commutator bound \(\Vert B_{Y}\Vert \le 2\vert Y\vert\)
proves

% DISPLAY a.9
% SOURCE ||K||<=2s,   s=|Y|/r_*,
% SOURCE (1+4s^2)^(-1) I_C <= g_C^(-1) <= I_C.                    (6)
\begin{gather*}
\begin{aligned}
& \Vert K\Vert \le 2s, \quad  s=\vert Y\vert /r_{*}, \\
& (1+4s^{2})^{-1} I_{C} \le  g_{C}^{-1} \le  I_{C}.
\end{aligned}\tag{6}
\end{gather*}

All normal-frame connections remain inside the covariant center
momentum. Their unbounded \(Y \partial _{Y}\) generators are never
estimated by oscillator energy on arbitrary vectors.

\subsection{3. Faddeev-Popov matrix, explicit noncircular invertibility,
and
density}\label{a-faddeev-popov-matrix-explicit-noncircular-invertibility-and-density}

The differential of the normalized slice constraint is
\(N^{*}=(-K,I_{m})\); thus

% DISPLAY a.10
% SOURCE N=(-K^*,I_m),   W=N^*N=I_m+K K^*.
\begin{gather*}
\begin{aligned}
& N=(-K^{*},I_{m}), \quad  W=N^{*}N=I_{m}+K K^{*}.
\end{aligned}\notag
\end{gather*}

The full transverse-orbit matrix is

% DISPLAY a.11
% SOURCE M=N^*O=E^*B_X-K O_C
% SOURCE  =E^*B_(Z+A)+E^*B_Y-R^(-1)O_C^*O_C.                     (7)
\begin{gather*}
\begin{aligned}
& M=N^{*}O=E^{*}B_{X}-K O_{C} \\
& =E^{*}B_{Z+A}+E^{*}B_{Y}-R^{-1}O_{C}^{*}O_{C}.
\end{aligned}\tag{7}
\end{gather*}

Put \(F=E^{*}B_{Z+A}\). \(F\) is pair-block diagonal. Under the
canonical complex identification of a pair, its block is unitarily
equivalent to

% DISPLAY a.12
% SOURCE r_ab I+T_(ab,n),
% SOURCE T_(ab,n) V=A_n^(a)V-V A_n^(b),  n=b_ab/r_ab.
\begin{gather*}
\begin{aligned}
& r_{ab} I+T_{ab,n}, \\
& T_{ab,n} V=A_{n}^{(a)}V-V A_{n}^{(b)}, \quad  n=b_{ab}/r_{ab}.
\end{aligned}\notag
\end{gather*}

Consequently

% DISPLAY a.13
% SOURCE (1-kappa)R <= F <= (1+kappa)R,   FR=RF.                  (8)
\begin{gather*}
\begin{aligned}
& (1-\kappa )R \le  F \le  (1+\kappa )R, \quad  FR=RF.
\end{aligned}\tag{8}
\end{gather*}

The transverse-frame convention \(E=B_{Z} R^{-1}\) has absorbed the
fixed real root-plane complex structure; that is why \(F\) is positive
here. In a fixed real root frame the corresponding matrix carries that
orthogonal complex structure.

Equations (5)-(8) give the explicit bound

% DISPLAY a.14
% SOURCE ||M R^(-1)-I|| <= kappa+2s+4s^2 <=0.0304.                (9)
\begin{gather*}
\begin{aligned}
& \Vert M R^{-1}-I\Vert  \le  \kappa +2s+4s^{2} \le 0.0304.
\end{aligned}\tag{9}
\end{gather*}

No \(r_{\max}/r_{*}\) or unrelated internal radius enters (9). The full
differential from \((h,\xi )\) to ambient coordinates has block matrix

% DISPLAY a.15
% SOURCE T=[ I_H   O_H ]
% SOURCE   [ K     E^*O].
\[T=\begin{bmatrix}I_H&O_H\\K&E^{*}O\end{bmatrix}.\]

Its determinant is det \(M\). Equivalently the induced slice volume is
\(\sqrt{\det  W}\), while the normalized-normal orbit contraction is
\(\vert \det (W^{-1/2}M)\vert\). They cancel. The physical measure after
the compact off-block gauge integration is therefore

% DISPLAY a.16
% SOURCE J dZ dA dY_N,   J=det M>0,                              (10)
\begin{gather*}
\begin{aligned}
& J dZ dA dY_{N}, \quad  J=\det  M>0,
\end{aligned}\tag{10}
\end{gather*}

up to one constant orbit normalization. This is physical measure
including orbit volume, not bare quotient Riemannian volume.

At \(Y=0\),

% DISPLAY a.17
% SOURCE J_0=det F=product_(a<b) det_C(r_ab I+T_(ab,n))^2,
% SOURCE d_0=sqrt(J_0)=product_(a<b) det_C(r_ab I+T_(ab,n)).        (11)
\begin{gather*}
\begin{aligned}
& J_{0}=\det  F=\prod _{a<b} \det _{C}(r_{ab} I+T_{ab,n})^{2}, \\
& d_{0}=\sqrt{J_{0}}=\prod _{a<b} \det _{C}(r_{ab} I+T_{ab,n}).
\end{aligned}\tag{11}
\end{gather*}

The determinant is positive by (1). Internal coincidences introduce no
singularity.

There is no linear-Y density term, even for nonzero internal A. Indeed
\(L=E^{*}B_{Y}\) has zero diagonal pair blocks: acting on an ab gauge
parameter, an off-block \(Y\) produces block-diagonal output or a
different pair. Since \(F^{-1}\) is pair-block diagonal,

% DISPLAY a.18
% SOURCE tr(F^(-1)L)=0.                                         (12)
\begin{gather*}
\begin{aligned}
& \operatorname{tr} (F^{-1}L)=0.
\end{aligned}\tag{12}
\end{gather*}

The remaining correction -K \(O_{C}\) is quadratic in \(Y\). In fact the
absolutely convergent logarithm series gives, for the stated \(\kappa\)
and \(\sigma\),

% DISPLAY a.19
% SOURCE |log(J/J_0)| <=8 q s^2.                                 (13)
\begin{gather*}
\begin{aligned}
& \vert \log (J/J_{0})\vert  \le 8 q s^{2}.
\end{aligned}\tag{13}
\end{gather*}

For verification:
\(\Vert F^{-1}L\Vert \le 2s/(1-\kappa ), \Vert F^{-1}K O_{C}\Vert \le 4s^{2}/(1-\kappa )\).
The linear trace vanishes, and the sum of powers \(\ge 2\) is bounded by
\(q e^{2}/[2(1-e)]\), where \(e=(2s+4s^{2})/(1-\kappa )\). Together with
the quadratic trace this is \(<8q s^{2}\).

\subsection{4. Intrinsic projectors and transition
law}\label{a-intrinsic-projectors-and-transition-law}

For an ordered orthogonal family of projectors \(P_{a}\) of ranks
\(n_{a}\) define the gauge-invariant function

% DISPLAY a.20
% SOURCE Phi_X(P)=sum_(a,i) n_a^(-1) (Tr(P_a X_i))^2.
\begin{gather*}
\begin{aligned}
& \Phi _{X}(P)=\sum _{a,i} n_{a}^{-1} (\operatorname{Tr} (P_{a} X_{i}))^{2}.
\end{aligned}\notag
\end{gather*}

Stationarity under off-block variations is precisely \(B_{Z}^{*}Y=0\).
At such a point its Hessian in a gauge parameter \(\xi\) is

% DISPLAY a.21
% SOURCE Hess Phi_X[xi,xi]
% SOURCE   =2||O_C xi||^2-2<B_Z xi,B_X xi>
% SOURCE   =-2<xi,R M xi>.                                      (14)
\begin{gather*}
\begin{aligned}
& \operatorname{Hess}  \Phi _{X}[\xi ,\xi ] \\
& =2\Vert O_{C} \xi \Vert ^{2}-2\langle B_{Z} \xi ,B_{X} \xi\rangle \\
& =-2\langle \xi ,R M \xi\rangle .
\end{aligned}\tag{14}
\end{gather*}

Thus RM is symmetric. In the normalized orbit variable \(\eta =R \xi\),
the negative half-Hessian is \(M R^{-1}\), which by (9) lies between
\(0.9696 I\) and \(1.0304\) I. Every such stationary decomposition is a
strict local maximum with a nonsingular transverse Hessian. The implicit
function theorem therefore gives a unique smooth nearby projector
branch, equivariant under SU(N), with no internal eigenvalue gap
assumption.

On overlap of charts describing the same branch, the ordered \(P_{a}\)
agree. Their block frames differ by \(K\) and, when the labels are
changed consistently, by a permutation of equal-rank blocks. Normal
frames differ by orthogonal maps on \(N_{Z}\). These are exact unitary
bundle changes on bosonic fibers and the fermionic module. \(J\) is
invariant; no extra nonunit half-density amplitude appears. All
expressions (3)-(14) and the kinetic identity below intertwine under
these changes. One may prove them in local frames without introducing
scalar angular IMS cutoffs.

This is a \textbf{same-branch} transition theorem. Comparing different
nested partitions, proving that their selectors agree with all
commuting-tree choices, and uniformly transporting every decision-tree
cutoff off the commuting cone remain separate obligations. Local
nonsingularity alone does not prove any of those global claims.

\subsubsection{4.1 Smooth fast projection on the same
branch}\label{a-smooth-fast-projection-on-the-same-branch}

The frozen pair matrices \(G_{ab}, K_{ab}\) and \(D_{ab}\) of the
general fast-energy lemma are obtained at \(Y=0\) from (15)-(16) and the
exact quadratic potential. Their bosonic frequencies and fermionic
spectral gaps are bounded below by fixed positive multiples of
\(r_{ab}\) under (1). Functional calculus therefore gives a smooth
positive Gaussian and a smooth negative-energy occupied subspace,
including at every internal eigenvalue collision. The resulting rank-one
fast-vacuum projection \(P\) is gauge equivariant. Local vacuum frames
may have a nonflat scalar Berry connection; none is declared parallel
here.

A same-branch chart change acts on every finite pair matrix by its
actual orthogonal/color/spin unitary representation. Spectral functional
calculus and the positive Gaussian construction commute with that
action. Hence \(P'=V P V^{*}\) exactly. The same is true after a smooth
invariant radial cutoff, for example with
\(h=(\sum _{a<b}r_{ab}^{-2})^{-1/2}\) and
\(\chi (\vert Y\vert /(\sigma  h))\), followed by fiberwise
normalization. There is no angular scalar partition in this
construction. Since \(r_{*}/\sqrt{k(k-1)/2}\le h\le r_{*}\), the cutoff
tail bounds differ only by fixed N-dependent constants.

For the uncut product \(U\), an internal derivative in \(A_{a}\)
differentiates only the pair factors incident to a. For \(U_{\sigma}\)
the derivative also differentiates its global scalar normalization; this
term is bounded by the same derivative norms, and its difference from
the uncut normalization is exponentially small in the indicated tail
scale. The resolvent formula for their spectral projections and the
Gaussian square-root formula give norm bounds \(C_{N,j} r_{ab}^{-j}\)
for each fixed number \(j\) of such derivatives, after the natural fast
dilation. These follow from external pair gaps, never from an internal
Hamiltonian resolvent. Mixed derivatives have the corresponding products
of inverse incident scales. Global phase choices are unnecessary because
the projection and its induced connection are intrinsic.

\subsection{5. Exact all-vector kinetic plus Gauss
form}\label{a-exact-all-vector-kinetic-plus-gauss-form}

Let \(p_{H}\) be the slice coordinate momentum stack, with the normal
connection in its center component and ordinary internal and fast
components. Let \(S_{f}\) denote the off-block fermionic Gauss
generator, in the same sign convention as \(O^{*}p=S_{f}\). Residual
K-equivariance remains imposed. The orthogonal tangential projection of
the orbit is

% DISPLAY a.22
% SOURCE O_T=S g_S^(-1) S^*O.
\begin{gather*}
\begin{aligned}
& O_{T}=S g_{S}^{-1} S^{*}O.
\end{aligned}\notag
\end{gather*}

For a raw equivariant section \(\phi\) define

% DISPLAY a.23
% SOURCE D phi=W^(1/2) M^(-*)[S_f phi-O^*S g_S^(-1)p_H phi].      (15)
\begin{gather*}
\begin{aligned}
& D \phi =W^{1/2} M^{-*}[S_{f} \phi -O^{*}S g_{S}^{-1}p_{H} \phi ].
\end{aligned}\tag{15}
\end{gather*}

Then the exact ambient reduced kinetic form, in coefficient-one units,
is

% DISPLAY a.24
% SOURCE t_raw[phi]=integral J {
% SOURCE    <p_H phi,g_S^(-1)p_H phi>+||D phi||^2 }.              (16)
\begin{gather*}
\begin{aligned}
& t_{\mathrm{raw}}[\phi ]=\int  J \{ \\
& \langle p_{H} \phi ,g_{S}^{-1}p_{H} \phi\rangle +\Vert D \phi \Vert ^{2} \}.
\end{aligned}\tag{16}
\end{gather*}

Proof: write an ambient covector as its slice-tangent projection plus
\(N \beta\). Its tangential restriction is \(p_{H}\); the Gauss equation
fixes \(M^{*} \beta =S_{f}-O^{*}S g_{S}^{-1}p_{H}\). Tangent and normal
parts are orthogonal, and
\(\Vert N \beta \Vert ^{2}=\langle \beta ,W \beta\rangle\). This proves
(16) for arbitrary covectors, arbitrary spinor values, and unrestricted
momenta. It is a first-order identity, not substitution into a squared
differential operator.

For the flat measure put \(d=\sqrt{J}\) and \(\psi =d \phi\). The exact
flattened formula is obtained from (15)-(16) by

% DISPLAY a.25
% SOURCE p_H -> p_H+i d_H log d.                                 (17)
\begin{gather*}
\begin{aligned}
& p_{H} \to  p_{H}+i d_{H} \log  d.
\end{aligned}\tag{17}
\end{gather*}

This fixes every ordering. Equivalently integration of the real density
cross terms gives the unshifted quadratic stack plus the scalar
multiplication

% DISPLAY a.26
% SOURCE V_d=d^(-1) div_H(a d_H d),                              (18)
\begin{gather*}
\begin{aligned}
& V_{d}=d^{-1} \operatorname{div} _{H}(a d_{H} d),
\end{aligned}\tag{18}
\end{gather*}

where a is the complete scalar principal co-metric, including the
orbital part of \(D\). Divergence is that of the flat bundle measure.
The finite Hermitian spin terms have zero real cross pairing with the
purely imaginary scalar density shift. No separate spin connection is
appended. Formula (17) remains the definitive expression if a convention
for momenta is changed.

On a fixed-shape regular region \(r_{*}\ge \delta  \rho _{C}\) with
\(\rho _{C}^{2}=\sum _{a} n_{a}\vert z_{a}\vert ^{2}\), and a closed
smaller tube, \(V_{d}\) and all coefficient derivatives needed at any
fixed order are bounded by constants times their homogeneous powers. In
particular
\(\vert V_{d}\vert \le C_{N,\delta ,\kappa ,\sigma} \rho _{C}^{-2}\).
These are bounds from explicit smooth finite matrices and compact shape
sets; they do not use a full Hamiltonian inequality.

Combining the exact internal kinetic coefficient with the exact internal
potential and Yukawa pieces retains

% DISPLAY a.27
% SOURCE sum_a [t_(A_a)+V_a+Yukawa_a]=sum_a h_(n_a)
\begin{gather*}
\begin{aligned}
& \sum _{a} [t_{A_{a}}+V_{a}+\mathrm{Yukawa}_{a}]=\sum _{a} h_{n_{a}}
\end{aligned}\notag
\end{gather*}

with coefficient one, also on colored internal sections. The
averaged-supercharge nonnegativity of these internal forms is the usual
algebraic statement, not a physical-sector Hardy inequality or an
internal spectral gap.

Equation (16) does \textbf{not} justify adding a separately extracted
adapted fast kinetic matrix and the whole same Gauss square. The frozen
adapted kinetic matrix is already part of that square.

\subsection{6. Incident-edge reserve for unrestricted low internal
derivatives}\label{a-incident-edge-reserve-for-unrestricted-low-internal-derivatives}

Let U(Z,A) be the normalized product of the pairwise kinetic-aware
Gaussian and occupied fermionic determinant from the frozen-pair lemma.
Interpret it as an isometric injection of its vacuum line; no flat Berry
connection is assumed. On each pair its covariance obeys

% DISPLAY a.28
% SOURCE E_U |Y_ab|^2 <= C n_a n_b ell^3/r_ab,
\begin{gather*}
\begin{aligned}
& E_{U} \vert Y_{ab}\vert ^{2} \le  C n_{a} n_{b} \ell ^{3}/r_{ab},
\end{aligned}\notag
\end{gather*}

with an absolute \(C\) on the stated \(\kappa\) tube. Use a smooth
invariant cutoff strictly inside (2), normalize fiberwise, and call the
resulting injection \(U_{\sigma}\). At fixed finite \(N\) and
sufficiently large \(r_{*}\), the same covariance bound holds with twice
\(C\); cutoff errors are bounded by polynomial factors times
\(\exp (-c \sigma ^{2} r_{*}^{3}/\ell ^{3})\), uniformly in unrelated A
and \(r_{\max}\). Derivatives can have fixed shape-dependent prefactors
if expressed in \(\rho _{C}\) units.

The internal orbital map in (15) is exactly

% DISPLAY a.29
% SOURCE O_(T,A)=O_A=P_I B_Y.                                    (19)
\begin{gather*}
\begin{aligned}
& O_{T,A}=O_{A}=P_{I} B_{Y}.
\end{aligned}\tag{19}
\end{gather*}

There is no Y-independent internal orbital derivative and no
contribution from an unrelated pair. For an internal covector
\(p=(p_{a})\), the ab component satisfies

% DISPLAY a.30
% SOURCE |(O_A^*p)_ab| <= C |Y_ab| (|p_a|+|p_b|).                 (20)
\begin{gather*}
\begin{aligned}
& \vert (O_{A}^{*}p)_{ab}\vert  \le  C \vert Y_{ab}\vert  (\vert p_{a}\vert +\vert p_{b}\vert ).
\end{aligned}\tag{20}
\end{gather*}

Write \(M=F+P\). The elementary bound

% DISPLAY a.31
% SOURCE ||P F^(-1)|| <=(2s+4s^2)/(1-kappa)<0.021
\begin{gather*}
\begin{aligned}
& \Vert P F^{-1}\Vert  \le (2s+4s^{2})/(1-\kappa )<0.021
\end{aligned}\notag
\end{gather*}

shows that \(M^{-*}=(I+F^{-*}P^{*})^{-1}F^{-*}\) has a uniformly bounded
left factor. \(W^{1/2}\) is uniformly bounded too. Thus (20) gives the
pointwise estimate

% DISPLAY a.32
% SOURCE ||W^(1/2)M^(-*)O_A^*p||^2
% SOURCE     <=C sum_(a<b) |Y_ab|^2/r_ab^2 (|p_a|^2+|p_b|^2).    (21)
\begin{gather*}
\begin{aligned}
& \Vert W^{1/2}M^{-*}O_{A}^{*}p\Vert ^{2} \\
& \le C \sum _{a<b} \vert Y_{ab}\vert ^{2}/r_{ab}^{2} (\vert p_{a}\vert ^{2}+\vert p_{b}\vert ^{2}).
\end{aligned}\tag{21}
\end{gather*}

This is valid for every internal covector, not just bounded internal
momenta. Compressing only its coefficients in \(U_{\sigma}\) yields

% DISPLAY a.33
% SOURCE ||D_A(U_sigma nabla_A f)||^2
% SOURCE     <=sum_a w_a |nabla_(A_a) f|^2,
% SOURCE w_a=C sum_(b!=a) n_a n_b ell^3 r_ab^(-3).                (22)
\begin{gather*}
\begin{aligned}
& \Vert D_{A}(U_{\sigma} \nabla _{A} f)\Vert ^{2} \\
& \le \sum _{a} w_{a} \vert \nabla _{A_{a}} f\vert ^{2}, \\
& w_{a}=C \sum _{b\ne a} n_{a} n_{b} \ell ^{3} r_{ab}^{-3}.
\end{aligned}\tag{22}
\end{gather*}

The derivatives here are the covariant derivatives on the possibly
nonflat vacuum line. The expression in (22) denotes the part where the
derivative acts on \(f\); derivatives acting on \(U_{\sigma}\) are
separate coefficient sources. Estimate (22) records the
incident-weighted low-leg coefficient bound. Its use in the mixed form
must follow the complete-charge and postcompletion estimates of Appendix
I, Sections \(4\) and 5.2-5.4; it does not authorize retaining or
spending a second copy of the original Gauss square.

Most importantly, the loss is incident-edge weighted. For an ordinary
untwisted internal derivative, both
\(t_{A_{a}}\le h_{a}+C_{n_{a}} \ell ^{-3} \alpha _{a}\) and
\(V_{a}\le h_{a}+C_{n_{a}} \ell ^{-3} \alpha _{a}\) follow from the
finite-Clifford Yukawa bound. Independently of that form comparison, (1)
gives

% DISPLAY a.34
% SOURCE w_a alpha_a <=C kappa sum_(b!=a)n_a n_b ell^3 r_ab^(-2). (23)
\begin{gather*}
\begin{aligned}
& w_{a} \alpha _{a} \le C \kappa  \sum _{b\ne a}n_{a} n_{b} \ell ^{3} r_{ab}^{-2}.
\end{aligned}\tag{23}
\end{gather*}

The weighted derivative error has a direct payment that does not require
a new weighted full-supercharge theorem. Use the internal-radial gauge
constructed in source Section \(4.2\) in Appendix \(F\): each pair
connection obeys \(\vert a_{ab}\vert \le C_{N} \kappa ^{2}/r_{ab}\), and
its \(A_{a}\) component vanishes unless that pair is incident to a. In
the resulting identification with the flat \(A=0\) center line,

% DISPLAY a.35
% SOURCE sum_a w_a ||nabla_(A_a)f||^2
% SOURCE    <=2 sum_a w_a t_(A_a)[f]
% SOURCE       +C_N kappa^4 ell^3 r_*^(-3)
% SOURCE                      integral sum_(a<b)r_ab^(-2)|f|^2
% SOURCE    <=2 sum_a w_a h_a[f]
% SOURCE       +C_N [kappa+kappa^4 ell^3 r_*^(-3)]
% SOURCE                      integral sum_(a<b)r_ab^(-2)|f|^2.   (26)
\begin{gather*}
\begin{aligned}
& \sum _{a} w_{a} \Vert \nabla _{A_{a}}f\Vert ^{2} \\
& \le 2 \sum _{a} w_{a} t_{A_{a}}[f] \\
& +C_{N} \kappa ^{4} \ell ^{3} r_{*}^{-3} \\
& \int  \sum _{a<b}r_{ab}^{-2}\vert f\vert ^{2} \\
& \le 2 \sum _{a} w_{a} h_{a}[f] \\
& +C_{N} [\kappa +\kappa ^{4} \ell ^{3} r_{*}^{-3}] \\
& \int  \sum _{a<b}r_{ab}^{-2}\vert f\vert ^{2}.
\end{aligned}\tag{26}
\end{gather*}

All rank-dependent Clifford and endpoint factors are included in
\(C_{N}\). The first inequality is the elementary
\(\vert \partial  f+i a f\vert ^{2}\le 2\vert \partial  f\vert ^{2}+2\vert a\vert ^{2}\vert f\vert ^{2}\);
the second uses (23). The weights depend only on centers, so multiplying
the internal form creates no omitted internal weight derivative. At
large \(r_{*}\) their \(h_{a}\) coefficients vanish. The inverse-square
term is arbitrarily small in \(\kappa\) and has a shape-independent
pair-center Hardy payment. This argument pays only the small derivative
error (22); it does not replace the full-supercharge comparison needed
for the principal Berry-coupled slow form. A cruder replacement
\(w_{a}\le C r_{*}^{-3}\), followed by \(r_{*}^{-3} \alpha _{a}\), would
destroy this conclusion for a far extended block and is not used.

The variable induced center metric is also paid as a genuine
derivative-form pairing against a retained covariant center kinetic
square. Its correction is exactly

% DISPLAY a.36
% SOURCE g_C^(-1)-I_C=-K^*(I+K K^*)^(-1)K.
\begin{gather*}
\begin{aligned}
& g_{C}^{-1}-I_{C}=-K^{*}(I+K K^{*})^{-1}K.
\end{aligned}\notag
\end{gather*}

Its Gaussian coefficient norm is \(O_{N}(\ell ^{3} r_{*}^{-3})\), so the
squared cost in that direct pairing is \(O_{N}(\ell ^{6} r_{*}^{-6})\)
times the low center derivative form. This step involves only center
derivatives; it cannot create an internal kinetic loss. Derivatives of
\(U_{\sigma}\) supply further coefficient sources, which still require
their own full source inventory.

\clearpage
\section{Frozen fast energy defect}\label{b-frozen-fast-energy-defect}

Technical derivation for review. Equation and subsection numbers in this
appendix are local to the source derivation. Historical local labels are
retained, including numbering gaps or nonsequential labels. The
accompanying source notes retain the original expressions for
comparison.

This appendix writes the frozen energy in physical units. Multiply
\(E_{0}\) by \(2/R_{\mathrm{BFSS}}\) to use the kinetic-one form of the
main text; set \(\ell =1\) after the stated fixed dilation.

\subsection{1. Statement and
normalization}\label{b-statement-and-normalization}

Let n,m be positive integers, \(d=nm\), and let \(A_{i}\) and
\(B_{i}, i=1,\ldots ,9\), be Hermitian internal matrices of the two
blocks. Rotate the spatial coordinates so that their center difference
is \(b=r e_{9}, r>0\). On \(\operatorname{Hom} (C^{m},C^{n})\), with its
Hilbert--Schmidt metric, put

% DISPLAY b.1
% SOURCE T_i Z=A_i Z−Z B_i,
% SOURCE Q_i=r δ_i9 I_d+T_i,
% SOURCE H=sum_i Q_i²,     C_ij=[Q_i,Q_j]=[T_i,T_j].
\begin{gather*}
\begin{aligned}
& T_{i} Z=A_{i} Z-Z B_{i}, \\
& Q_{i}=r \delta _{i9} I_{d}+T_{i}, \\
& H=\sum _{i} Q_{i}^{2}, \quad  C_{ij}=[Q_{i},Q_{j}]=[T_{i},T_{j}].
\end{aligned}\notag
\end{gather*}

Assume each \(\Vert T_{i}\Vert _{\mathrm{op}}\le r/100\). A
coordinate-invariant sufficient condition is

% DISPLAY b.2
% SOURCE (sum_i ||A_i||F²)^(1/2)+(sum_i ||B_i||F²)^(1/2)≤r/100.
\begin{gather*}
\begin{aligned}
& (\sum _{i} \Vert A_{i}\Vert _{F}^{2})^{1/2}+(\sum _{i} \Vert B_{i}\Vert _{F}^{2})^{1/2}\le r/100.
\end{aligned}\notag
\end{gather*}

Thus it suffices that each endpoint radius is at most r/200. No relation
to any unrelated pair\textquotesingle s separation is imposed. Each pair
may use its own orthonormal spatial frame; a common global axial
direction for all pairs is unnecessary. All traces below are finite
ordinary, unnormalized traces.

On the slice \(Z_{9}=0\) define 8-by-8 matrices with d-by-d blocks:

% DISPLAY b.3
% SOURCE G_ij=δ_ij I_d+Q_i Q_9^(-2) Q_j,
% SOURCE K_ij=δ_ij H−Q_i Q_j+2 C_ij,       i,j≤8.
\begin{gather*}
\begin{aligned}
& G_{ij}= \delta _{ij} I_{d}+Q_{i} Q_{9}^{-2} Q_{j}, \\
& K_{ij}= \delta _{ij} H-Q_{i} Q_{j}+2 C_{ij}, \quad  i,j\le 8.
\end{aligned}\notag
\end{gather*}

The kinetic-aware bosonic frequency matrix is \(B=G^{1/2} K G^{1/2}\).
For Hermitian Spin(9) matrices satisfying
\(\{ \Gamma _{i}, \Gamma _{j}\}=2 \delta _{ij} I_{16}\) put

% DISPLAY b.4
% SOURCE D=sum_i Γ_i⊗Q_i.
\begin{gather*}
\begin{aligned}
& D=\sum _{i} \Gamma _{i} \otimes  Q_{i}.
\end{aligned}\notag
\end{gather*}

Both \(K\) and \(B\) are positive, and \(D\) has precisely 8d negative
and 8d positive eigenvalues. The frozen ground energy, in physical
units, is

% DISPLAY b.5
% SOURCE E_0=(R/ell³)[Tr_(8d) sqrt(B)−(1/2)Tr_(16d)|D|].
\begin{gather*}
\begin{aligned}
& E_{0}=(R/\ell ^{3})[\operatorname{Tr} _{8d} \sqrt{B}-(1/2)\operatorname{Tr} _{16d}\vert D\vert ].
\end{aligned}\notag
\end{gather*}

The following explicit, deliberately non-sharp constant is valid for
every n,m:

% DISPLAY b.6
% SOURCE |E_0| ≤ 150002 (R/ell³) r^(-3)
% SOURCE               sum_(i<j)||[T_i,T_j]||F².                    (1)
\begin{gather*}
\begin{aligned}
& \vert E_{0}\vert  \le  150002 (R/\ell ^{3}) r^{-3} \\
& \sum _{i<j}\Vert [T_{i},T_{j}]\Vert _{F}^{2}.
\end{aligned}\tag{1}
\end{gather*}

The constant in (1) is independent of \(d\); endpoint weights below
carry the finite-rank dependence. The proof neither diagonalizes
internal matrices nor invokes an internal Hamiltonian gap, internal
eigenvalue gap, regularity of the commuting variety, or a
nearest-commuting-tuple estimate.

\subsection{2. Why these are the correct frozen
matrices}\label{b-why-these-are-the-correct-frozen-matrices}

The eliminated off-block Gauss principal equation is

% DISPLAY b.7
% SOURCE Q_9 p_9+sum_(i≤8)Q_i p_i=0.
\begin{gather*}
\begin{aligned}
& Q_{9} p_{9}+\sum _{i\le 8}Q_{i} p_{i}=0.
\end{aligned}\notag
\end{gather*}

Its positive kinetic square gives \(G\) exactly. It must not be dropped
or treated as a separate small bounded perturbation.

For
\(V(X)=(R/(2\ell ^{6}))\sum _{i<j}\Vert [X_{i},X_{j}]\Vert _{F}^{2}\),
the quadratic part in one off-block pair is \((R/\ell ^{6}) Z^{*} K Z\).
The square of the linear off-block commutator gives
\(\delta _{ij} H-Q_{j} Q_{i}\). The internal-background commutator
paired with \([Y_{i},Y_{j}]\) contributes \(C_{ij}\). Their sum is

% DISPLAY b.8
% SOURCE δ_ij H−Q_j Q_i+C_ij=δ_ij H−Q_i Q_j+2C_ij.
\begin{gather*}
\begin{aligned}
& \delta _{ij} H-Q_{j} Q_{i}+C_{ij}= \delta _{ij} H-Q_{i} Q_{j}+2C_{ij}.
\end{aligned}\notag
\end{gather*}

This works for two arbitrary internal blocks. The second block
contributes with the minus sign in \(T_{i}\), including in its induced
commutator. At quadratic order other block pairs do not couple into this
pair\textquotesingle s Hessian. Interactions involving three distinct
blocks begin in the higher-degree cross terms and are not estimated
here.

The 8d complex bosonic coordinates are 16d real oscillators. Their
ground energy is \((R/\ell ^{3})\operatorname{Tr}  \sqrt{B}\). The fast
fermionic sector has 16d complex annihilators; filling the 8d negative
eigenmodes gives \(-(R/(2\ell ^{3}))\operatorname{Tr}  \vert D\vert\).
The reverse block is the adjoint, not an independent second copy.

\subsection{3. Uniform cross gaps}\label{b-uniform-cross-gaps}

Scale \(r\) to one and write \(\kappa  =1/100\). Then
\(Q_{9}=I+T_{9}, Q_{i}=T_{i}\) for \(i\le 8\). Block norm row sums give

% DISPLAY b.9
% SOURCE ||K−I|| ≤ p(κ),        p(t)=2t+29t²,
% SOURCE 0≤G−I,   ||G−I||≤g(κ), g(t)=8t²/(1−t)².
\begin{gather*}
\begin{aligned}
& \Vert K-I\Vert  \le  p( \kappa  ), \quad  p(t)=2t+29t^{2}, \\
& 0\le G-I, \quad  \Vert G-I\Vert \le g( \kappa  ), g(t)=8t^{2}/(1-t)^{2}.
\end{aligned}\notag
\end{gather*}

For \(K\), a diagonal block has \(2T_{9}\) plus eight squares, and each
of its seven off-diagonal blocks is \(T_{i} T_{j}-2T_{j} T_{i}\). This
explains the coefficient \(29\). Consequently

% DISPLAY b.10
% SOURCE K≥0.9771 I,     G≥I,     B≥0.9771 I.
\begin{gather*}
\begin{aligned}
& K\ge 0.9771 I, \quad  G\ge I, \quad  B\ge 0.9771 I.
\end{aligned}\notag
\end{gather*}

Also
\(\Vert D- \Gamma _{9} \otimes  I\Vert \le \sum _{i}\Vert T_{i}\Vert \le 9 \kappa\),
so \(\min \vert D\vert \ge 0.91\). Continuous scaling of \(T\) to zero
preserves its signature. \(H\ge Q_{9}^{2}\ge 0.9801\) I. Restoring \(r\)
multiplies \(K,B,H,D^{2}\) by \(r^{2}\) and \(D\) by \(r\).

\subsection{4. The real-trace two-commutator
lemma}\label{b-the-real-trace-two-commutator-lemma}

Let \(X_{1},\ldots ,X_{s}\) be Hermitian matrices with
\(\Vert X_{i}\Vert \le  \kappa\). Set
\(F^{2}=\sum _{i<j}\Vert [X_{i},X_{j}]\Vert _{F}^{2}\). For any words
\(w\) and w\textquotesingle{} of length \(q\) having the same multiset
of letters,

% DISPLAY b.11
% SOURCE |Re Tr w−Re Tr w'|
% SOURCE    ≤ ((q)_4/8) κ^(q−4) F²,                  q≥4,          (2)
\begin{gather*}
\begin{aligned}
& \vert \operatorname{Re}  \operatorname{Tr}  w-\operatorname{Re}  \operatorname{Tr}  w'\vert \\
& \le  ((q)_{4}/8) \kappa ^{q-4} F^{2}, \quad  q\ge 4,
\end{aligned}\tag{2}
\end{gather*}

where \((q)_{4}=q(q-1)(q-2)(q-3)\). The left side is zero for
\(q\le 3\).

Proof. At most \(q(q-1)/2\) adjacent swaps transform one ordering into
another. The trace difference from one swap is
\(\operatorname{Tr}  U[X_{i},X_{j}]V\). By cyclicity this is Tr \(C W\),
with \(C=[X_{i},X_{j}]\) anti-Hermitian and \(W=VU\) a word of length
\(q-2\). Exactly,

% DISPLAY b.12
% SOURCE Re Tr(CW)=(1/2)Tr(C(W−W*)).                              (3)
\begin{gather*}
\begin{aligned}
& \operatorname{Re}  \operatorname{Tr} (CW)=(1/2)\operatorname{Tr} (C(W-W^{*})).
\end{aligned}\tag{3}
\end{gather*}

Reversing \(W\) takes at most \((q-2)(q-3)/2\) adjacent swaps. Each
resulting term contains a second commutator and \(q-4\) other letters.
Hilbert--Schmidt Cauchy--Schwarz bounds its trace by
\(\kappa ^{q-4}F^{2}\). Equation (3), followed by the first sequence of
swaps, proves (2). If \(q\le 3, W\) is constant or Hermitian and (3) is
zero.

In particular, take a real-coefficient homogeneous noncommutative trace
polynomial whose commutative polynomial is zero. Sort every word to the
same canonical word for its multiset. The sorted contributions cancel
coefficient by coefficient; (2) bounds the original real trace by the
sum of absolute coefficients times \((q)_{4} \kappa ^{q-4}F^{2}/8\).
This is an explicit algebraic estimate; it does not divide by an ideal
on a singular algebraic variety.

\subsection{5. Bosonic functional calculus and an explicit coefficient
majorant}\label{b-bosonic-functional-calculus-and-an-explicit-coefficient-majorant}

Continue at \(r=1\). Put \(L=GK-I\). Since GK is similar to \(B\),

% DISPLAY b.13
% SOURCE Tr sqrt(B)=Tr sqrt(I+L)
% SOURCE           =8d+sum_(j≥1) binom(1/2,j) Tr L^j.            (4)
\begin{gather*}
\begin{aligned}
& \operatorname{Tr}  \sqrt{B}=\operatorname{Tr}  \sqrt{I+L} \\
& =8d+\sum _{j\ge 1} \binom{1/2}{j} \operatorname{Tr}  L^{j}.
\end{aligned}\tag{4}
\end{gather*}

The binomial series converges in operator norm:
\(\Vert L\Vert \le l( \kappa  )<1\), where

% DISPLAY b.14
% SOURCE l(t)=p(t)+g(t)(1+p(t)).
\begin{gather*}
\begin{aligned}
& l(t)=p(t)+g(t)(1+p(t)).
\end{aligned}\notag
\end{gather*}

The expansion \(Q_{9}^{-2}=\sum _{j\ge 0}(-1)^{j}(j+1)T_{9}^{j}\)
converges as well. Thus (4), and the corresponding expansion of 8Tr
\(\sqrt{H}\), define absolutely convergent noncommutative power series
with real coefficients.

Here is a coefficientwise majorant, not just a bound on evaluations.
Replace every letter \(T_{i}\) by a common nonnegative variable \(t\),
every scalar coefficient by its absolute value, and take maximum block
row sums before taking the eight spatial diagonal traces. The series for
\(K-I\) is majorized by p(t), for \(G-I\) by g(t), for \(L\) by l(t),
and for \(H-I\) by

% DISPLAY b.15
% SOURCE h(t)=2t+9t².
\begin{gather*}
\begin{aligned}
& h(t)=2t+9t^{2}.
\end{aligned}\notag
\end{gather*}

Since \(\sum _{j\ge 1}\vert \binom{1/2}{j}\vert z^{j}=1-\sqrt{1-z}\),
the sum of absolute trace-word coefficients of

% DISPLAY b.16
% SOURCE F_B(T)=Tr sqrt(B)−8Tr sqrt(H)
\begin{gather*}
\begin{aligned}
& F_{B}(T)=\operatorname{Tr}  \sqrt{B}-8\operatorname{Tr}  \sqrt{H}
\end{aligned}\notag
\end{gather*}

is majorized degree by degree by

% DISPLAY b.17
% SOURCE A(t)=8[2−sqrt(1−l(t))−sqrt(1−h(t))].                     (5)
\begin{gather*}
\begin{aligned}
& A(t)=8[2-\sqrt{1-l(t)}-\sqrt{1-h(t)}].
\end{aligned}\tag{5}
\end{gather*}

The eight is the spatial trace multiplicity; there is no hidden factor
\(d\), because the color trace is retained inside each word and
estimated by (2).

For scalar commuting variables, \(G K=H I_{8}\) exactly, by the rank-one
identity

% DISPLAY b.18
% SOURCE (I+vv*/q_9²)(H I−vv*)=H I,      H=q_9²+v*v.
\begin{gather*}
\begin{aligned}
& (I+vv^{*}/q_{9}^{2})(H I-vv^{*})=H I, \quad  H=q_{9}^{2}+v^{*}v.
\end{aligned}\notag
\end{gather*}

It follows that the commutative specialization of \(F_{B}\) vanishes
identically as a convergent power series. Each homogeneous polynomial
therefore has zero commutative polynomial. Apply (2) termwise, which is
allowed by the absolute majorant.

For a wholly explicit bound use \(t_{0}=1/10\). Direct arithmetic gives

% DISPLAY b.19
% SOURCE p(t_0)=49/100,
% SOURCE g(t_0)=8/81,
% SOURCE l(t_0)=5161/8100<1,
% SOURCE h(t_0)=29/100,
% SOURCE A(t_0)<5.
\begin{gather*}
\begin{aligned}
& p(t_{0})=49/100, \\
& g(t_{0})=8/81, \\
& l(t_{0})=5161/8100<1, \\
& h(t_{0})=29/100, \\
& A(t_{0})<5.
\end{aligned}\notag
\end{gather*}

If \(A(t)=\sum  a_{q} t^{q}\), all \(a_{q}\) are nonnegative and

% DISPLAY b.20
% SOURCE sum_(q≥4) a_q (q)_4 κ^(q−4)
% SOURCE    ≤24 t_0^(-4) A(t_0),             κ≤1/100.
\begin{gather*}
\begin{aligned}
& \sum _{q\ge 4} a_{q} (q)_{4} \kappa ^{q-4} \\
& \le 24 t_{0}^{-4} A(t_{0}), \quad  \kappa  \le 1/100.
\end{aligned}\notag
\end{gather*}

Indeed \((q)_{4}( \kappa  /t_{0})^{q-4}\) is maximized at \(q=4\): the
consecutive-term ratio is at most \(1/2\). Combining with (2) gives

% DISPLAY b.21
% SOURCE |Tr sqrt(B)−8Tr sqrt(H)|
% SOURCE     ≤150000 sum_(i<j)||C_ij||F².                        (6)
\begin{gather*}
\begin{aligned}
& \vert \operatorname{Tr}  \sqrt{B}-8\operatorname{Tr}  \sqrt{H}\vert \\
& \le 150000 \sum _{i<j}\Vert C_{ij}\Vert _{F}^{2}.
\end{aligned}\tag{6}
\end{gather*}

Restoring \(r\) adds \(r^{-3}\). This establishes the needed bosonic
quadratic-commutator estimate without a claim that vanishing on
commuting matrices alone implies it.

\subsection{6. Fermionic resolvent
cancellation}\label{b-fermionic-resolvent-cancellation}

Exactly,

% DISPLAY b.22
% SOURCE D²=I_16⊗H+S,
% SOURCE S=sum_(i<j)Γ_iΓ_j⊗C_ij.
\begin{gather*}
\begin{aligned}
& D^{2}=I_{16} \otimes  H+S, \\
& S=\sum _{i<j} \Gamma _{i} \Gamma _{j} \otimes  C_{ij}.
\end{aligned}\notag
\end{gather*}

The Clifford trace gives

% DISPLAY b.23
% SOURCE Tr[(I⊗H)^(-1/2)S]=0,
% SOURCE ||S||F²=16 sum_(i<j)||C_ij||F².                          (7)
\begin{gather*}
\begin{aligned}
& \operatorname{Tr} [(I \otimes  H)^{-1/2}S]=0, \\
& \Vert S\Vert _{F}^{2}=16 \sum _{i<j}\Vert C_{ij}\Vert _{F}^{2}.
\end{aligned}\tag{7}
\end{gather*}

For positive matrices \(A,A+S\ge  \lambda  I\), the square-root
resolvent identity gives

% DISPLAY b.24
% SOURCE Tr sqrt(A+S)−Tr sqrt(A)−(1/2)Tr(A^(-1/2)S)
% SOURCE   =−(1/π)integral_0^infinity t^(1/2)
% SOURCE      Tr[(A+t)^(-1)S(A+t)^(-1)S(A+S+t)^(-1)] dt.
\begin{gather*}
\begin{aligned}
& \operatorname{Tr}  \sqrt{A+S}-\operatorname{Tr}  \sqrt{A}-(1/2)\operatorname{Tr} (A^{-1/2}S) \\
& =-(1/ \pi  )\int _{0}^{\infty} t^{1/2} \\
& \operatorname{Tr} [(A+t)^{-1}S(A+t)^{-1}S(A+S+t)^{-1}] dt.
\end{aligned}\notag
\end{gather*}

The absolute value is bounded by

% DISPLAY b.25
% SOURCE ||S||F²/(8λ^(3/2)),
\begin{gather*}
\begin{aligned}
& \Vert S\Vert _{F}^{2}/(8 \lambda ^{3/2}),
\end{aligned}\notag
\end{gather*}

using
\(\int _{0}^{\infty} t^{1/2}( \lambda  +t)^{-3}dt= \pi  /(8 \lambda ^{3/2})\).
Take \(A=I \otimes  H\) and \(\lambda  =(0.91r)^{2}\). Equations (7)
prove

% DISPLAY b.26
% SOURCE |(1/2)Tr |D|−8Tr sqrt(H)|
% SOURCE    ≤(0.91)^(-3) r^(-3)sum_(i<j)||C_ij||F²
% SOURCE    <2 r^(-3)sum_(i<j)||C_ij||F².                        (8)
\begin{gather*}
\begin{aligned}
& \vert (1/2)\operatorname{Tr}  \vert D\vert -8\operatorname{Tr}  \sqrt{H}\vert \\
& \le (0.91)^{-3} r^{-3}\sum _{i<j}\Vert C_{ij}\Vert _{F}^{2} \\
& <2 r^{-3}\sum _{i<j}\Vert C_{ij}\Vert _{F}^{2}.
\end{aligned}\tag{8}
\end{gather*}

Combining (6) and (8) proves (1). The first-order fermionic term cancels
exactly under the Clifford trace; the bosonic first-order-in-commutators
term cancels by (3) and commutative specialization.

\subsection{7. Endpoint commutators, pair weights, and
absorption}\label{b-endpoint-commutators-pair-weights-and-absorption}

In a matrix representation of \(\operatorname{Hom} (C^{m},C^{n})\),

% DISPLAY b.27
% SOURCE T_i=A_i⊗I_m−I_n⊗B_i^T,
% SOURCE [T_i,T_j]=[A_i,A_j]⊗I_m−I_n⊗[B_i,B_j]^T.
\begin{gather*}
\begin{aligned}
& T_{i}=A_{i} \otimes  I_{m}-I_{n} \otimes  B_{i}^{T}, \\
& [T_{i},T_{j}]=[A_{i},A_{j}] \otimes  I_{m}-I_{n} \otimes  [B_{i},B_{j}]^{T}.
\end{aligned}\notag
\end{gather*}

There is no mixed term \([A_{i},B_{j}]\): the two endpoint actions
commute. Since traces of commutators vanish, exactly

% DISPLAY b.28
% SOURCE ||[T_i,T_j]||F²
% SOURCE    =m||[A_i,A_j]||F²+n||[B_i,B_j]||F².                 (9)
\begin{gather*}
\begin{aligned}
& \Vert [T_{i},T_{j}]\Vert _{F}^{2} \\
& =m\Vert [A_{i},A_{j}]\Vert _{F}^{2}+n\Vert [B_{i},B_{j}]\Vert _{F}^{2}.
\end{aligned}\tag{9}
\end{gather*}

For a partition with block sizes \(n_{a}\) and separations \(r_{ab}\),
summing (1) gives the pairwise-relative statement

% DISPLAY b.29
% SOURCE |sum_(a<b) E_0,ab|
% SOURCE   ≤C (R/ell³)sum_a F_a² sum_(b≠a)n_b r_ab^(-3),
% SOURCE F_a²=sum_(i<j)||[A_i^(a),A_j^(a)]||F²,   C=150002.      (10)
\begin{gather*}
\begin{aligned}
& \vert \sum _{a<b} E_{0,ab}\vert \\
& \le C (R/\ell ^{3})\sum _{a} F_{a}^{2} \sum _{b\ne a}n_{b} r_{ab}^{-3}, \\
& F_{a}^{2}=\sum _{i<j}\Vert [A_{i}^{(a)},A_{j}^{(a)}]\Vert _{F}^{2}, \quad  C=150002.
\end{aligned}\tag{10}
\end{gather*}

The sum of absolute pair defects obeys the same bound. Defining the
internal bosonic potential \(V_{a}=R F_{a}^{2}/(2\ell ^{6})\), its
multiplier is

% DISPLAY b.30
% SOURCE 2C ell³ sum_(b≠a)n_b r_ab^(-3).                         (11)
\begin{gather*}
\begin{aligned}
& 2C \ell ^{3} \sum _{b\ne a}n_{b} r_{ab}^{-3}.
\end{aligned}\tag{11}
\end{gather*}

This tends to zero once all retained pair separations tend to
\(\infty\). Its constants have no dependence on a ratio involving the
smallest unrelated center gap. At fixed finite \(N, r_{ab}\ge r_{*}\)
gives the explicit upper bound 2C \(\ell ^{3}(N-n_{a})r_{*}^{-3}\).

For clarity about full-form absorption, \(V_{a}\) by itself is not
bounded by \(H_{a}\) with coefficient one without a Yukawa remainder. If
the elementary finite-Clifford estimate is written as

% DISPLAY b.31
% SOURCE ||Y_a(A)||≤c_(n_a) (R/ell³) alpha_a,
% SOURCE alpha_a²=sum_i ||A_i^(a)||F²,
\begin{gather*}
\begin{aligned}
& \Vert Y_{a}(A)\Vert \le c_{n_{a}} (R/\ell ^{3}) \alpha _{a}, \\
& \alpha _{a}^{2}=\sum _{i} \Vert A_{i}^{(a)}\Vert _{F}^{2},
\end{aligned}\notag
\end{gather*}

then \(V_{a}\le H_{a}+c_{n_{a}}(R/\ell ^{3})\alpha _{a}\), since the
internal kinetic form is nonnegative. One explicit choice in these
conventions is \(c_{n}=24(n^{2}-1)\), with \(c_{1}=0\). Indeed, on the
\(16(n^{2}-1)\) real Majoranas, the internal mass is
\(D_{\operatorname{ad}}=\sum _{i} \Gamma _{i} \otimes  \operatorname{ad} (A_{i})\),
and \(Y_{a}=(R/(2\ell ^{3}))\theta ^{T} D_{\operatorname{ad}} \theta\).
In a real Hermitian generator basis \(D_{\operatorname{ad}}\) is purely
imaginary antisymmetric. Pairing its \(\pm  \lambda\) eigenvalues gives
\(\Vert Y_{a}\Vert =(R/(4\ell ^{3}))\operatorname{Tr} \vert D_{\operatorname{ad}}\vert\),
while

% DISPLAY b.32
% SOURCE ||D_ad||≤2sum_i||A_i||op≤6alpha_a.
\begin{gather*}
\begin{aligned}
& \Vert D_{\operatorname{ad}}\Vert \le 2\sum _{i}\Vert A_{i}\Vert _{\mathrm{op}}\le 6\alpha _{a}.
\end{aligned}\notag
\end{gather*}

This proves the stated \(c_{n}\) without an internal spectral gap. Hence
(10) is bounded by an arbitrarily small multiple of \(\sum _{a} H_{a}\),
after taking \(r_{*}\) large, plus

% DISPLAY b.33
% SOURCE 2C R sum_a c_(n_a)alpha_a sum_(b≠a)n_b r_ab^(-3).
\begin{gather*}
\begin{aligned}
& 2C R \sum _{a} c_{n_{a}}\alpha _{a} \sum _{b\ne a}n_{b} r_{ab}^{-3}.
\end{aligned}\notag
\end{gather*}

Under the endpoint condition \(\alpha _{a}\le  \kappa  r_{ab}\), this
last expression is at most

% DISPLAY b.34
% SOURCE 2C R κ sum_(a<b)(c_(n_a)n_b+c_(n_b)n_a)r_ab^(-2).
\begin{gather*}
\begin{aligned}
& 2C R \kappa  \sum _{a<b}(c_{n_{a}}n_{b}+c_{n_{b}}n_{a})r_{ab}^{-2}.
\end{aligned}\notag
\end{gather*}

It is an arbitrarily small center-Hardy cost by reducing \(\kappa\)
before choosing the final exterior radius. This is a precise way to use
the defect; replacing \(V_{a}\) by \(H_{a}\) with no remainder would be
incorrect.

For an explicit finite-N allocation, write
\(T_{\mathrm{cent}}=(R/2)t_{\mathrm{cent}}\) in the mass metric
\(\sum _{a} n_{a}\vert dz_{a}\vert ^{2}\). The nine-dimensional relative
Hardy inequality gives

% DISPLAY b.35
% SOURCE R integral r_ab^(-2)|f|²
% SOURCE   ≤(8/49)(1/n_a+1/n_b)^(-1) T_cent[f].
\begin{gather*}
\begin{aligned}
& R \int  r_{ab}^{-2}\vert f\vert ^{2} \\
& \le (8/49)(1/n_{a}+1/n_{b})^{-1} T_{\mathrm{cent}}[f].
\end{aligned}\notag
\end{gather*}

The preceding residual multiplier is therefore bounded by

% DISPLAY b.36
% SOURCE (16Cκ/49) sum_(a<b)
% SOURCE   (c_(n_a)n_b+c_(n_b)n_a)(1/n_a+1/n_b)^(-1) T_cent[f]
% SOURCE   ≤(192C/49)N^4 κ T_cent[f].
\begin{gather*}
\begin{aligned}
& (16C \kappa  /49) \sum _{a<b} \\
& (c_{n_{a}}n_{b}+c_{n_{b}}n_{a})(1/n_{a}+1/n_{b})^{-1} T_{\mathrm{cent}}[f] \\
& \le (192C/49)N^{4} \kappa  T_{\mathrm{cent}}[f].
\end{aligned}\notag
\end{gather*}

The last bound uses \(c_{n}=24(n^{2}-1)\) and
\(\sum _{a<b}n_{a}^{2}n_{b}^{2}\le N^{4}/2\). Thus, for any
\(\epsilon  >0\), the choices

% DISPLAY b.37
% SOURCE κ≤min(1/200, 49ε/(192C N^4)),
% SOURCE r_*≥[2C ell³ N/ε]^(1/3),
% SOURCE alpha_a≤κ r_ab for every incident pair,
% SOURCE r_ab≥r_* for every pair,
\begin{gather*}
\begin{aligned}
& \kappa  \le \min (1/200, 49 \epsilon  /(192C N^{4})), \\
& r_{*}\ge [2C \ell ^{3} N/ \epsilon  ]^{1/3}, \\
& \alpha _{a}\le  \kappa  r_{ab} \text{ for every incident pair }, \\
& r_{ab}\ge r_{*} \text{ for every pair },
\end{aligned}\notag
\end{gather*}

give the form bound

% DISPLAY b.38
% SOURCE integral sum_(a<b)|E_0,ab| |f|²
% SOURCE   ≤ε sum_a H_a[f]+ε T_cent[f].                         (12)
\begin{gather*}
\begin{aligned}
& \int  \sum _{a<b}\vert E_{0,ab}\vert  \vert f\vert ^{2} \\
& \le  \epsilon  \sum _{a} H_{a}[f]+ \epsilon  T_{\mathrm{cent}}[f].
\end{aligned}\tag{12}
\end{gather*}

Here the standard nonnegative internal supersymmetric forms are used on
their common core, and the center Hardy estimate is valid also for
Hilbert-valued coefficients and unitary connections. The multiplier of
each \(H_{a}\) depends only on the centers, so no derivative of an
internal-dependent weight has been omitted. Equation (12) is a bound on
this defect multiplier, not a claim that the full Hamiltonian has
already been compared with the slow forms on its right.

\clearpage
\section{Centered sources and graph
estimates}\label{c-centered-sources-and-graph-estimates}

Technical derivation for review. Equation and subsection numbers in this
appendix are local to the source derivation. Historical local labels are
retained, including numbering gaps or nonsequential labels. The
accompanying source notes retain the original expressions for
comparison.

\subsection{0. Result and scope}\label{c-result-and-scope}

For any fixed partition \((n_{1},\ldots ,n_{k})\), the centered
coefficient calculation has an exact root-graph resolution:

\begin{itemize}
\tightlist
\item
  Each two-block charge source occupies two quanta of one root pair and
  has fast energy \(4r_{ab}\).
\item
  Each genuine three-block charge source occupies one quantum on every
  edge of a triangle and has fast energy \(2(r_{ab}+r_{bc}+r_{ca})\).
\item
  On residual block-center-neutral states, the leading inverse-square
  Schur multiplication cancels as an operator on the entire internal
  Clifford module. Internal derivatives of the adapted vacuum are
  included in this statement; freezing the internal coordinates is not
  permission to discard those derivatives.
\item
  The associated finite dressing has pair weight \(C_{N} r_{ab}^{-3}\)
  and triangle weight
  \(C_{N} [r_{ab}^{-2}+r_{bc}^{-2}+r_{ca}^{-2}]/(r_{ab}+r_{bc}+r_{ca})\).
  The latter summed denominator is essential beside a far large block
  and a close opposite pair.
\item
  A graph-local multiplier lemma turns these weights into a small
  multiple of the complete slow supersymmetric form plus a small
  pairwise center-Hardy multiplier, without expanding an inverse around
  an internal Hamiltonian and without applying physical internal Hardy
  to colored vectors.
\end{itemize}

The nonzero internal coefficient continuation is supplied in Appendix
\(D\); analyticity of the oscillator matrices alone would not supply
that continuation.

\subsection{1. Fibers, neutral sector, and the correct fast
inverse}\label{c-fibers-neutral-sector-and-the-correct-fast-inverse}

Write \(b_{ab}=z_{a}-z_{b}\), \(r_{ab}=\vert b_{ab}\vert\), and
\(d_{ab}=n_{a} n_{b}\). Centers have mass metric
\(\sum _{a} n_{a} \vert dz_{a}\vert ^{2}\). Internal matrices
\(A^{(a)}\) are traceless Hermitian, with radii

% DISPLAY c.1
% SOURCE alpha_a² = sum_i ||A_i^(a)||_F².
\begin{gather*}
\begin{aligned}
& \alpha _{a}^{2} = \sum _{i} \Vert A_{i}^{(a)}\Vert _{F}^{2}.
\end{aligned}\notag
\end{gather*}

The hypotheses used in graph estimates are genuinely pairwise:

% DISPLAY c.2
% SOURCE alpha_a <= kappa r_ab  for every b != a,
% SOURCE r_ab >= r_* > 0.
\begin{gather*}
\begin{aligned}
& \alpha _{a} \le  \kappa  r_{ab} \quad  \text{ for every } b \ne  a, \\
& r_{ab} \ge  r_{*} > 0.
\end{aligned}\notag
\end{gather*}

No comparison of one gap with an unrelated gap is made.

At \(A=0\) the fast vacuum is a product
\(U_{0}= \otimes _{a<b} U_{ab}\). Each pair has \(16d_{ab}\) real
transverse bosonic coordinates and \(16d_{ab}\) complex fermionic modes,
with \(8d_{ab}\) occupied modes. The reverse block is the adjoint, not a
second independent pair. The vacuum is invariant under the residual
\(S(\prod  U(n_{a}))\) group and is even in boson parity and fast
fermion parity.

For any one normalized real supercharge let \(q_{0}\) be its centered
fast part. Before imposing the block-center Gauss constraints, its
square contains the center-weighted gauge term:

% DISPLAY c.3
% SOURCE q_0² = H_f + sum_i (Gamma_i)_(alpha alpha)
% SOURCE                      sum_a z_a^i G_(a,center,fast),
\begin{gather*}
\begin{aligned}
& q_{0}^{2} = H_{f} + \sum _{i} (\Gamma _{i})_{\alpha  \alpha} \\
& \sum _{a} z_{a}^{i} G_{a,\mathrm{center},\mathrm{fast}},
\end{aligned}\notag
\end{gather*}

up to the common gauge-generator normalization. The background
\(z_{a} I_{n_{a}}\) annihilates all traceless internal generators in
this term. Consequently

% DISPLAY c.4
% SOURCE q_0² = H_f
\begin{gather*}
\begin{aligned}
& q_{0}^{2} = H_{f}
\end{aligned}\notag
\end{gather*}

on the residual block-center-neutral fast space. This includes fast
states carrying nontrivial internal \(SU(n_{a})\) representations. It
does not assert the identity on arbitrary block-center-charged vectors.
Physical total residual equivariance implies this center neutrality
because traceless internal coordinates and internal adjoint fermions
carry no block-center charge.

Every source below is center neutral: a pair contracts opposite root
charges, while a triangle follows an oriented root cycle. On this space
\(P=U_{0} U_{0}^{*}\) is the kernel of \(H_{f}\) and of \(q_{0}\), and

% DISPLAY c.5
% SOURCE q_0^-1 = q_0 H_f^-1
\begin{gather*}
\begin{aligned}
& q_{0}^{-1} = q_{0} H_{f}^{-1}
\end{aligned}\notag
\end{gather*}

is a genuine inverse on \(1-P\). This is not the inverse of an
interacting internal Hamiltonian.

\subsection{2. The internal vacuum jet cannot be
omitted}\label{c-the-internal-vacuum-jet-cannot-be-omitted}

Let \(U(z,A)\) be the normalized adapted product vacuum from the gapped
quadratic matrices in Appendix \(B\). At one pair, rotate \(n=b/r\) to
direction \(9\) and write

% DISPLAY c.6
% SOURCE T_i = A_i^(a) tensor I - I tensor (A_i^(b))^T,
% SOURCE F = rI + T_9,
% SOURCE D = Gamma_9 r + sum_i Gamma_i T_i.
\begin{gather*}
\begin{aligned}
& T_{i} = A_{i}^{(a)} \otimes  I - I \otimes  (A_{i}^{(b)})^{T}, \\
& F = rI + T_{9}, \\
& D = \Gamma _{9} r + \sum _{i} \Gamma _{i} T_{i}.
\end{aligned}\notag
\end{gather*}

At \(A=0\), the bosonic kinetic correction \(G-I\) has no linear term
and the bosonic frequency has first variation
\(\delta  \Omega _{ij}=\delta _{ij} \delta  T_{9}\). The fermionic
negative projector has first variation

% DISPLAY c.7
% SOURCE delta P_- = -(1/(2r)) sum_(i perpendicular n) Gamma_i tensor delta T_i.
\begin{gather*}
\begin{aligned}
& \delta  P_{-} = -(1/(2r)) \sum _{i \perp  n} \Gamma _{i} \otimes  \delta  T_{i}.
\end{aligned}\notag
\end{gather*}

In a local unitary frame with zero scalar Berry connection at this
point, the vacuum derivative is orthogonal to \(U_{0}\) and obeys the
exact norm identity

% DISPLAY c.8
% SOURCE ||delta U_ab||² = (2/r_ab²) sum_i Tr_(d_ab)(delta Q_i²),       (2.1)
\begin{gather*}
\begin{aligned}
& \Vert \delta  U_{ab}\Vert ^{2} = (2/r_{ab}^{2}) \sum _{i} \operatorname{Tr} _{d_{ab}}(\delta  Q_{i}^{2}),
\end{aligned}\tag{2.1}
\end{gather*}

where \(\delta  Q_{i}=\delta  b_{i} I+\delta  T_{i}\). The longitudinal
contribution is the Gaussian variation; the transverse contributions are
the occupied-determinant variation. Formula (2.1) is a norm in the fast
fiber, leaving the entire slow Clifford module untouched.

For example, summing over all canonical center and endpoint-internal
coordinate directions gives

% DISPLAY c.9
% SOURCE sum_mu ||partial_mu U_ab||²
% SOURCE    = 18 n_a n_b (n_a+n_b) / r_ab².                         (2.2)
\begin{gather*}
\begin{aligned}
& \sum _{\mu} \Vert \partial _{\mu} U_{ab}\Vert ^{2} \\
& = 18 n_{a} n_{b} (n_{a}+n_{b}) / r_{ab}^{2}.
\end{aligned}\tag{2.2}
\end{gather*}

Indeed the center contribution is \(18(n_{a}+n_{b})/r^{2}\), while the
internal contributions are
\(18[n_{b}(n_{a}^{2}-1)+n_{a}(n_{b}^{2}-1)]/r^{2}\). Thus the adapted
vacuum has nonzero internal derivatives even at \(A=0\).

\subsubsection{2.1 Zero first jet of the frozen-charge
defect}\label{c-zero-first-jet-of-the-frozen-charge-defect}

Let \(q_{f}(z,A)\) be the frozen linear fast charge, with the exact
frozen eliminated-Gauss kinetic coefficient. It need not annihilate the
adapted vacuum at nonzero internal commutators. Nevertheless

% DISPLAY c.10
% SOURCE R(z,A)=q_f(z,A)U(z,A),
% SOURCE R(z,0)=0,   partial_A R(z,0)=0.                            (2.3)
\begin{gather*}
\begin{aligned}
& R(z,A)=q_{f}(z,A)U(z,A), \\
& R(z,0)=0, \quad  \partial _{A} R(z,0)=0.
\end{aligned}\tag{2.3}
\end{gather*}

Here is the algebra behind the first derivative. In the frozen
superalgebra the discrepancy between \(q_{f}^{2}\) and the quadratic
Hamiltonian consists of the residual internal gauge term and
background-commutator terms. The latter are quadratic in \(A\). The
derivative of the former annihilates \(U_{0}\), since \(U_{0}\) is
internal-gauge invariant. Differentiating the adapted ground equation
and using \(E_{0}'(0)=0\) therefore gives

% DISPLAY c.11
% SOURCE q_0 [q_f'(0)U_0 + q_0 U'(0)] = 0.
\begin{gather*}
\begin{aligned}
& q_{0} [q_{f}'(0)U_{0} + q_{0} U'(0)] = 0.
\end{aligned}\notag
\end{gather*}

The vector in brackets is center neutral and odd in boson parity: a
linear fast charge maps an even Gaussian to an odd Gaussian-polynomial
vector. The kernel of \(q_{0}\) is the even vacuum line, so the bracket
is zero. This parity step is required; the differentiated square alone
does not rule out a vacuum component.

There is also a direct commuting-axis proof. Vary one canonical internal
coordinate \(A_{i}^{(a)}=t T\), with all other internal matrices zero.
This whole one-parameter background is commuting. A unitary
diagonalizing the single Hermitian \(T\) splits the adapted fast system
into scalar color-pair oscillators at shifted relative vectors. The
kinetic-aware Gaussian is the scalar transverse vacuum written in the
chosen oblique slice, and the filled fermion vacuum is neutral for each
color-pair phase. The frozen charge annihilates this product exactly.
Thus every canonical internal partial derivative of \(R\) at zero
vanishes, and linearity gives the full first jet. This argument is only
along commuting coordinate axes; it does not incorrectly extend
annihilation to noncommuting internal tuples.

Center derivatives of \(R\) vanish as well because \(R(z,0)=0\)
identically in \(z\). Equation (2.3), rather than a false assertion that
the adapted fast charge has an exact kernel for all \(A\), is what the
centered homological calculation uses.

\subsection{3. Exact centered coefficient identity, including internal
derivatives}\label{c-exact-centered-coefficient-identity-including-internal-derivatives}

Under the simultaneous coefficient dilation

% DISPLAY c.12
% SOURCE z=L zeta,  A=L a,  Y=L^(-1/2) xi,
\begin{gather*}
\begin{aligned}
& z=L \zeta , \quad  A=L a, \quad  Y=L^{-1/2} \xi ,
\end{aligned}\notag
\end{gather*}

the exact reduced charge has the formal homogeneous inventory

% DISPLAY c.13
% SOURCE q(L) = L² Q_I^V(a) + L^(1/2) q_0(zeta,a)
% SOURCE        + L^-1 q_1 + L^(-5/2) q_2 + L^-4 q_3 + ... .      (3.1)
\begin{gather*}
\begin{aligned}
& q(L) = L^{2} Q_{I}^{V}(a) + L^{1/2} q_{0}(\zeta ,a) \\
& + L^{-1} q_{1} + L^{-5/2} q_{2} + L^{-4} q_{3} + \ldots  .
\end{aligned}\tag{3.1}
\end{gather*}

This is a coefficient expansion near the zero section, not a claim that
the internal interacting charge is bounded or small. In an intrinsic
multiblock chart the inverse Gauss matrix can create further terms; one
must not copy the finite three-term axial expansion without checking
this. All charge coefficients are first order in slow momenta. Since
\(Q_{I}^{V}(a)\) and its first derivative vanish at \(a=0\), its
anticommutators with \(q_{2}\) and \(q_{3}\) vanish there. Thus the
relevant centered Hamiltonian coefficients are exactly

% DISPLAY c.14
% SOURCE H_1 = {q_0,q_1},
% SOURCE H_2 = q_1² + {q_0,q_2},                                  (3.2)
\begin{gather*}
\begin{aligned}
& H_{1} = \{q_{0},q_{1}\}, \\
& H_{2} = q_{1}^{2} + \{q_{0},q_{2}\},
\end{aligned}\tag{3.2}
\end{gather*}

when evaluated at the centered internal jet. Terms from the quadratic
density of the intrinsic Gauss determinant belong to this inventory. The
fact that its linear-in-Y trace vanishes does not permit omitting its
quadratic triangle terms.

Put

% DISPLAY c.15
% SOURCE D_s = U* q_1 U,
% SOURCE B = (1-P) q_1 U.
\begin{gather*}
\begin{aligned}
& D_{s} = U^{*} q_{1} U, \\
& B = (1-P) q_{1} U.
\end{aligned}\notag
\end{gather*}

The derivative of a coefficient section \(f\) in \(q_{1}(Uf)\) lies in
the vacuum line. Therefore \(B\) is a multiplication map, but it
includes the internal derivatives in (2.1):

% DISPLAY c.16
% SOURCE q_1(Uf) = U D_s f + Bf.
\begin{gather*}
\begin{aligned}
& q_{1}(Uf) = U D_{s} f + Bf.
\end{aligned}\notag
\end{gather*}

Equation (2.3) gives the centered off-diagonal source and diagonal
compression

% DISPLAY c.17
% SOURCE (1-P) H_1 U = q_0 B,
% SOURCE U* H_2 U = D_s² + B*B.                                  (3.3)
\begin{gather*}
\begin{aligned}
& (1-P) H_{1} U = q_{0} B, \\
& U^{*} H_{2} U = D_{s}^{2} + B^{*}B.
\end{aligned}\tag{3.3}
\end{gather*}

These are differential-operator/form identities on the common smooth
coefficient core. There is no projection onto internal bound states or
restriction of the spectator Clifford module.

The exact neutral fast inverse then gives

% DISPLAY c.18
% SOURCE B* q_0 H_f^-1 q_0 B = B*B.                               (3.4)
\begin{gather*}
\begin{aligned}
& B^{*} q_{0} H_{f}^{-1} q_{0} B = B^{*}B.
\end{aligned}\tag{3.4}
\end{gather*}

Consequently the centered effective coefficient is \(D_{s}^{2}\), and
the entire additional multiplication coefficient cancels. It is
important to identify \(D_{s}\) with its first jet before replacing
\(D_{s}^{2}\) by a free slow Laplacian.

\subsubsection{3.1 Compression and its first
jet}\label{c-compression-and-its-first-jet}

The following checks supply that last identification at \(A=0\).

\begin{enumerate}
\def\labelenumi{\arabic{enumi}.}
\item
  The scalar Berry curvature is zero at \(A=0\), including pure internal
  and mixed center/internal components. A local gauge therefore has zero
  connection and zero first jet there.
\item
  The expectation of the quadratic-Y commutator charge is zero to first
  order in \(A\). The centered covariance and its first variation are
  spatial-diagonal, so contraction with \(\Gamma _{ij}\) vanishes.
\item
  The triangle orbital term and the fast-fast spin-Gauss term have an
  odd number of fast fermions and have zero vacuum compression. The same
  remains true after one internal derivative.
\item
  The \(\delta\) part of the mixed slow-fast fermion covariance cancels
  the half-density derivative. The remaining sign-covariance term is, up
  to the fixed charge sign and factors of \(i\), a linear combination of

  % DISPLAY c.19
  % SOURCE Gamma_n Im Tr_color[T_c F^-1 sign(D)],                 (3.5)
  \begin{gather*}
  \begin{aligned}
  & \Gamma _{n} \operatorname{Im}  \operatorname{Tr} _{\mathrm{color}}[T_{c} F^{-1} \operatorname{sign} (D)],
  \end{aligned}\tag{3.5}
  \end{gather*}

  where \(T_{c}\) is the Hermitian bifundamental action of a slow color
  generator. The trace leaves the spin matrix untraced. At the centered
  point it vanishes. Its first variation is a linear combination of

  % DISPLAY c.20
  % SOURCE r^-2 Im Tr(T_c delta T_n) Gamma_n,
  % SOURCE r^-2 Im Tr(T_c delta T_i) Gamma_i.
  \begin{gather*}
  \begin{aligned}
  & r^{-2} \operatorname{Im}  \operatorname{Tr} (T_{c} \delta  T_{n}) \Gamma _{n}, \\
  & r^{-2} \operatorname{Im}  \operatorname{Tr} (T_{c} \delta  T_{i}) \Gamma _{i}.
  \end{aligned}\notag
  \end{gather*}

  Both are zero because the trace of a product of two Hermitian matrices
  is real. This is why the compression must be checked as a jet, not
  only by its value at \(A=0\).
\end{enumerate}

In a realification of the complex pair space, the covariance is \(I/2\)
plus the complex structure times the realification of
\(\operatorname{sign} (D)\), multiplied by \(i/2\). The contracted mixed
Gauss coefficient is the realification of \(F^{-1} T_{c}\). Its
\(\delta\) trace gives the density derivative, and its complex-structure
trace is exactly the imaginary color trace in (3.5). This verifies the
color mechanism without assuming that \(F\) commutes with the internal
generators.

It follows that \(D_{s}\) has the centered free slow Dirac value and
first jet, so the coefficient in (3.3)-(3.4) is the full free
center-plus-internal kinetic operator. At nonzero internal coordinates,
the complete internal potential charge must of course be retained; (3.1)
never authorizes discarding it from an all-vector estimate.

\subsection{4. Root graph grading and exact virtual
denominators}\label{c-root-graph-grading-and-exact-virtual-denominators}

Use a root edge \(e=\{a,b\}\) and a triangle \(t=\{a,b,c\}\). For a root
oscillator let \(N_{e}\) count bosonic excitations plus fermionic
particles and holes relative to its vacuum. At the centered point

% DISPLAY c.21
% SOURCE H_f = 2 sum_e r_e N_e.
\begin{gather*}
\begin{aligned}
& H_{f} = 2 \sum _{e} r_{e} N_{e}.
\end{aligned}\notag
\end{gather*}

The leading complementary charge map has the decomposition

% DISPLAY c.22
% SOURCE B = sum_e B_e + sum_t B_t.                               (4.1)
\begin{gather*}
\begin{aligned}
& B = \sum _{e} B_{e} + \sum _{t} B_{t}.
\end{aligned}\tag{4.1}
\end{gather*}

Every \(B_{e}\) has \(N_{e}=2\) and every other \(N_{f}=0\). Its terms
are the Gaussian and determinant derivatives, the two-fast-variable
commutator with slow fermion output, and the complementary mixed
spin-Gauss contraction. Removing \(P\) removes their only possible
zero-quantum part.

Every \(B_{t}\) has one excitation on each of its three edges. Its
possible charge monomials are exactly of these types, with color
contractions along an oriented cycle:

% DISPLAY c.23
% SOURCE Y_e Y_f theta_g,
% SOURCE r_g^-1 Y_e p_f theta_g,
% SOURCE r_h^-1 theta_e theta_f theta_g.                          (4.2)
\begin{gather*}
\begin{aligned}
& Y_{e} Y_{f} \theta _{g}, \\
& r_{g}^{-1} Y_{e} p_{f} \theta _{g}, \\
& r_{h}^{-1} \theta _{e} \theta _{f} \theta _{g}.
\end{aligned}\tag{4.2}
\end{gather*}

Here \(e,f,g\) are the three distinct triangle edges, while \(h\) is one
of them. A fast Majorana applied once to the filled vacuum creates one
particle or hole. Different edges have independent Gaussian and Fock
factors, so no within-edge vacuum contraction occurs in a triangle
monomial.

Thus

% DISPLAY c.24
% SOURCE H_f B_e = 4r_e B_e,
% SOURCE H_f B_t = 2s_t B_t,       s_t=sum_(e in t) r_e.           (4.3)
\begin{gather*}
\begin{aligned}
& H_{f} B_{e} = 4r_{e} B_{e}, \\
& H_{f} B_{t} = 2s_{t} B_{t}, \quad  s_{t}=\sum _{e \in  t} r_{e}.
\end{aligned}\tag{4.3}
\end{gather*}

Distinct graph summands are orthogonal as fast-space maps. Pair sources
are nonvacuum on only their own edge. Triangle sources are odd in
combined excitation parity on exactly their three edges. The centered
charge preserves each combined excitation parity and annihilates vacuum
factors, so the orthogonality survives application of \(q_{0}\) and
\(H_{f}^{-1}\).

The finite dressing has the explicit graph resolution

% DISPLAY c.25
% SOURCE Z_e = -(4r_e)^-1 q_0 B_e,
% SOURCE Z_t = -(2s_t)^-1 q_0 B_t.                               (4.4)
\begin{gather*}
\begin{aligned}
& Z_{e} = -(4r_{e})^{-1} q_{0} B_{e}, \\
& Z_{t} = -(2s_{t})^{-1} q_{0} B_{t}.
\end{aligned}\tag{4.4}
\end{gather*}

In particular, as positive operators on the full slow Clifford module,

% DISPLAY c.26
% SOURCE Z_e* Z_e = (4r_e)^-1 B_e*B_e,
% SOURCE Z_t* Z_t = (2s_t)^-1 B_t*B_t.                           (4.5)
\begin{gather*}
\begin{aligned}
& Z_{e}^{*} Z_{e} = (4r_{e})^{-1} B_{e}^{*}B_{e}, \\
& Z_{t}^{*} Z_{t} = (2s_{t})^{-1} B_{t}^{*}B_{t}.
\end{aligned}\tag{4.5}
\end{gather*}

Replacing the second denominator by one arbitrarily selected root
frequency is incorrect. Replacing every denominator by the unrelated
global minimum throws away exactly the information needed for large
far-away blocks.

\subsection{5. Explicit coefficient seminorms and shape-independent
graph
estimates}\label{c-explicit-coefficient-seminorms-and-shape-independent-graph-estimates}

All constants below are finite at fixed \(N\). They can be calculated
from the normalized Lie structure constants, \(\gamma\) matrices, and
the displayed source monomials; no compact global gap-ratio chart occurs
in their definition.

One precise conservative convention is the following. Write each pair
source as a sum of elementary maps after factoring \(1/r_{e}\); evaluate
the elementary Hermite/Fock monomials on unit-frequency normalized vacua
and include their operator norms and absolute scalar coefficients in a
sum \(K_{e}\). Include the derivative contribution using (2.1). For a
triangle, group the monomials (4.2) into the three types, factoring
respectively

% DISPLAY c.27
% SOURCE (r_e r_f)^-1/2,
% SOURCE r_g^-1 (r_f/r_e)^1/2,
% SOURCE r_h^-1.
\begin{gather*}
\begin{aligned}
& (r_{e} r_{f})^{-1/2}, \\
& r_{g}^{-1} (r_{f}/r_{e})^{1/2}, \\
& r_{h}^{-1}.
\end{aligned}\notag
\end{gather*}

Let \(K_{t}\) be the sum of their unit-frequency Hermite/Fock norms
times absolute coefficients. \(\Gamma\) matrices have norm one, a real
Majorana has norm \(1/\sqrt{2}\), and degree-two normalized Gaussian
monomials have norms at most \(\sqrt{3}/2\). The remaining coefficients
are normalized color brackets and unit spatial-direction contractions.
This is a finite, explicit coefficient seminorm, not an unspecified
supremum over gap ratios. For example (2.2) and Cauchy-Schwarz give a
derivative-source contribution no larger than

% DISPLAY c.28
% SOURCE K_(e,derivative)² <= 162 N² n_a n_b(n_a+n_b).
\begin{gather*}
\begin{aligned}
& K_{e,\mathrm{derivative}}^{2} \le  162 N^{2} n_{a} n_{b}(n_{a}+n_{b}).
\end{aligned}\notag
\end{gather*}

Enlarge \(K_{t}\), if necessary, by a factor two to account for the
following weight reduction. Triangle geometry gives
\(r_{f}\le r_{e}+r_{g}\), and hence

% DISPLAY c.29
% SOURCE r_g^-1 sqrt(r_f/r_e)
% SOURCE    <= r_g^-1 + (r_e r_g)^-1/2
% SOURCE    <= (3/2)r_g^-1+(1/2)r_e^-1.
\begin{gather*}
\begin{aligned}
& r_{g}^{-1} \sqrt{r_{f}/r_{e}} \\
& \le  r_{g}^{-1} + (r_{e} r_{g})^{-1/2} \\
& \le  (3/2)r_{g}^{-1}+(1/2)r_{e}^{-1}.
\end{aligned}\notag
\end{gather*}

Also \((r_{e} r_{f})^{-1/2} \le  (r_{e}^{-1}+r_{f}^{-1})/2\). Therefore

% DISPLAY c.30
% SOURCE ||B_e|| <= K_e/r_e,
% SOURCE ||B_t|| <= K_t sum_(e in t) r_e^-1.                     (5.1)
\begin{gather*}
\begin{aligned}
& \Vert B_{e}\Vert  \le  K_{e}/r_{e}, \\
& \Vert B_{t}\Vert  \le  K_{t} \sum _{e \in  t} r_{e}^{-1}.
\end{aligned}\tag{5.1}
\end{gather*}

Set \(w_{t}=\sum _{e \in  t} r_{e}^{-2}\). Equations (4.5) and (5.1)
imply

% DISPLAY c.31
% SOURCE ||Z_e||² <= K_e²/(4r_e³),
% SOURCE ||Z_t||² <= 3 K_t² w_t/(2s_t).                          (5.2)
\begin{gather*}
\begin{aligned}
& \Vert Z_{e}\Vert ^{2} \le  K_{e}^{2}/(4r_{e}^{3}), \\
& \Vert Z_{t}\Vert ^{2} \le  3 K_{t}^{2} w_{t}/(2s_{t}).
\end{aligned}\tag{5.2}
\end{gather*}

The centered map by itself has only center derivatives; an internal
derivative requires an extension. For the following internal derivative
bounds, choose smooth gapped pairwise Gaussian/Fock transport of the
centered graph coefficient from \(A=0\), for example transport along
internal radial rays. This gives a concrete graph-factorized extension,
but it is not asserted to equal the actual nonzero-internal homological
source. The center derivative statement is unconditional. For these
derivatives, the same graph coefficients and one differentiated vacuum
factor give, with another computable coefficient seminorm \(K'_{g}\),

% DISPLAY c.32
% SOURCE ||nabla Z_e||² <= K'_e² r_e^-5,
% SOURCE ||nabla Z_t||² <= K'_t² (sum_(e in t) r_e^-1)² w_t/s_t. (5.3)
\begin{gather*}
\begin{aligned}
& \Vert \nabla  Z_{e}\Vert ^{2} \le  {K'_{e}}^{2} r_{e}^{-5}, \\
& \Vert \nabla  Z_{t}\Vert ^{2} \le  {K'_{t}}^{2} (\sum _{e \in  t} r_{e}^{-1})^{2} w_{t}/s_{t}.
\end{aligned}\tag{5.3}
\end{gather*}

In (5.3), \(\nabla\) differentiates the active graph coefficient and its
active oscillator factors. The full map also contains a vacuum factor on
every spectator edge. Its full derivative has the corrected bound

% DISPLAY c.33
% SOURCE ||nabla Z_g||² <= C_N[d_(g,active)+m_g W_2],
% SOURCE W_p=sum_e r_e^-p,                                       (5.4)
\begin{gather*}
\begin{aligned}
& \Vert \nabla  Z_{g}\Vert ^{2} \le  C_{N}[d_{g,\mathrm{active}}+m_{g} W_{2}], \\
& W_{p}=\sum _{e} r_{e}^{-p},
\end{aligned}\tag{5.4}
\end{gather*}

where \(m_{g}\) is the right side of (5.2). The extra term is necessary
when a spectator edge is much shorter. For adapted vacua it includes
internal derivatives of spectator factors as well as center derivatives.
The active triangle right side in (5.3) is at most a fixed numerical
multiple of \({K'_{t}}^{2} \sum _{e \in  t} r_{e}^{-5}\), by Holder and
\(s_{t}\ge \max  r_{e}\). Moreover

% DISPLAY c.34
% SOURCE W_3 W_2 <= (#edges) W_5.
\begin{gather*}
\begin{aligned}
& W_{3} W_{2} \le  (\#\text{edges}) W_{5}.
\end{aligned}\notag
\end{gather*}

Thus the full derivatives still have a summed \(C_{N} W_{5}\) majorant.
Graph-resolved weights must nevertheless be kept when an internal Yukawa
multiplier is to be paid.

\subsubsection{5.1 Perturbed finite source
spaces}\label{c-perturbed-finite-source-spaces}

The Gaussian/Fock statement also has a useful conditional continuation.
On the adapted pairwise-relative tube, unitary/symplectic transport
identifies the finite degree-at-most-three source spaces with their
centered spaces. Frequencies have fixed positive bounds
\(c r_{e} \le  \omega _{e,j} \le  C r_{e}\). A source odd in each
triangle edge therefore has complementary excitation energy at least
\(c s_{t}\), even when modes within one pair split or coincide.

If the exact leading Hamiltonian graph source \(S_{g}(A)\) has the
coefficient bounds

% DISPLAY c.35
% SOURCE ||S_e(A)||² <= C_N/r_e,
% SOURCE ||S_t(A)||² <= C_N s_t w_t,
\begin{gather*}
\begin{aligned}
& \Vert S_{e}(A)\Vert ^{2} \le  C_{N}/r_{e}, \\
& \Vert S_{t}(A)\Vert ^{2} \le  C_{N} s_{t} w_{t},
\end{aligned}\notag
\end{gather*}

then the inverse of the shifted excitation operator

% DISPLAY c.36
% SOURCE H_exc = H_fast - E_(0,total)
\begin{gather*}
\begin{aligned}
& H_{\mathrm{exc}} = H_{\mathrm{fast}} - E_{0,\mathrm{total}}
\end{aligned}\notag
\end{gather*}

on the relevant excited sector satisfies precisely the graph dressing
bounds (5.2), with enlarged constants. The vacuum-energy defect is
handled separately by the general frozen-energy lemma. An unshifted
inverse of the full fast Hamiltonian need not have these denominators,
because unrelated pair vacuum-energy defects would enter it. This is
solely the finite-source fast excitation inverse. No internal momentum
appears in its denominator.

To deduce a small adapted inverse-square remainder from the centered
identity, one additionally needs the exact diagonal \(H_{2}\) inventory,
including its density, ordering, and internal-potential cross terms.
Under coefficientwise perturbation bounds of order \(\kappa\), the
finite resolvent identity gives a defect at most
\(C_{N} \kappa  W_{2}\). Establishing all those hypotheses from the full
reduced Hamiltonian is separate from the graph calculation.

\subsection{6. A graph-local complete slow form
lemma}\label{c-a-graph-local-complete-slow-form-lemma}

This is an all-vector statement, not a finite-Gaussian approximation to
an internal wavefunction.

Let

% DISPLAY c.37
% SOURCE H_s = T_cent + sum_a H_a,
% SOURCE H_a = T_a + V_a + Y_a,
% SOURCE V_a >= 0,     ||Y_a(A)|| <= c_(n_a) alpha_a.
\begin{gather*}
\begin{aligned}
& H_{s} = T_{\mathrm{cent}} + \sum _{a} H_{a}, \\
& H_{a} = T_{a} + V_{a} + Y_{a}, \\
& V_{a} \ge  0, \quad  \Vert Y_{a}(A)\Vert  \le  c_{n_{a}} \alpha _{a}.
\end{aligned}\notag
\end{gather*}

The averaged sixteen-charge identity makes \(H_{a}\) nonnegative also on
colored internal vectors. The ordinary physical lower-rank Hardy
estimate is not being invoked. For a graph-local multiplication map
\(Z_{g}\), suppose

\begin{itemize}
\tightlist
\item
  it is an active-graph coefficient, depending on internal coordinates
  at graph vertices only, tensored with the vacua on spectator edges;
\item
  it acts on internal Clifford factors only at these vertices;
\item
  it is even in total fermion parity, as is \(q_{0}^{-1} B_{g}\);
\item
  \(\Vert Z_{g}\Vert ^{2} \le  m_{g}(z)\) and the full derivative obeys
  the center-only majorant \(d_{g}\) in (5.4).
\end{itemize}

For a block \(a\) outside \(g\), its internal potential charge commutes
exactly with the map. The complete internal charge satisfies

% DISPLAY c.38
% SOURCE Q_(a,alpha) Z_g = Z_g Q_(a,alpha)+C_(a,alpha,g),
\begin{gather*}
\begin{aligned}
& Q_{a,\alpha} Z_{g} = Z_{g} Q_{a,\alpha}+C_{a,\alpha ,g},
\end{aligned}\notag
\end{gather*}

where \(C_{a,\alpha ,g}\) is bounded multiplication coming only from
derivatives of spectator vacua. Its squared norm is bounded by the
corresponding derivative majorant, with the finite Clifford contraction
factor. Averaging the exact charge squares and applying Young gives

% DISPLAY c.39
% SOURCE H_a[Z_g f] <= 2 integral m_g(z) [H_a-density of f]
% SOURCE               + C_N integral d_(a,g)(z)|f|².            (6.1)
\begin{gather*}
\begin{aligned}
& H_{a}[Z_{g} f] \le  2 \int  m_{g}(z) [H_{a}\text{-density of }f] \\
& + C_{N} \int  d_{a,g}(z)\vert f\vert ^{2}.
\end{aligned}\tag{6.1}
\end{gather*}

In particular no unrelated \(\alpha _{a} m_{g}\) loss appears. It would
be incorrect to assert exact commutation while differentiating a moving
spectator vacuum. For a graph vertex use the elementary uncompressed
estimate, with no internal spectral assumptions,

% DISPLAY c.40
% SOURCE H_a[Z_g f]
% SOURCE   <= 2 integral m_g(z) [H_a-density of f]
% SOURCE      + 3c_(n_a) integral alpha_a m_g(z)|f|²
% SOURCE      + 2 integral |partial_Aa Z_g|² |f|².               (6.2)
\begin{gather*}
\begin{aligned}
& H_{a}[Z_{g} f] \\
& \le  2 \int  m_{g}(z) [H_{a}\text{-density of }f] \\
& + 3c_{n_{a}} \int  \alpha _{a} m_{g}(z)\vert f\vert ^{2} \\
& + 2 \int  \vert \partial _{Aa} Z_{g}\vert ^{2} \vert f\vert ^{2}.
\end{aligned}\tag{6.2}
\end{gather*}

The notation for a weighted internal density is harmless here because
its weight depends only on centers; it is the nonnegative internal form
on every fixed-center fiber. The center kinetic estimate has the same
derivative-product bound with no Yukawa term.

For pair and triangle dressings from (5.2),

% DISPLAY c.41
% SOURCE sup m_g <= C_N r_*^-3.
\begin{gather*}
\begin{aligned}
& \sup  m_{g} \le  C_{N} r_{*}^{-3}.
\end{aligned}\notag
\end{gather*}

At a pair endpoint, \(\alpha _{a}\le \kappa  r_{e}\) gives

% DISPLAY c.42
% SOURCE alpha_a m_e <= C_N kappa r_e^-2.
\begin{gather*}
\begin{aligned}
& \alpha _{a} m_{e} \le  C_{N} \kappa  r_{e}^{-2}.
\end{aligned}\notag
\end{gather*}

At any triangle vertex,
\(\alpha _{a}\le \kappa  \min (\text{ incident } r_{e})\le \kappa  s_{t}\),
and hence

% DISPLAY c.43
% SOURCE alpha_a m_t <= C_N kappa w_t.                           (6.3)
\begin{gather*}
\begin{aligned}
& \alpha _{a} m_{t} \le  C_{N} \kappa  w_{t}.
\end{aligned}\tag{6.3}
\end{gather*}

This remains true when the opposite edge is arbitrarily shorter than the
two incident edges. The factor \(s_{t}\) from the virtual denominator is
doing essential work.

Combining (6.1)-(6.3), then using positivity of \(H_{s}\) to sum the
finitely many graph maps, proves

% DISPLAY c.44
% SOURCE H_s[sum_g Z_g f]
% SOURCE   <= C_N r_*^-3 H_s[f]
% SOURCE      + C_N kappa integral W_2 |f|²
% SOURCE      + C_N integral W_5 |f|².                           (6.4)
\begin{gather*}
\begin{aligned}
& H_{s}[\sum _{g} Z_{g} f] \\
& \le  C_{N} r_{*}^{-3} H_{s}[f] \\
& + C_{N} \kappa  \int  W_{2} \vert f\vert ^{2} \\
& + C_{N} \int  W_{5} \vert f\vert ^{2}.
\end{aligned}\tag{6.4}
\end{gather*}

The constants can be chosen from the finite sums of \(K_{g}^{2}\),
\({K'_{g}}^{2}\), vertex counts, and \(c_{n_{a}}\). For instance a
direct Cauchy-Schwarz sum multiplies their sum by the number of graphs
\(\binom{k}{2}+\binom{k}{3}\). This gives an explicit finite-N recipe,
without a global ratio parameter.

The full derivative terms in this argument include spectator factors;
their product weight is controlled by
\(W_{3} W_{2} \le  (\#\text{edges}) W_{5}\). The same argument applies
to a genuinely adapted dressing if its active-graph coefficients and
spectator factors satisfy the hypotheses. In particular, (6.4) retains
arbitrary internal momenta. No absolute-value inverse expansion of
\(H_{f}+H_{I}\) has been used. It also explains why the coarse estimate
\(\Vert Z\Vert ^{2}\le C W_{3}\) should not first be multiplied by
\(\sum  \alpha _{a}\): that would incorrectly charge a distant large
block against an unrelated short root.

\subsubsection{6.1 Direct center-Hardy
payment}\label{c-direct-center-hardy-payment}

For each pair, in the mass center metric,

% DISPLAY c.45
% SOURCE integral r_ab^-2 |f|²
% SOURCE    <= (4/49)(1/n_a+1/n_b)^-1 T_cent[f].                 (6.5)
\begin{gather*}
\begin{aligned}
& \int  r_{ab}^{-2} \vert f\vert ^{2} \\
& \le  (4/49)(1/n_{a}+1/n_{b})^{-1} T_{\mathrm{cent}}[f].
\end{aligned}\tag{6.5}
\end{gather*}

The inequality holds for Hilbert-valued coefficients and unitary
connections. Also \(W_{5} \le  r_{*}^{-3} W_{2}\). Hence the last two
terms in (6.4) are bounded by an arbitrarily small center kinetic
fraction by first choosing \(\kappa\) small at fixed \(N\) and then
\(r_{*}\) large. The first term is a small fraction of the entire
interacting slow form, not a fixed internal kinetic loss.

\clearpage
\section{Actual adapted source
inventory}\label{d-actual-adapted-source-inventory}

Technical derivation for review. Equation and subsection numbers in this
appendix are local to the source derivation. Historical local labels are
retained, including numbering gaps or nonsequential labels. The
accompanying source notes retain the original expressions for
comparison.

\subsection{1. Notation and exact
origin}\label{d-notation-and-exact-origin}

Use the kinetic-one normalization and the intrinsic slice in Appendix A.
The only internal thinness hypothesis is

% DISPLAY d.1
% SOURCE alpha_a+alpha_b <= kappa r_ab,
\begin{gather*}
\begin{aligned}
& \alpha _{a}+\alpha _{b} \le  \kappa  r_{ab},
\end{aligned}\notag
\end{gather*}

with fixed sufficiently small \(\kappa\). Set \(r_{*}=\min  r_{ab}\).
All finite constants below depend on \(N\) but not on a ratio of
unrelated gaps.

To avoid overloading center \(C\), use these coefficient maps:

% DISPLAY d.2
% SOURCE F = E*B_(Z+A),             C = P_N B_A,
% SOURCE L = E*B_Y,                 N_Y = P_N B_Y,
% SOURCE I_Y = P_I B_Y,             C_Y = P_C B_Y,
% SOURCE K = R^-1 C_Y*,             J_2 = K C_Y.
\begin{gather*}
\begin{aligned}
& F = E^{*}B_{Z+A}, \quad  C = P_{N} B_{A}, \\
& L = E^{*}B_{Y}, \quad  N_{Y} = P_{N} B_{Y}, \\
& I_{Y} = P_{I} B_{Y}, \quad  C_{Y} = P_{C} B_{Y}, \\
& K = R^{-1} C_{Y}^{*}, \quad  J_{2} = K C_{Y}.
\end{aligned}\notag
\end{gather*}

Thus \(F\) is pair-block diagonal, \(C\) is pair-block diagonal as a map
from gauge parameters into normal fast coordinates, and

% DISPLAY d.3
% SOURCE M=F+L-J_2,   W=I+KK*.
\begin{gather*}
\begin{aligned}
& M=F+L-J_{2}, \quad  W=I+KK^{*}.
\end{aligned}\notag
\end{gather*}

The maps \(L\) and \(N_{Y}\) between distinct fast root pairs have
triangle support. They have no same-pair block. The maps \(I_{Y}\) and
\(C_{Y}\) use one root and its endpoints. The center metric is
\(g_{C}=I+K^{*}K\); the internal and fast tangent metrics are exactly
identity.

The definitive flattened kinetic form is

% DISPLAY d.4
% SOURCE ||(p_H+i rho)u||_(g_S^-1)^2 + ||D u||²,
% SOURCE D=W^(1/2)M^-*[S_f-O*Sg_S^-1(p_H+i rho)],
% SOURCE rho=d_H log d,       d=sqrt(det M).                     (1.1)
\begin{gather*}
\begin{aligned}
& \Vert (p_{H}+i \rho )u\Vert _{g_{S}^{-1}}^{2} + \Vert D u\Vert ^{2}, \\
& D=W^{1/2}M^{-*}[S_{f}-O^{*}Sg_{S}^{-1}(p_{H}+i \rho )], \\
& \rho =d_{H} \log  d, \quad  d=\sqrt{\det  M}.
\end{aligned}\tag{1.1}
\end{gather*}

All products below inherit this ordering. Fast adjoints are taken in the
flat fiber measure. The normal connection is retained in center
derivatives. Nothing is obtained by substituting a momentum constraint
into a squared operator.

\subsection{2. Exact coefficients needed through
H2}\label{d-exact-coefficients-needed-through-h2}

Use homogeneous weights

% DISPLAY d.5
% SOURCE (Z,A) -> L_scale (Z,A),  Y -> L_scale^-1/2 Y,
% SOURCE p_s -> L_scale^-1 p_s,  p_F -> L_scale^1/2 p_F.
\begin{gather*}
\begin{aligned}
& (Z,A) \to  L_{\mathrm{scale}} (Z,A), \quad  Y \to  L_{\mathrm{scale}}^{-1/2} Y, \\
& p_{s} \to  L_{\mathrm{scale}}^{-1} p_{s}, \quad  p_{F} \to  L_{\mathrm{scale}}^{1/2} p_{F}.
\end{aligned}\notag
\end{gather*}

Here \(s\) comprises centers and all internal coordinates. Keep the
internal potential and internal Yukawa terms in their complete internal
Hamiltonians throughout.

Write

% DISPLAY d.6
% SOURCE log d = log d_0 + phi_2 + terms of Y-degree >=3,
% SOURCE d_0=sqrt(det F),
% SOURCE phi_2=-(1/2)tr(F^-1 J_2)
% SOURCE         -(1/4)tr(F^-1 L F^-1 L),                       (2.1)
% SOURCE rho_0=d_s log d_0,    rho_F2=d_F phi_2,
% SOURCE u_s=p_s+i rho_0.
\begin{gather*}
\begin{aligned}
& \log  d = \log  d_{0} + \phi _{2} + \text{ terms of Y-degree } \ge 3, \\
& d_{0}=\sqrt{\det  F}, \\
& \phi _{2}=-(1/2)\operatorname{tr} (F^{-1} J_{2}) \\
& -(1/4)\operatorname{tr} (F^{-1} L F^{-1} L),
\end{aligned}\tag{2.1} \\
\begin{aligned}
& \rho _{0}=d_{s} \log  d_{0}, \quad  \rho _{F2}=d_{F} \phi _{2}, \\
& u_{s}=p_{s}+i \rho _{0}.
\end{aligned}\notag
\end{gather*}

There is no degree-one term because \(\operatorname{tr} (F^{-1} L)=0\)
for every A, not only at \(A=0\). The triangle term in (2.1) must be
retained.

The Gauss derivative stack has coefficients

% DISPLAY d.7
% SOURCE D = D_0 + D_1 + D_2 + higher homogeneous orders,
\begin{gather*}
\begin{aligned}
& D = D_{0} + D_{1} + D_{2} + \text{ higher homogeneous orders },
\end{aligned}\notag
\end{gather*}

of scaling degrees \(1/2\), \(-1\), and \(-5/2\), respectively. They are

% DISPLAY d.8
% SOURCE D_0 = -F^-1 C* p_F,                                     (2.2)
% SOURCE 
% SOURCE D_1 = F^-1(S_f-N_Y* p_F)
% SOURCE          +F^-1 L* F^-1 C* p_F,                         (2.3)
\begin{gather*}
\begin{aligned}
& D_{0} = -F^{-1} C^{*} p_{F},
\end{aligned}\tag{2.2} \\
\begin{aligned}
& D_{1} = F^{-1}(S_{f}-N_{Y}^{*} p_{F}) \\
& +F^{-1} L^{*} F^{-1} C^{*} p_{F},
\end{aligned}\tag{2.3}
\end{gather*}

and

% DISPLAY d.9
% SOURCE D_2 = -F^-1[I_Y* u_I+(C_Y*+F K)u_C+i C* rho_F2]
% SOURCE       -F^-1 L* F^-1(S_f-N_Y* p_F)
% SOURCE       -F^-1 L* F^-1 L* F^-1 C* p_F
% SOURCE       -F^-1 J_2* F^-1 C* p_F
% SOURCE       -(1/2)KK* F^-1 C* p_F.                           (2.4)
\begin{gather*}
\begin{aligned}
& D_{2} = -F^{-1}[I_{Y}^{*} u_{I}+(C_{Y}^{*}+F K)u_{C}+i C^{*} \rho _{F2}] \\
& -F^{-1} L^{*} F^{-1}(S_{f}-N_{Y}^{*} p_{F}) \\
& -F^{-1} L^{*} F^{-1} L^{*} F^{-1} C^{*} p_{F} \\
& -F^{-1} J_{2}^{*} F^{-1} C^{*} p_{F} \\
& -(1/2)KK^{*} F^{-1} C^{*} p_{F}.
\end{aligned}\tag{2.4}
\end{gather*}

These follow by expanding

% DISPLAY d.10
% SOURCE M^-*=F^-1-F^-1L*F^-1
% SOURCE          +F^-1L*F^-1L*F^-1+F^-1J_2*F^-1+...,
% SOURCE W^(1/2)=I+(1/2)KK*+...,
% SOURCE O*Sg_S^-1p_H
% SOURCE    =C*p_F+N_Y*p_F+I_Y*p_I+(C_Y*+F K)p_C+... .          (2.5)
\begin{gather*}
\begin{aligned}
& M^{-*}=F^{-1}-F^{-1}L^{*}F^{-1} \\
& +F^{-1}L^{*}F^{-1}L^{*}F^{-1}+F^{-1}J_{2}^{*}F^{-1}+\ldots , \\
& W^{1/2}=I+(1/2)KK^{*}+\ldots , \\
& O^{*}Sg_{S}^{-1}p_{H} \\
& =C^{*}p_{F}+N_{Y}^{*}p_{F}+I_{Y}^{*}p_{I}+(C_{Y}^{*}+F K)p_{C}+\ldots  .
\end{aligned}\tag{2.5}
\end{gather*}

For example \(L^{*}K p_{C}\) first enters the next Gauss order, and
\(g_{C}^{-1}-I\) first contributes to direct center kinetic energy at
scaling degree -5; neither belongs to H2. The \(F K\) in (2.4) is
essential and becomes \(C_{Y}^{*}\) at \(A=0\), so the centered
center-orbital coefficient is twice \(C_{Y}^{*}\).

The exact leading fast Hamiltonian comprises the direct \(p_{F}^{2}\),
\(D_{0}^{*}D_{0}\), the full adapted quadratic potential, and the fast
fermion mass. Its product vacuum is \(U(A)\), with energy
\(E_{0,\mathrm{total}}\). Define

% DISPLAY d.11
% SOURCE H_exc = H_fast-E_0,total.
\begin{gather*}
\begin{aligned}
& H_{\mathrm{exc}} = H_{\mathrm{fast}}-E_{0,\mathrm{total}}.
\end{aligned}\notag
\end{gather*}

The vacuum-energy multiplier is handled by Appendix \(B\). It is not
inserted into a virtual excitation denominator.

The next two homogeneous form coefficients are exactly

% DISPLAY d.12
% SOURCE H_1 = V_3+Y_Y +D_0*D_1+D_1*D_0,                       (2.6)
% SOURCE 
% SOURCE H_2 = ||u_s (.)||² +V_4
% SOURCE         +2 Re <p_F(.),i rho_F2(.)>
% SOURCE         +||D_1(.)||²+2 Re <D_0(.),D_2(.)>.             (2.7)
\begin{gather*}
\begin{aligned}
& H_{1} = V_{3}+Y_{Y} +D_{0}^{*}D_{1}+D_{1}^{*}D_{0},
\end{aligned}\tag{2.6} \\
\begin{aligned}
& H_{2} = \Vert u_{s} (.)\Vert ^{2} +V_{4} \\
& +2 \operatorname{Re}  \langle p_{F}(.),i \rho _{F2}(.)\rangle \\
& +\Vert D_{1}(.)\Vert ^{2}+2 \operatorname{Re}  \langle D_{0}(.),D_{2}(.)\rangle .
\end{aligned}\tag{2.7}
\end{gather*}

Here \(Y_{Y}\) is the Yukawa term linear in \(Y\), including both
slow-fast and fast-fast fermion terms. \(V_{3}\) is the cross term
between \([Z+A,Y]\) and \([Y,Y]\); \(V_{4}\) is the pure quartic fast
potential. The internal potential term paired with \([Y,Y]\) is already
in the adapted quadratic potential. Equations (2.6)-(2.7) therefore
include the terms which appear as anticommutators of the internal
potential charge with higher inverse-Gauss charge coefficients; there is
no permission to omit them in a charge derivation. This Hamiltonian
derivation fixes them directly from the original scalar potential and
the exact kinetic form.

\subsection{3. Actual H1 source and its graph
weights}\label{d-actual-h1-source-and-its-graph-weights}

Define the actual leading source

% DISPLAY d.13
% SOURCE S(A) = H_1 U(A).                                        (3.1)
\begin{gather*}
\begin{aligned}
& S(A) = H_{1} U(A).
\end{aligned}\tag{3.1}
\end{gather*}

It is complementary by boson parity. Each summand is a finite
Gaussian-polynomial/Fock map on the full slow Clifford module. It has
the graph decomposition

% DISPLAY d.14
% SOURCE S=sum_e S_e+sum_t S_t,
\begin{gather*}
\begin{aligned}
& S=\sum _{e} S_{e}+\sum _{t} S_{t},
\end{aligned}\notag
\end{gather*}

where an edge source has exactly two excitations on that edge
(\(N_{e}=2\)) and a triangle source has exactly one excitation on every
triangle edge (\(N_{e}=N_{f}=N_{g}=1\)). In the pair case these are one
bosonic and one fast-fermionic excitation. Graph subspaces remain
orthogonal and invariant under the adapted quadratic oscillator. Parity
alone would not be the complete finite-sector statement.

The inventory is explicit:

\begin{itemize}
\tightlist
\item
  \(Y_{Y} U\): a pair has \(Y_{e} \theta _{e} \theta _{s}\); a triangle
  has \(Y_{e} \theta _{f} \theta _{g}\).
\item
  \(V_{3} U\): only three distinct triangle edges, one \(Y\) on each.
\item
  \(D_{0}^{*}F^{-1} S_{f} U\) and its adjoint: pair terms have one
  \(p_{e}\) and one fast fermion; triangle terms have
  \(p_{e} \theta _{f} \theta _{g}\).
\item
  \(D_{0}^{*}F^{-1} N_{Y}^{*}p_{F} U\) and its adjoint: a triangle has
  \(p_{e} Y_{f} p_{g}\).
\item
  The last term of \(D_{1}\) yields the same triangle support with one
  additional factor \(F_{f}^{-1} C_{f}^{*}\).
\end{itemize}

A slow internal derivative does not occur in H1. Only the pair source
contains internal/center Clifford generators; the triangle source acts
trivially on the internal Clifford factors.

For a triangle let \(s_{t}=\sum _{e \in  t}r_{e}\),
\(w_{t}=\sum _{e \in  t}r_{e}^{-2}\). Pairwise thinness gives

% DISPLAY d.15
% SOURCE ||F_e^-1|| <= C/r_e,    ||F_e^-1 C_e*|| <= C kappa.
\begin{gather*}
\begin{aligned}
& \Vert F_{e}^{-1}\Vert  \le  C/r_{e}, \quad  \Vert F_{e}^{-1} C_{e}^{*}\Vert  \le  C \kappa .
\end{aligned}\notag
\end{gather*}

On a fixed finite Gaussian degree, \(Y_{e}\) costs \(C r_{e}^{-1/2}\),
and \(p_{e}\) costs \(C r_{e}^{1/2}\). Thus

% DISPLAY d.16
% SOURCE ||S_e||² <= C_N/r_e,
% SOURCE ||S_t||² <= C_N s_t w_t.                               (3.2)
\begin{gather*}
\begin{aligned}
& \Vert S_{e}\Vert ^{2} \le  C_{N}/r_{e}, \\
& \Vert S_{t}\Vert ^{2} \le  C_{N} s_{t} w_{t}.
\end{aligned}\tag{3.2}
\end{gather*}

For the cubic potential, use

% DISPLAY d.17
% SOURCE s_t/(r_e r_f r_g)
% SOURCE    =1/(r_e r_f)+1/(r_f r_g)+1/(r_g r_e) <= w_t.
\begin{gather*}
\begin{aligned}
& s_{t}/(r_{e} r_{f} r_{g}) \\
& =1/(r_{e} r_{f})+1/(r_{f} r_{g})+1/(r_{g} r_{e}) \le  w_{t}.
\end{aligned}\notag
\end{gather*}

For the Gauss triangle \(p_{e} Y_{f} p_{g}/r_{e}\), its norm is bounded
by

% DISPLAY d.18
% SOURCE C kappa sqrt(r_g/(r_e r_f))
% SOURCE    <= C kappa sqrt(s_t) (r_e r_f)^-1/2.
\begin{gather*}
\begin{aligned}
& C \kappa  \sqrt{r_{g}/(r_{e} r_{f})} \\
& \le  C \kappa  \sqrt{s_{t}} (r_{e} r_{f})^{-1/2}.
\end{aligned}\notag
\end{gather*}

The inverse-L correction has one additional bounded factor \(\kappa\).
These are shape-independent estimates, including when one edge is much
shorter than the other two.

Transporting only the finite oscillator coefficient spaces to their
centered spaces by the actual pair functional calculus gives the
coefficientwise perturbation estimates

% DISPLAY d.19
% SOURCE ||S_e(A)-S_e(0)|| <= C_N kappa r_e^-1/2,
% SOURCE ||S_t(A)-S_t(0)|| <= C_N kappa sqrt(s_t w_t).             (3.3)
\begin{gather*}
\begin{aligned}
& \Vert S_{e}(A)-S_{e}(0)\Vert  \le  C_{N} \kappa  r_{e}^{-1/2}, \\
& \Vert S_{t}(A)-S_{t}(0)\Vert  \le  C_{N} \kappa  \sqrt{s_{t} w_{t}}.
\end{aligned}\tag{3.3}
\end{gather*}

This is an estimate for the actual source (3.1), not a definition by
transport. Every new Gauss H1 term has an explicit factor
\(F^{-1} C^{*}\), and the remaining actual coefficients have relative
\(O(\kappa )\) variation in the pair matrices and vacuum factors.

\subsection{4. Compressed H2 and its actual scalar/Clifford
multiplier}\label{d-compressed-h2-and-its-actual-scalarclifford-multiplier}

Separate the complete slow Berry-covariant form

% DISPLAY d.20
% SOURCE H_s^B=T_cent^B+sum_a(T_a^B+V_a+Y_a)
\begin{gather*}
\begin{aligned}
& H_{s}^{B}=T_{\mathrm{cent}}^{B}+\sum _{a}(T_{a}^{B}+V_{a}+Y_{a})
\end{aligned}\notag
\end{gather*}

from the vacuum compression of (2.7). The remainder is a multiplication
operator \(W(A)\) on the entire slow Clifford fiber. The following
procedure is an exact specification of \(W\):

\begin{enumerate}
\def\labelenumi{\arabic{enumi}.}
\item
  In \(\Vert u_{s}(Uf)\Vert ^{2}\), retain
  \(\Vert \nabla _{s}^{B} f\Vert ^{2}\); put the Born-Huang multiplier
  and \(\vert \rho _{0}\vert ^{2}+\operatorname{div}  \rho _{0}\) into
  \(W\).
\item
  Put the vacuum expectations of \(V_{4}\), the displayed fast-density
  cross term, and \(D_{1}^{*}D_{1}\) into \(W\).
\item
  In \(2 \operatorname{Re} \langle D_{0}(Uf),D_{2}(Uf)\rangle\), keep
  the written ordering. For the slow-orbital part put

  % DISPLAY d.21
  % SOURCE L_mu = component_mu{F^-1[I_Y*,C_Y*+F K]},
  % SOURCE <D_0 U,-L_mu U> = i j_mu,       j_mu real.            (4.1)
  \begin{gather*}
  \begin{aligned}
  & L_{\mu} = \operatorname{component} _{\mu}\{F^{-1}[I_{Y}^{*},C_{Y}^{*}+F K]\}, \\
  & \langle D_{0} U,-L_{\mu} U\rangle  = i j_{\mu}, \quad  j_{\mu} \text{ real }.
  \end{aligned}\tag{4.1}
  \end{gather*}

  This coefficient is purely imaginary: the bosonic Gaussian and the
  orbital coefficient maps are real, while \(D_{0}\) contains one fast
  momentum. The fermionic factor is its norm. Its first-order
  coefficient on \(f\) therefore integrates to the scalar multiplier
  \(-\operatorname{div}  j\). All remaining terms, including derivatives
  on \(U\) and \(i \rho _{0}\), are placed in \(W\). A scalar Berry term
  in \(p_{\mu} U\) has zero real pairing with \(i j_{\mu}\), so this
  step is gauge covariant.
\end{enumerate}

An equivalent exact scalar formula removes even the implicit derivative
bookkeeping. Write \(U=g(Y;s)v(s)\), with \(g\) real, put

% DISPLAY d.22
% SOURCE A_op=F^-1 C*,    a=A_op d_F log g,
% SOURCE B_mu=component_mu{F^-1[I_Y*,C_Y*+F K]}.
\begin{gather*}
\begin{aligned}
& A_{\mathrm{op}}=F^{-1} C^{*}, \quad  a=A_{\mathrm{op}} d_{F} \log  g, \\
& B_{\mu}=\operatorname{component} _{\mu}\{F^{-1}[I_{Y}^{*},C_{Y}^{*}+F K]\}.
\end{aligned}\notag
\end{gather*}

Then \(D_{0} U=i a U\), and the entire slow-orbital compression is

% DISPLAY d.23
% SOURCE V_orb=-sum_mu <partial_mu(a dot B_mu)>_U
% SOURCE          -2 sum_mu rho_(0,mu)<a dot B_mu>_U.            (4.3)
\begin{gather*}
\begin{aligned}
& V_{\mathrm{orb}}=-\sum _{\mu} \langle \partial _{\mu}(a \cdot  B_{\mu})\rangle _{U} \\
& -2 \sum _{\mu} \rho _{0,\mu}\langle a \cdot  B_{\mu}\rangle _{U}.
\end{aligned}\tag{4.3}
\end{gather*}

The derivative acts on the displayed fast coefficient; Gaussian-density
differentiation has already canceled the derivative of the fiber
expectation. Fock/Berry derivative terms have zero real part. The direct
fast-density term and the \(\rho _{F2}\) term in \(D_{2}\) combine
exactly to

% DISPLAY d.24
% SOURCE V_fast-density=tr_F[(I+C F^-2 C*) Hess_Y phi_2].          (4.4)
\begin{gather*}
\begin{aligned}
& V_{\mathrm{fast\text{-}density}}=\operatorname{tr} _{F}[(I+C F^{-2} C^{*}) \operatorname{Hess} _{Y} \phi _{2}].
\end{aligned}\tag{4.4}
\end{gather*}

Since \(\phi _{2}\) is quadratic in \(Y\), its Hessian is a coefficient
matrix independent of \(Y\). These formulas follow by one integration by
parts and retain all ordering terms.

No slow derivative has been discarded or bounded by a momentum cutoff.
In particular differentiating \(j\), or the coefficient in (4.3),
includes internal derivatives at \(A=0\). Those terms participate in the
centered cancellation.

\subsubsection{4.1 Complete root support and coefficient
bounds}\label{d-complete-root-support-and-coefficient-bounds}

Every surviving vacuum contraction in \(W\) is supported on one edge or
the edges of a triangle. The following list accounts for all terms in
(2.7):

\begin{itemize}
\tightlist
\item
  Born-Huang derivatives on two different pair vacua are orthogonal
  after removing the scalar Berry connection. Each remaining term is
  bounded by \(C_{N}/r_{e}^{2}\).
\item
  \(\rho _{0}\) is a sum of pair gradients. Their inner product vanishes
  for disjoint endpoint sets. A cross product on two incident edges has
  weight \(1/(r_{e} r_{f})\) and belongs to their triangle. The
  divergence is pairwise.
\item
  Wick contractions in \(V_{4}\) involve either one repeated root or two
  roots sharing an endpoint. Their weights are \(r_{e}^{-2}\) or
  \((r_{e} r_{f})^{-1}\).
\item
  The first term of \(\phi _{2}\) is a sum of \(Y_{e}^{2}/r_{e}^{2}\)
  contractions. Its second term follows a root-space path
  \(e \to  f \to  e\) and has \(Y_{g}^{2}/(r_{e} r_{f})\), where e,f,g
  form a triangle. A fast derivative and the momentum pairing remove the
  \(Y_{g}\) scale; the resulting weights are again \(r_{e}^{-2}\) or
  \((r_{e} r_{f})^{-1}\).
\item
  \(D_{1}\) contains a pair spin term with weight \(r_{e}^{-1}\), and
  triangle terms with weights \(r_{e}^{-1}\) or
  \(r_{e}^{-1} \sqrt{r_{g}/r_{f}}\). Squaring on the vacuum preserves
  graph orthogonality. The triangle inequality reduces the latter square
  to \(C w_{t}\).
\item
  The slow-orbital part of \(D_{2}\) paired with \(D_{0}\) gives
  \(j_{e}=O(\kappa /r_{e})\) on endpoint coordinates. Its derivative is
  \(O(r_{e}^{-2})\) and varies from its centered value by
  \(O(\kappa /r_{e}^{2})\). Derivatives of spectator vacua are
  orthogonal in this compression; the products with \(\rho _{0}\) have
  pair/triangle support.
\item
  In \(D_{0}^{*}F^{-1}L^{*}F^{-1} S_{f}\), the two bosonic factors lie
  on distinct roots, so its vacuum compression is zero.
\item
  In \(D_{0}^{*}F^{-1}L^{*}F^{-1} N_{Y}^{*}p_{F}\), nonzero Wick
  contractions force the second root path to close the same triangle.
  The two possible pairings have weights \(\kappa /(r_{f} r_{g})\) and
  \(\kappa /(r_{e} r_{f})\).
\item
  In the two-L term of \(D_{2}\) the same closure gives the preceding
  weights with \(\kappa\) squared.
\item
  In the \(J_{2}\) and \(KK^{*}\) terms, the two orbital center maps
  vanish against each other for disjoint endpoint sets. Surviving
  weights are \(\kappa ^{2}/(r_{e} r_{f})\), with the repeated-edge case
  included.
\item
  The \(\rho _{F2}\) part of \(D_{2}\) has the density supports already
  listed and a prefactor bounded by \(\kappa\) squared after pairing
  with \(D_{0}\).
\end{itemize}

For estimates on finite excited states, a root path need not close under
Wick contraction. It still has a uniform inverse-square coefficient
bound. For example a two-L path has normalized weight

% DISPLAY d.25
% SOURCE sqrt(r_h)/(sqrt(r_e) r_f sqrt(r_g r_j)).
\begin{gather*}
\begin{aligned}
& \sqrt{r_{h}}/(\sqrt{r_{e}} r_{f} \sqrt{r_{g} r_{j}}).
\end{aligned}\notag
\end{gather*}

The triangle relation \(r_{h}\le r_{f}+r_{j}\) bounds this by a sum of

% DISPLAY d.26
% SOURCE 1/sqrt(r_e r_f r_g r_j),
% SOURCE 1/(sqrt(r_e) r_f sqrt(r_g)),
\begin{gather*}
\begin{aligned}
& 1/\sqrt{r_{e} r_{f} r_{g} r_{j}}, \\
& 1/(\sqrt{r_{e}} r_{f} \sqrt{r_{g}}),
\end{aligned}\notag
\end{gather*}

both bounded by a numerical multiple of the sum of inverse-square edge
weights by weighted arithmetic-geometric mean. This observation is
useful for mixed source terms and prevents an illicit global-frequency
comparison.

All of these are finite coefficient expressions in pair matrices with
external gaps comparable to their own \(r_{e}\). Normalize each
occurrence of an internal matrix by the incident \(r_{e}\) in that
factor. Its norm is \(O(\kappa )\). Functional calculus and its required
derivatives have uniform convergent resolvent bounds on that fixed pair
ball. Termwise product estimates in the above inventory therefore prove

% DISPLAY d.27
% SOURCE ||W(A)|| <= C_N W_2,
% SOURCE ||W(A)-W(0)|| <= C_N kappa W_2,
% SOURCE W_2=sum_e r_e^-2.                                      (4.2)
\begin{gather*}
\begin{aligned}
& \Vert W(A)\Vert  \le  C_{N} W_{2}, \\
& \Vert W(A)-W(0)\Vert  \le  C_{N} \kappa  W_{2}, \\
& W_{2}=\sum _{e} r_{e}^{-2}.
\end{aligned}\tag{4.2}
\end{gather*}

The derivative of \(j\) is included in the second estimate. It is not
estimated merely by the small value of \(j\), which would miss a nonzero
centered internal derivative.

\subsection{5. Actual shifted-fast correction and leading operator
defect}\label{d-actual-shifted-fast-correction-and-leading-operator-defect}

On the pair source space the adapted excitation energy is at least
\(c r_{e}\); on the triangle source space it is at least \(c s_{t}\).
The finite matrix resolvent identity, applied to the actual sources
(3.1), gives

% DISPLAY d.28
% SOURCE ||S_e*H_exc^-1 S_e-S_e(0)*H_f(0)^-1 S_e(0)||
% SOURCE       <= C_N kappa r_e^-2,
% SOURCE 
% SOURCE ||S_t*H_exc^-1 S_t-S_t(0)*H_f(0)^-1 S_t(0)||
% SOURCE       <= C_N kappa w_t.                                (5.1)
\begin{gather*}
\begin{aligned}
& \Vert S_{e}^{*}H_{\mathrm{exc}}^{-1} S_{e}-S_{e}(0)^{*}H_{f}(0)^{-1} S_{e}(0)\Vert \\
& \le  C_{N} \kappa  r_{e}^{-2},
\end{aligned}\notag \\
\begin{aligned}
& \Vert S_{t}^{*}H_{\mathrm{exc}}^{-1} S_{t}-S_{t}(0)^{*}H_{f}(0)^{-1} S_{t}(0)\Vert \\
& \le  C_{N} \kappa  w_{t}.
\end{aligned}\tag{5.1}
\end{gather*}

Cross terms between distinct graph spaces vanish. The centered source
lemma, including all internal vacuum derivatives and compression jets,
gives

% DISPLAY d.29
% SOURCE W(0)=S(0)*H_f(0)^-1 S(0).
\begin{gather*}
\begin{aligned}
& W(0)=S(0)^{*}H_{f}(0)^{-1} S(0).
\end{aligned}\notag
\end{gather*}

Consequently the actual leading operator defect satisfies

% DISPLAY d.30
% SOURCE ||W(A)-S(A)*H_exc(A)^-1 S(A)||
% SOURCE       <= C_N kappa W_2.                                (5.2)
\begin{gather*}
\begin{aligned}
& \Vert W(A)-S(A)^{*}H_{\mathrm{exc}}(A)^{-1} S(A)\Vert \\
& \le  C_{N} \kappa  W_{2}.
\end{aligned}\tag{5.2}
\end{gather*}

This is an operator-norm estimate on the complete internal Clifford
module. Its actual fast dressing is

% DISPLAY d.31
% SOURCE Z(A)=-H_exc(A)^-1 S(A),                                 (5.3)
\begin{gather*}
\begin{aligned}
& Z(A)=-H_{\mathrm{exc}}(A)^{-1} S(A),
\end{aligned}\tag{5.3}
\end{gather*}

and obeys the root-graph bounds and spectator-derivative correction from
Appendix \(C\). No interacting internal Hamiltonian occurs in (5.3). The
fast energy E0 remains a separate multiplier with its proved
commutator-square payment.

\subsection{6. All-vector use and the remaining reserve
requirement}\label{d-all-vector-use-and-the-remaining-reserve-requirement}

The graph-local lemma now applies to the actual dressing (5.3): its
coefficients are the explicit graph-local functions above. Smooth pair
functional calculus gives active derivative weights \(C r_{e}^{-5}\) or
\(C(\sum _{t} r_{e}^{-1})^{2}w_{t}/s_{t}\); spectator derivatives add
\(C m_{g} W_{2}\). Therefore

% DISPLAY d.32
% SOURCE H_s[Zf] <= C_N r_*^-3 H_s[f]
% SOURCE              +C_N kappa W_2[f]+C_N W_5[f].              (6.1)
\begin{gather*}
\begin{aligned}
& H_{s}[Zf] \le  C_{N} r_{*}^{-3} H_{s}[f] \\
& +C_{N} \kappa  W_{2}[f]+C_{N} W_{5}[f].
\end{aligned}\tag{6.1}
\end{gather*}

In fact triangle sources have no internal Clifford generator, so outside
derivative commutators are the only issue for their internal potential
charges. The more general graph bound is valid without using this
improvement. All internal momenta in (6.1) are unrestricted.

The identity \(H_{\mathrm{exc}} Z+S=0\) cancels the actual leading
complementary source exactly. The complete internal/center form pairs
involving \(Z\) can be controlled by Cauchy-Schwarz in \(H_{s}\), using
(6.1). The remaining derivative source from \(U\) is incident-edge
weighted and has the same payment as Section \(6\) of the geometric
lemma. These steps are valid form estimates on arbitrary complementary
vectors in the corresponding control norm; they are not a finite
oscillator truncation of those vectors.

What is not supplied here is an inequality extracting that control norm
from the original complementary form while retaining the same Gauss
square. In particular, the expression

% DISPLAY d.33
% SOURCE H_s[v]+H_exc[v]+||D_full v||²
\begin{gather*}
\begin{aligned}
& H_{s}[v]+H_{\mathrm{exc}}[v]+\Vert D_{\mathrm{full}} v\Vert ^{2}
\end{aligned}\notag
\end{gather*}

cannot be added to the right side of the original Hamiltonian merely
because its three terms are nonnegative. \(H_{\mathrm{exc}}\) already
uses the frozen piece of \(D_{\mathrm{full}}\). The desired actual
all-vector comparison still needs a non-double-counted reserve,
compatible with unrestricted slow momenta, and a complete
mixed/remainder estimate in that reserve. Finite coefficient
cancellation (5.2) cannot prove this coercivity.

Thus Sections 2--5 identify and estimate the actual leading H1/H2 source
coefficients. Section \(6\) supplies their internal-momentum payment and
states the precise remaining all-vector obstruction rather than silently
assuming it away.

\clearpage
\section{Rootwise all vector connection
estimates}\label{e-rootwise-all-vector-connection-estimates}

Technical derivation for review. Equation and subsection numbers in this
appendix are local to the source derivation. Historical local labels are
retained, including numbering gaps or nonsequential labels. The
accompanying source notes retain the original expressions for
comparison.

\subsection{1. Assumptions and the excitation
energy}\label{e-assumptions-and-the-excitation-energy}

Fix the rank and a partition. For every root/pair \(e=(a,b)\), let
\(r_{e}=\vert z_{a}-z_{b}\vert\) and assume
\(\alpha _{a}+\alpha _{b} \le  \kappa  r_{e}\), with \(\kappa\)
sufficiently small. Write \(r_{*}=\min _{e} r_{e}\) and
\(W_{2}=\sum _{e} r_{e}^{-2}\). No ratio \(r_{\max}/r_{*}\) is
restricted.

Freeze the slow variables. Let

% DISPLAY e.1
% SOURCE H_fast,e = H_B,e + H_F,e,
% SOURCE E_0,e = E_B,e + E_F,min,e,
% SOURCE h_e = H_fast,e - E_0,e
% SOURCE     = (H_B,e-E_B,e) + (H_F,e-E_F,min,e).
\begin{gather*}
\begin{aligned}
& H_{\mathrm{fast},e} = H_{B,e} + H_{F,e}, \\
& E_{0,e} = E_{B,e} + E_{F,\mathrm{min},e}, \\
& h_{e} = H_{\mathrm{fast},e} - E_{0,e} \\
& = (H_{B,e}-E_{B,e}) + (H_{F,e}-E_{F,\mathrm{min},e}).
\end{aligned}\notag
\end{gather*}

Here \(H_{B,e}\) is the exact kinetic-aware bosonic oscillator. The
uniform pair-matrix bounds give

% DISPLAY e.2
% SOURCE c(p_e^2+r_e^2 Y_e^2) <= H_B,e <= C(p_e^2+r_e^2 Y_e^2),
% SOURCE E_B,e <= C r_e,
% SOURCE h_e >= c r_e Q_e,
\begin{gather*}
\begin{aligned}
& c(p_{e}^{2}+r_{e}^{2} Y_{e}^{2}) \le  H_{B,e} \le  C(p_{e}^{2}+r_{e}^{2} Y_{e}^{2}), \\
& E_{B,e} \le  C r_{e}, \\
& h_{e} \ge  c r_{e} Q_{e},
\end{aligned}\notag
\end{gather*}

where \(P_{e}\) is its full boson/fermion vacuum projection and
\(Q_{e}=1-P_{e}\). In particular

% DISPLAY e.3
% SOURCE ||p_e Q_e v||^2 <= C h_e[v],
% SOURCE ||Y_e v||^2 <= C r_e^-2 h_e[v] + C r_e^-1 ||v||^2,
% SOURCE ||p_e P_e v||^2 <= C r_e ||P_e v||^2,
% SOURCE ||Y_e p_e P_e v||^2 <= C ||P_e v||^2.                 (1)
\begin{gather*}
\begin{aligned}
& \Vert p_{e} Q_{e} v\Vert ^{2} \le  C h_{e}[v], \\
& \Vert Y_{e} v\Vert ^{2} \le  C r_{e}^{-2} h_{e}[v] + C r_{e}^{-1} \Vert v\Vert ^{2}, \\
& \Vert p_{e} P_{e} v\Vert ^{2} \le  C r_{e} \Vert P_{e} v\Vert ^{2}, \\
& \Vert Y_{e} p_{e} P_{e} v\Vert ^{2} \le  C \Vert P_{e} v\Vert ^{2}.
\end{aligned}\tag{1}
\end{gather*}

The last two statements also hold with any uniformly bounded finite
color/spatial matrices between the displayed operators. Constants depend
only on rank and the fixed small pair-matrix neighborhood. \(h_{e}\),
rather than \(H_{B,e}\) or a centered oscillator, is essential.

Let \(\chi\) be any multiplication cutoff with
\(\vert \chi \vert \le 1\) supported in
\(\vert Y\vert \le \sigma  r_{*}\). Derivatives of \(\chi\) are not used
below. Every estimate is a fast-fiber estimate for arbitrary vectors,
with arbitrary spectator fast modes and arbitrary slow/Clifford
coefficient vectors. It can be integrated over the slow variables. No
internal momentum or oscillator-degree cutoff is imposed. No individual
internal singlet condition is imposed; the estimates remain valid on the
total residual-equivariant space.

\subsection{2. Same-root and mixed-root
lemmas}\label{e-same-root-and-mixed-root-lemmas}

For any pair \(e\) and any positive denominator scale \(r_{g}\) drawn
from the retained roots,

% DISPLAY e.4
% SOURCE ||chi Y_e p_e v/r_g||^2
% SOURCE   <= C sigma^2 (r_*^2/r_g^2) h_e[v]
% SOURCE        + C r_g^-2 ||v||^2.                           (2)
\begin{gather*}
\begin{aligned}
& \Vert \chi  Y_{e} p_{e} v/r_{g}\Vert ^{2} \\
& \le  C \sigma ^{2} (r_{*}^{2}/r_{g}^{2}) h_{e}[v] \\
& + C r_{g}^{-2} \Vert v\Vert ^{2}.
\end{aligned}\tag{2}
\end{gather*}

Proof: split the INPUT as \(P_{e} v+Q_{e} v\). On \(Q_{e}\) use
\(\vert \chi  Y_{e}\vert \le \sigma  r_{*}\) and the first inequality in
(1). On \(P_{e}\) use the last inequality in (1), dropping \(\chi\). The
factor two from the split is harmless. \(P_{e} v\) need not be supported
in the tube. No argument asserts that \(P_{e}\) preserves tube support.

For distinct roots \(e\) and \(f\),

% DISPLAY e.5
% SOURCE ||chi Y_f p_e v/r_g||^2
% SOURCE   <= C sigma^2 (r_*^2/r_g^2) h_e[v]
% SOURCE    + C r_e/(r_g^2 r_f^2) h_f[v]
% SOURCE    + C r_e/(r_g^2 r_f) ||v||^2.                      (3)
\begin{gather*}
\begin{aligned}
& \Vert \chi  Y_{f} p_{e} v/r_{g}\Vert ^{2} \\
& \le  C \sigma ^{2} (r_{*}^{2}/r_{g}^{2}) h_{e}[v] \\
& + C r_{e}/(r_{g}^{2} r_{f}^{2}) h_{f}[v] \\
& + C r_{e}/(r_{g}^{2} r_{f}) \Vert v\Vert ^{2}.
\end{aligned}\tag{3}
\end{gather*}

Indeed the \(Q_{e}\) input is as above. On the \(P_{e}\) input, drop
\(\chi\), commute \(Y_{f}\) with \(P_{e}\), and use
\(\Vert p_{e} P_{e}\Vert ^{2}\le C r_{e}\) followed by the \(Y_{f}\)
bound in (1). Equivalently split the remaining \(f\) input into
\(P_{f}\) and \(Q_{f}\). This is the step that cannot be omitted from a
mixed-root claim.

If e,f,g are the three edges of a triangle, \(r_{e}\le r_{f}+r_{g}\), so

% DISPLAY e.6
% SOURCE r_e/(r_g^2 r_f^2) <= 2 r_*^-3,
% SOURCE r_e/(r_g^2 r_f) <= r_g^-2+(r_g r_f)^-1 <= C W_2,t.
\begin{gather*}
\begin{aligned}
& r_{e}/(r_{g}^{2} r_{f}^{2}) \le  2 r_{*}^{-3}, \\
& r_{e}/(r_{g}^{2} r_{f}) \le  r_{g}^{-2}+(r_{g} r_{f})^{-1} \le  C W_{2,t}.
\end{aligned}\notag
\end{gather*}

Consequently

% DISPLAY e.7
% SOURCE ||chi Y_f p_e v/r_g||^2
% SOURCE   <= C (sigma^2+r_*^-3) (h_e[v]+h_f[v])
% SOURCE        + C W_2,t ||v||^2.                            (4)
\begin{gather*}
\begin{aligned}
& \Vert \chi  Y_{f} p_{e} v/r_{g}\Vert ^{2} \\
& \le  C (\sigma ^{2}+r_{*}^{-3}) (h_{e}[v]+h_{f}[v]) \\
& + C W_{2,t} \Vert v\Vert ^{2}.
\end{aligned}\tag{4}
\end{gather*}

For the input-root denominator \(g=e\), the same conclusion follows
directly, without any triangle hypothesis, with residual
\((r_{e} r_{f})^{-1}\). In (2)-(4), \(Y\) and \(p\) can be contracted
with any bounded root-preserving finite matrices. Finite sums of such
sources obey the same bounds after enlarging the rank-dependent
constant.

A triangle hypothesis (or another restriction on \(r_{e}\) relative to
\(r_{f},r_{g}\)) is genuinely necessary to replace (3) by (4): the
\(P_{e} P_{f}\) input has squared size comparable to
\(r_{e}/(r_{g}^{2} r_{f})\). An arbitrary unrelated \(r_{e}\) cannot
disappear from this expression.

\subsection{3. Exact inverse-M factorization for the fast
column}\label{e-exact-inverse-m-factorization-for-the-fast-column}

Use the notation of the intrinsic geometry and adapted source inventory:

% DISPLAY e.8
% SOURCE M=F+P,   P=L-J_2,   J_2=K C_Y=R^-1 C_Y* C_Y,
% SOURCE W=I+K K*,   K=R^-1 C_Y*,
% SOURCE T=(I+F^-1 P*)^-1,
% SOURCE M^-*=T F^-1,
% SOURCE d_0(v)=F^-1 C* p_F v.
\begin{gather*}
\begin{aligned}
& M=F+P, \quad  P=L-J_{2}, \quad  J_{2}=K C_{Y}=R^{-1} C_{Y}^{*} C_{Y}, \\
& W=I+K K^{*}, \quad  K=R^{-1} C_{Y}^{*}, \\
& T=(I+F^{-1} P^{*})^{-1}, \\
& M^{-*}=T F^{-1}, \\
& d_{0}(v)=F^{-1} C^{*} p_{F} v.
\end{aligned}\notag
\end{gather*}

\(F\) is self-adjoint pair-block diagonal;
\(\Vert F_{e}^{-1} C_{e}^{*}\Vert \le C \kappa\). On the tube,

% DISPLAY e.9
% SOURCE ||P F^-1||+||F^-1 P*|| <= C sigma,
% SOURCE ||T||+||W^(1/2)|| <= C,
% SOURCE ||K|| <= C sigma.
\begin{gather*}
\begin{aligned}
& \Vert P F^{-1}\Vert +\Vert F^{-1} P^{*}\Vert  \le  C \sigma , \\
& \Vert T\Vert +\Vert W^{1/2}\Vert  \le  C, \\
& \Vert K\Vert  \le  C \sigma .
\end{aligned}\notag
\end{gather*}

Let \(B p_{F}=W^{1/2} M^{-*}(C^{*}+N_{Y}^{*})p_{F}\) be the exact fast
column of the Gauss stack. Algebraically, with the indicated operator
ordering,

% DISPLAY e.10
% SOURCE B p_F-d_0
% SOURCE   =(W^(1/2)-I)d_0
% SOURCE    +W^(1/2)T[
% SOURCE       F^-1 N_Y* p_F
% SOURCE       -F^-1 L* d_0
% SOURCE       +F^-1 J_2* d_0].                               (5)
\begin{gather*}
\begin{aligned}
& B p_{F}-d_{0} \\
& =(W^{1/2}-I)d_{0} \\
& +W^{1/2}T[ \\
& F^{-1} N_{Y}^{*} p_{F} \\
& -F^{-1} L^{*} d_{0} \\
& +F^{-1} J_{2}^{*} d_{0}].
\end{aligned}\tag{5}
\end{gather*}

There are no derivatives of inverse matrices in this first-order form
identity.

The first two bracketed expressions are finite triangle sources
\(Y_{f} p_{e}/r_{g}\), with respectively bounded and \(O(\kappa )\)
root-preserving factors. The last expression factors as

% DISPLAY e.11
% SOURCE F^-1 J_2* d_0
% SOURCE   =F^-1 C_Y* [C_Y R^-1 d_0],                         (6)
\begin{gather*}
\begin{aligned}
& F^{-1} J_{2}^{*} d_{0} \\
& =F^{-1} C_{Y}^{*} [C_{Y} R^{-1} d_{0}],
\end{aligned}\tag{6}
\end{gather*}

where \(\Vert F^{-1} C_{Y}^{*}\Vert \le C \sigma\) and
\(C_{Y} R^{-1} d_{0}\) is a sum of same-root terms \(Y_{e} p_{e}/r_{e}\)
with \(O(\kappa )\) coefficients. Also

% DISPLAY e.12
% SOURCE (W^(1/2)-I)d_0
% SOURCE   =(W^(1/2)+I)^-1 K [K* d_0],
% SOURCE K* d_0=C_Y R^-1 d_0.                                 (7)
\begin{gather*}
\begin{aligned}
& (W^{1/2}-I)d_{0} \\
& =(W^{1/2}+I)^{-1} K [K^{*} d_{0}], \\
& K^{*} d_{0}=C_{Y} R^{-1} d_{0}.
\end{aligned}\tag{7}
\end{gather*}

Thus all potentially troublesome inverse-M matrices remain bounded left
factors. Equations (2)-(7) prove

% DISPLAY e.13
% SOURCE ||chi (B p_F-d_0)v||^2
% SOURCE   <= C (sigma^2+r_*^-3) sum_e h_e[v]
% SOURCE        + C W_2 ||v||^2.                              (8)
\begin{gather*}
\begin{aligned}
& \Vert \chi  (B p_{F}-d_{0})v\Vert ^{2} \\
& \le  C (\sigma ^{2}+r_{*}^{-3}) \sum _{e} h_{e}[v] \\
& + C W_{2} \Vert v\Vert ^{2}.
\end{aligned}\tag{8}
\end{gather*}

The estimate holds for arbitrary fast vectors. It is not an estimate
restricted to vacuum or a finite excited coefficient space.

The spin column obeys separately

% DISPLAY e.14
% SOURCE ||chi W^(1/2)M^-* S_f v||^2 <= C W_2 ||v||^2,          (9)
\begin{gather*}
\begin{aligned}
& \Vert \chi  W^{1/2}M^{-*} S_{f} v\Vert ^{2} \le  C W_{2} \Vert v\Vert ^{2},
\end{aligned}\tag{9}
\end{gather*}

because the finite fermionic Gauss matrices are bounded and \(F^{-1}\)
retains the individual output-root denominators.

\subsection{4. Exact slow-square coefficient and its Schur
loss}\label{e-exact-slow-square-coefficient-and-its-schur-loss}

Let \(G_{s}=\operatorname{diag} (g_{C}^{-1},I_{I})\), so the direct slow
kinetic stack is \(w=G_{s}^{1/2}p_{s}\). The exact normalized slow
column in the Gauss square is

% DISPLAY e.15
% SOURCE A=W^(1/2)M^-* Psi,
% SOURCE Psi=[(C_Y*+(F+L)*K)g_C^-1/2, I_Y*].                  (10)
\begin{gather*}
\begin{aligned}
& A=W^{1/2}M^{-*} \Psi , \\
& \Psi =[(C_{Y}^{*}+(F+L)^{*}K)g_{C}^{-1/2}, I_{Y}^{*}].
\end{aligned}\tag{10}
\end{gather*}

This follows directly from \(O^{*}Sg_{S}^{-1}\). In particular the
\(F K\) term must be retained. One has

% DISPLAY e.16
% SOURCE ||Psi*F^-1|| <= C sigma,
% SOURCE ||A|| <= C sigma.                                    (11)
\begin{gather*}
\begin{aligned}
& \Vert \Psi ^{*}F^{-1}\Vert  \le  C \sigma , \\
& \Vert A\Vert  \le  C \sigma .
\end{aligned}\tag{11}
\end{gather*}

The adjoint coefficient acting on the frozen fast column has the exact
expression

% DISPLAY e.17
% SOURCE A* d_0=Psi*F^-1 T*W^(1/2)d_0.                        (12)
\begin{gather*}
\begin{aligned}
& A^{*} d_{0}=\Psi ^{*}F^{-1} T^{*}W^{1/2}d_{0}.
\end{aligned}\tag{12}
\end{gather*}

Expand

% DISPLAY e.18
% SOURCE T*W^(1/2)-I=(T*-I)+T*(W^(1/2)-I),
% SOURCE T*-I=-T* P F^-1.                                    (13)
\begin{gather*}
\begin{aligned}
& T^{*}W^{1/2}-I=(T^{*}-I)+T^{*}(W^{1/2}-I), \\
& T^{*}-I=-T^{*} P F^{-1}.
\end{aligned}\tag{13}
\end{gather*}

The leading term \(\Psi ^{*}F^{-1} d_{0}\) is the stack of

% DISPLAY e.19
% SOURCE I_Y F^-1 d_0,
% SOURCE g_C^-1/2[C_Y F^-1 d_0+K* d_0+K* L F^-1 d_0].         (14)
\begin{gather*}
\begin{aligned}
& I_{Y} F^{-1} d_{0}, \\
& g_{C}^{-1/2}[C_{Y} F^{-1} d_{0}+K^{*} d_{0}+K^{*} L F^{-1} d_{0}].
\end{aligned}\tag{14}
\end{gather*}

The first three contractions are same-root \(Y_{e} p_{e}/r_{e}\).
\(L F^{-1} d_{0}\) is a triangle source with INPUT-root denominator
\(r_{e}\), so (3) applies even without using a frequency ratio. The
\(J_{2}\) part of \(P F^{-1} d_{0}\) factors as

% DISPLAY e.20
% SOURCE J_2 F^-1 d_0=R^-1 C_Y* [C_Y F^-1 d_0],               (15)
\begin{gather*}
\begin{aligned}
& J_{2} F^{-1} d_{0}=R^{-1} C_{Y}^{*} [C_{Y} F^{-1} d_{0}],
\end{aligned}\tag{15}
\end{gather*}

again a bounded \(O(\sigma )\) left factor times a same-root
contraction. Equations (7), (11), and (13)-(15) therefore prove the
sharper estimate

% DISPLAY e.21
% SOURCE ||chi A* d_0(v)||^2
% SOURCE   <= C kappa^2 (sigma^2+r_*^-3) sum_e h_e[v]
% SOURCE        + C kappa^2 W_2 ||v||^2.                      (16)
\begin{gather*}
\begin{aligned}
& \Vert \chi  A^{*} d_{0}(v)\Vert ^{2} \\
& \le  C \kappa ^{2} (\sigma ^{2}+r_{*}^{-3}) \sum _{e} h_{e}[v] \\
& + C \kappa ^{2} W_{2} \Vert v\Vert ^{2}.
\end{aligned}\tag{16}
\end{gather*}

Combining (8), (9), and \(\Vert A\Vert \le C \sigma\) also controls
\(A^{*}\) applied to the full fast-plus-spin stack, with a coefficient
\(C(\kappa ^{2}+\sigma ^{2})\) multiplying the right side of (8), and a
harmless \(C \sigma ^{2} W_{2}\) spin term.

For completeness, for any vectors w,b and \(\eta >0\),

% DISPLAY e.22
% SOURCE eta||w||^2+||b-Aw||^2
% SOURCE   =||(eta I+A*A)^(1/2)
% SOURCE       [w-(eta I+A*A)^-1 A*b]||^2
% SOURCE       +<b,(I+eta^-1 A A*)^-1 b>.
\begin{gather*}
\begin{aligned}
& \eta \Vert w\Vert ^{2}+\Vert b-Aw\Vert ^{2} \\
& =\Vert (\eta  I+A^{*}A)^{1/2} \\
& [w-(\eta  I+A^{*}A)^{-1} A^{*}b]\Vert ^{2} \\
& +\langle b,(I+\eta ^{-1} A A^{*})^{-1} b\rangle .
\end{aligned}\notag
\end{gather*}

The difference between \(\Vert b\Vert ^{2}\) and the final term is at
most \(\eta ^{-1}\Vert A^{*}b\Vert ^{2}\). Thus (16) controls the
frozen-column Schur loss by exact excitation energies and \(W_{2}\),
with no \(\delta  \sum _{e} r_{e}\) zero-point cost. This is a
first-order form identity: variable coefficients create no omitted
ordering terms at this stage.

\subsection{5. What this lemma does and does not
close}\label{e-what-this-lemma-does-and-does-not-close}

It closes the all-vector rootwise CONNECTION-NORM sublemma, including
mixed roots, exact inverse-M left factors, and the coefficient occurring
after slow-square completion. It remains valid with unrestricted
internal momenta because none of its steps estimates those momenta or
differentiates a fast projection in a slow direction.

At the connection-norm stage, two additional arguments are necessary.
Section \(6\) below supplies the first for the fast kinetic column:

\begin{enumerate}
\def\labelenumi{\arabic{enumi}.}
\tightlist
\item
  Bounding a squared connection norm is not a bound for the bilinear
  pairing with \(d_{0}\). The direct Young estimate against
  \(\Vert d_{0} v\Vert ^{2}\) introduces its frozen vacuum energy
  \(C \kappa ^{2} \sum _{e} r_{e}\). The odd H1 source / normal-ordering
  or graph-local form argument must be used instead; (8) alone does not
  perform it.
\item
  Completing all slow squares leaves shifted slow derivatives. Splitting
  off \(\eta\) times ordinary internal kinetic energy is not the same as
  splitting \(\eta\) times the colored internal supersymmetric form: an
  unweighted Yukawa remainder \(\eta  C \alpha _{a}\) may be large for a
  far block. One must retain a suitable full slow-supercharge form or
  prove an incident-weighted conversion of the shifted slow squares.
  Equations (10)-(16) do not silently supply that conversion.
\end{enumerate}

Accordingly these results are compatible with, but do not establish, a
near-sharp complementary lower bound in the original Hamiltonian. They
provide the rootwise estimate that avoids a global-frequency
substitution at the connection stage.

\subsection{6. Addendum: signed fast kinetic terms without a
vacuum-frequency
loss}\label{e-addendum-signed-fast-kinetic-terms-without-a-vacuum-frequency-loss}

The first obstruction listed in Section \(5\) can also be removed for
the exact fast kinetic column. The following statements concern a smooth
core vector \(v\) supported in the fast tube; an outer cutoff equal to
one near its support may always be used in Sections 2-4. Only that
support requirement, not a restriction on the oscillator content of
\(v\), is imposed.

\subsubsection{6.1 All-vector triangle
multiplication}\label{e-all-vector-triangle-multiplication}

For distinct triangle roots e,f,g, write \(s_{t}=r_{e}+r_{f}+r_{g}\).
Any uniformly bounded real coefficient tensor that may depend on the
slow variables but is independent of the fast variables \(Y\) obeys

% DISPLAY e.23
% SOURCE |<v,s_t Y_e Y_f Y_g v>|
% SOURCE    <= C sigma sum_(j in t) h_j[v]
% SOURCE       + C sigma^-1 W_2,t ||v||^2.                   (17)
\begin{gather*}
\begin{aligned}
& \vert \langle v,s_{t} Y_{e} Y_{f} Y_{g} v\rangle \vert \\
& \le  C \sigma  \sum _{j \in  t} h_{j}[v] \\
& + C \sigma ^{-1} W_{2,t} \Vert v\Vert ^{2}.
\end{aligned}\tag{17}
\end{gather*}

First, splitting inputs into \(Q_{f}, P_{f} Q_{g}\), and \(P_{f} P_{g}\)
gives

% DISPLAY e.24
% SOURCE ||chi Y_f Y_g v||^2
% SOURCE    <= C sigma^2 r_*^2 [r_f^-2 h_f[v]+r_g^-2 h_g[v]]
% SOURCE       + C/(r_f r_g) ||v||^2.                        (18)
\begin{gather*}
\begin{aligned}
& \Vert \chi  Y_{f} Y_{g} v\Vert ^{2} \\
& \le  C \sigma ^{2} r_{*}^{2} [r_{f}^{-2} h_{f}[v]+r_{g}^{-2} h_{g}[v]] \\
& + C/(r_{f} r_{g}) \Vert v\Vert ^{2}.
\end{aligned}\tag{18}
\end{gather*}

On the first component bound \(\chi  Y_{g}\) by \(\sigma  r_{*}\) and
use the excited \(Y_{f}\) oscillator bound; on the second use
\(\chi  Y_{f}\) and the excited \(Y_{g}\) bound; on the last use both
Gaussian widths. The splitting is for estimation only; the projected
pieces need not have tube support.

Choose \(e\) to be a largest triangle root. If
\(U_{B,e}=\text{ constant } \exp (-Y_{e} \cdot  \Omega _{e} Y_{e}/2)\),
put \(a_{e}=\partial _{e}+\Omega _{e} Y_{e}\). The exact bosonic
excitation form controls \(\Vert a_{e} v\Vert ^{2}\). Since the other
two factors do not depend on \(Y_{e}\), integration by parts gives,
componentwise,

% DISPLAY e.25
% SOURCE <v,Y_e Y_f Y_g v>
% SOURCE    =Re<a_e v,Omega_e^-1(Y_f Y_g)v>.
\begin{gather*}
\begin{aligned}
& \langle v,Y_{e} Y_{f} Y_{g} v\rangle \\
& =\operatorname{Re} \langle a_{e} v,\Omega _{e}^{-1}(Y_{f} Y_{g})v\rangle .
\end{aligned}\notag
\end{gather*}

Using \(\Vert \Omega _{e}^{-1}\Vert \le C/r_{e}, s_{t}/r_{e}\le 3\), and
(18) gives (17) by Young with parameter \(\sigma\). The proof is valid
for arbitrary fermionic and slow vector coefficients.

\subsubsection{6.2 Leading triangle differential
forms}\label{e-leading-triangle-differential-forms}

Let a(Y) be a real symmetric fast coefficient matrix supported only in
the off-diagonal root blocks \((e,g)\) and \((g,e)\), with each entry a
bounded linear contraction of \(Y_{f}\) divided by \(r_{e}\), multiplied
by a coefficient \(O(\kappa )\) that is independent of the fast
variables \(Y\) (it may depend on the slow variables). For the product
bosonic Gaussian \(U_{B}\), writing \(v=U_{B} u\) gives the exact
identity

% DISPLAY e.26
% SOURCE integral (partial v)* a (partial v)
% SOURCE   =integral U_B^2 (partial u)* a (partial u)
% SOURCE    +integral [div(a Omega Y)-(Omega Y)·a(Omega Y)] |v|^2.
% SOURCE                                                                 (19)
\begin{gather*}
\begin{aligned}
& \int  (\partial  v)^{*} a (\partial  v) \\
& =\int  U_{B}^{2} (\partial  u)^{*} a (\partial  u) \\
& +\int  [\operatorname{div} (a \Omega  Y)-(\Omega  Y) \cdot  a(\Omega  Y)] \vert v\vert ^{2}.
\end{aligned}\tag{19}
\end{gather*}

This identity does not project onto the fermionic vacuum. The excitation
form satisfies

% DISPLAY e.27
% SOURCE sum_e (H_B,e-E_B,e)[v]
% SOURCE    =integral U_B^2 (partial u)*G(partial u),
\begin{gather*}
\begin{aligned}
& \sum _{e} (H_{B,e}-E_{B,e})[v] \\
& =\int  U_{B}^{2} (\partial  u)^{*}G(\partial  u),
\end{aligned}\notag
\end{gather*}

with \(G\) uniformly positive. On tube-supported \(v\), the first term
in (19) is bounded in absolute value by
\(C \kappa  \sigma  \sum _{e} h_{e}[v]\). The divergence term in (19)
vanishes: a connects distinct roots, depends on the third root, and
\(\Omega\) is pair-block diagonal. The remaining scalar is a triangle
cubic with coefficient \(O(\kappa  r_{g})\), hence is bounded by (17).
Consequently

% DISPLAY e.28
% SOURCE |integral (partial v)* a (partial v)|
% SOURCE   <= C kappa sigma sum_(j in t) h_j[v]
% SOURCE       + C kappa sigma^-1 W_2,t ||v||^2.              (20)
\begin{gather*}
\begin{aligned}
& \vert \int  (\partial  v)^{*} a (\partial  v)\vert \\
& \le  C \kappa  \sigma  \sum _{j \in  t} h_{j}[v] \\
& + C \kappa  \sigma ^{-1} W_{2,t} \Vert v\Vert ^{2}.
\end{aligned}\tag{20}
\end{gather*}

The same proof applies to \(O(\kappa ^{2})\) coefficients. Equation (20)
applies to both leading triangle forms from

% DISPLAY e.29
% SOURCE 2 Re<d_0,F^-1 N_Y* p_F-F^-1 L* d_0>.
\begin{gather*}
\begin{aligned}
& 2 \operatorname{Re} \langle d_{0},F^{-1} N_{Y}^{*} p_{F}-F^{-1} L^{*} d_{0}\rangle .
\end{aligned}\notag
\end{gather*}

No derivative cutoff is needed: (19) is applied to the original
polynomial coefficient and \(v\) itself has tube support. If one instead
wants a globally cut-off form on arbitrary unsupported inputs, cutoff
derivatives must be included; an even radial cutoff generates additional
triangle cubics, with coefficient smaller by
\(O((\sigma ^{2} r_{*}^{3})^{-1})\).

\subsubsection{6.3 Full inverse-M
reduction}\label{e-full-inverse-m-reduction}

There is a short exact reduction avoiding an infinite expansion. Set

% DISPLAY e.30
% SOURCE Z=F^-1 P*,   T=(I+Z)^-1,
% SOURCE n=F^-1 N_Y* p_F,
% SOURCE X=n-Z d_0.
\begin{gather*}
\begin{aligned}
& Z=F^{-1} P^{*}, \quad  T=(I+Z)^{-1}, \\
& n=F^{-1} N_{Y}^{*} p_{F}, \\
& X=n-Z d_{0}.
\end{aligned}\notag
\end{gather*}

Then

% DISPLAY e.31
% SOURCE B p_F=W^(1/2)(d_0+T X).
\begin{gather*}
\begin{aligned}
& B p_{F}=W^{1/2}(d_{0}+T X).
\end{aligned}\notag
\end{gather*}

Because \(W\ge I\) and \(T-I=-ZT\),

% DISPLAY e.32
% SOURCE ||B p_F v||^2
% SOURCE   >=||d_0 v||^2+2 Re<d_0 v,Xv>
% SOURCE       -2 |<Z* d_0 v,T Xv>|.                         (21)
\begin{gather*}
\begin{aligned}
& \Vert B p_{F} v\Vert ^{2} \\
& \ge \Vert d_{0} v\Vert ^{2}+2 \operatorname{Re} \langle d_{0} v,Xv\rangle \\
& -2 \vert \langle Z^{*} d_{0} v,T Xv\rangle \vert .
\end{aligned}\tag{21}
\end{gather*}

Both \(X\) and \(Z^{*} d_{0}=P F^{-1} d_{0}\) are covered by the
rootwise connection estimates. The last pairing is therefore bounded by

% DISPLAY e.33
% SOURCE C(sigma^2+r_*^-3) sum_e h_e[v]+C W_2||v||^2.
\begin{gather*}
\begin{aligned}
& C(\sigma ^{2}+r_{*}^{-3}) \sum _{e} h_{e}[v]+C W_{2}\Vert v\Vert ^{2}.
\end{aligned}\notag
\end{gather*}

Inside \(\langle d_{0},X\rangle\), the \(J_{2}\) term factors exactly as

% DISPLAY e.34
% SOURCE <d_0,F^-1 J_2* d_0>
% SOURCE    =<C_Y F^-1 d_0,C_Y R^-1 d_0>,                    (22)
\begin{gather*}
\begin{aligned}
& \langle d_{0},F^{-1} J_{2}^{*} d_{0}\rangle \\
& =\langle C_{Y} F^{-1} d_{0},C_{Y} R^{-1} d_{0}\rangle ,
\end{aligned}\tag{22}
\end{gather*}

and has the same bound. Only the two linear triangle forms in Section
\(6.2\) remain. Combining (20)-(22) gives

% DISPLAY e.35
% SOURCE ||B p_F v||^2
% SOURCE   >=||d_0 v||^2
% SOURCE      -C[kappa sigma+sigma^2+r_*^-3] sum_e h_e[v]
% SOURCE      -C(1+kappa/sigma) W_2||v||^2.                  (23)
\begin{gather*}
\begin{aligned}
& \Vert B p_{F} v\Vert ^{2} \\
& \ge \Vert d_{0} v\Vert ^{2} \\
& -C[\kappa  \sigma +\sigma ^{2}+r_{*}^{-3}] \sum _{e} h_{e}[v] \\
& -C(1+\kappa /\sigma ) W_{2}\Vert v\Vert ^{2}.
\end{aligned}\tag{23}
\end{gather*}

Thus direct \(p_{F}^{2}\) plus the exact Gauss FAST column, the adapted
quadratic potential, and its fast fermion mass have a lower bound

% DISPLAY e.36
% SOURCE E_0,total ||v||^2
% SOURCE    +(1-C[kappa sigma+sigma^2+r_*^-3]) H_exc[v]
% SOURCE    -C(1+kappa/sigma) W_2||v||^2,                    (24)
\begin{gather*}
\begin{aligned}
& E_{0,\mathrm{total}} \Vert v\Vert ^{2} \\
& +(1-C[\kappa  \sigma +\sigma ^{2}+r_{*}^{-3}]) H_{\mathrm{exc}}[v] \\
& -C(1+\kappa /\sigma ) W_{2}\Vert v\Vert ^{2},
\end{aligned}\tag{24}
\end{gather*}

with the frozen \(E_{0}\) multiplier understood pointwise in the slow
variables. This is a fast-column comparison only; the slow derivatives
in the full Gauss square have not been dropped or independently counted.

\subsubsection{6.4 Frozen spin-linear
term}\label{e-frozen-spin-linear-term}

The leading spin cross is a finite \(\sum  A_{e} p_{e}/r_{e}\), with
Y-independent Hermitian Clifford matrices
\(\Vert A_{e}\Vert \le C \kappa\). Since the \(\Omega _{e} Y_{e}\) part
of \(a_{e}=\partial _{e}+\Omega _{e} Y_{e}\) contributes a real
expectation against \(A_{e}\),

% DISPLAY e.37
% SOURCE |<v,A_e p_e v/r_e>|
% SOURCE   <= C kappa r_e^-1 sqrt(h_B,e[v]) ||v||
% SOURCE   <= epsilon h_e[v]+C kappa^2/(epsilon r_e^2)||v||^2. (25)
\begin{gather*}
\begin{aligned}
& \vert \langle v,A_{e} p_{e} v/r_{e}\rangle \vert \\
& \le  C \kappa  r_{e}^{-1} \sqrt{h_{B,e}[v]} \Vert v\Vert \\
& \le  \epsilon  h_{e}[v]+C \kappa ^{2}/(\epsilon  r_{e}^{2})\Vert v\Vert ^{2}.
\end{aligned}\tag{25}
\end{gather*}

Here \(h_{B,e}=H_{B,e}-E_{B,e}\). This also avoids any bosonic
vacuum-frequency cost. The inverse-M corrections to this frozen spin
cross pair rootwise connection errors with the bounded spin column (9),
and are bounded by \(\epsilon  H_{\mathrm{exc}}+C_{\epsilon} W_{2}\).

\subsection{7. Final scope after the
addendum}\label{e-final-scope-after-the-addendum}

Equations (2)-(16) prove the connection and Schur-loss norm statements
for arbitrary fast vectors. Equations (17)-(25) additionally prove the
signed fast-column comparison on an unrestricted oscillator core
supported in the fast tube. These estimates remove the need for a
\(\delta  \sum _{e} r_{e}\) loss at the fast-connection stage.

They still do not convert the completed, shifted slow derivative form
into a nonnegative colored internal supersymmetric form. One must
control its connection curvature/full-supercharge terms or use another
incident-weighted argument. In particular a constant loss of internal
kinetic energy cannot simply be identified with the same loss of
\(H_{a}\). That slow-form issue, the other physical potential/Yukawa
terms, and the final allocation of one non-double-counted reserve belong
to the full complementary theorem, not to this helper lemma.

\clearpage
\section{Pairwise Berry connection
comparison}\label{f-pairwise-berry-connection-comparison}

Technical derivation for review. Equation and subsection numbers in this
appendix are local to the source derivation. Historical local labels are
retained, including numbering gaps or nonsequential labels. The
accompanying source notes retain the original expressions for
comparison.

\subsection{4.1 Complete charge
comparison}\label{f-complete-charge-comparison}

A naive kinetic comparison

% DISPLAY f.1
% SOURCE |(partial+iA)f|² >= (1−epsilon)|partial f|²−epsilon^(-1)|A|²|f|²
\begin{gather*}
\begin{aligned}
& \vert (\partial +iA)f\vert ^{2} \ge  (1-\epsilon )\vert \partial  f\vert ^{2}-\epsilon ^{-1}\vert A\vert ^{2}\vert f\vert ^{2}
\end{aligned}\notag
\end{gather*}

loses a fixed part of internal kinetic energy. Replacing it by \(h_{I}\)
incurs the unwanted \(O(\epsilon  \kappa  r)\) internal Yukawa mass.
This is the wrong comparison.

Instead, normalize the sixteen internal charges so \(h_{I}\) is their
averaged squared norm on the core. Replace each momentum by its scalar
Berry-covariant momentum. Direct Clifford expansion gives

% DISPLAY f.2
% SOURCE h_I(A)= average_alpha ||Q_(I,alpha)+C_alpha(A)||²
% SOURCE          − a finite Clifford contraction of F_A.
\begin{gather*}
\begin{aligned}
& h_{I}(A)= \operatorname{average} _{\alpha} \Vert Q_{I,\alpha}+C_{\alpha}(A)\Vert ^{2} \\
& - \text{ a finite Clifford contraction of } F_{A}.
\end{aligned}\notag
\end{gather*}

The difference \(C_{\alpha}(A)\) is a bounded multiplication matrix,
with norm \(\le C_{N}\vert A\vert\). Consequently

% DISPLAY f.3
% SOURCE h_I(A) >= (1−epsilon) h_I
% SOURCE            −C_N[epsilon^(-1)|A|²+|F_A|].
\begin{gather*}
\begin{aligned}
& h_{I}(A) \ge  (1-\epsilon ) h_{I} \\
& -C_{N}[\epsilon ^{-1}\vert A\vert ^{2}+\vert F_{A}\vert ].
\end{aligned}\notag
\end{gather*}

This retains the complete interacting \(h_{I}\), rather than its kinetic
part, and uses no gap. Spinor averaging eliminates the ordinary
gauge-term obstruction to positivity even on the unrestricted internal
module; the low coefficient additionally has the physical residual
equivariance required by the inductive Hardy input.

On a two-block small-ratio tube the gapped finite matrix is analytic.
All determinant curvatures vanish at internal \(A=0\): the two-projector
derivative calculation contains a Spin(9) trace of at most three
Clifford matrices. Homogeneity and smoothness therefore give

% DISPLAY f.4
% SOURCE |F_A| <= C_N kappa/r².
\begin{gather*}
\begin{aligned}
& \vert F_{A}\vert  \le  C_{N} \kappa /r^{2}.
\end{aligned}\notag
\end{gather*}

Use equivariant radial gauge in the internal coordinates, based on the
flat center vacuum at \(A=0\). The radial-gauge curvature formula gives
\(\vert A_{A}\vert \le C_{N} \kappa ^{2}/r\) for internal components.
Mixed relative components must also be kept: in this gauge they are the
flat \(A=0\) base connection plus the integral of \(F_{A,b}\) contracted
with A. The same vanishing-at-\(A=0\) argument and homogeneity give
their bound \(C_{N} \kappa ^{2}/r\). Thus, after this separate
mixed-component check,

% DISPLAY f.5
% SOURCE |A_Berry| <= C_N kappa²/r.
\begin{gather*}
\begin{aligned}
& \vert A_{\mathrm{Berry}}\vert  \le  C_{N} \kappa ^{2}/r.
\end{aligned}\notag
\end{gather*}

It follows that the preceding comparison costs only

% DISPLAY f.6
% SOURCE C_N[kappa+epsilon^(-1) kappa^4] r^(-2),
\begin{gather*}
\begin{aligned}
& C_{N}[\kappa +\epsilon ^{-1} \kappa ^{4}] r^{-2},
\end{aligned}\notag
\end{gather*}

which can be made arbitrarily small after fixing \(\epsilon\). A
gauge-equivariant radial trivialization preserves ordinary residual
SU(n) physicality of the coefficient. Weyl/center flat identifications
remain part of the bundle.

This establishes a viable repair of the flatness step provided the local
Hamiltonian comparison really identifies the covariant internal form and
retains its curvature terms. The following rootwise construction avoids
an artificial dependence on the global minimum-gap ratio.

\subsection{4.2 Pairwise relative multiblock Berry
bound}\label{f-pairwise-relative-multiblock-berry-bound}

At quadratic order the fast fermionic space is a direct sum over
bifundamental block pairs. Its one-particle matrices are

% DISPLAY f.7
% SOURCE D_ab=Gamma dot b_ab tensor I
% SOURCE        +sum_i Gamma_i tensor (A_i^(a) tensor I−I tensor (A_i^(b))^T).
\begin{gather*}
\begin{aligned}
& D_{ab}=\Gamma  \cdot  b_{ab} \otimes  I \\
& +\sum _{i} \Gamma _{i} \otimes  (A_{i}^{(a)} \otimes  I-I \otimes  (A_{i}^{(b)})^{T}).
\end{aligned}\notag
\end{gather*}

If each endpoint internal norm is at most \(\kappa\) times
\(r_{ab}=\vert b_{ab}\vert\), the gap is a positive fixed multiple of
\(r_{ab}\) after reducing \(\kappa\) by the fixed nine-dimensional norm
factor. The adapted determinant line is the tensor product of the
occupied pair determinants. On each pair, choose the preceding
internal-radial gauge from the flat \(A=0\) center line. The Clifford
trace calculation at \(A=0\) is independent of bifundamental dimension,
so every pure/mixed curvature component vanishes there.

Smoothness on the compact dimensionless pair tube and homogeneity give,
in canonical block/center coordinates,

% DISPLAY f.8
% SOURCE |F_ab|<=C_N kappa r_ab^(-2),
% SOURCE |a_ab|<=C_N kappa² r_ab^(-1).
\begin{gather*}
\begin{aligned}
& \vert F_{ab}\vert \le C_{N} \kappa  r_{ab}^{-2}, \\
& \vert a_{ab}\vert \le C_{N} \kappa ^{2} r_{ab}^{-1}.
\end{aligned}\notag
\end{gather*}

The relative-component bound follows from its own mixed-curvature
integral, not from simply declaring it gauged away. Product connections
add, and Cauchy--Schwarz over the finitely many pairs gives

% DISPLAY f.9
% SOURCE |F_total|<=C_N kappa sum_(a<b) r_ab^(-2),
% SOURCE |a_total|²<=C_N kappa^4 sum_(a<b) r_ab^(-2).
\begin{gather*}
\begin{aligned}
& \vert F_{\mathrm{total}}\vert \le C_{N} \kappa  \sum _{a<b} r_{ab}^{-2}, \\
& \vert a_{\mathrm{total}}\vert ^{2}\le C_{N} \kappa ^{4} \sum _{a<b} r_{ab}^{-2}.
\end{aligned}\notag
\end{gather*}

The full slow averaged charge stack includes all internal interacting
charges and free center charges. Its comparison therefore loses only
\(\epsilon\) times the entire slow supersymmetric form and the
inverse-square multiplier

% DISPLAY f.10
% SOURCE C_N(kappa+epsilon^(-1)kappa^4) sum_(a<b) r_ab^(-2).
\begin{gather*}
\begin{aligned}
& C_{N}(\kappa +\epsilon ^{-1}\kappa ^{4}) \sum _{a<b} r_{ab}^{-2}.
\end{aligned}\notag
\end{gather*}

This multiplier has a direct shape-independent center-kinetic payment.
In the mass metric \(\sum  n_{a}\vert dz_{a}\vert ^{2}\), the linear map
\(z\to z_{a}-z_{b}\) has squared norm \(l_{ab}^{2}=1/n_{a}+1/n_{b}\).
Let \(T_{ab}\) be kinetic energy in its nine-dimensional orthogonal
relative subspace. Then

% DISPLAY f.11
% SOURCE T_ab >= (49/4) l_ab² integral r_ab^(-2)|f|²,
% SOURCE T_ab <= T_centers.
\begin{gather*}
\begin{aligned}
& T_{ab} \ge  (49/4) l_{ab}^{2} \int  r_{ab}^{-2}\vert f\vert ^{2}, \\
& T_{ab} \le  T_{\mathrm{centers}}.
\end{aligned}\notag
\end{gather*}

Thus sum
\(\int  r_{ab}^{-2}\vert f\vert ^{2} \le (4/49) \sum  l_{ab}^{-2} T_{\mathrm{centers}}\).
The finite rank-dependent coefficient contains no inverse shape gap
\(\delta _{N}\). Choose \(\epsilon\), then \(\kappa\), before \(L\) and
the global radius. Nonflatness therefore causes no parameter-order
circle in this Berry comparison.

This is a lemma about the finite adapted vacuum connection and the full
slow charge forms. It does not prove that all remaining full-Hamiltonian
compression and Schur terms match that covariant slow model, nor does it
prove the adapted bosonic/fermionic energy-defect estimate.

\clearpage
\section{Differential connections and colored
forms}\label{g-differential-connections-and-colored-forms}

Technical derivation for review. Equation and subsection numbers in this
appendix are local to the source derivation. Historical local labels are
retained, including numbering gaps or nonsequential labels. The
accompanying source notes retain the original expressions for
comparison.

\subsection{1. Setting and result}\label{g-setting-and-result}

Fix the rank and block partition. Assume only

% DISPLAY g.1
% SOURCE alpha_a+alpha_b <= kappa r_ab,     |Y| <= sigma r_*,
% SOURCE r_* = min_e r_e >= 1,             W_2 = sum_e r_e^-2.
\begin{gather*}
\begin{aligned}
& \alpha _{a}+\alpha _{b} \le  \kappa  r_{ab}, \quad  \vert Y\vert  \le  \sigma  r_{*}, \\
& r_{*} = \min _{e} r_{e} \ge  1, \quad  W_{2} = \sum _{e} r_{e}^{-2}.
\end{aligned}\notag
\end{gather*}

Constants depend on fixed rank and fixed sufficiently small upper bounds
for \(\kappa\) and \(\sigma\), never on \(r_{\max}/r_{*}\). Work first
on the compactly supported smooth residual-equivariant core in a
normal-frame chart. No individual block-singlet assumption is imposed.
Let

% DISPLAY g.2
% SOURCE H_exc = sum_e (H_fast,e-E_0,e) >= 0,
% SOURCE lambda(Z) = sigma^2+r_*^-3,
% SOURCE E[v] = (lambda H_exc)[v]+W_2[v].                       (1.1)
\begin{gather*}
\begin{aligned}
& H_{\mathrm{exc}} = \sum _{e} (H_{\mathrm{fast},e}-E_{0,e}) \ge  0, \\
& \lambda (Z) = \sigma ^{2}+r_{*}^{-3}, \\
& E[v] = (\lambda  H_{\mathrm{exc}})[v]+W_{2}[v].
\end{aligned}\tag{1.1}
\end{gather*}

The weights are center-only and the excitation form acts in fast fibers.
For the INTERNAL part of the completion connection write

% DISPLAY g.3
% SOURCE beta_j = b_j(Y,A,Z) dot p_F,
% SOURCE beta = A_s^* B_F p_F,     A_s = B_s G_s^-1/2.
\begin{gather*}
\begin{aligned}
& \beta _{j} = b_{j}(Y,A,Z) \cdot  p_{F}, \\
& \beta  = A_{s}^{*} B_{F} p_{F}, \quad  A_{s} = B_{s} G_{s}^{-1/2}.
\end{aligned}\notag
\end{gather*}

Here \(A_{s}\) is the normalized Gauss slow coefficient, not the
internal matrix tuple. The internal block of \(G_{s}\) is exactly
identity. Define the ordered shifted form inherited from the exact
kinetic identity:

% DISPLAY g.4
% SOURCE H_I^raw(beta)[v]
% SOURCE   = sum_j ||(p_j+beta_j)v||^2+(V_I+Y_I)[v].             (1.2)
\begin{gather*}
\begin{aligned}
& H_{I}^{\mathrm{raw}}(\beta )[v] \\
& = \sum _{j} \Vert (p_{j}+\beta _{j})v\Vert ^{2}+(V_{I}+Y_{I})[v].
\end{aligned}\tag{1.2}
\end{gather*}

For every \(\epsilon\) in (0,1),

% DISPLAY g.5
% SOURCE H_I^raw(beta)[v]
% SOURCE   >= (1-epsilon)H_I[v]
% SOURCE        -C_(N,epsilon)(lambda H_exc)[v]
% SOURCE        -C_(N,epsilon)W_2[v].                         (1.3)
\begin{gather*}
\begin{aligned}
& H_{I}^{\mathrm{raw}}(\beta )[v] \\
& \ge  (1-\epsilon )H_{I}[v] \\
& -C_{N,\epsilon}(\lambda  H_{\mathrm{exc}})[v] \\
& -C_{N,\epsilon}W_{2}[v].
\end{aligned}\tag{1.3}
\end{gather*}

This retains the ENTIRE interacting internal form, including its
original potential and Yukawa terms. There is no \(\epsilon\) times bare
internal kinetic loss and no \(\epsilon  \alpha _{a}\) mass remainder.
No internal gap or physical-sector lower-rank Hardy estimate is used.

The proof gives the actual curvature expectation bound

% DISPLAY g.6
% SOURCE |R_beta[v]| <= C_N(lambda H_exc)[v]+C_N W_2[v].       (1.4)
\begin{gather*}
\begin{aligned}
& \vert R_{\beta}[v]\vert  \le  C_{N}(\lambda  H_{\mathrm{exc}})[v]+C_{N} W_{2}[v].
\end{aligned}\tag{1.4}
\end{gather*}

The stronger possible coefficient \(C_{N} r_{*}^{-3}\) with
\(\sigma ^{2}\) absent is not needed and is not claimed. Fixing
\(\sigma\) sufficiently small and then taking the exterior radius large
makes the excitation coefficient as small as desired. On fast-complement
inputs \(W_{2}\) is also paid by the gap.

\subsection{2. Exact completion}\label{g-exact-completion}

Remove the finite fermionic Gauss source by the separate estimate chosen
for that source. Its estimate is not part of this algebra. Let

% DISPLAY g.7
% SOURCE w=G_s^1/2 p_s v,   a=B_F p_F v,   A_s=B_s G_s^-1/2.
\begin{gather*}
\begin{aligned}
& w=G_{s}^{1/2} p_{s} v, \quad  a=B_{F} p_{F} v, \quad  A_{s}=B_{s} G_{s}^{-1/2}.
\end{aligned}\notag
\end{gather*}

Exactly, as a first-order form identity,

% DISPLAY g.8
% SOURCE ||w||^2+||a+A_s w||^2
% SOURCE   =||w+A_s^*a||^2+||a||^2+||A_s w||^2-||A_s^*a||^2. (2.1)
\begin{gather*}
\begin{aligned}
& \Vert w\Vert ^{2}+\Vert a+A_{s} w\Vert ^{2} \\
& =\Vert w+A_{s}^{*}a\Vert ^{2}+\Vert a\Vert ^{2}+\Vert A_{s} w\Vert ^{2}-\Vert A_{s}^{*}a\Vert ^{2}.
\end{aligned}\tag{2.1}
\end{gather*}

Coefficients multiply the differentiated section in their original
order. No integration by parts and no coefficient derivative occurs in
this completion.

The internal components of \(w+A_{s}^{*}a\) are \(p_{j}+\beta _{j}\).
Its center components use the exact center metric and NORMAL-BUNDLE
covariant momentum. The center square and \(\Vert A_{s} w\Vert ^{2}\)
may be retained or discarded as positive forms. They are not additional
copies of the fast Gauss energy. The negative term is the connection
norm already bounded in the rootwise lemma.

If the half-density was separated as in the intrinsic geometry note, its
scalar \(V_{d}\) must still be included separately. The symmetrization
below neither replaces nor duplicates that calculation.

\subsection{3. Coefficients and differentiated rootwise
bounds}\label{g-coefficients-and-differentiated-rootwise-bounds}

Use intrinsic notation

% DISPLAY g.9
% SOURCE F=E^*B_(Z+A),   C=P_N B_A,   N_Y=P_N B_Y,
% SOURCE I_Y=P_I B_Y,    M=F+L-J_2,   W=I+KK^*.
\begin{gather*}
\begin{aligned}
& F=E^{*}B_{Z+A}, \quad  C=P_{N} B_{A}, \quad  N_{Y}=P_{N} B_{Y}, \\
& I_{Y}=P_{I} B_{Y}, \quad  M=F+L-J_{2}, \quad  W=I+KK^{*}.
\end{aligned}\notag
\end{gather*}

The internal coefficient is, with the common Gauss sign convention,

% DISPLAY g.10
% SOURCE beta_I=I_Y M^-1 W M^-* (C+N_Y)^*p_F.                   (3.1)
\begin{gather*}
\begin{aligned}
& \beta _{I}=I_{Y} M^{-1} W M^{-*} (C+N_{Y})^{*}p_{F}.
\end{aligned}\tag{3.1}
\end{gather*}

All coefficients are real in real trace-orthonormal frames. Normal
frames depend on centers only; an internal derivative differentiates
neither \(p_{F}\) nor the frame. The connection commutes with internal
Clifford matrices and every internal potential-supercharge
multiplication coefficient.

For a unit internal direction \(h\),

% DISPLAY g.11
% SOURCE partial_h M=partial_h F,
% SOURCE partial_h F and partial_h C are pair-block diagonal,
% SOURCE ||partial_h F_e||+||partial_h C_e|| <= C_N,
% SOURCE partial_h W=partial_h I_Y=partial_h N_Y=0.             (3.2)
\begin{gather*}
\begin{aligned}
& \partial _{h} M=\partial _{h} F, \\
& \partial _{h} F \text{ and } \partial _{h} C \text{ are pair-block diagonal }, \\
& \Vert \partial _{h} F_{e}\Vert +\Vert \partial _{h} C_{e}\Vert  \le  C_{N}, \\
& \partial _{h} W=\partial _{h} I_{Y}=\partial _{h} N_{Y}=0.
\end{aligned}\tag{3.2}
\end{gather*}

Only pairs incident to the differentiated block contribute. The coarse
\(r_{*}\) estimate below does not convert any kinetic or Yukawa term, so
it creates no unrelated \(\alpha _{a}/r_{*}\) error.

The cited rootwise connection proof gives

% DISPLAY g.12
% SOURCE sum_j ||beta_jv||^2 <= C_N E[v].                     (3.3)
\begin{gather*}
\begin{aligned}
& \sum _{j} \Vert \beta _{j} v\Vert ^{2} \le  C_{N} E[v].
\end{aligned}\tag{3.3}
\end{gather*}

Its explicit first-order decomposition also gives

% DISPLAY g.13
% SOURCE sum_(j,k)|| (partial_j beta_k)v ||^2
% SOURCE    <= C_N [ (r_*^-2 lambda H_exc)[v]
% SOURCE                  +(r_*^-2 W_2)[v] ].                 (3.4)
\begin{gather*}
\begin{aligned}
& \sum _{j,k}\Vert  (\partial _{j} \beta _{k})v \Vert ^{2} \\
& \le  C_{N} [ (r_{*}^{-2} \lambda  H_{\mathrm{exc}})[v] \\
& +(r_{*}^{-2} W_{2})[v] ].
\end{aligned}\tag{3.4}
\end{gather*}

Here \(\partial _{j} \beta _{k}\) differentiates its coefficients only.
The following proves (3.4); it is not inferred from the abstract bound
(3.3).

\subsubsection{3.1 Proof of the differentiated
estimate}\label{g-proof-of-the-differentiated-estimate}

In the rootwise lemma put

% DISPLAY g.14
% SOURCE d_0=F^-1 C^*p_F,    T=(I+F^-1 P^*)^-1,    P=L-J_2.
\begin{gather*}
\begin{aligned}
& d_{0}=F^{-1} C^{*}p_{F}, \quad  T=(I+F^{-1} P^{*})^{-1}, \quad  P=L-J_{2}.
\end{aligned}\notag
\end{gather*}

The exact decompositions (5)-(7) and (12)-(15) in that lemma express
\(\beta\) as finite sums of bounded multiplication matrices times the
source classes

% DISPLAY g.15
% SOURCE Y_e p_e/r_e,
% SOURCE Y_f p_e/r_g,       when (e,f,g) is a triangle,
% SOURCE Y_f p_e/r_e,       with the input-root denominator,
\begin{gather*}
\begin{aligned}
& Y_{e} p_{e}/r_{e}, \\
& Y_{f} p_{e}/r_{g}, \quad  \text{ when } (e,f,g) \text{ is a triangle }, \\
& Y_{f} p_{e}/r_{e}, \quad  \text{ with the input-root denominator },
\end{aligned}\notag
\end{gather*}

with bounded root-preserving factors between the displayed operators.
Non-root-preserving inverse matrices remain bounded LEFT factors. Thus
no unweighted spectator momentum is hidden.

Differentiate those exact decompositions.
\(Y, r_{e}, W, K, I_{Y}, C_{Y}, L, J_{2}\) and \(g_{C}\) have zero
internal derivatives. The differentiated factors satisfy

% DISPLAY g.16
% SOURCE ||partial_j F_e^-1|| <= C_N/r_e^2,
% SOURCE ||partial_j(F_e^-1 C_e^*)|| <= C_N/r_e,
% SOURCE ||partial_j T||+||partial_j T^*|| <= C_N/r_*.
\begin{gather*}
\begin{aligned}
& \Vert \partial _{j} F_{e}^{-1}\Vert  \le  C_{N}/r_{e}^{2}, \\
& \Vert \partial _{j}(F_{e}^{-1} C_{e}^{*})\Vert  \le  C_{N}/r_{e}, \\
& \Vert \partial _{j} T\Vert +\Vert \partial _{j} T^{*}\Vert  \le  C_{N}/r_{*}.
\end{aligned}\notag
\end{gather*}

The inverse derivative formula and (3.2) prove these inequalities. All
other bounded left factors in the displayed decompositions consequently
have derivative at most \(C_{N}/r_{*}\). Differentiating
\(F_{e}^{-1} C_{e}^{*}\) replaces its \(O(\kappa )\) factor by
\(O(1/r_{e})\); use the source bound with its optional \(\kappa\) gain
omitted.

Every resulting term is therefore the SAME source class with one extra
factor at most \(C_{N}/r_{*}\). Applying the proved same-root and
mixed-root inequalities (2)-(4) of the cited lemma, and finite-sum
Cauchy, proves (3.4). The argument is all-vector, does not bound
internal momenta, and does not require a fast projection to preserve
tube support.

\subsubsection{3.2 Coarse pointwise derivatives for
divergence}\label{g-coarse-pointwise-derivatives-for-divergence}

Only for the scalar divergence correction we need

% DISPLAY g.17
% SOURCE |b| <= C_N sigma(kappa+sigma),
% SOURCE |partial_Y b| <= C_N/r_*,
% SOURCE |partial_A partial_Y b|+|partial_Y^2 b| <= C_N/r_*^2. (3.5)
\begin{gather*}
\begin{aligned}
& \vert b\vert  \le  C_{N} \sigma (\kappa +\sigma ), \\
& \vert \partial _{Y} b\vert  \le  C_{N}/r_{*}, \\
& \vert \partial _{A} \partial _{Y} b\vert +\vert \partial _{Y}^{2} b\vert  \le  C_{N}/r_{*}^{2}.
\end{aligned}\tag{3.5}
\end{gather*}

Indeed the internal normalized slow column is \(O(\sigma )\), the full
fast Gauss column is \(O(\kappa +\sigma )\), and
\(\Vert M^{-1}\Vert \le C_{N}/r_{*}\). First derivatives of \(M, I_{Y}\)
and \(C+N_{Y}\) are \(O(1)\); second derivatives of \(M\) are
\(O(1/r_{*})\), coming from \(J_{2}\). First and second derivatives of
\(W\) and its positive square root are \(O(1/r_{*})\) and
\(O(1/r_{*}^{2})\). Differentiate the product \(B_{I}^{*}B_{F}\) using
the inverse derivative formula. The potentially large \(C\) is always
kept inside \(F^{-1} C^{*}\) or the bounded full fast column, rather
than estimated by \(\kappa  r_{\max}\). This proves (3.5); fixed
coordinate dimensions absorb component sums.

\subsection{4. Exact self-adjoint
symmetrization}\label{g-exact-self-adjoint-symmetrization}

In flat fast-fiber measure, with \(p_{F}=-i \partial _{Y}\), put

% DISPLAY g.18
% SOURCE d_j=div_F b_j,
% SOURCE beta_hat_j=(beta_j+beta_j^*)/2
% SOURCE           =b_j dot p_F-(i/2)d_j,
% SOURCE P_j=p_j+beta_hat_j,
% SOURCE X_j=partial_j+b_j dot partial_Y.
\begin{gather*}
\begin{aligned}
& d_{j}=\operatorname{div} _{F} b_{j}, \\
& \widehat{\beta}_{j}=(\beta _{j}+\beta _{j}^{*})/2 \\
& =b_{j} \cdot  p_{F}-(i/2)d_{j}, \\
& P_{j}=p_{j}+\widehat{\beta}_{j}, \\
& X_{j}=\partial _{j}+b_{j} \cdot  \partial _{Y}.
\end{aligned}\notag
\end{gather*}

Then \(\beta _{j}=\widehat{\beta}_{j}+(i/2)d_{j}\) and
\(P_{j}=-i(X_{j}+d_{j}/2)\) is symmetric on the core. Direct integration
by parts gives the EXACT identity

% DISPLAY g.19
% SOURCE ||(p_j+beta_j)v||^2
% SOURCE   =||P_jv||^2
% SOURCE       + integral [(1/2)X_j d_j+(1/4)d_j^2]|v|^2.     (4.1)
\begin{gather*}
\begin{aligned}
& \Vert (p_{j}+\beta _{j})v\Vert ^{2} \\
& =\Vert P_{j} v\Vert ^{2} \\
& + \int  [(1/2)X_{j} d_{j}+(1/4)d_{j}^{2}]\vert v\vert ^{2}.
\end{aligned}\tag{4.1}
\end{gather*}

The cross term is \((1/2) \int  X_{j}(d_{j})\vert v\vert ^{2}\); div
\(X_{j}=d_{j}\). Thus simply treating \(b \cdot  p_{F}\) as symmetric
would miss a scalar correction.

By (3.5),

% DISPLAY g.20
% SOURCE |d_j| <= C_N/r_*,
% SOURCE |partial_A d_j|+|partial_Y d_j| <= C_N/r_*^2,
% SOURCE sum_j |(1/2)X_j d_j+(1/4)d_j^2| <= C_N W_2.           (4.2)
\begin{gather*}
\begin{aligned}
& \vert d_{j}\vert  \le  C_{N}/r_{*}, \\
& \vert \partial _{A} d_{j}\vert +\vert \partial _{Y} d_{j}\vert  \le  C_{N}/r_{*}^{2}, \\
& \sum _{j} \vert (1/2)X_{j} d_{j}+(1/4)d_{j}^{2}\vert  \le  C_{N} W_{2}.
\end{aligned}\tag{4.2}
\end{gather*}

Hence the full correction is inverse-square multiplication. Also

% DISPLAY g.21
% SOURCE sum_j ||beta_hat_jv||^2 <= C_N E[v],                  (4.3)
\begin{gather*}
\begin{aligned}
& \sum _{j} \Vert \widehat{\beta}_{j} v\Vert ^{2} \le  C_{N} E[v],
\end{aligned}\tag{4.3}
\end{gather*}

and (3.4) holds with \(\widehat{\beta}\) replacing \(\beta\): its extra
multiplication derivative is \(O(r_{*}^{-2})\), whose squared norm is
bounded by \(r_{*}^{-2} W_{2}\).

\subsection{5. Actual curvature
expectation}\label{g-actual-curvature-expectation}

Define the symmetric curvature operator

% DISPLAY g.22
% SOURCE F_jk=i[P_j,P_k]
% SOURCE     =partial_j beta_hat_k-partial_k beta_hat_j
% SOURCE                      +i[beta_hat_j,beta_hat_k].       (5.1)
\begin{gather*}
\begin{aligned}
& F_{jk}=i[P_{j},P_{k}] \\
& =\partial _{j} \widehat{\beta}_{k}-\partial _{k} \widehat{\beta}_{j} \\
& +i[\widehat{\beta}_{j},\widehat{\beta}_{k}].
\end{aligned}\tag{5.1}
\end{gather*}

Equivalently \(F_{jk}\) is the symmetrized momentum of the fast vector
field

% DISPLAY g.23
% SOURCE partial_j b_k-partial_k b_j
% SOURCE   +(b_j dot partial_Y)b_k-(b_k dot partial_Y)b_j.
\begin{gather*}
\begin{aligned}
& \partial _{j} b_{k}-\partial _{k} b_{j} \\
& +(b_{j} \cdot  \partial _{Y})b_{k}-(b_{k} \cdot  \partial _{Y})b_{j}.
\end{aligned}\notag
\end{gather*}

It is first order, not bounded multiplication.

The charge identity contracts (5.1) with fixed bounded internal-Clifford
matrices. If \(L\) is any such matrix, it commutes with
\(\widehat{\beta}\) and its coefficient derivatives. Symmetry gives
exactly

% DISPLAY g.24
% SOURCE |<v,L i[beta_hat_j,beta_hat_k]v>|
% SOURCE    <=2||L|| ||beta_hat_jv|| ||beta_hat_kv||.           (5.2)
\begin{gather*}
\begin{aligned}
& \vert \langle v,L i[\widehat{\beta}_{j},\widehat{\beta}_{k}]v\rangle \vert \\
& \le 2\Vert L\Vert  \Vert \widehat{\beta}_{j} v\Vert  \Vert \widehat{\beta}_{k} v\Vert .
\end{aligned}\tag{5.2}
\end{gather*}

Thus (4.3) pays all beta-beta curvature terms. There is no need to bound
the norm of a fast-differentiated curvature symbol.

For the derivative part, (3.4) and Cauchy give

% DISPLAY g.25
% SOURCE |<v,L(partial_j beta_hat_k)v>|
% SOURCE    <= C_N ||r_*^-1 v|| E[v]^1/2
% SOURCE    <= C_N E[v],                                     (5.3)
\begin{gather*}
\begin{aligned}
& \vert \langle v,L(\partial _{j} \widehat{\beta}_{k})v\rangle \vert \\
& \le  C_{N} \Vert r_{*}^{-1} v\Vert  E[v]^{1/2} \\
& \le  C_{N} E[v],
\end{aligned}\tag{5.3}
\end{gather*}

since \(\Vert r_{*}^{-1}v\Vert ^{2}\le W_{2}[v]\le E[v]\). To justify
the first inequality with varying centers, apply the coefficient
derivative bound at each center and then apply weighted Cauchy over
centers. Summing the finitely many Clifford components proves (1.4).

This is an all-vector curvature EXPECTATION estimate. It derives the
extra coefficient derivative from the actual rootwise expansion, not
from a generic relative-norm assertion. It makes no oscillator-degree
cutoff and loses no internal kinetic energy.

\subsection{6. Averaged sixteen-supercharge
comparison}\label{g-averaged-sixteen-supercharge-comparison}

Normalize the internal charges by

% DISPLAY g.26
% SOURCE H_I[v]=(1/16)sum_alpha ||Q_alpha v||^2,
% SOURCE Q_alpha=sum_j c_(alpha,j)p_j+V_alpha(A).
\begin{gather*}
\begin{aligned}
& H_{I}[v]=(1/16)\sum _{\alpha} \Vert Q_{\alpha} v\Vert ^{2}, \\
& Q_{\alpha}=\sum _{j} c_{\alpha ,j}p_{j}+V_{\alpha}(A).
\end{aligned}\notag
\end{gather*}

The \(c_{\alpha ,j}\) are fixed internal Clifford matrices. The ordinary
colored Gauss terms cancel in the spinor average. An individual
\(Q_{\alpha}^{2}\) must not be identified with \(H_{I}\) on the colored
core.

Replace \(p_{j}\) by \(P_{j}\):

% DISPLAY g.27
% SOURCE Q_alpha^beta=Q_alpha+C_alpha(beta_hat),
% SOURCE C_alpha(beta_hat)=sum_j c_(alpha,j) beta_hat_j.
\begin{gather*}
\begin{aligned}
& Q_{\alpha}^{\beta}=Q_{\alpha}+C_{\alpha}(\widehat{\beta}), \\
& C_{\alpha}(\widehat{\beta})=\sum _{j} c_{\alpha ,j} \widehat{\beta}_{j}.
\end{aligned}\notag
\end{gather*}

The exact averaged expansion is

% DISPLAY g.28
% SOURCE (1/16)sum_alpha ||Q_alpha^beta v||^2
% SOURCE    =sum_j||P_jv||^2+(V_I+Y_I)[v]+R_beta[v].           (6.1)
\begin{gather*}
\begin{aligned}
& (1/16)\sum _{\alpha} \Vert Q_{\alpha}^{\beta} v\Vert ^{2} \\
& =\sum _{j}\Vert P_{j} v\Vert ^{2}+(V_{I}+Y_{I})[v]+R_{\beta}[v].
\end{aligned}\tag{6.1}
\end{gather*}

Here \(R_{\beta}\) is a finite-Clifford linear contraction of
\(F_{jk}\). The symmetric momentum products give the kinetic form,
antisymmetric momentum products give curvature, and

% DISPLAY g.29
% SOURCE [P_j,V_alpha]=[p_j,V_alpha]
\begin{gather*}
\begin{aligned}
& [P_{j},V_{\alpha}]=[p_{j},V_{\alpha}]
\end{aligned}\notag
\end{gather*}

because \(\widehat{\beta}\) is fast-bosonic with scalar bosonic
coefficients. The original potential and Yukawa terms are consequently
unchanged. Spinor-averaged cancellation of the original Gauss term is
unchanged too.

Finite-dimensional Clifford bounds and (4.3) imply

% DISPLAY g.30
% SOURCE (1/16)sum_alpha ||C_alpha(beta_hat)v||^2 <= C_N E[v].  (6.2)
\begin{gather*}
\begin{aligned}
& (1/16)\sum _{\alpha} \Vert C_{\alpha}(\widehat{\beta})v\Vert ^{2} \le  C_{N} E[v].
\end{aligned}\tag{6.2}
\end{gather*}

Young\textquotesingle s inequality on the COMPLETE charge stack now
gives

% DISPLAY g.31
% SOURCE (1/16)sum_alpha ||Q_alpha^beta v||^2
% SOURCE    >=(1-epsilon)H_I[v]-C_N epsilon^-1 E[v].           (6.3)
\begin{gather*}
\begin{aligned}
& (1/16)\sum _{\alpha} \Vert Q_{\alpha}^{\beta} v\Vert ^{2} \\
& \ge (1-\epsilon )H_{I}[v]-C_{N} \epsilon ^{-1} E[v].
\end{aligned}\tag{6.3}
\end{gather*}

Subtract curvature using (1.4), then apply (4.1)-(4.2). This proves
(1.3). The loss is \(\epsilon  H_{I}\), not \(\epsilon  t_{I}\); that is
why no large unrelated internal Yukawa mass is introduced.

\subsection{7. Normal-bundle center connection and chart
changes}\label{g-normal-bundle-center-connection-and-chart-changes}

The normal fast bundle \(N_{Z}\) depends on centers alone. Its
orthogonal connection is already in \(p_{C}\) in the exact intrinsic
formula. In a local frame its center generator may contain unbounded
\(Y \partial _{Y}\) operators. This argument never estimates those by
\(H_{\mathrm{exc}}\).

Internal derivatives act in a center-dependent but A-independent frame,
so (3.2) has no frame derivative. The internal supercharges contain no
center momentum. Neither center normal-bundle curvature nor mixed
center/internal curvature occurs in the INTERNAL square (6.1).

The completed center square remains a genuine positive covariant
derivative form. Any subsequent center Hardy argument must use that
actual connection with its domain justified. It may not replace its
unbounded normal-frame generator by a bounded vacuum Berry coefficient.
That separate comparison is not claimed here.

All identities intertwine under center-dependent orthogonal changes of
fast normal frame and residual gauge changes. \(\beta\) is an actual
vector field on the fast fiber, \(\widehat{\beta}\) is its half-density
action, and curvature is its commutator. The internal comparison is
therefore intrinsic on a fixed projector branch.

\subsection{8. Parameter order and exact
scope}\label{g-parameter-order-and-exact-scope}

Choose the requested internal loss \(\epsilon\) first, then \(\sigma\)
so \(C_{N,\epsilon}\sigma ^{2}\) fits the available fast reserve, then
\(r_{*}\) large enough for \(C_{N,\epsilon}r_{*}^{-3}\) and the
inverse-square multiplier.

On the global fast complement \(H_{\mathrm{exc}}\ge c_{N} r_{*}\) and
\(W_{2}\le m_{N} r_{*}^{-2}\), so

% DISPLAY g.32
% SOURCE W_2[v] <= C_N (r_*^-3 H_exc)[v].
\begin{gather*}
\begin{aligned}
& W_{2}[v] \le  C_{N} (r_{*}^{-3} H_{\mathrm{exc}})[v].
\end{aligned}\notag
\end{gather*}

This final payment is restricted to complementary inputs; (1.3) itself
holds on arbitrary colored inputs. No internal Hamiltonian is inverted.

Still separate are the finite spin source, density multiplier, signed
fast kinetic remainder, potential/Yukawa remainders, and branch/cutoff
compatibility. The present lemma closes the complete internal-form and
curvature step only.

\clearpage
\section{Colored complementary
reserve}\label{h-colored-complementary-reserve}

Technical derivation for review. Equation and subsection numbers in this
appendix are local to the source derivation. Historical local labels are
retained, including numbering gaps or nonsequential labels. The
accompanying source notes retain the original expressions for
comparison.

\subsection{1. Coordinate conventions and exact
matrices}\label{h-coordinate-conventions-and-exact-matrices}

Use the intrinsic geometry note and the notation of the actual adapted
source inventory:

% DISPLAY h.1
% SOURCE F=E*B_(Z+A), C=P_NB_A, L=E*B_Y, N_Y=P_NB_Y,
% SOURCE I_Y=P_IB_Y, C_Y=P_CB_Y, K=R^−1C_Y*, J2=KC_Y,
% SOURCE M=F+L−J2, W=I+KK*, g_C=I+K*K.
\begin{gather*}
\begin{aligned}
& F=E^{*}B_{Z+A}, C=P_{N} B_{A}, L=E^{*}B_{Y}, N_{Y}=P_{N} B_{Y}, \\
& I_{Y}=P_{I} B_{Y}, C_{Y}=P_{C} B_{Y}, K=R^{-1}C_{Y}^{*}, J_{2}=KC_{Y}, \\
& M=F+L-J_{2}, W=I+KK^{*}, g_{C}=I+K^{*}K.
\end{aligned}\notag
\end{gather*}

Let \(p_{s}=(g_{C}^{-1/2} p_{C},p_{I})\), where \(p_{C}\) includes the
exact normal connection. Define

% DISPLAY h.2
% SOURCE Ψ=[(C_Y*+(F+L)*K)g_C^−1/2,I_Y*],
% SOURCE A=W^1/2M^−*Ψ,
% SOURCE a=Bp_F=W^1/2M^−*(C+N_Y)*p_F,
% SOURCE s=−W^1/2M^−*S_f.
\begin{gather*}
\begin{aligned}
& \Psi  =[(C_{Y}^{*}+(F+L)^{*}K)g_{C}^{-1/2},I_{Y}^{*}], \\
& A=W^{1/2}M^{-*} \Psi  , \\
& a=Bp_{F}=W^{1/2}M^{-*}(C+N_{Y})^{*}p_{F}, \\
& s=-W^{1/2}M^{-*}S_{f}.
\end{aligned}\notag
\end{gather*}

The minus sign in \(s\) matches the original Gauss expression
\(S_{f}-\mathrm{fast}-\mathrm{slow}\) after multiplying that whole stack
by \(-1\). The unshifted scalar kinetic stack is exactly

% DISPLAY h.3
% SOURCE |p_s v|²+|p_Fv|²+|a v+A p_s v+s v|².                 (1.1)
\begin{gather*}
\begin{aligned}
& \vert p_{s} v\vert ^{2}+\vert p_{F} v\vert ^{2}+\vert a v+A p_{s} v+s v\vert ^{2}.
\end{aligned}\tag{1.1}
\end{gather*}

The scalar half-density term is added after this first-order identity.
We never add the same Gauss square a second time to
\(H_{\mathrm{fast}}\).

Throughout, \(T_{\mathrm{cent}}[v]=\Vert p_{C}v\Vert ^{2}\) with the
exact normal connection. Its direct induced-metric coefficient obeys
\(g_{C}^{-1}\ge (1+C \sigma ^{2})^{-1}I\); this small center-only loss
is included in the final allocation.

Set
\(\mathrm{D0}=F^{-1}C^{*}p_{F}, P=L-J_{2}, Z=F^{-1}P^{*}, T=(I+Z)^{-1}, n=F^{-1}N_{Y}^{*}p_{F}\).
Thus

% DISPLAY h.4
% SOURCE a=W^1/2 T(D0+n),   ||T||+||T^−1||≤C,   ||A||≤Cσ.   (1.2)
\begin{gather*}
\begin{aligned}
& a=W^{1/2} T(\mathrm{D0}+n), \quad  \Vert T\Vert +\Vert T^{-1}\Vert \le C, \quad  \Vert A\Vert \le C \sigma  .
\end{aligned}\tag{1.2}
\end{gather*}

The frozen kinetic contribution in \(H_{\mathrm{fast}}\) is
\(p_{F}^{2}+\Vert \mathrm{D0}\Vert ^{2}\).

\subsection{2. Oscillator estimates with the vacuum kept
explicitly}\label{h-oscillator-estimates-with-the-vacuum-kept-explicitly}

The adapted bosonic Gaussian on edge \(e\) is
\(\exp (-Y_{e} \cdot  \Omega _{e}Y_{e}/2)\), with
\(c r_{e} I\le  \Omega _{e}\le C r_{e}\) I. Let
\(a_{e}= \partial  _{e}+ \Omega _{e}Y_{e}\) be its annihilation
gradient. Then \(\Vert a_{e} v\Vert ^{2}\le C h_{e}[v]\), including
arbitrary fermionic and spectator components. The shifted fermionic
excitation form is nonnegative and only improves this inequality.

The following estimates are consequences of splitting each edge into its
adapted full vacuum \(P_{e}\) and \(Q_{e}=1-P_{e}\); see Appendix \(E\)
for the first two:

% DISPLAY h.5
% SOURCE ||χY_ep_e/r_e v||²≤Cσ²h_e[v]+Cr_e^−2||v||²,         (2.1)
% SOURCE 
% SOURCE ||χY_fp_e/r_g v||²≤C(σ²+r_*^−3)(h_e+h_f)[v]
% SOURCE                          +CΣ_(u in triangle)r_u^−2||v||². (2.2)
\begin{gather*}
\begin{aligned}
& \Vert  \chi  Y_{e} p_{e}/r_{e} v\Vert ^{2}\le C \sigma ^{2}h_{e}[v]+Cr_{e}^{-2}\Vert v\Vert ^{2},
\end{aligned}\tag{2.1} \\
\begin{aligned}
& \Vert  \chi  Y_{f} p_{e}/r_{g} v\Vert ^{2}\le C( \sigma ^{2}+r_{*}^{-3})(h_{e}+h_{f})[v] \\
& +C \sum _{u \in  \text{ triangle }}r_{u}^{-2}\Vert v\Vert ^{2}.
\end{aligned}\tag{2.2}
\end{gather*}

A useful sharpened two-coordinate version is

% DISPLAY h.6
% SOURCE ||χY_fY_gv||²≤Cσ²r_*²(h_f[v]/r_f²+h_g[v]/r_g²)
% SOURCE                          +C/(r_fr_g)||v||²,               (2.3)
\begin{gather*}
\begin{aligned}
& \Vert  \chi  Y_{f} Y_{g} v\Vert ^{2}\le C \sigma ^{2}r_{*}^{2}(h_{f}[v]/r_{f}^{2}+h_{g}[v]/r_{g}^{2}) \\
& +C/(r_{f} r_{g})\Vert v\Vert ^{2},
\end{aligned}\tag{2.3}
\end{gather*}

for distinct edges f,g. To prove it, decompose into
\(Q_{f}, P_{f}Q_{g}\) and \(P_{f}P_{g}\). On the first two pieces use
\(\vert  \chi  Y_{g}\vert \le C \sigma  r_{*}\) or
\(\vert  \chi  Y_{f}\vert \le C \sigma  r_{*}\), and
\(\Vert Y_{f}Q_{f} v\Vert ^{2}\le C h_{f}[v]/r_{f}^{2}\) (similarly g).
On the last use the exact Gaussian second moments. Finite-sum
Cauchy--Schwarz supplies the displayed constant. The cutoff acts only on
the left; it need not commute with \(P_{f}\).

In every triangle e,f,g,

% DISPLAY h.7
% SOURCE r_*r_g/(r_er_f)≤2,    r_g/(r_e²r_f)≤C(r_e^−2+r_f^−2). (2.4)
\begin{gather*}
\begin{aligned}
& r_{*}r_{g}/(r_{e} r_{f})\le 2, \quad  r_{g}/(r_{e}^{2}r_{f})\le C(r_{e}^{-2}+r_{f}^{-2}).
\end{aligned}\tag{2.4}
\end{gather*}

These elementary triangle inequalities are what remove remote zero-point
energies.

\subsection{3. Direct cubic potential, with arbitrary fast
vectors}\label{h-direct-cubic-potential-with-arbitrary-fast-vectors}

Let e,f,g be a triangle, \(s_{t}=r_{e}+r_{f}+r_{g}\), and let \(B_{t}\)
be any fixed bounded trilinear contraction of the three real fast
coordinate spaces, also allowing a bounded spectator matrix. Then

% DISPLAY h.8
% SOURCE |<v,s_t B_t(Y_e,Y_f,Y_g)v>|
% SOURCE   ≤Cσ(h_e+h_f+h_g)[v]+Cσ^−1Σ_(u in t)r_u^−2||v||². (3.1)
\begin{gather*}
\begin{aligned}
& \vert \langle v,s_{t} B_{t}(Y_{e},Y_{f},Y_{g})v\rangle \vert \\
& \le C \sigma  (h_{e}+h_{f}+h_{g})[v]+C \sigma ^{-1} \sum _{u \in  t}r_{u}^{-2}\Vert v\Vert ^{2}.
\end{aligned}\tag{3.1}
\end{gather*}

Choose \(e\) with largest \(r_{e}\). For a scalar component, integration
by parts gives exactly

% DISPLAY h.9
% SOURCE < v,Y_eY_fY_g v >
% SOURCE  =Re< a_ev, Ω_e^−1(Y_fY_g)v >.                           (3.2)
\begin{gather*}
\begin{aligned}
& \langle v,Y_{e} Y_{f} Y_{g} v\rangle \\
& =\operatorname{Re} \langle a_{e} v, \Omega _{e}^{-1}(Y_{f} Y_{g})v\rangle .
\end{aligned}\tag{3.2}
\end{gather*}

There is no derivative of \(Y_{f} Y_{g}\) in the \(e\) coordinate.
Equation (3.2) applies componentwise to the finite contraction, with
\(\Omega _{e}^{-1}\) acting on its \(e\) index. Use
\(s_{t}/r_{e}\le 3\), (2.3), \(r_{*}\le r_{f},r_{g}\) and Young with
parameter \(\sigma\). The source term is at most
\(C \sigma ^{-1}/(r_{f} r_{g})\), which is bounded by the triangle
\(W_{2}\). This proves (3.1). No finite oscillator truncation is made.

Every actual cubic potential term has a bounded coefficient times
\(s_{t}\): its coefficient is a commutator with \(Z+A\), whose norm is
\(\le C(r_{e}+r_{f}+r_{g})\) under the pairwise-relative internal tube.
Therefore (3.1) applies to the full \(V_{3}\). \(V_{4}\) remains
positive and is not spent.

\subsection{4. Signed triangle kinetic
forms}\label{h-signed-triangle-kinetic-forms}

For any bounded real coefficient tensor independent of the fast
variables \(Y\) (but possibly dependent on the slow variables), and
three distinct triangle edges,

% DISPLAY h.10
% SOURCE |2Re<p_ev, (Y_f/r_e)p_gv>|
% SOURCE  ≤Cσ(h_e+h_f+h_g)[v]+Cσ^−1Σ_(u in t)r_u^−2||v||². (4.1)
\begin{gather*}
\begin{aligned}
& \vert 2\operatorname{Re} \langle p_{e} v, (Y_{f}/r_{e})p_{g} v\rangle \vert \\
& \le C \sigma  (h_{e}+h_{f}+h_{g})[v]+C \sigma ^{-1} \sum _{u \in  t}r_{u}^{-2}\Vert v\Vert ^{2}.
\end{aligned}\tag{4.1}
\end{gather*}

The clean proof is the Gaussian ground-state transform. The real
symmetric principal coefficient \(a_{eg}(Y)=Y_{f}/r_{e}\) has no
divergence in either active derivative coordinate e,g. For the product
Gaussian \(g_{B}\),

% DISPLAY h.11
% SOURCE ∫(∂v)*a∂v
% SOURCE  =∫g_B²(∂(v/g_B))*a∂(v/g_B)
% SOURCE   +∫[div(aΩY)−(ΩY)·a(ΩY)]|v|².                         (4.2)
\begin{gather*}
\begin{aligned}
& \int  ( \partial  v)^{*}a \partial  v \\
& = \int  g_{B}^{2}( \partial  (v/g_{B}))^{*}a \partial  (v/g_{B}) \\
& + \int  [\operatorname{div} (a \Omega  Y)-( \Omega  Y) \cdot  a( \Omega  Y)]\vert v\vert ^{2}.
\end{aligned}\tag{4.2}
\end{gather*}

The first term is \(\le C \sigma  \Sigma  h_{e}\) because
\(\vert Y_{f}\vert /r_{e}\le C \sigma\) on the support. The divergence
term is zero since the three roots are distinct. The last term is a
triangle cubic with coefficient bounded by \(C r_{g}\), hence by
\(C s_{t}\), and (3.1) proves (4.1). An outer cutoff equal to one near
the support contributes no derivative term there.

One may instead expand \(p=-i a+i \Omega  Y\). Formula (2.3), together
with (2.4), controls the apparently dangerous factor \(r_{g}/r_{e}\) in
the mixed annihilator/quadratic term. This gives the same bound without
ever estimating \(p_{e}^{2}\) by \(h_{e}\) plus an unassigned \(r_{e}\)
vacuum mass.

\subsection{5. Exact inverse-M reduction: higher paths are products of
errors}\label{h-exact-inverse-m-reduction-higher-paths-are-products-of-errors}

Let \(X=n-Z\mathrm{D0}\). The connection bounds imply

% DISPLAY h.12
% SOURCE ||Xv||²+||Z*D0v||²
% SOURCE  ≤C(σ²+r_*^−3)H_exc[v]+CW2[v].                          (5.1)
\begin{gather*}
\begin{aligned}
& \Vert Xv\Vert ^{2}+\Vert Z^{*}\mathrm{D0}v\Vert ^{2} \\
& \le C( \sigma ^{2}+r_{*}^{-3})H_{\mathrm{exc}}[v]+CW_{2}[v].
\end{aligned}\tag{5.1}
\end{gather*}

For \(X\) this is the displayed helper lemma. For
\(Z^{*}\mathrm{D0}=P F^{-1}\mathrm{D0}\) the linear \(L\) term is a
triangle map \(Y_{f} p_{e}/r_{e}\); the \(J_{2}\) term factors through
\(C_{Y} F^{-1}\mathrm{D0}\), a same-edge map, followed by bounded \(K\).
Thus the identical proof applies. Denominators stay on their actual
input or output root.

Since \(W\ge I\) and \(T-I=-ZT\) exactly,

% DISPLAY h.13
% SOURCE ||av||²≥||D0v+TXv||²
% SOURCE  ≥||D0v||²+2Re<D0v,Xv>−2|<Z*D0v,TXv>|.                 (5.2)
\begin{gather*}
\begin{aligned}
& \Vert av\Vert ^{2}\ge \Vert \mathrm{D0}v+TXv\Vert ^{2} \\
& \ge \Vert \mathrm{D0}v\Vert ^{2}+2\operatorname{Re} \langle \mathrm{D0}v,Xv\rangle -2\vert \langle Z^{*}\mathrm{D0}v,TXv\rangle \vert .
\end{aligned}\tag{5.2}
\end{gather*}

The last term is bounded by (5.1). The \(J_{2}\) part of the remaining
cross factors as

% DISPLAY h.14
% SOURCE <D0,F^−1J2*D0>=<C_YF^−1D0,K*D0>,                       (5.3)
\begin{gather*}
\begin{aligned}
& \langle \mathrm{D0},F^{-1}J_{2}^{*}\mathrm{D0}\rangle =\langle C_{Y} F^{-1}\mathrm{D0},K^{*}\mathrm{D0}\rangle ,
\end{aligned}\tag{5.3}
\end{gather*}

and both factors satisfy the same-edge bound. The two other terms are
precisely

% DISPLAY h.15
% SOURCE 2Re<D0,F^−1N_Y*p_F>−2Re<D0,F^−1L*D0>.                   (5.4)
\begin{gather*}
\begin{aligned}
& 2\operatorname{Re} \langle \mathrm{D0},F^{-1}N_{Y}^{*}p_{F}\rangle -2\operatorname{Re} \langle \mathrm{D0},F^{-1}L^{*}\mathrm{D0}\rangle .
\end{aligned}\tag{5.4}
\end{gather*}

Their root supports are triangles; their coefficients are bounded by
\(C \kappa  /r_{e}\) and \(C \kappa ^{2}/r_{e}\). They obey (4.1).
Consequently

% DISPLAY h.16
% SOURCE ||av||²≥||D0v||²−C(κσ+σ²+r_*^−3)H_exc[v]
% SOURCE                          −C(1+κ/σ)W2[v].                 (5.5)
\begin{gather*}
\begin{aligned}
& \Vert av\Vert ^{2}\ge \Vert \mathrm{D0}v\Vert ^{2}-C( \kappa  \sigma  + \sigma ^{2}+r_{*}^{-3})H_{\mathrm{exc}}[v] \\
& -C(1+ \kappa  / \sigma  )W_{2}[v].
\end{aligned}\tag{5.5}
\end{gather*}

This is an actual all-vector signed-form comparison for the coupled
inverse-M fast column. It does not require a separate estimate of every
higher Neumann path. Importantly, it never bounds the large
\(\mathrm{D0}\) vacuum norm by the nearest complementary gap.

\subsection{6. Exact slow completion and finite-spin
payments}\label{h-exact-slow-completion-and-finite-spin-payments}

With \(s\) temporarily omitted, the identity

% DISPLAY h.17
% SOURCE ||p_s v||²+||a v+A p_s v||²
% SOURCE  =||(p_s+A*a)v||²+||av||²+||A p_sv||²−||A*a v||²         (6.1)
\begin{gather*}
\begin{aligned}
& \Vert p_{s} v\Vert ^{2}+\Vert a v+A p_{s} v\Vert ^{2} \\
& =\Vert (p_{s}+A^{*}a)v\Vert ^{2}+\Vert av\Vert ^{2}+\Vert A p_{s} v\Vert ^{2}-\Vert A^{*}a v\Vert ^{2}
\end{aligned}\tag{6.1}
\end{gather*}

holds pointwise on the first-order stack. The final negative term is
controlled by

% DISPLAY h.18
% SOURCE ||A*a v||²≤C(σ²+r_*^−3)H_exc[v]+CW2[v],                  (6.2)
\begin{gather*}
\begin{aligned}
& \Vert A^{*}a v\Vert ^{2}\le C( \sigma ^{2}+r_{*}^{-3})H_{\mathrm{exc}}[v]+CW_{2}[v],
\end{aligned}\tag{6.2}
\end{gather*}

because
\(A^{*}a=A^{*}\mathrm{D0}+A^{*}(a-\mathrm{D0}), \Vert A\Vert \le C \sigma\),
and both helper norms are established. Equation (6.1) leaves the
explicit nonnegative reserve \(\Vert A p_{s} v\Vert ^{2}\). It does not
reserve another copy of the full Gauss square.

The center component of the shifted square is compared to the original
exact normal-covariant center kinetic by ordinary Young, costing
\(\epsilon  T_{\mathrm{cent}}\) plus \(\epsilon ^{-1}\) times (6.2).
Internal components are treated by the complete averaged charges in
Section \(7\); one must not apply this ordinary Young inequality to the
internal kinetic square alone.

Restore \(s\). Drop its positive norm square and expand its two cross
terms. Put \(s_{0}=-F^{-1}S_{f}\), so
\(\Vert s_{0}\Vert ^{2}\le CW_{2}\). The cross \(\langle s,a\rangle\)
differs from \(\langle s_{0},\mathrm{D0}\rangle\) by a quantity bounded
by \(\epsilon  H_{\mathrm{exc}}+C_{\epsilon} W_{2}\). To verify this,
expand \(T^{*}WT-I\) once: every factor moved onto \(\mathrm{D0}\) is
\(Z\mathrm{D0}, Z^{*}\mathrm{D0}\), or \(KK^{*}\mathrm{D0}\), all
controlled by (2.1), (2.2), (5.1). The remaining factors are uniformly
bounded. The \(n\) terms use (2.2). The leading
\(\langle s_{0},\mathrm{D0}\rangle\) is a finite \(\sum  p_{e}\) times a
Hermitian fast/slow Clifford matrix of norm \(C \kappa  /r_{e}\). Its
real form pairs with \(a_{e}\), since the Gaussian drift has purely
imaginary expectation against that Hermitian matrix. Hence it is bounded
by \(\epsilon  h_{e}+C_{\epsilon} \kappa ^{2}r_{e}^{-2}\).

For the internal part of \(\langle s,A p_{s}\rangle\), define
\(w_{a}= \sum _{e \text{ incident } a}r_{e}^{-3}\). Factoring
\(M^{-1}=R^{-1}(MR^{-1})^{-1}\) gives, for its coefficient
\(\beta _{s,a}\),

% DISPLAY h.19
% SOURCE ||β_s,av||²/w_a
% SOURCE  ≤CW2 Σ_(e incident a)r_e||Y_ev||²
% SOURCE  ≤C r_*^−3H_exc[v]+CW2[v].                              (6.3)
\begin{gather*}
\begin{aligned}
& \Vert  \beta _{s,a} v\Vert ^{2}/w_{a} \\
& \le CW_{2} \sum _{e \text{ incident } a}r_{e}\Vert Y_{e} v\Vert ^{2} \\
& \le C r_{*}^{-3}H_{\mathrm{exc}}[v]+CW_{2}[v].
\end{aligned}\tag{6.3}
\end{gather*}

Weighted Young loses \(\gamma  \sum _{a} w_{a} t_{a}\), plus
\(\gamma ^{-1}\) times the right side. This is a vanishing loss of the
complete colored internal forms because

% DISPLAY h.20
% SOURCE t_a≤H_a+c_a alpha_a,
% SOURCE Σ_a w_a alpha_a≤CκW2.                                   (6.4)
\begin{gather*}
\begin{aligned}
& t_{a}\le H_{a}+c_{a} \alpha _{a}, \\
& \sum _{a} w_{a} \alpha _{a}\le C \kappa  W_{2}.
\end{aligned}\tag{6.4}
\end{gather*}

The weights depend only on centers, so no internal derivative of a
weight occurs. Center spin cross terms use ordinary center Young. The
direct linear-Y Yukawa terms obey
\(\epsilon  H_{\mathrm{exc}}+C_{\epsilon} W_{2}\) by (3.2) with one
\(Y\) and a bounded Hermitian Clifford matrix.

All spin estimates remain on the whole colored internal module. They
make no physical-sector assumption.

\subsection{7. Complete differential-connection curvature
payment}\label{h-complete-differential-connection-curvature-payment}

Write the internal component of \(\beta  =A^{*}a\) as
\(\beta _{j}=b_{j}(Y;Z,A) \cdot  p_{F}\). This is real and fast-bosonic,
hence commutes with internal coordinates and internal Clifford matrices.
Its symmetric version is

% DISPLAY h.21
% SOURCE βhat_j=β_j−(i/2)div_F b_j.
\begin{gather*}
\begin{aligned}
& \widehat{\beta}_{j}= \beta _{j}-(i/2)\operatorname{div} _{F} b_{j}.
\end{aligned}\notag
\end{gather*}

The averaged sixteen internal charges can therefore be used with
\(p_{A}+ \widehat{\beta}\). There is no additional internal-potential
anticommutator; the only algebraic correction is a finite Clifford
contraction of the fast-bosonic curvature

% DISPLAY h.22
% SOURCE F_jk=i[P_j,P_k]=p_sym(∂_j b_k−∂_k b_j+[b_j,b_k]),
% SOURCE P_j=p_Aj+βhat_j.                                         (7.1)
\begin{gather*}
\begin{aligned}
& F_{jk}=i[P_{j},P_{k}]=p_{\mathrm{sym}}( \partial  _{j} b_{k}- \partial  _{k} b_{j}+[b_{j},b_{k}]), \\
& P_{j}=p_{Aj}+ \widehat{\beta}_{j}.
\end{aligned}\tag{7.1}
\end{gather*}

The nonsymmetric-to-symmetric kinetic conversion also has the exact
scalar correction

% DISPLAY h.23
% SOURCE (1/2)X_j(div_F b_j)+(1/4)(div_F b_j)²,
% SOURCE X_j=∂_Aj+b_j·∂Y.                                        (7.2)
\begin{gather*}
\begin{aligned}
& (1/2)X_{j}(\operatorname{div} _{F} b_{j})+(1/4)(\operatorname{div} _{F} b_{j})^{2}, \\
& X_{j}= \partial  _{Aj}+b_{j} \cdot  \partial  Y.
\end{aligned}\tag{7.2}
\end{gather*}

The companion proof Appendix \(G\) establishes the actual rootwise bound

% DISPLAY h.24
% SOURCE |<v,F_jkv>|+|<(7.2)v,v>|
% SOURCE  ≤C(σ²+r_*^−3)H_exc[v]+C W2[v].                          (7.3)
\begin{gather*}
\begin{aligned}
& \vert \langle v,F_{jk} v\rangle \vert +\vert \langle (7.2)v,v\rangle \vert \\
& \le C( \sigma ^{2}+r_{*}^{-3})H_{\mathrm{exc}}[v]+C W_{2}[v].
\end{aligned}\tag{7.3}
\end{gather*}

The proof does not infer (7.3) from a norm bound alone. Internal
differentiation of the explicit rootwise source decomposition preserves
the same-edge and triangle source classes and adds one factor at most
\(C/r_{*}\). Thus
\(\Vert r_{*} \partial  A \widehat{\beta} v\Vert ^{2}\le C[( \sigma ^{2}+r_{*}^{-3})H_{\mathrm{exc}}[v]+W_{2}[v]]\).
Weighted Cauchy with \(\Vert r_{*}^{-1}v\Vert ^{2}\le W_{2}[v]\) pays
the derivative curvature. Symmetry bounds each commutator-curvature
expectation by
\(2\Vert  \widehat{\beta}_{j} v\Vert \Vert  \widehat{\beta}_{k} v\Vert\).
The scalar divergence correction is bounded by \(CW_{2}\) from the
actual first and second coefficient derivatives. Averaged-charge
Cauchy--Schwarz therefore yields

% DISPLAY h.25
% SOURCE H_I(p+β)≥(1−ε)H_I−C_ε(σ²+r_*^−3)H_exc−C_εW2,           (7.4)
\begin{gather*}
\begin{aligned}
& H_{I}(p+ \beta  )\ge (1- \epsilon  )H_{I}-C_{\epsilon} ( \sigma ^{2}+r_{*}^{-3})H_{\mathrm{exc}}-C_{\epsilon} W_{2},
\end{aligned}\tag{7.4}
\end{gather*}

retaining the entire interacting \(H_{I}\), not a fixed fraction of
\(t_{A}\) alone.

\subsection{8. Density, vacuum defect, and the final
allocation}\label{h-density-vacuum-defect-and-the-final-allocation}

The density term is bounded by \(CW_{2}\) uniformly on the pairwise
tube. Indeed \(M R^{-1}\) and its inverse are uniformly bounded; each
slow or fast derivative of their explicit normalized entries is
\(\le C/r_{*}\), and a second derivative is \(\le C/r_{*}^{2}\). The
principal co-metric and the required normal-covariant coefficient
derivatives obey the same bounds. In
\(d^{-1}\operatorname{div} (a \partial  d)\), these give
\(C/r_{*}^{2}\le CW_{2}\). This argument uses actual normalized matrices
and does not compactify a set of gap ratios.

The frozen \(E_{0,\mathrm{total}}\) is paid separately by the incident
commutator-square result Appendix \(B\). Directly,
\(\vert E_{0,\mathrm{total}}\vert \le C \sum _{a} w_{a}V_{a}\le Cr_{*}^{-3}H_{I}+C \kappa  W_{2}\),
with fixed endpoint multiplicities absorbed in \(C\). This uses
\(V_{a}\le H_{a}+c_{a} \alpha _{a}\) and the incident bound
\(w_{a} \alpha _{a}\le C \kappa  W_{2}\). The first term is a vanishing
complete-internal-form loss; the second remains in the allowed
inverse-square remainder. Thus this local reserve does not require a
separate center-Hardy argument. \(E_{0,\mathrm{total}}\) is never
included in \(H_{\mathrm{exc}}\) or a virtual excitation denominator.

Equations (7.3), (3.1), (5.5), (6.1)--(6.4), the density bound, and the
exact Hamiltonian potential/Yukawa inventory give, for every desired
\(\delta  >0\), constants
\(\kappa _{\delta} , \sigma _{\delta} , R_{\delta}\) and \(C_{\delta}\),
independent of all gap ratios,

% DISPLAY h.26
% SOURCE H_full[v]≥(1−δ)(H_I[v]+T_cent[v]+H_exc[v])−C_δW2[v],     (8.1)
\begin{gather*}
\begin{aligned}
& H_{\mathrm{full}}[v]\ge (1- \delta  )(H_{I}[v]+T_{\mathrm{cent}}[v]+H_{\mathrm{exc}}[v])-C_{\delta} W_{2}[v],
\end{aligned}\tag{8.1}
\end{gather*}

on the tube with
\(\kappa  \le  \kappa _{\delta} , \sigma  \le  \sigma _{\delta}\) and
\(r_{*}\ge R_{\delta}\). Uniformity for nested smaller fast tubes
follows by proving the estimate at the fixed positive width
\(\sigma _{\delta}\) and restricting its core; \(C_{\delta}\) is not
obtained by substituting a varying \(\sigma\) into \(\sigma ^{-1}\). A
nonnegative fraction of \(\Vert A p_{s} v\Vert ^{2}\) and \(V_{4}\) may
also be retained by allocating their positive terms explicitly. This is
not a reserve of the original full Gauss square.

For \(v\) fiberwise orthogonal to the adapted product vacuum,
\(H_{\mathrm{exc}}[v]\ge c r_{*}\Vert v\Vert ^{2}\). Since
\(W_{2}\le (\#\text{edges})r_{*}^{-2}\), increasing \(R_{\delta}\)
absorbs \(C_{\delta} W_{2}\) into \(\delta  H_{\mathrm{exc}}\). The
resulting complementary reserve is near-unit in the full colored
\(H_{I}\), center kinetic, and shifted excitation form. Starting (8.1)
with \(\delta  /2\) gives a final displayed coefficient \(1- \delta\)
after this absorption.

For a projection-stable tube realization first fix a positive cutoff
width \(\sigma _{0}\le  \sigma _{\delta}\), then choose the exterior
radius so \(\sigma _{0}^{2}r_{*}^{3}\) is large. Choose a smooth
\(\chi _{0}\) with support strictly inside the validity tube and equal
to one on a smaller tube, using the smooth harmonic minimum of pair
separations. Put \(U_{\sigma} = \chi _{0}U/\Vert  \chi _{0}U\Vert\). The
fixed-rank Gaussian tail gives
\(\Vert U_{\sigma} -U\Vert \le  \tau  \le C \exp (-c \sigma _{0}^{2}r_{*}^{3})\),
uniformly in all gap ratios. For \(v\) fiberwise orthogonal to
\(U_{\sigma} , \vert \langle U,v\rangle \vert \le  \tau  \Vert v\Vert\),
hence \(H_{\mathrm{exc}}[v]\ge cr_{*}(1- \tau ^{2})\Vert v\Vert ^{2}\).
Both \(U_{\sigma} f\) and \(v=u-U_{\sigma} f\) have Dirichlet form
support. Thus the same complementary reserve holds for this cutoff
projection once the exterior radius is enlarged. Gaussian derivative
tails needed in a subsequent mixed-source calculation are separate
estimates.

The reserve retains the actual normal-bundle center derivative, with no
replacement by a bounded Berry coefficient. On the low vacuum line the
scalar pure-internal and mixed center/internal Berry comparison is the
rootwise full-charge lemma in source Section \(4.2\) in Appendix \(F\):
its costs are \(C_{N}( \kappa  + \epsilon ^{-1} \kappa ^{4})W_{2}\).
This does not require a flat determinant line.

No claim about mixed low/complement source estimates follows merely from
(8.1), especially if those estimates spend the full original Gauss
square instead of the actual reserve left by (6.1). The full
inverse-square Feshbach conclusion remains separate from this
complementary theorem.

\clearpage
\section{Local Schur comparison on the form
domain}\label{i-local-schur-comparison-on-the-form-domain}

Technical derivation for review. Equation and subsection numbers in this
appendix are local to the source derivation. Historical local labels are
retained, including numbering gaps or nonsequential labels. The
accompanying source notes retain the original expressions for
comparison.

\subsection{1. Statement and inputs}\label{i-statement-and-inputs}

Fix a partition of a finite \(N\). Put

% DISPLAY i.1
% SOURCE r_* = min_e r_e,    W_p = sum_e r_e^(-p),
% SOURCE H_s = H_I + T_C,    H_exc = sum_e(H_fast,e-E_0,e).
\begin{gather*}
\begin{aligned}
& r_{*} = \min _{e} r_{e}, \quad  W_{p} = \sum _{e} r_{e}^{-p}, \\
& H_{s} = H_{I} + T_{C}, \quad  H_{\mathrm{exc}} = \sum _{e}(H_{\mathrm{fast},e}-E_{0,e}).
\end{aligned}\notag
\end{gather*}

The internal tube is \(\alpha _{a}+\alpha _{b} \le  \kappa  r_{ab}\).
The fast tube has \(\vert Y\vert  < \sigma  r_{*}\). There is no bound
on ratios of unrelated \(r_{e}\) and no internal spectral-gap
assumption. \(H_{I}\) is the complete interacting colored internal form.
\(T_{C}\) on the full fiber retains the intrinsic normal connection. On
the slow bundle, \(H_{s}\) denotes the center-plus-internal model
identified by internal radial transport with the \(A=0\) vacuum line.
Its unperturbed internal forms are the colored ones; no lower-rank
physical-sector Hardy theorem is used.

For every \(\epsilon >0\), there are positive \(\kappa , \sigma , R\)
and a smooth even normalized cutoff vacuum \(U_{\sigma}\) supported
strictly inside the fast tube such that, for \(r_{*}\ge R\),

% DISPLAY i.2
% SOURCE (1-epsilon) H_s[f] - epsilon W_2[f]
% SOURCE    <= F_sigma[f]
% SOURCE    <= (1+epsilon) H_s[f] + epsilon W_2[f],              (1.1)
\begin{gather*}
\begin{aligned}
& (1-\epsilon ) H_{s}[f] - \epsilon  W_{2}[f] \\
& \le  F_{\sigma}[f] \\
& \le  (1+\epsilon ) H_{s}[f] + \epsilon  W_{2}[f],
\end{aligned}\tag{1.1}
\end{gather*}

where the actual zero-energy variational Hamiltonian Schur form is

% DISPLAY i.3
% SOURCE F_sigma[f] = inf_{v in Q(H), U_sigma^*v=0}
% SOURCE                       H[U_sigma f+v].                 (1.2)
\begin{gather*}
\begin{aligned}
& F_{\sigma}[f] = \inf _{v \in  Q(H), U_{\sigma}^{*}v=0} \\
& H[U_{\sigma} f+v].
\end{aligned}\tag{1.2}
\end{gather*}

Here Q(H) denotes the CLOSED QUADRATIC-FORM DOMAIN of the Hamiltonian,
not its operator domain D(H). Orthogonality is fiberwise over the slow
variables. For every \(f\) in \(Q(H_{s})\), the infimum is attained by a
unique \(v_{f}\) in Q(H) with \(U_{\sigma}^{*}v_{f}=0\). The form is
closed on the slow model\textquotesingle s form domain, interpreted with
Dirichlet support on the selected branch. The constants depend on the
fixed partition and \(\epsilon\), never on \(r_{\max}/r_{*}\).

Inputs proved in the companion notes are:

\begin{enumerate}
\def\labelenumi{\arabic{enumi}.}
\item
  The exact intrinsic kinetic identity and normalized matrices from
  Appendix A.
\item
  The rootwise all-vector connection estimates in Appendix \(E\) and
  differential-curvature identities in Appendix \(G\).
\item
  The actual H1/H2 inventory, finite source graph grading, and operator
  identity

  % DISPLAY i.4
  % SOURCE ||W-S^*H_exc^(-1)S|| <= C_N kappa W_2,              (1.3)
  \begin{gather*}
  \begin{aligned}
  & \Vert W-S^{*}H_{\mathrm{exc}}^{-1}S\Vert  \le  C_{N} \kappa  W_{2},
  \end{aligned}\tag{1.3}
  \end{gather*}

  from Appendix \(D\) . Here \(S=H1 U\) is the actual source, not a
  transported replacement.
\item
  The graph-local complete-slow-form bound for
  \(Z=-H_{\mathrm{exc}}^{-1}S\):

  % DISPLAY i.5
  % SOURCE H_s[Zf] <= C R^(-3)H_s[f]+C kappa W_2[f]+C W_5[f], (1.4)
  \begin{gather*}
  \begin{aligned}
  & H_{s}[Zf] \le  C R^{-3}H_{s}[f]+C \kappa  W_{2}[f]+C W_{5}[f],
  \end{aligned}\tag{1.4}
  \end{gather*}

  including spectator-vacuum derivatives, from Appendix \(C\) and the
  actual-source continuation.
\item
  The colored complementary reserve. A fixed positive part of the actual
  square \(J[v]=\Vert A p_{s} v\Vert ^{2}\) can be retained:

  % DISPLAY i.6
  % SOURCE H[v] >= (1-delta)(H_s[v]+H_exc[v])+c_0 J[v]         (1.5)
  \begin{gather*}
  \begin{aligned}
  & H[v] \ge  (1-\delta )(H_{s}[v]+H_{\mathrm{exc}}[v])+c_{0} J[v]
  \end{aligned}\tag{1.5}
  \end{gather*}

  on the cutoff complement, after choosing the parameters and exterior
  radius. For example \(c_{0}=1/2\) is allowed. This is NOT another copy
  of the original full Gauss square.
\end{enumerate}

This note supplies the mixed and diagonal-remainder estimates needed to
assemble these inputs. It never inverts \(H_{I}\) or expands a resolvent
in internal momenta.

\subsection{2. Exact error-column factorization, including
spin}\label{i-exact-error-column-factorization-including-spin}

Use the common positive-fast/positive-slow convention

% DISPLAY i.7
% SOURCE H_kin = ||p_s u||^2+||p_F u||^2+||a u+A p_s u+s u||^2,
% SOURCE p_s=(g_C^(-1/2)p_C,p_I),
% SOURCE a=W^(1/2)M^(-*) (C+N_Y)^*p_F,
% SOURCE s=-W^(1/2)M^(-*)S_f.
\begin{gather*}
\begin{aligned}
& H_{\mathrm{kin}} = \Vert p_{s} u\Vert ^{2}+\Vert p_{F} u\Vert ^{2}+\Vert a u+A p_{s} u+s u\Vert ^{2}, \\
& p_{s}=(g_{C}^{-1/2}p_{C},p_{I}), \\
& a=W^{1/2}M^{-*} (C+N_{Y})^{*}p_{F}, \\
& s=-W^{1/2}M^{-*}S_{f}.
\end{aligned}\notag
\end{gather*}

The scalar density term is added separately. Set

% DISPLAY i.8
% SOURCE d_0=F^(-1)C^*p_F,  P=L-J_2,  Z_m=F^(-1)P^*,
% SOURCE T=(I+Z_m)^(-1),  n=F^(-1)N_Y^*p_F,
% SOURCE X=n-Z_m d_0,  s_0=-F^(-1)S_f,
% SOURCE q=X+s_0,  X_1=n-F^(-1)L^*d_0,  X_2=F^(-1)J_2^*d_0.
\begin{gather*}
\begin{aligned}
& d_{0}=F^{-1}C^{*}p_{F}, \quad  P=L-J_{2}, \quad  Z_{m}=F^{-1}P^{*}, \\
& T=(I+Z_{m})^{-1}, \quad  n=F^{-1}N_{Y}^{*}p_{F}, \\
& X=n-Z_{m} d_{0}, \quad  s_{0}=-F^{-1}S_{f}, \\
& q=X+s_{0}, \quad  X_{1}=n-F^{-1}L^{*}d_{0}, \quad  X_{2}=F^{-1}J_{2}^{*}d_{0}.
\end{aligned}\notag
\end{gather*}

The symbol \(Z_{m}\) is an inverse-matrix error, distinct from the
dressing \(Z\). Exactly,

% DISPLAY i.9
% SOURCE a+s = W^(1/2)(d_0+Tq),
\begin{gather*}
\begin{aligned}
& a+s = W^{1/2}(d_{0}+Tq),
\end{aligned}\notag
\end{gather*}

and, as ordered first-order forms,

% DISPLAY i.10
% SOURCE ||(a+s)u||^2-||d_0u||^2
% SOURCE   =2 Re<d_0u,(X_1+s_0)u> + 2 Re<d_0u,X_2u>
% SOURCE    -2 Re<Z_m^*d_0u,Tq u> +||Tq u||^2
% SOURCE    +||K^*(d_0+Tq)u||^2.                              (2.1)
\begin{gather*}
\begin{aligned}
& \Vert (a+s)u\Vert ^{2}-\Vert d_{0} u\Vert ^{2} \\
& =2 \operatorname{Re} \langle d_{0} u,(X_{1}+s_{0})u\rangle  + 2 \operatorname{Re} \langle d_{0} u,X_{2} u\rangle \\
& -2 \operatorname{Re} \langle Z_{m}^{*}d_{0} u,Tq u\rangle  +\Vert Tq u\Vert ^{2} \\
& +\Vert K^{*}(d_{0}+Tq)u\Vert ^{2}.
\end{aligned}\tag{2.1}
\end{gather*}

The first line\textquotesingle s first term is exactly the
fast-orbital/spin part of H1. No derivative has been moved through an
inverse matrix. The remaining apparently dangerous \(d_{0}-X_{2}\)
product factors exactly:

% DISPLAY i.11
% SOURCE <d_0u,X_2u>
% SOURCE    =<C_Y F^(-1)d_0u,C_Y R^(-1)d_0u>.                 (2.2)
\begin{gather*}
\begin{aligned}
& \langle d_{0} u,X_{2} u\rangle \\
& =\langle C_{Y} F^{-1}d_{0} u,C_{Y} R^{-1}d_{0} u\rangle .
\end{aligned}\tag{2.2}
\end{gather*}

Every factor on the right of (2.1)-(2.2), after removing the displayed
H1 term, belongs to the finite error-column list

% DISPLAY i.12
% SOURCE X_1, X_2, Z_m^*d_0, K^*d_0,
% SOURCE C_Y F^(-1)d_0, C_Y R^(-1)d_0, s_0,
\begin{gather*}
\begin{aligned}
& X_{1}, X_{2}, Z_{m}^{*}d_{0}, K^{*}d_{0}, \\
& C_{Y} F^{-1}d_{0}, C_{Y} R^{-1}d_{0}, s_{0},
\end{aligned}\notag
\end{gather*}

followed by uniformly bounded left multiplication matrices. Their
all-vector estimates are

% DISPLAY i.13
% SOURCE ||E u||^2 <= C[(lambda H_exc)[u]+W_2[u]],
% SOURCE lambda=sigma^2+r_*^(-3).                              (2.3)
\begin{gather*}
\begin{aligned}
& \Vert E u\Vert ^{2} \le  C[(\lambda  H_{\mathrm{exc}})[u]+W_{2}[u]], \\
& \lambda =\sigma ^{2}+r_{*}^{-3}.
\end{aligned}\tag{2.3}
\end{gather*}

For the spin column the stronger bound
\(\Vert s_{0}u\Vert ^{2}\le CW_{2}[u]\) holds. On the uncut vacuum, and
on each fixed-degree Gaussian/Fock coefficient vector normalized in its
own roots,

% DISPLAY i.14
% SOURCE ||E U f||^2 <= C W_2[f].                              (2.4)
\begin{gather*}
\begin{aligned}
& \Vert E U f\Vert ^{2} \le  C W_{2}[f].
\end{aligned}\tag{2.4}
\end{gather*}

For (2.4), a same-edge \(Y_{e} p_{e}/r_{e}\) costs \(C/r_{e}\); a
triangle \(Y_{f} p_{e}/r_{g}\) costs \(C \sqrt{r_{e}/r_{f}}/r_{g}\),
whose square is bounded by \(C W_{2,t}\) by \(r_{e}\le r_{f}+r_{g}\).
Root-preserving pair matrices have uniformly bounded normalized
coefficients. All inverse matrices in the exact expressions remain
bounded LEFT factors. The same proof therefore works on a fixed finite
excited degree, with a degree-dependent constant, and does not introduce
\(r_{\max}\).

For a complementary \(v, W_{2}[v]\le C R^{-3}H_{\mathrm{exc}}[v]\).
Polarization of each product in (2.1)-(2.2) consequently gives, for any
\(\eta >0\),

% DISPLAY i.15
% SOURCE |R_fast[Uf,v]| <= eta H_exc[v]
% SOURCE                      +C_eta(sigma^2+R^(-3)) W_2[f].   (2.5)
\begin{gather*}
\begin{aligned}
& \vert R_{\mathrm{fast}}[Uf,v]\vert  \le  \eta  H_{\mathrm{exc}}[v] \\
& +C_{\eta}(\sigma ^{2}+R^{-3}) W_{2}[f].
\end{aligned}\tag{2.5}
\end{gather*}

This includes spin-fast and spin-square remainders. It is the central
reason a finite inverse-square constant is not left unpaid: no Taylor
remainder is ever paired directly with an unassigned \(d_{0}\) vacuum
norm.

For diagonal matching, retain the quadratic terms in (2.1), replace
\(T\) by I, \(q\) by \(X_{1}+s_{0}, Z_{m}\) by \(F^{-1}L^{*}\), and
retain (2.2) and \(\Vert K^{*}d_{0}\Vert ^{2}\). These are precisely the
corresponding H2 terms. Each discarded product has one additional
uniformly \(O(\sigma )\) factor and two low error columns. Hence

% DISPLAY i.16
% SOURCE |R_fast[Uf]-R_fast,H2[Uf]| <= C sigma W_2[f].          (2.6)
\begin{gather*}
\begin{aligned}
& \vert R_{\mathrm{fast}}[Uf]-R_{\mathrm{fast},H2}[Uf]\vert  \le  C \sigma  W_{2}[f].
\end{aligned}\tag{2.6}
\end{gather*}

One may obtain an additional inverse-radius improvement from Gaussian
moments, but it is unnecessary. The order of the two operations matters:
first perform the exact error-column factorization, then bound the
inverse difference. Doing them in reverse would expose the forbidden
remote vacuum mass.

\subsection{3. Other nondifferential fast
terms}\label{i-other-nondifferential-fast-terms}

The actual \(V_{3}\) and linear-Y Yukawa are exactly the remaining H1
terms; there is no omitted higher polynomial potential. They are
boson-odd, so their diagonal expectation on \(U\) or an even cutoff
\(U_{\sigma}\) is zero.

The quartic potential is a sum of squares of quadratic fast-coordinate
contractions. In addition to the distinct-root Q/P estimate already
proved, the repeated-root estimate is

% DISPLAY i.17
% SOURCE ||chi Y_e^2 u||^2 <= C sigma^2 h_e[u]+C r_e^(-2)||u||^2. (3.1)
\begin{gather*}
\begin{aligned}
& \Vert \chi  Y_{e}^{2} u\Vert ^{2} \le  C \sigma ^{2} h_{e}[u]+C r_{e}^{-2}\Vert u\Vert ^{2}.
\end{aligned}\tag{3.1}
\end{gather*}

Split the input into \(Q_{e}\) and \(P_{e}\), bound one \(Y_{e}\) by
\(\sigma  r_{*}\) on \(Q_{e}\), use
\(\Vert Y_{e} Q_{e} u\Vert ^{2}\le Cr_{e}^{-2}h_{e}[u]\), and use the
exact fourth Gaussian moment on \(P_{e}\). This gives

% DISPLAY i.18
% SOURCE 0<=V4[u]<=C sigma^2 H_exc[u]+C W_2[u],
% SOURCE V4[Uf]<=C W_2[f].                                    (3.2)
\begin{gather*}
\begin{aligned}
& 0\le V_{4}[u]\le C \sigma ^{2} H_{\mathrm{exc}}[u]+C W_{2}[u], \\
& V_{4}[Uf]\le C W_{2}[f].
\end{aligned}\tag{3.2}
\end{gather*}

Thus its low/complement cross obeys the same bound as (2.5). Its full
diagonal expectation belongs to \(W\) and is retained there.

The density multiplier satisfies \(\vert V_{d}\vert \le CW_{2}\). Its
mixed term is paid by the complementary L2 gap:

% DISPLAY i.19
% SOURCE |<Uf,V_d v>|<=eta H_exc[v]+C_eta R^(-3)W_2[f].         (3.3)
\begin{gather*}
\begin{aligned}
& \vert \langle Uf,V_{d} v\rangle \vert \le \eta  H_{\mathrm{exc}}[v]+C_{\eta} R^{-3}W_{2}[f].
\end{aligned}\tag{3.3}
\end{gather*}

The leading coefficient \(V_{d,2}\) is the one specified by the exact H2
inventory, including the slow density, fast Hessian, and orbital
ordering terms. Differentiating the normalized matrices in the exact
density formula and subtracting their zero-fast-coordinate jets gives

% DISPLAY i.20
% SOURCE |V_d-V_d,2|<=C sigma W_2                              (3.4)
\begin{gather*}
\begin{aligned}
& \vert V_{d}-V_{d,2}\vert \le C \sigma  W_{2}
\end{aligned}\tag{3.4}
\end{gather*}

on the tube. Each normalized inverse and its required derivatives are
bounded in its own root scale. More explicitly, \(V_{d}\) uses two
coefficient derivatives, and one additional fast derivative of that
finite rational expression is bounded by \(C/r_{*}^{3}\); integrating it
from \(Y=0\) gives \(C\vert Y\vert /r_{*}^{3}\le C \sigma  W_{2}\). The
leading zero-fast-coordinate coefficient includes the quadratic
determinant jet before fast differentiation. This is a
density-coefficient estimate, not a bound on \(d_{0}^{2}\). Equations
(3.3)-(3.4) settle both density uses.

The frozen \(E_{0}\) is a separate scalar multiplier. It has zero U-v
cross because it commutes with the vacuum projection. Its incident
payment is

% DISPLAY i.21
% SOURCE |E_0[u]|<=C R^(-3)H_I[u]+C kappa W_2[u].              (3.5)
\begin{gather*}
\begin{aligned}
& \vert E_{0}[u]\vert \le C R^{-3}H_{I}[u]+C \kappa  W_{2}[u].
\end{aligned}\tag{3.5}
\end{gather*}

It is never included in \(H_{\mathrm{exc}}\) or in a virtual
denominator.

\subsection{\texorpdfstring{4. Scalar fast-slow orbit cross: complete
charges instead of bare
\(t_{I}\)}{4. Scalar fast-slow orbit cross: complete charges instead of bare t\_\{I\}}}\label{i-scalar-fast-slow-orbit-cross-complete-charges-instead-of-bare-t_i}

Write \(\beta =A^{*}a\), with internal components
\(\beta _{j}=b_{j}(Y;s) \cdot  p_{F}\). The \(b_{j}\) are real, scalar
in internal Clifford factors, and commute with all internal
potential-charge coefficients. Let

% DISPLAY i.22
% SOURCE d_j=div_F b_j,    beta_hat,j=beta_j-(i/2)d_j,
% SOURCE C_alpha(beta_hat)=sum_j c_(alpha,j) beta_hat,j.
\begin{gather*}
\begin{aligned}
& d_{j}=\operatorname{div} _{F} b_{j}, \quad  \widehat{\beta}_{j}=\beta _{j}-(i/2)d_{j}, \\
& C_{\alpha}(\widehat{\beta})=\sum _{j} c_{\alpha ,j} \widehat{\beta}_{j}.
\end{aligned}\notag
\end{gather*}

The differentiated connection lemma and its explicit rootwise
decomposition imply

% DISPLAY i.23
% SOURCE sum_j ||beta_hat,j v||^2 <= C(lambda H_exc[v]+W_2[v]),
% SOURCE sum_j ||beta_hat,j U f||^2 <= C(kappa^2+sigma^2)W_2[f], (4.1)
\begin{gather*}
\begin{aligned}
& \sum _{j} \Vert \widehat{\beta}_{j} v\Vert ^{2} \le  C(\lambda  H_{\mathrm{exc}}[v]+W_{2}[v]), \\
& \sum _{j} \Vert \widehat{\beta}_{j} U f\Vert ^{2} \le  C(\kappa ^{2}+\sigma ^{2})W_{2}[f],
\end{aligned}\tag{4.1}
\end{gather*}

and

% DISPLAY i.24
% SOURCE sum_jk ||(partial_Aj beta_hat,k)U f||^2
% SOURCE                        <= C r_*^(-2)W_2[f].          (4.2)
\begin{gather*}
\begin{aligned}
& \sum _{jk} \Vert (\partial _{Aj} \widehat{\beta}_{k})U f\Vert ^{2} \\
& \le  C r_{*}^{-2}W_{2}[f].
\end{aligned}\tag{4.2}
\end{gather*}

The SMALL low-leg divergence estimate in (4.1) is essential. Directly,
\(b=A^{*}B\) with
\(A=O(\sigma ), B=O(\kappa +\sigma ), \partial _{Y} A=O(r_{*}^{-1})\)
and \(\partial _{Y} B=O(r_{*}^{-1})\). Hence

% DISPLAY i.25
% SOURCE |div_F b|<=C(kappa+sigma)/r_*.
\begin{gather*}
\begin{aligned}
& \vert \operatorname{div} _{F} b\vert \le C(\kappa +\sigma )/r_{*}.
\end{aligned}\notag
\end{gather*}

Together with the stronger rootwise estimate for \(A^{*}d_{0}\) and the
\(O(\sigma )\) factor multiplying \(a-d_{0}\), this proves the second
line of (4.1). It is not inferred from a coarse \(O(1/r_{*})\)
divergence bound.

\subsubsection{4.1 Exact linear charge
identity}\label{i-exact-linear-charge-identity}

Polarize the following identity, initially on the smooth core:

% DISPLAY i.26
% SOURCE 2 Re sum_j <p_j u,beta_j u>
% SOURCE   =2 Re (1/16)sum_alpha <Q_alpha u,C_alpha(beta_hat)u>
% SOURCE      -R_der[u]+(1/2)sum_j <u,(partial_Aj d_j)u>.      (4.3)
\begin{gather*}
\begin{aligned}
& 2 \operatorname{Re}  \sum _{j} \langle p_{j} u,\beta _{j} u\rangle \\
& =2 \operatorname{Re}  (1/16)\sum _{\alpha} \langle Q_{\alpha} u,C_{\alpha}(\widehat{\beta})u\rangle \\
& -R_{\mathrm{der}}[u]+(1/2)\sum _{j} \langle u,(\partial _{Aj} d_{j})u\rangle .
\end{aligned}\tag{4.3}
\end{gather*}

Here \(R_{\mathrm{der}}\) is the fixed bounded-Clifford contraction of
\(\partial _{j} \widehat{\beta}_{k}-\partial _{k} \widehat{\beta}_{j}\).
To verify it, expand the exact shifted averaged-charge identity and
subtract the beta-square. The beta-beta curvature in
\(C_{\alpha}(\widehat{\beta})^{2}\) cancels the beta-beta curvature in
the full shifted charge identity. Only the derivative curvature remains.
The raw-versus-symmetric cross contributes \(+(1/2)\partial _{A} d\),
with this sign. No anticommutator with the internal potential survives
the spinor average.

The first term in the polarized identity contains

% DISPLAY i.27
% SOURCE <Q(Uf),C(beta_hat)v>  and  <C(beta_hat)Uf,Qv>.
\begin{gather*}
\begin{aligned}
& \langle Q(Uf),C(\widehat{\beta})v\rangle  \quad  \text{ and } \quad  \langle C(\widehat{\beta})Uf,Qv\rangle .
\end{aligned}\notag
\end{gather*}

The first is bounded using (4.1), the complementary gap, and

% DISPLAY i.28
% SOURCE (1/16)sum ||Q_alpha(Uf)||^2 <= C(H_s[f]+W_2[f]).
\begin{gather*}
\begin{aligned}
& (1/16)\sum  \Vert Q_{\alpha}(Uf)\Vert ^{2} \le  C(H_{s}[f]+W_{2}[f]).
\end{aligned}\notag
\end{gather*}

The second uses the full \(H_{I}[v]\) reserve and the small low-leg
estimate in (4.1). No constant fraction of bare \(t_{I}\) is spent.
Since each derivative-curvature operator is symmetric, it can be moved
onto \(U\) in a mixed pairing. Equation (4.2) and the complementary gap
then give a \(C_{\eta} R^{-3}W_{2}\) cost. The scalar term has the same
payment.

Consequently the complete internal scalar-orbit mixed form obeys

% DISPLAY i.29
% SOURCE |R_orb,I[Uf,v]| <= eta(H_I[v]+H_exc[v])
% SOURCE    +C_eta lambda H_s[f]
% SOURCE    +C_eta(kappa^2+sigma^2+R^(-3))W_2[f].              (4.4)
\begin{gather*}
\begin{aligned}
& \vert R_{\mathrm{orb},I}[Uf,v]\vert  \le  \eta (H_{I}[v]+H_{\mathrm{exc}}[v]) \\
& +C_{\eta} \lambda  H_{s}[f] \\
& +C_{\eta}(\kappa ^{2}+\sigma ^{2}+R^{-3})W_{2}[f].
\end{aligned}\tag{4.4}
\end{gather*}

This is uniform over unrestricted internal momenta and colored internal
states. The center component obeys the same inequality by ordinary
center-kinetic Cauchy-Schwarz, with its actual normal connection
retained. Normal-frame generators are not estimated by a vacuum Berry
scalar on arbitrary \(v\).

\subsubsection{4.2 Independent incident-row
verification}\label{i-independent-incident-row-verification}

There is a second direct check of the high-momentum issue. With \(U=g\)
times its fermion vacuum and real Gaussian \(g\),

% DISPLAY i.30
% SOURCE (beta_j+beta_j^*)U
% SOURCE     =i(2 b_j dot Omega Y-div_F b_j)U.
\begin{gather*}
\begin{aligned}
& (\beta _{j}+\beta _{j}^{*})U \\
& =i(2 b_{j} \cdot  \Omega  Y-\operatorname{div} _{F} b_{j})U.
\end{aligned}\notag
\end{gather*}

Ground-state integration by parts therefore gives

% DISPLAY i.31
% SOURCE ||H_exc^(-1/2)(beta_j+beta_j^*)U||^2
% SOURCE                           <=C||b_jU||^2.            (4.5)
\begin{gather*}
\begin{aligned}
& \Vert H_{\mathrm{exc}}^{-1/2}(\beta _{j}+\beta _{j}^{*})U\Vert ^{2} \\
& \le C\Vert b_{j}U\Vert ^{2}.
\end{aligned}\tag{4.5}
\end{gather*}

For an internal row belonging to block a, factor the exact coefficient
as

% DISPLAY i.32
% SOURCE b_a=I_Y,a F^(-1)[T^* W M^(-*)(C+N_Y)^*].
\begin{gather*}
\begin{aligned}
& b_{a}=I_{Y,a} F^{-1}[T^{*} W M^{-*}(C+N_{Y})^{*}].
\end{aligned}\notag
\end{gather*}

The bracket is uniformly \(O(\kappa +\sigma )\), while
\(I_{Y,a} F^{-1}\) contains only \(Y_{e}/r_{e}\) on edges incident to a.
Thus

% DISPLAY i.33
% SOURCE ||b_aU||^2<=C(kappa+sigma)^2 w_a,
% SOURCE w_a=sum_(e incident a)r_e^(-3).                       (4.6)
\begin{gather*}
\begin{aligned}
& \Vert b_{a}U\Vert ^{2}\le C(\kappa +\sigma )^{2} w_{a}, \\
& w_{a}=\sum _{e \text{ incident } a}r_{e}^{-3}.
\end{aligned}\tag{4.6}
\end{gather*}

This proves directly that integrating the single slow derivative in the
orbit cross cannot create an unrelated internal-radius loss. It also
accounts for the adjoint fast divergence, rather than dropping it.
Either (4.3) or (4.5)-(4.6) supplies the needed first-derivative
estimate; the proof above uses (4.3).

\subsubsection{4.3 Diagonal orbit
remainder}\label{i-diagonal-orbit-remainder}

Let \(\beta _{1}\) be the first homogeneous slow-orbit coefficient: use
the linear slow column with \(F^{-1}\), the leading \(C_{Y}^{*}+FK\)
center term, and \(d_{0}\). Its raw diagonal compression, including
derivatives of its coefficient at \(A=0\), is retained in \(W\).
Together with the corresponding \(V_{d,2}\) orbital-density term, it
equals the \(V_{\mathrm{orb}}\) in the density-shifted inventory. In
particular the inventory term -2 \(\sum _{\mu} \rho _{0,\mu} j_{\mu}\)
belongs to \(V_{d,2}\) in the present unshifted-kinetic convention; it
is counted exactly once, not inserted a second time in the raw
\(\beta _{1}\) compression. Only the SUM of these contributions is
identified with the H2 scalar. The raw \(\beta _{1}\) contribution is
NOT discarded on the grounds that \(\beta _{1}\) vanishes at \(A=0\).

The exact identities for \(A^{*}d_{0}\) in the rootwise helper express
\(\beta -\beta _{1}\) as a sum with one extra \(O(\sigma )\) factor
multiplying a controlled error column. They give

% DISPLAY i.34
% SOURCE ||(beta-beta_1)U||^2+||(div(b-b_1))U||^2<=C sigma^2 W_2,
% SOURCE ||partial_A(beta-beta_1)U||<=C sigma r_*^(-1) W_2^(1/2).
\begin{gather*}
\begin{aligned}
& \Vert (\beta -\beta _{1})U\Vert ^{2}+\Vert (\operatorname{div} (b-b_{1}))U\Vert ^{2}\le C \sigma ^{2} W_{2}, \\
& \Vert \partial _{A}(\beta -\beta _{1})U\Vert \le C \sigma  r_{*}^{-1} W_{2}^{1/2}.
\end{aligned}\notag
\end{gather*}

Differentiate those identities, not an abstract norm bound: the extra
fast-degree factor remains when \(F\) or \(F^{-1}C^{*}\) is
differentiated. Applying (4.3) to this difference gives, for any
\(\eta >0\),

% DISPLAY i.35
% SOURCE |R_orb,diag[f]|<=eta H_s[f]
% SOURCE                          +C_eta(sigma+sigma^2)W_2[f]. (4.7)
\begin{gather*}
\begin{aligned}
& \vert R_{\mathrm{orb},\mathrm{diag}}[f]\vert \le \eta  H_{s}[f] \\
& +C_{\eta}(\sigma +\sigma ^{2})W_{2}[f].
\end{aligned}\tag{4.7}
\end{gather*}

This retains the exact leading coefficient while paying only its true
higher-order remainder.

\subsection{5. Spin-slow, slow-orbit square, and induced center
metric}\label{i-spin-slow-slow-orbit-square-and-induced-center-metric}

\subsubsection{5.1 Spin-slow cross}\label{i-spin-slow-cross}

Put \(c=A^{*}s\). Its internal row a satisfies the exact incident
coefficient bound

% DISPLAY i.36
% SOURCE |c_a|^2<=C W_2 sum_(e incident a)|Y_e|^2/r_e^2,
% SOURCE ||c_a U||^2<=C w_a W_2.                              (5.1)
\begin{gather*}
\begin{aligned}
& \vert c_{a}\vert ^{2}\le C W_{2} \sum _{e \text{ incident } a}\vert Y_{e}\vert ^{2}/r_{e}^{2}, \\
& \Vert c_{a} U\Vert ^{2}\le C w_{a} W_{2}.
\end{aligned}\tag{5.1}
\end{gather*}

For the mixed form, integrate its single slow derivative onto the smooth
coefficient cU \(f\). This is legitimate for the complete matrix \(c\),
with no commutation claim about internal Clifford matrices. The
coefficient of \(\partial _{Aa} f\) has norm squared bounded by
\(C w_{a} W_{2}\). The complementary gap therefore gives the weighted
derivative cost

% DISPLAY i.37
% SOURCE C(W_2/r_*) sum_a w_a t_a[f]
% SOURCE    <=C R^(-3)[R^(-3)H_I[f]+kappa W_2[f]].             (5.2)
\begin{gather*}
\begin{aligned}
& C(W_{2}/r_{*}) \sum _{a} w_{a} t_{a}[f] \\
& \le C R^{-3}[R^{-3}H_{I}[f]+\kappa  W_{2}[f]].
\end{aligned}\tag{5.2}
\end{gather*}

The exact inverse derivative formula, incident left factor
\(I_{Y,a} F^{-1}\), and finite Gaussian derivative moments give

% DISPLAY i.38
% SOURCE sum_a ||partial_Aa(c_aU)||^2
% SOURCE    <=C r_*^(-2)(sum_a w_a)W_2+C(sum_a w_a)W_2^2
% SOURCE    <=C R^(-5)W_2.                                   (5.3)
\begin{gather*}
\begin{aligned}
& \sum _{a} \Vert \partial _{Aa}(c_{a} U)\Vert ^{2} \\
& \le C r_{*}^{-2}(\sum _{a} w_{a})W_{2}+C(\sum _{a} w_{a})W_{2}^{2} \\
& \le C R^{-5}W_{2}.
\end{aligned}\tag{5.3}
\end{gather*}

Spectator-vacuum derivatives are included in the \(W_{2}^{2}\) term.
Center derivatives have the same or better bound and are paid by
\(T_{C}\). After the complementary gap, (5.2)-(5.3) are smaller than the
mixed budget in (4.4).

For the DIAGONAL, (5.1) alone would leave a fixed \(W_{2}\) constant and
is insufficient. The leading coefficient \(c_{1}\) is boson-odd: its
internal part is \(-I_{Y} F^{-2}S_{f}\), with the corresponding linear
center row. \(U\), its slow derivatives, and the even cutoff are
boson-even. Its diagonal cross is exactly zero. Every term of
\(c-c_{1}\) has an additional inverse/metric difference of size
\(O(\sigma )\), so

% DISPLAY i.39
% SOURCE ||(c_a-c_1,a)U||^2<=C sigma^2 w_a W_2.
\begin{gather*}
\begin{aligned}
& \Vert (c_{a}-c_{1,a})U\Vert ^{2}\le C \sigma ^{2} w_{a} W_{2}.
\end{aligned}\notag
\end{gather*}

Weighted Young and \(\sum _{a} w_{a} \alpha _{a}\le C \kappa  W_{2}\)
yield only

% DISPLAY i.40
% SOURCE eta R^(-3)H_s[f]+C_eta(kappa+sigma^2)W_2[f].          (5.4)
\begin{gather*}
\begin{aligned}
& \eta  R^{-3}H_{s}[f]+C_{\eta}(\kappa +\sigma ^{2})W_{2}[f].
\end{aligned}\tag{5.4}
\end{gather*}

This is the specific parity use that prevents an unpaid diagonal
constant.

\subsubsection{5.2 The positive square actually retained by
completion}\label{i-the-positive-square-actually-retained-by-completion}

Let \(J[u]=\Vert A p_{s} u\Vert ^{2}\). Its low-leg estimate is

% DISPLAY i.41
% SOURCE J[Uf]<=C R^(-3)H_s[f]+C kappa W_2[f]+C W_5[f].        (5.5)
\begin{gather*}
\begin{aligned}
& J[Uf]\le C R^{-3}H_{s}[f]+C \kappa  W_{2}[f]+C W_{5}[f].
\end{aligned}\tag{5.5}
\end{gather*}

For derivatives on \(f\), the exact incident factorization gives
\(\sum _{a} w_{a} t_{a}\) and the analogous center coefficient. Convert
only these incident weights using

% DISPLAY i.42
% SOURCE sum_a w_a t_a<=C R^(-3)H_I+C kappa W_2.
\begin{gather*}
\begin{aligned}
& \sum _{a} w_{a} t_{a}\le C R^{-3}H_{I}+C \kappa  W_{2}.
\end{aligned}\notag
\end{gather*}

Derivatives on \(U\) are finite degree-two Gaussian/Fock vectors. Their
rootwise moments give \(W_{5}\), including spectator products
\(W_{3}W_{2}\le C_{N}W_{5}\). The Berry connection adds its already
proved smaller incident contribution. Thus

% DISPLAY i.43
% SOURCE |<A p_sUf,A p_sv>|
% SOURCE    <=eta J[v]+C_eta[R^(-3)H_s[f]+kappa W_2[f]+W_5[f]]. (5.6)
\begin{gather*}
\begin{aligned}
& \vert \langle A p_{s} U f,A p_{s} v\rangle \vert \\
& \le \eta  J[v]+C_{\eta}[R^{-3}H_{s}[f]+\kappa  W_{2}[f]+W_{5}[f]].
\end{aligned}\tag{5.6}
\end{gather*}

Equation (5.6) uses only the true postcompletion reserve \(J\). No
original full Gauss square is used.

\subsubsection{5.3 Center metric}\label{i-center-metric}

The direct center coefficient differs from identity by

% DISPLAY i.44
% SOURCE g_C^(-1)-I=-K^*(I+KK^*)^(-1)K.
\begin{gather*}
\begin{aligned}
& g_{C}^{-1}-I=-K^{*}(I+KK^{*})^{-1}K.
\end{aligned}\notag
\end{gather*}

It is \(O(\sigma ^{2})\) on arbitrary tube vectors. On a Gaussian low
leg its squared coefficient norm is \(O(R^{-6})\); derivatives on \(U\)
have the corresponding \(W_{5}\)-or-smaller moments. Ordinary
center-form Young controls the mixed term by
\(\eta  T_{C}[v]+C_{\eta} R^{-6}T_{C}[f]+C_{\eta}W_{5}[f]\). Its
diagonal has bound \(C R^{-3}T_{C}[f]+CW_{5}[f]\). No internal
derivative is involved.

\subsubsection{5.4 Direct slow kinetic
cross}\label{i-direct-slow-kinetic-cross}

The preceding orbit estimates do not replace the cross from the DIRECT
center and internal kinetic forms. For \(v\) perpendicular to \(U\),
differentiate \(U^{*}v=0\) and integrate the remaining slow derivative
once. In a covariant vacuum-line frame the complementary source is

% DISPLAY i.45
% SOURCE -2 sum_j (1-P)(nabla_j U) nabla_j f
% SOURCE       -(1-P)sum_j nabla_j^2 U f.                    (5.7)
\begin{gather*}
\begin{aligned}
& -2 \sum _{j} (1-P)(\nabla _{j} U) \nabla _{j} f \\
& -(1-P)\sum _{j} \nabla _{j}^{2} U f.
\end{aligned}\tag{5.7}
\end{gather*}

This is a form identity and contains no second derivative of \(f\).
Internal potential and Yukawa terms commute with the vacuum injection
and have zero mixed compression. The same is true of \(E_{0}\).
Normal-bundle derivatives are kept covariantly in (5.7).

An internal derivative in block a differentiates only incident pair
vacua. After removing its scalar Berry part, each term has a fixed
excited sector on that pair, norm \(C/r_{e}\), and gap at least
\(c r_{e}\). Its squared \(H_{\mathrm{exc}}\)-dual norm is therefore at
most \(C w_{a}\). The center derivative satisfies the analogous sum of
pair weights. Second derivatives have active-pair and two-pair finite
excitation terms, of norm at most \(C/r_{e}^{2}\) or
\(C/(r_{e} r_{f})\). Their dual squared norm is bounded by \(C W_{5}\);
spectator products are included using \(W_{3}W_{2}\le C_{N}W_{5}\). Thus
ordinary source/gap Young gives

% DISPLAY i.46
% SOURCE 2|H_s[Uf,v]| <= eta H_exc[v]
% SOURCE        +C_eta[R^(-3)H_s[f]+kappa W_2[f]+W_5[f]].    (5.8)
\begin{gather*}
\begin{aligned}
& 2\vert H_{s}[Uf,v]\vert  \le  \eta  H_{\mathrm{exc}}[v] \\
& +C_{\eta}[R^{-3}H_{s}[f]+\kappa  W_{2}[f]+W_{5}[f]].
\end{aligned}\tag{5.8}
\end{gather*}

The conversion on the right is incident:
\(\sum _{a} w_{a} t_{a}\le C R^{-3}H_{I}+C \kappa  W_{2}\). It is
performed before replacing \(w_{a}\) by its maximum. The cutoff version
gains only the source-specific exponentially small errors in Section
\(7\). This explicitly includes the direct slow derivative source in the
mixed ledger.

\subsection{6. Summary of the mixed and diagonal
estimates}\label{i-summary-of-the-mixed-and-diagonal-estimates}

Let \(B[v]=H_{s}[v]+H_{\mathrm{exc}}[v]+J[v]\). After a fixed finite
allocation of \(\eta\) to the preceding terms, for every \(\eta >0\) the
exact mixed form has the decomposition

% DISPLAY i.47
% SOURCE H[Uf,v]=<Sf,v>+R_mix[f,v],
\begin{gather*}
\begin{aligned}
& H[Uf,v]=\langle Sf,v\rangle +R_{\mathrm{mix}}[f,v],
\end{aligned}\notag
\end{gather*}

with

% DISPLAY i.48
% SOURCE 2|R_mix[f,v]|
% SOURCE    <=eta B[v]+a_(eta)(kappa,sigma,R)H_s[f]
% SOURCE                      +b_(eta)(kappa,sigma,R)W_2[f],  (6.1)
\begin{gather*}
\begin{aligned}
& 2\vert R_{\mathrm{mix}}[f,v]\vert \\
& \le \eta  B[v]+a_{\eta}(\kappa ,\sigma ,R)H_{s}[f] \\
& +b_{\eta}(\kappa ,\sigma ,R)W_{2}[f],
\end{aligned}\tag{6.1}
\end{gather*}

where one can choose finite constants such that

% DISPLAY i.49
% SOURCE a_(eta)<=C_eta(sigma^2+R^(-3)),
% SOURCE b_(eta)<=C_eta(kappa+kappa^2+sigma^2+R^(-3)).          (6.2)
\begin{gather*}
\begin{aligned}
& a_{\eta}\le C_{\eta}(\sigma ^{2}+R^{-3}), \\
& b_{\eta}\le C_{\eta}(\kappa +\kappa ^{2}+\sigma ^{2}+R^{-3}).
\end{aligned}\tag{6.2}
\end{gather*}

The exact diagonal is

% DISPLAY i.50
% SOURCE H[Uf]=H_s^B[f]+E_0[f]+W[f]+R_diag[f],                (6.3)
\begin{gather*}
\begin{aligned}
& H[Uf]=H_{s}^{B}[f]+E_{0}[f]+W[f]+R_{\mathrm{diag}}[f],
\end{aligned}\tag{6.3}
\end{gather*}

and for an independently chosen \(\theta >0\),

% DISPLAY i.51
% SOURCE |R_diag[f]|<=theta H_s[f]
% SOURCE      +C_theta[kappa+sigma+sigma^2+R^(-3)]W_2[f]
% SOURCE      +C R^(-3)H_s[f].                               (6.4)
\begin{gather*}
\begin{aligned}
& \vert R_{\mathrm{diag}}[f]\vert \le \theta  H_{s}[f] \\
& +C_{\theta}[\kappa +\sigma +\sigma ^{2}+R^{-3}]W_{2}[f] \\
& +C R^{-3}H_{s}[f].
\end{aligned}\tag{6.4}
\end{gather*}

The \(H_{s}^{B}\) comparison is the full averaged-charge Berry
comparison from source Section \(4.2\) of Appendix \(F\):

% DISPLAY i.52
% SOURCE |H_s^B[f]-H_s[f]|
% SOURCE    <=theta H_s[f]+C(theta^(-1)kappa^4+kappa)W_2[f].  (6.5)
\begin{gather*}
\begin{aligned}
& \vert H_{s}^{B}[f]-H_{s}[f]\vert \\
& \le \theta  H_{s}[f]+C(\theta ^{-1}\kappa ^{4}+\kappa )W_{2}[f].
\end{aligned}\tag{6.5}
\end{gather*}

Equations (3.5), (6.3)-(6.5) keep every internal potential and Yukawa
term. No frozen slow-momentum resolvent has appeared.

\subsection{7. Cutoff vacuum, derivative tails, and source
parity}\label{i-cutoff-vacuum-derivative-tails-and-source-parity}

Fix \(\sigma >0\) before increasing \(R\). Let
\(h=(\sum _{e} r_{e}^{-2})^{-1/2}\), and choose a smooth even radial
\(\chi =\chi _{0}(\vert Y\vert ^{2}/(\sigma _{0}^{2}h^{2}))\), with
\(0<\sigma _{0}<c_{N} \sigma\) fixed, equal to one on an inner tube and
supported strictly inside the validity tube. Put

% DISPLAY i.53
% SOURCE U_sigma=chi U/n,    n=||chi U||.
\begin{gather*}
\begin{aligned}
& U_{\sigma}=\chi  U/n, \quad  n=\Vert \chi  U\Vert .
\end{aligned}\notag
\end{gather*}

This is intrinsic and residual-equivariant. Its support is
projection-stable. The normalized Gaussian coordinates
\(\xi _{e}=r_{e}^{1/2}Y_{e}\) have uniformly positive bounded covariance
matrices. On the cutoff transition region,

% DISPLAY i.54
% SOURCE sum_e |xi_e|^2 >= c sigma_0^2 r_*^3.
\begin{gather*}
\begin{aligned}
& \sum _{e} \vert \xi _{e}\vert ^{2} \ge  c \sigma _{0}^{2} r_{*}^{3}.
\end{aligned}\notag
\end{gather*}

Consequently any fixed-degree root-normalized Gaussian polynomial has a
tail bounded by

% DISPLAY i.55
% SOURCE tau_R=C_(N,sigma_0) exp(-c_N sigma_0^2 R^3),           (7.1)
\begin{gather*}
\begin{aligned}
& \tau _{R}=C_{N,\sigma _{0}} \exp (-c_{N} \sigma _{0}^{2} R^{3}),
\end{aligned}\tag{7.1}
\end{gather*}

after absorbing fixed powers of \(R\) into the exponential. These
estimates do not contain \(r_{\max}\). They apply after the rootwise
error-column factorization, and to finite graph source/dressing
coefficients with their own graph weights.

One must not claim that \(\Vert p_{F}(U_{\sigma}-U)\Vert\) has a
root-independent bound: a far vacuum derivative can indeed carry an
arbitrarily large zero-point norm. Instead use the excitation
ground-state transform. Since

% DISPLAY i.56
% SOURCE H_exc U=0,
% SOURCE H_exc[chi U]=integral |grad_F chi|_G^2 |U|^2,
\begin{gather*}
\begin{aligned}
& H_{\mathrm{exc}} U=0, \\
& H_{\mathrm{exc}}[\chi  U]=\int  \vert \nabla _{F} \chi \vert _{G}^{2} \vert U\vert ^{2},
\end{aligned}\notag
\end{gather*}

no spectator vacuum energy occurs. In the operator cutoff commutator,
\(\nabla _{e} \chi\) is a scalar multiple of \(Y_{e}/h^{2}\); its
product with \(\Omega _{e}Y_{e}\) has a root-normalized quadratic
coefficient of order \(h^{-2}\), because \(\Omega _{e}/r_{e}\) is
bounded. The same observation applies to the active finite graph states
in \(Z\) and to H1: replacing a \(p_{e}\) by \(\nabla _{e} \chi\) adds a
factor bounded by \(C/(r_{e} h^{2})\), with a Gaussian tail. Thus the
only cutoff errors required here satisfy, uniformly in all root ratios,

% DISPLAY i.57
% SOURCE |delta diagonal|<=tau_R(H_s[f]+W_2[f]),
% SOURCE |delta mixed|<=eta B[v]+C_eta tau_R(H_s[f]+W_2[f]),
% SOURCE H_exc[chi Zf/n]=H_exc[Zf]+O(tau_R W_2[f]),             (7.2)
\begin{gather*}
\begin{aligned}
& \vert \delta  \text{ diagonal }\vert \le \tau _{R}(H_{s}[f]+W_{2}[f]), \\
& \vert \delta  \text{ mixed }\vert \le \eta  B[v]+C_{\eta} \tau _{R}(H_{s}[f]+W_{2}[f]), \\
& H_{\mathrm{exc}}[\chi  Zf/n]=H_{\mathrm{exc}}[Zf]+O(\tau _{R} W_{2}[f]),
\end{aligned}\tag{7.2}
\end{gather*}

with the analogous graph-local \(H_{s}\) and \(J\) estimates. Slow
cutoff derivatives use \(\vert \partial  h\vert /h\le C/r_{*}\);
internal derivatives of \(h\) vanish. Derivatives of \(n\) are
controlled by the same normalized tail calculation. The normal
connection preserves \(\vert Y\vert ^{2}\), so it creates no extra
radial-cutoff generator.

\(S\) and \(Z\) are boson-odd. The cutoff is even, hence \(\chi\) Z/n is
EXACTLY orthogonal to \(U_{\sigma}\). The leading cutoff H1 source is
also odd and is orthogonal to the uncut \(U\), so its shifted-fast
inverse never encounters the vacuum kernel. Replacing it by \(S\) in the
mixed identity costs only (7.2). In the lower bound the source
completion with the uncut \(H_{\mathrm{exc}}\) is valid for every \(v\),
whether or not \(v\) is exactly orthogonal to \(U\): \(S\) is orthogonal
to its kernel.

Finally \(v\) perp \(U_{\sigma}\) has
\(\Vert Pv\Vert \le \tau _{R}\Vert v\Vert\) and therefore

% DISPLAY i.58
% SOURCE H_exc[v]>=c r_*(1-tau_R^2)||v||^2.                   (7.3)
\begin{gather*}
\begin{aligned}
& H_{\mathrm{exc}}[v]\ge c r_{*}(1-\tau _{R}^{2})\Vert v\Vert ^{2}.
\end{aligned}\tag{7.3}
\end{gather*}

All estimates (6.1)-(6.4) consequently hold for \(U_{\sigma}\) and its
genuine Dirichlet complement after adding the vanishing \(\tau _{R}\)
terms.

\subsection{8. Lower Schur comparison}\label{i-lower-schur-comparison}

Choose the reserve loss \(\delta\) and the mixed allocation \(\eta\)
small. A fixed positive fraction of \(J\) is available, so \(\eta\) can
be chosen below \(c_{0}\). Equations (1.5) and (6.1) leave

% DISPLAY i.59
% SOURCE H[U_sigma f+v]
% SOURCE   >=H[U_sigma f]+c_s H_s[v]+c_J J[v]
% SOURCE        +(1-delta-eta)H_exc[v]+2Re<Sf,v>
% SOURCE        -a_eta H_s[f]-b_eta W_2[f],                  (8.1)
\begin{gather*}
\begin{aligned}
& H[U_{\sigma} f+v] \\
& \ge H[U_{\sigma} f]+c_{s} H_{s}[v]+c_{J} J[v] \\
& +(1-\delta -\eta )H_{\mathrm{exc}}[v]+2\operatorname{Re} \langle Sf,v\rangle \\
& -a_{\eta} H_{s}[f]-b_{\eta} W_{2}[f],
\end{aligned}\tag{8.1}
\end{gather*}

where \(c_{s}\) and \(c_{J}\) are positive after the finite allocation.
Completing only the shifted fast form gives

% DISPLAY i.60
% SOURCE (1-delta-eta)H_exc[v]+2Re<Sf,v>
% SOURCE   >=-(1-delta-eta)^(-1)<f,S^*H_exc^(-1)Sf>.          (8.2)
\begin{gather*}
\begin{aligned}
& (1-\delta -\eta )H_{\mathrm{exc}}[v]+2\operatorname{Re} \langle Sf,v\rangle \\
& \ge -(1-\delta -\eta )^{-1}\langle f,S^{*}H_{\mathrm{exc}}^{-1}Sf\rangle .
\end{aligned}\tag{8.2}
\end{gather*}

The virtual multiplier is \(\le CW_{2}\). Therefore replacing its
coefficient by one costs only \(C(\delta +\eta )W_{2}\). Use (1.3),
(3.5), and (6.3)-(6.5). Their total error is made smaller than
\(\epsilon (H_{s}+W_{2})\) in the parameter order below. Taking the
infimum proves the lower half of (1.1). One may retain a further small
positive fraction of \(H_{s}[v]+H_{\mathrm{exc}}[v]+J[v]\) for the
form-domain argument, at an additional \(O(\eta  W_{2})\) cost.

\subsection{9. Upper trial and the complete internal-source graph
bound}\label{i-upper-trial-and-the-complete-internal-source-graph-bound}

Use the genuine complementary trial

% DISPLAY i.61
% SOURCE v_f=chi Zf/n,   Z=-H_exc^(-1)S.
\begin{gather*}
\begin{aligned}
& v_{f}=\chi  Zf/n, \quad  Z=-H_{\mathrm{exc}}^{-1}S.
\end{aligned}\notag
\end{gather*}

The actual graph components obey

% DISPLAY i.62
% SOURCE ||Z_e||^2<=C r_e^(-3),
% SOURCE ||Z_t||^2<=C w_t/s_t,
% SOURCE w_t=sum_(e in t)r_e^(-2),  s_t=sum_(e in t)r_e,
\begin{gather*}
\begin{aligned}
& \Vert Z_{e}\Vert ^{2}\le C r_{e}^{-3}, \\
& \Vert Z_{t}\Vert ^{2}\le C w_{t}/s_{t}, \\
& w_{t}=\sum _{e \in  t}r_{e}^{-2}, \quad  s_{t}=\sum _{e \in  t}r_{e},
\end{aligned}\notag
\end{gather*}

and their full derivatives include the spectator term \(m_{g}W_{2}\).
The complete interacting internal estimate is (1.4). At a graph vertex,
\(\alpha _{a} m_{g}\le C \kappa  W_{2,g}\); outside the graph the
potential charge commutes with the even graph map, and only
differentiated spectator vacua remain. This is why no \(\alpha _{a}\)
from an unrelated far block can multiply a nearest-root dressing norm.
The rootwise graph proof covers arbitrary momenta of \(f\).

There is also a direct incident estimate for the retained orbit square:

% DISPLAY i.63
% SOURCE J[Zf]<=C R^(-6)H_s[f]
% SOURCE                   +C kappa R^(-3)W_2[f]+C sigma^2 W_5[f]. (9.1)
\begin{gather*}
\begin{aligned}
& J[Zf]\le C R^{-6}H_{s}[f] \\
& +C \kappa  R^{-3}W_{2}[f]+C \sigma ^{2} W_{5}[f].
\end{aligned}\tag{9.1}
\end{gather*}

Indeed each fixed-degree graph state has the same \(Y_{e}\)
second-moment scale \(C/r_{e}\) on active and spectator roots. The
coefficient of \(t_{a}[f]\) is at most \(C w_{a} \sum _{g} m_{g}\).
Convert \(w_{a} t_{a}\) BEFORE using \(\sum _{g} m_{g}\le CR^{-3}\).
Derivatives on \(Z\) use the full active-plus-spectator derivative bound
and the pointwise coefficient \(\Vert A\Vert \le C \sigma\), giving
\(C \sigma ^{2} W_{5}\). This proves (9.1) without a global internal
kinetic loss.

The same absolute estimates used in the reserve proof, with the upper
Young inequality for the complete shifted charges, give

% DISPLAY i.64
% SOURCE H[w]<= (1+delta)(H_s[w]+H_exc[w])
% SOURCE                     +C_delta W_2[w]+C J[w]+V4[w].   (9.2)
\begin{gather*}
\begin{aligned}
& H[w]\le  (1+\delta )(H_{s}[w]+H_{\mathrm{exc}}[w]) \\
& +C_{\delta} W_{2}[w]+C J[w]+V_{4}[w].
\end{aligned}\tag{9.2}
\end{gather*}

For the fast part this is immediate from the exact identity (2.1), whose
leading signed terms have two-sided bounds. The internal shifted-charge
comparison is also two-sided; density, E0 and spin estimates are
absolute. If desired \(V_{4}\) can be absorbed using (3.2). Thus (9.2)
uses no unproved upper control of the full original Gauss square.

Apply (9.2) to \(v_{f}\). Since
\(H_{\mathrm{exc}}[Zf]=\langle f,S^{*}H_{\mathrm{exc}}^{-1}Sf\rangle\),
(1.4), (9.1), (3.2), and the cutoff estimates give

% DISPLAY i.65
% SOURCE H[v_f]<= (1+delta+C sigma^2)H_exc[Zf]
% SOURCE    +C R^(-3)H_s[f]+C kappa W_2[f]
% SOURCE    +C_delta R^(-3)W_2[f]+O(tau_R(H_s+W_2)).          (9.3)
\begin{gather*}
\begin{aligned}
& H[v_{f}]\le  (1+\delta +C \sigma ^{2})H_{\mathrm{exc}}[Zf] \\
& +C R^{-3}H_{s}[f]+C \kappa  W_{2}[f] \\
& +C_{\delta} R^{-3}W_{2}[f]+O(\tau _{R}(H_{s}+W_{2})).
\end{aligned}\tag{9.3}
\end{gather*}

The mixed remainder (6.1) applied to this trial is small by the same
graph estimates. Meanwhile

% DISPLAY i.66
% SOURCE 2Re<Sf,Zf>=-2<f,S^*H_exc^(-1)Sf>.
\begin{gather*}
\begin{aligned}
& 2\operatorname{Re} \langle Sf,Zf\rangle =-2\langle f,S^{*}H_{\mathrm{exc}}^{-1}Sf\rangle .
\end{aligned}\notag
\end{gather*}

Together with the diagonal expansion, this leaves
\(W-S^{*}H_{\mathrm{exc}}^{-1}S\) and only the displayed arbitrarily
small errors. Equation (1.3) therefore proves the upper half of (1.1).
The graph trial is used only for this upper bound; arbitrary
complementary vectors remain unrestricted in the lower bound.

\subsection{10. Simultaneous parameter
order}\label{i-simultaneous-parameter-order}

There is no circular order and no unattached constant \(C W_{2}\):

\begin{enumerate}
\def\labelenumi{\arabic{enumi}.}
\tightlist
\item
  Fix the partition and desired \(\epsilon\). Choose the finite Young
  allocations \(\theta , \eta\) and reserve loss \(\delta\) so the
  coefficients they alone contribute, including \(C(\delta +\eta )\)
  from (8.2), use less than epsilon/10 each.
\item
  Choose \(\kappa\) small enough for the fixed pair functional-calculus
  neighborhood, \(C_{\eta} \kappa , C \kappa\) from the source defect
  and E0 payment, and \(C \theta ^{-1}\kappa ^{4}\) to fit their
  assigned \(W_{2}\) budgets.
\item
  Choose a FIXED positive \(\sigma\) small enough for the reserve,
  \(C_{\theta} \sigma , C_{\eta} \sigma ^{2}\) and the complete-charge
  mixed costs to fit their budgets. Choose a fixed smaller
  \(\sigma _{0}\) for the radial cutoff.
\item
  Increase \(R\) until every constant now fixed times \(R^{-3}\), all
  complementary \(W_{2}\) absorptions, and all cutoff-tail quantities
  fit the remaining budgets.
\end{enumerate}

Constants such as \(C_{\delta}\) or \(1/\sigma\) are finite before \(R\)
is chosen. No estimate uses a later inverse power of a parameter to undo
an earlier smallness requirement without allowing that earlier parameter
to be chosen in the stated order. All estimates are uniform on nested
smaller supports at the selected fixed widths.

\subsection{11. Closed domains and variational
meaning}\label{i-closed-domains-and-variational-meaning}

First work on compact slow supports strictly inside the chosen branch
and tube. The exact first-order coefficients and density are smooth
there, with uniformly positive principal kinetic form and bounded local
matrix potentials. The associated Dirichlet form is closed after a
scalar shift, with its ordinary local H1 core. \(U_{\sigma}\) has smooth
compact fast support; \(P_{\sigma}\) and \(Q_{\sigma}\) map this core
into the same form domain. No assertion is made that the uncut \(P\)
preserves Dirichlet support.

The estimates above imply, after adding a bounded scalar multiple of
\(\Vert u\Vert ^{2}\),

% DISPLAY i.67
% SOURCE H[u]+C_R||u||^2 >= c H_s[f]+c B[v],
% SOURCE u=U_sigma f+v,
\begin{gather*}
\begin{aligned}
& H[u]+C_{R}\Vert u\Vert ^{2} \ge  c H_{s}[f]+c B[v], \\
& u=U_{\sigma} f+v,
\end{aligned}\notag
\end{gather*}

with \(c>0\); \(W_{2}\le C R^{-2}\) is bounded. Conversely the diagonal
upper estimate bounds \(H[U_{\sigma} f]\) by
\(C(H_{s}[f]+\Vert f\Vert ^{2})\). Hence \(P_{\sigma}\) extends
boundedly in the closed quadratic-form norm, equivalently the graph norm
of the averaged supercharge stack for the stated supersymmetric
realization. This proves \(P_{\sigma}\) Q(H) subset Q(H); preservation
of the Hamiltonian operator domain D(H), or boundedness in its operator
graph norm, is not asserted. Exhaustion by compact slow supports gives
the same identification on an unbounded selected branch, using these
UNIFORM estimates rather than local ellipticity constants that might
depend on \(r_{\max}\).

The restriction of the closed quadratic form of \(H\) to Q(H) intersect
\(\ker (U_{\sigma}^{*})\) is closed and strictly L2-coercive by (1.5)
and (7.3). For \(f\) in the slow form domain, the mixed functional is
continuous in that complementary form norm by (6.1) and the exact source
dual estimate. Riesz representation gives a unique minimizing \(v_{f}\)
in Q(H) intersect \(\ker (U_{\sigma}^{*})\) for each \(f\) in
\(Q(H_{s})\). This is form-domain attainment only; no operator-domain
regularity of \(v_{f}\) is inferred. The resulting Schur form is
continuous in the shifted \(H_{s}\) form norm, and (1.1) makes the
shifted quadratic-form norms equivalent after a bounded scalar shift. It
is therefore closed.

These are form statements. They do not claim a bounded operator \(B\) or
a second-slow-derivative operator identity on arbitrary \(f\), and do
not assume an inverse for \(H_{I}\). All cutoffs, projections, graph
dressings, and estimates intertwine under the same-branch
residual/color/normal-frame unitaries. Different cluster branches and
off-cone selector transport remain outside this local theorem.

\subsection{13. Explicit module and differentiated-tail
checks}\label{i-explicit-module-and-differentiated-tail-checks}

The following coefficient identities make the retained-module and
differentiated-tail steps explicit.

\subsubsection{13.1 Internal potential/Yukawa commute with the actual
injection}\label{i-internal-potentialyukawa-commute-with-the-actual-injection}

At a fixed intrinsic branch, the orthogonal Lie-algebra splitting into
off-block and retained block/center directions gives the graded tensor
product of the fast and slow Clifford modules. The injection is

% DISPLAY i.68
% SOURCE U(A,Z)f = g_(A,Z)(Y) v_F(A,Z) tensor f,
\begin{gather*}
\begin{aligned}
& U(A,Z)f = g_{A,Z}(Y) v_{F}(A,Z) \otimes  f,
\end{aligned}\notag
\end{gather*}

with the identity map on the ENTIRE slow Clifford module. There is no
projection onto an internal singlet or internal eigenstate. Each complex
bifundamental fast mass has \(8 n_{a} n_{b}\) occupied negative modes,
an even number. This remains constant under the gapped continuation, so
the product determinant \(v_{F}\) is even. In particular its graded
tensor injection intertwines the internal odd Clifford generators with
the same retained representation. The internal
Hamiltonian\textquotesingle s Yukawa term is even in any case: it is a
sum of bilinears in retained internal fermions only.

Pointwise in \((A,Z), V_{I}\) is scalar multiplication and \(Y_{I}(A)\)
acts only on those retained Clifford factors. Therefore, as actual maps,

% DISPLAY i.69
% SOURCE V_I U=U V_I,    Y_I U=U Y_I,
% SOURCE U^*Y_I v=Y_I U^*v.
\begin{gather*}
\begin{aligned}
& V_{I} U=U V_{I}, \quad  Y_{I} U=U Y_{I}, \\
& U^{*}Y_{I} v=Y_{I} U^{*}v.
\end{aligned}\notag
\end{gather*}

This proves their zero mixed compression for \(U^{*}v=0\). The cutoff is
a scalar bosonic multiplier, so the identities hold exactly for
\(U_{\sigma}\) as well.

The pair negative-projector radial transport is implemented on the fast
Fock module by a fast-fermion even unitary; its bosonic covariance
transport also acts only on the fast module. Both commute with the
retained Clifford action. The line may carry a nontrivial scalar Berry
connection, but this does not rotate the slow Clifford representation.
Differentiation gives

% DISPLAY i.70
% SOURCE p_A(Uf)=U nabla_A^B f+(1-P)(p_A U)f.
\begin{gather*}
\begin{aligned}
& p_{A}(Uf)=U \nabla _{A}^{B} f+(1-P)(p_{A} U)f.
\end{aligned}\notag
\end{gather*}

The transverse term is retained in (5.7); the Berry term is retained in
\(H_{s}^{B}\) and compared by the full charge estimate (6.5). Nothing in
the radial identification sets the full derivative of \(U\) to zero.
Residual color transformations act simultaneously on the fast
determinant line and the retained module, and \(U\) is the corresponding
equivariant intertwiner. These pointwise tensor identities therefore
descend to the residual-equivariant/physical subspace without requiring
individual block singlets.

\subsubsection{13.2 Exact beta-tail decomposition before internal
differentiation}\label{i-exact-beta-tail-decomposition-before-internal-differentiation}

Use the notation of Section \(2\), and set

% DISPLAY i.71
% SOURCE B=Psi^*F^(-1),
% SOURCE B_1=[C_Y F^(-1)+K^* ; I_Y F^(-1)],
% SOURCE R_m=T^*WT-I.
\begin{gather*}
\begin{aligned}
& B=\Psi ^{*}F^{-1}, \\
& B_{1}=[C_{Y} F^{-1}+K^{*} ; I_{Y} F^{-1}], \\
& R_{m}=T^{*}WT-I.
\end{aligned}\notag
\end{gather*}

The center row of \(B\) is

% DISPLAY i.72
% SOURCE B_C=g_C^(-1/2)[C_Y+K^*(F+L)]F^(-1),
\begin{gather*}
\begin{aligned}
& B_{C}=g_{C}^{-1/2}[C_{Y}+K^{*}(F+L)]F^{-1},
\end{aligned}\notag
\end{gather*}

and the internal row is \(B_{I}=I_{Y} F^{-1}\). Direct substitution into
\(\beta =A^{*}a\) gives the exact identities

% DISPLAY i.73
% SOURCE beta=B T^*WT(d_0+n),      beta_1=B_1d_0,
% SOURCE 
% SOURCE beta-beta_1=(B-B_1)d_0+B R_m d_0+B T^*WT n.           (13.1)
\begin{gather*}
\begin{aligned}
& \beta =B T^{*}WT(d_{0}+n), \quad  \beta _{1}=B_{1} d_{0},
\end{aligned}\notag \\
\begin{aligned}
& \beta -\beta _{1}=(B-B_{1})d_{0}+B R_{m} d_{0}+B T^{*}WT n.
\end{aligned}\tag{13.1}
\end{gather*}

The internal part of \(B-B_{1}\) is zero, and its center part has the
exact factorization

% DISPLAY i.74
% SOURCE (B-B_1)_C d_0
% SOURCE   =(g_C^(-1/2)-I)(C_Y F^(-1)d_0+K^*d_0)
% SOURCE      +g_C^(-1/2)K^* L F^(-1)d_0.                    (13.2)
\begin{gather*}
\begin{aligned}
& (B-B_{1})_{C} d_{0} \\
& =(g_{C}^{-1/2}-I)(C_{Y} F^{-1}d_{0}+K^{*}d_{0}) \\
& +g_{C}^{-1/2}K^{*} L F^{-1}d_{0}.
\end{aligned}\tag{13.2}
\end{gather*}

The metric difference is \(O(\sigma ^{2})\), and the last term has the
explicit \(O(\sigma )\) left factor \(K^{*}\). The other factors in
(13.2) are same-root or input-root-denominator triangle columns.

To expand the middle term of (13.1) without exposing a bare frozen
vacuum norm, use \(T-I=-Z_{m} T\) and the fact that \(T\) commutes with
\(Z_{m}\). Exactly,

% DISPLAY i.75
% SOURCE R_m d_0
% SOURCE   =-T^* Z_m^*d_0 -T^*T Z_m d_0
% SOURCE      +T^*K(K^*d_0)-T^*KK^*T(Z_m d_0).              (13.3)
\begin{gather*}
\begin{aligned}
& R_{m} d_{0} \\
& =-T^{*} Z_{m}^{*}d_{0} -T^{*}T Z_{m} d_{0} \\
& +T^{*}K(K^{*}d_{0})-T^{*}KK^{*}T(Z_{m} d_{0}).
\end{aligned}\tag{13.3}
\end{gather*}

Thus it contains only uniformly bounded left factors multiplying
\(Z_{m}^{*}d_{0}, Z_{m} d_{0}\), or \(K^{*}d_{0}\). These are the
controlled rootwise columns. For example

% DISPLAY i.76
% SOURCE Z_m d_0=F^(-1)L^*d_0
% SOURCE              -F^(-1)C_Y^*[C_Y R^(-1)d_0],
% SOURCE Z_m^*d_0=L F^(-1)d_0-K[C_Y F^(-1)d_0],
% SOURCE K^*d_0=C_Y R^(-1)d_0.
\begin{gather*}
\begin{aligned}
& Z_{m} d_{0}=F^{-1}L^{*}d_{0} \\
& -F^{-1}C_{Y}^{*}[C_{Y} R^{-1}d_{0}], \\
& Z_{m}^{*}d_{0}=L F^{-1}d_{0}-K[C_{Y} F^{-1}d_{0}], \\
& K^{*}d_{0}=C_{Y} R^{-1}d_{0}.
\end{aligned}\notag
\end{gather*}

The last term of (13.1) has the same explicit \(O(\sigma )\) outer \(B\)
multiplying the triangle column \(n\). This proves the extra small
factor by an exact finite identity, not by a formal infinite series.

For an INTERNAL derivative
\(\partial _{j}, Y, K, L, N_{Y}, I_{Y}, C_{Y}, W\) and \(g_{C}\) are
independent of A. Consequently

% DISPLAY i.77
% SOURCE ||partial_j B||<=C sigma/r_*,
% SOURCE partial_j Z_m=-F^(-1)(partial_j F)Z_m,
% SOURCE ||partial_j T||<=C sigma/r_*.
\begin{gather*}
\begin{aligned}
& \Vert \partial _{j} B\Vert \le C \sigma /r_{*}, \\
& \partial _{j} Z_{m}=-F^{-1}(\partial _{j} F)Z_{m}, \\
& \Vert \partial _{j} T\Vert \le C \sigma /r_{*}.
\end{aligned}\notag
\end{gather*}

Every derivative of a root-preserving coefficient \(F_{e}^{-1}\) adds
\(1/r_{e}\). A derivative of \(F_{e}^{-1}C_{e}^{*}\) can remove its
optional \(O(\kappa )\) factor, but replaces it by \(O(1/r_{e})\). The
same-root or triangle column therefore still has an extra
\(1/r_{e}\le 1/r_{*}\) in its low Gaussian norm. Crucially the outer
\(B, K^{*}\), or \(g_{C}^{-1/2}-I\) displayed in (13.1)-(13.3) remains,
or its derivative has the SAME \(\sigma\) gain and another \(1/r_{*}\).

Applying the finite Gaussian rootwise bounds to these differentiated
identities gives

% DISPLAY i.78
% SOURCE ||(beta-beta_1)U||<=C sigma W_2^(1/2),
% SOURCE ||[partial_A(beta-beta_1)]U||
% SOURCE                     <=C sigma r_*^(-1)W_2^(1/2).    (13.4)
\begin{gather*}
\begin{aligned}
& \Vert (\beta -\beta _{1})U\Vert \le C \sigma  W_{2}^{1/2}, \\
& \Vert [\partial _{A}(\beta -\beta _{1})]U\Vert \\
& \le C \sigma  r_{*}^{-1}W_{2}^{1/2}.
\end{aligned}\tag{13.4}
\end{gather*}

The derivative in the second line acts on the operator coefficient,
exactly as required by the curvature term, not on \(U\). Derivatives of
\(U\) are already present in the complete-charge cross. At
\(A=0, \partial _{A} d_{0}\) need not vanish; (13.1)-(13.3) explicitly
retain its contribution with the same extra fast factor. The non-small
derivative of \(\beta _{1}\) itself is not included in the tail estimate
and remains in the leading H2 scalar.

Finally the coefficients of \(\beta -\beta _{1}\) have at least two
fast-coordinate factors before \(p_{F}\). A fast divergence removes at
most one. Differentiating the displayed normalized coefficients
therefore gives
\(\vert \operatorname{div} _{F}(b-b_{1})\vert \le C \sigma /r_{*}\); an
additional internal derivative costs at most \(1/r_{*}\) and leaves this
\(\sigma\) factor. This proves the divergence portions of Section
\(4.3\) with the required gain as well.

\subsubsection{13.3 Form domain
convention}\label{i-form-domain-convention}

The infimum, projection bound, and Riesz minimizer are statements about
Q(H), the closed quadratic-form domain. They do not establish invariance
of the Hamiltonian operator domain D(H). The minimizing complement is
asserted only in Q(H), for \(f\) in \(Q(H_{s})\). The global lower bound
does not require attainment.

\clearpage
\section{Off cone incidence
localization}\label{j-off-cone-incidence-localization}

Technical derivation for review. Equation and subsection numbers in this
appendix are local to the source derivation. Historical local labels are
retained, including numbering gaps or nonsequential labels. The
accompanying source notes retain the original expressions for
comparison.

\subsection{1. A slightly regularized commuting
tree}\label{j-a-slightly-regularized-commuting-tree}

Write a traceless ordered commuting configuration as
\(\lambda =(\lambda _{1},\ldots ,\lambda _{N})\), with
\(\rho ^{2}=\sum  \vert \lambda _{i}\vert ^{2}\). For a proper subset
\(B\) with at least two members put

% DISPLAY j.1
% SOURCE c_B=|B|^(-1) sum_(i in B) lambda_i,
% SOURCE a_B^2=sum_(i in B)|lambda_i-c_B|^2,
% SOURCE d_B^(-2)=sum_(j outside B) (|lambda_j-c_B|^2+a_B^2)^(-1).
\begin{gather*}
\begin{aligned}
& c_{B}=\vert B\vert ^{-1} \sum _{i \in  B} \lambda _{i}, \\
& a_{B}^{2}=\sum _{i \in  B}\vert \lambda _{i}-c_{B}\vert ^{2}, \\
& d_{B}^{-2}=\sum _{j \text{ outside } B} (\vert \lambda _{j}-c_{B}\vert ^{2}+a_{B}^{2})^{-1}.
\end{aligned}\notag
\end{gather*}

The extra \(a_{B}^{2}\) in the denominator is deliberate. It removes
inverse-matrix singularities from the off-cone extension whenever
\(a_{B}>0\), without changing the clustering mechanism. For
\(0<\kappa <1/20, L>1\) and \(\delta =\kappa  \exp (-L)\), use

% DISPLAY j.2
% SOURCE theta_B=Theta(log(a_B/(delta d_B))/L),
% SOURCE C_B=cos(theta_B), S_B=sin(theta_B),
\begin{gather*}
\begin{aligned}
& \theta _{B}=\Theta (\log (a_{B}/(\delta  d_{B}))/L), \\
& C_{B}=\cos (\theta _{B}), S_{B}=\sin (\theta _{B}),
\end{aligned}\notag
\end{gather*}

where \(\Theta\) is a fixed smooth nondecreasing function, \(0\) on
(-infinity,0{]} and \(\pi /2\) on {[}1,infinity). Every derivative of
the trigonometric factors is flat at the constant-tail boundaries. If
\(a_{B}=0\) and all exterior points are separated, \(C_{B}=1\) locally.
Coincident unused nodes are handled by the same identically zero-prefix
mechanism as in the decision tree.

The regularization retains a uniform derivative bound

% DISPLAY j.3
% SOURCE |d a_B|<=1, |d d_B|<=C_N,
% SOURCE |d theta_B|^2 <= C_N ||Theta'||_infinity^2 L^(-2) a_B^(-2)
\begin{gather*}
\begin{aligned}
& \vert d a_{B}\vert \le 1, \vert d d_{B}\vert \le C_{N}, \\
& \vert d \theta _{B}\vert ^{2} \le  C_{N} \Vert \Theta '\Vert _{\infty}^{2} L^{-2} a_{B}^{-2}
\end{aligned}\notag
\end{gather*}

on \(\delta  d_{B}<a_{B}<\kappa  d_{B}\). The constant is independent of
\(\delta\) and \(L\). For example a conservative \(C_{N}\) can be
obtained by taking \(\vert d d_{B}\vert \le 4N\) and then squaring
\(1+4N \kappa\).

For crossing subsets B,C, the same three-point argument gives

% DISPLAY j.4
% SOURCE d_B <= 2a_B+2a_C, d_C <= 2a_B+2a_C.
\begin{gather*}
\begin{aligned}
& d_{B} \le  2a_{B}+2a_{C}, d_{C} \le  2a_{B}+2a_{C}.
\end{aligned}\notag
\end{gather*}

Thus \(C_{B} C_{C}=0\) if \(\kappa <1/4\). All simultaneously selected
clusters are laminar. The decreasing-cardinality decision tree and its
exact pathwise IMS identity therefore still apply.

\subsubsection{Symmetric raw leaf
formula}\label{j-symmetric-raw-leaf-formula}

For an unordered partition \(\pi\) into at least two final blocks, its
tree leaf multiplier is exactly

% DISPLAY j.5
% SOURCE q_pi = product_(B in pi, |B|>=2) C_B
% SOURCE        product_(B proper union of >=2 pi-blocks) S_B.       (1)
\begin{gather*}
\begin{aligned}
& q_{\pi} = \prod _{B \in  \pi , \vert B\vert \ge 2} C_{B} \\
& \prod _{B \text{ proper union of } \ge 2 \pi\text{-blocks}} S_{B}.
\end{aligned}\tag{1}
\end{gather*}

This is not a claim that all visited nodes are unions on their entire
domain. Rather, any omitted rejected crossing node has \(S_{B}=1\) on a
neighborhood of the nonzero final-block selection support, by
laminarity. Nodes strictly inside a selected final block are skipped.
Thus (1) equals the original leaf product as a smooth function,
including its derivatives. In particular

% DISPLAY j.6
% SOURCE sum_pi q_pi^2=1,
% SOURCE sum_pi |d q_pi|^2 = sum_pi q_pi^2 E_pi,                  (2)
\begin{gather*}
\begin{aligned}
& \sum _{\pi} q_{\pi}^{2}=1, \\
& \sum _{\pi} \vert d q_{\pi}\vert ^{2} = \sum _{\pi} q_{\pi}^{2} E_{\pi},
\end{aligned}\tag{2}
\end{gather*}

where \(E_{\pi}\) sums \(\vert d \theta _{B}\vert ^{2}\) for the final
nontrivial blocks and the proper unions in (1). Terms that were skipped
or crossing contribute zero on that leaf. No arbitrary ordering of
equal-cardinality clusters remains in (1).

\subsection{2. Intrinsic matrix extension on one stationary
flag}\label{j-intrinsic-matrix-extension-on-one-stationary-flag}

Use the notation and associated-bundle coordinates of the intrinsic
multiblock geometry note. On an ordered flag of ranks
\(n_{1},\ldots ,n_{k}\), write

% DISPLAY j.7
% SOURCE X=Z+A+Y, B_Z^*Y=0, alpha_a^2=sum_i ||A_i^(a)||_F^2,
% SOURCE r_ab=|z_a-z_b|, r_*=min r_ab.
\begin{gather*}
\begin{aligned}
& X=Z+A+Y, B_{Z}^{*}Y=0, \alpha _{a}^{2}=\sum _{i} \Vert A_{i}^{(a)}\Vert _{F}^{2}, \\
& r_{ab}=\vert z_{a}-z_{b}\vert , r_{*}=\min  r_{ab}.
\end{aligned}\notag
\end{gather*}

For a union I of flag blocks define

% DISPLAY j.8
% SOURCE n_I=sum_(a in I)n_a,
% SOURCE c_I=n_I^(-1)sum_(a in I)n_a z_a,
% SOURCE a_I^2=sum_(a in I)alpha_a^2 + sum_(a in I)n_a|z_a-c_I|^2,
% SOURCE H_(b,I)=sum_i(A_i^(b)+(z_b-c_I)_i I)^2+a_I^2 I,
% SOURCE d_I^(-2)=sum_(b outside I)Tr H_(b,I)^(-1).                (3)
\begin{gather*}
\begin{aligned}
& n_{I}=\sum _{a \in  I}n_{a}, \\
& c_{I}=n_{I}^{-1}\sum _{a \in  I}n_{a} z_{a}, \\
& a_{I}^{2}=\sum _{a \in  I}\alpha _{a}^{2} + \sum _{a \in  I}n_{a}\vert z_{a}-c_{I}\vert ^{2}, \\
& H_{b,I}=\sum _{i}(A_{i}^{(b)}+(z_{b}-c_{I})_{i} I)^{2}+a_{I}^{2} I, \\
& d_{I}^{-2}=\sum _{b \text{ outside } I}\operatorname{Tr}  H_{b,I}^{-1}.
\end{aligned}\tag{3}
\end{gather*}

For a single nontrivial final block, \(a_{I}=\alpha _{a}\). Formula (3)
ignores \(Y\), as it must for the desired leading slow localization. It
uses traces, positive matrix inverses, and projector compressions; it
chooses no internal eigenbasis. If \(a_{I}>0\) every \(H\) is positive.
If \(a_{I}=0\) for a selected final block, exterior separation keeps
every \(H\) positive and its angle has a constant select tail. A union
of at least two blocks has \(a_{I}>0\) when \(r_{*}>0\).

The inverse-trace differentiation estimate remains uniform. If
\(H=\sum _{i} T_{i}^{2}+a^{2}\) I, then Cauchy-Schwarz and
\(\sum _{i}\Vert T_{i} H^{-2}\Vert _{F}^{2}\le \operatorname{Tr}  H^{-3}\)
bound the derivative of \(\operatorname{Tr}  H^{-1}\). The inequality
\(\sum  \operatorname{eigenvalues} (H)^{-3}\le (\sum  \operatorname{eigenvalues} (H)^{-1})^{3}\)
gives a dimension-controlled Lipschitz bound for \(d\). Center
subtraction only introduces the fixed factors from \(n_{I}\) and \(N\).
Hence the \(\theta\) derivative estimate of Section \(1\) holds for the
canonical \((Z,A)\) metric.

For any gauge-invariant scalar f(Z,A), the exact co-metric identity
gives

% DISPLAY j.9
% SOURCE |grad_X f|^2 <= (1+C_N sigma^2)|d_(Z,A)f|^2              (4)
\begin{gather*}
\begin{aligned}
& \vert \nabla _{X} f\vert ^{2} \le  (1+C_{N} \sigma ^{2})\vert d_{Z,A}f\vert ^{2}
\end{aligned}\tag{4}
\end{gather*}

on a thin fast tube. Indeed put \(p=(d_{Z} f,d_{A} f,0)\). Solving the
full coordinate differential and scalar gauge constraint gives the
ambient covector \((p-K^{*} \beta ,\beta )\), where
\(\beta =-M^{-*}O_{H}^{*}p\) and
\(O_{H}^{*}p=O_{C}^{*}d_{Z} f+O_{A}^{*}d_{A} f\). Both \(O_{C}\) and
\(O_{A}\) are projections of \(B_{Y}\). Factoring \(M^{-*}\) with
\(F^{-*}\) on the right gives
\(\vert \beta \vert \le C\vert Y\vert  \vert p\vert /r_{*}\). Since
\(\Vert K\Vert \le 2\vert Y\vert /r_{*}\), the full squared covector
norm is at most \((1+C_{N} \sigma ^{2})\vert p\vert ^{2}\). This
argument includes the full orbital contribution and has no unrelated
internal-radius loss.

\subsection{3. Support margins and an explicit center-gap
bound}\label{j-support-margins-and-an-explicit-center-gap-bound}

If \(q_{\pi}\) is nonzero, all nontrivial final blocks satisfy
\(\alpha _{a}<\kappa  d_{a}\). From

% DISPLAY j.10
% SOURCE d_a^2 <= r_ab^2+alpha_b^2+alpha_a^2
\begin{gather*}
\begin{aligned}
& d_{a}^{2} \le  r_{ab}^{2}+\alpha _{b}^{2}+\alpha _{a}^{2}
\end{aligned}\notag
\end{gather*}

and the analogous inequality with a,b interchanged, obtain

% DISPLAY j.11
% SOURCE alpha_a+alpha_b <= beta r_ab,
% SOURCE beta=2 kappa/sqrt(1-2 kappa^2).                          (5)
\begin{gather*}
\begin{aligned}
& \alpha _{a}+\alpha _{b} \le  \beta  r_{ab}, \\
& \beta =2 \kappa /\sqrt{1-2 \kappa ^{2}}.
\end{aligned}\tag{5}
\end{gather*}

Singleton radii are zero. Choose the tree \(\kappa\) so \(\beta\) is
strictly smaller, for example by a factor four, than the endpoint
constant allowed by the geometric lemma.

Every proper union I of at least two final blocks has been rejected, so
\(a_{I}\ge \delta  d_{I}\) on \(q_{\pi}\)\textquotesingle s support.
Suppose the currently included centers have diameter \(D\). Equation (5)
implies

% DISPLAY j.12
% SOURCE a_I <= sqrt(N(1+beta^2)) D.
\begin{gather*}
\begin{aligned}
& a_{I} \le  \sqrt{N(1+\beta ^{2})} D.
\end{aligned}\notag
\end{gather*}

From \(d_{I}\le a_{I}/\delta\), at least one eigenvalue of some exterior
\(H_{b,I}\) is \(\le N a_{I}^{2}/\delta ^{2}\). Evaluating its unit
eigenvector and using the triangle inequality yields

% DISPLAY j.13
% SOURCE |z_b-c_I| <= sqrt(N) a_I/delta + alpha_b.
\begin{gather*}
\begin{aligned}
& \vert z_{b}-c_{I}\vert  \le  \sqrt{N} a_{I}/\delta  + \alpha _{b}.
\end{aligned}\notag
\end{gather*}

Using \(\alpha _{b}\le \beta \vert z_{b}-z_{a}\vert\) for an included a
and adjoining this center grows \(D\) by at most \(M_{N}/\delta\), where

% DISPLAY j.14
% SOURCE M_N=[N sqrt(1+beta^2)+1]/(1-beta).
\begin{gather*}
\begin{aligned}
& M_{N}=[N \sqrt{1+\beta ^{2}}+1]/(1-\beta ).
\end{aligned}\notag
\end{gather*}

Starting from a closest pair and making at most k-2 additions proves

% DISPLAY j.15
% SOURCE r_* >= Delta_N R,
% SOURCE R^2=r_pi^2+sum_a alpha_a^2,
% SOURCE Delta_N=[N(1+beta^2)]^(-1/2)(delta/M_N)^(N-2).            (6)
\begin{gather*}
\begin{aligned}
& r_{*} \ge  \Delta _{N} R, \\
& R^{2}=r_{\pi}^{2}+\sum _{a} \alpha _{a}^{2}, \\
& \Delta _{N}=[N(1+\beta ^{2})]^{-1/2}(\delta /M_{N})^{N-2}.
\end{aligned}\tag{6}
\end{gather*}

Since \(\rho _{X}^{2}=R^{2}+\vert Y\vert ^{2}\) and
\(\vert Y\vert \le \sigma  r_{*}\), this also gives

% DISPLAY j.16
% SOURCE r_* >= Delta_N rho_X/sqrt(1+2sigma^2).                  (7)
\begin{gather*}
\begin{aligned}
& r_{*} \ge  \Delta _{N} \rho _{X}/\sqrt{1+2\sigma ^{2}}.
\end{aligned}\tag{7}
\end{gather*}

All nonzero raw weights, and the closures of their derivative supports,
have this margin. This is the key quantitative substitute for an
unspecified finite cover.

For a selected-block angle transition,

% DISPLAY j.17
% SOURCE d_a >= (1-beta) r_*/sqrt(N),
% SOURCE alpha_a >= delta(1-beta)r_*/sqrt(N).                    (8)
\begin{gather*}
\begin{aligned}
& d_{a} \ge  (1-\beta ) r_{*}/\sqrt{N}, \\
& \alpha _{a} \ge  \delta (1-\beta )r_{*}/\sqrt{N}.
\end{aligned}\tag{8}
\end{gather*}

For a proper-union angle, \(a_{I}\) dominates its center-relative
radius. Thus all nonconstant cutoff factors live at a fixed positive
multiple of \(\rho\), once \(N,\kappa ,L\) have been fixed. Internal
coincidences never cause an unbounded derivative of a surviving
multiplier.

\subsection{4. All branches, finite count, and smooth zero
extension}\label{j-all-branches-finite-count-and-smooth-zero-extension}

For every ordered composition \(N=n_{1}+\ldots +n_{k}, k\ge 2\), take
the incidence set of stationary ordered orthogonal projector flags for
\(\Phi _{X}\). Admit branches inside a padded domain with:

\begin{itemize}
\tightlist
\item
  endpoint bound strictly larger than \(\beta\), but still inside the
  stated geometric regime
\item
  center gap strictly smaller than the lower bound (7), for example one
  fourth of that bound
\item
  fast width strictly larger than the support of an invariant smooth
  fast cutoff
\end{itemize}

Use a smooth cutoff \(\chi _{Y}=1\) for \(\vert Y\vert \le \sigma\) h/4
and \(\chi _{Y}=0\) for \(\vert Y\vert \ge \sigma\) h/2, with
\(h=(\sum _{a<b}r_{ab}^{-2})^{-1/2}\). The geometric domain can allow
\(\vert Y\vert <\sigma  r_{*}\). Put

% DISPLAY j.18
% SOURCE b_j=(k!)^(-1/2) chi_Y q_pi                             (9)
\begin{gather*}
\begin{aligned}
& b_{j}=(k!)^{-1/2} \chi _{Y} q_{\pi}
\end{aligned}\tag{9}
\end{gather*}

on each stationary flag sheet \(j\). Extend it by zero outside the
padded support. Equations (5)-(7), the flat cutoff tails, and the strict
fast margin make this a smooth finite-cover section. There is no cutoff
associated with a choice of coordinate chart or block frame.

The Hessian estimate from the geometry note makes each admitted
stationary flag a nondegenerate critical point. In particular it is
locally a smooth branch, and no branch can be born or disappear inside
the nonzero support. Branches can have monodromy; chart transitions
merely permute direct-sum components.

\subsubsection{Explicit N-only multiplicity
bound}\label{j-explicit-n-only-multiplicity-bound}

A crude bound is sufficient. A flag of type \((n_{a})\) has real
dimension \(m=N^{2}-\sum  n_{a}^{2}\le N^{2}\). It is covered by at most
\(N\)! pivot charts. In a complex block-lower-unipotent chart, every
cumulative orthogonal projector is

% DISPLAY j.19
% SOURCE V(V^*V)^(-1)V^*.
\begin{gather*}
\begin{aligned}
& V(V^{*}V)^{-1}V^{*}.
\end{aligned}\notag
\end{gather*}

A common denominator for all projectors has degree at most \(2N^{2}\) in
real chart coordinates. \(\Phi\) has numerator and denominator degree at
most \(4N^{2}\); clearing denominators from each first derivative gives
polynomial equations of degree at most \(8N^{2}\). Nondegenerate real
critical points are isolated complex zeros as well. The elementary
Bezout isolated-zero bound therefore gives at most \((8N^{2})^{m}\) per
chart. Summing the at most \(2^{N-1}\) ordered compositions gives the
conservative bound

% DISPLAY j.20
% SOURCE D_N=2^(N-1) N! (8N^2)^(N^2).                           (10)
\begin{gather*}
\begin{aligned}
& D_{N}=2^{N-1} N! (8N^{2})^{N^{2}}.
\end{aligned}\tag{10}
\end{gather*}

The bound is deliberately huge and only used to ensure constants do not
secretly depend on the logarithmic inner scale. It counts ordered
sheets; (9) corrects their actual multiplicity.

\subsubsection{Why commuting admitted sheets are exactly cluster
partitions}\label{j-why-commuting-admitted-sheets-are-exactly-cluster-partitions}

Let \(B(X)=\sum _{i<j}\Vert [X_{i},X_{j}]\Vert _{F}^{2}\). On a
stationary thin flag, the transverse linear commutator has norm

% DISPLAY j.21
% SOURCE ||L_Z Y||^2=sum_(a<b) r_ab^2 |Y_ab|^2
\begin{gather*}
\begin{aligned}
& \Vert L_{Z} Y\Vert ^{2}=\sum _{a<b} r_{ab}^{2} \vert Y_{ab}\vert ^{2}
\end{aligned}\notag
\end{gather*}

with the corresponding fixed real-block normalization. Endpoint
perturbations are bounded by \(C\) times the padded geometric endpoint
constant times this norm (and by \(C \beta\) on the raw-weight support),
and the quadratic \(Y\) term by \(C \sigma\) times it. With the fixed
small geometric constants,

% DISPLAY j.22
% SOURCE B(X)^(1/2) >= c_N (sum r_ab^2|Y_ab|^2)^(1/2).          (11)
\begin{gather*}
\begin{aligned}
& B(X)^{1/2} \ge  c_{N} (\sum  r_{ab}^{2}\vert Y_{ab}\vert ^{2})^{1/2}.
\end{aligned}\tag{11}
\end{gather*}

Consequently a commuting admitted sheet has \(Y=0\). Its projectors
group common eigenspaces; a repeated joint eigenvalue cannot be split
between two selected thin blocks, by (5). Each unordered particle
partition appears exactly \(k\)! times among all ordered compositions
and flags. Thus on the commuting cone

% DISPLAY j.23
% SOURCE Q=sum_j b_j^2=1.                                      (12)
\begin{gather*}
\begin{aligned}
& Q=\sum _{j} b_{j}^{2}=1.
\end{aligned}\tag{12}
\end{gather*}

No extra \(N\)! multiplier belongs in (9).

Equation (11) also ensures that on a sufficiently small near-commuting
region the fast cutoff is identically one on every contributing sheet.
Using (7) and \(h\ge r_{*}/\sqrt{k(k-1)/2}\), it suffices to take

% DISPLAY j.24
% SOURCE t=B(X)/rho_X^4 <= c_N sigma^2 Delta_N^4.                (13)
\begin{gather*}
\begin{aligned}
& t=B(X)/\rho _{X}^{4} \le  c_{N} \sigma ^{2} \Delta _{N}^{4}.
\end{aligned}\tag{13}
\end{gather*}

The constant in (13) can be reduced to include every fixed geometric
margin. There is then no fast-cutoff derivative in the near-cone
localization; outside it, its usual cost could instead be paid by a
retained fast reserve.

\subsection{5. The exact cone identity and the small off-cone
defect}\label{j-the-exact-cone-identity-and-the-small-off-cone-defect}

Where \(Q>0\) define \(w_{j}=b_{j}/\sqrt{Q}\). Then
\(\sum _{j} w_{j}^{2}=1\). Let

% DISPLAY j.25
% SOURCE E(X)=sum_j |grad_X w_j|^2,
% SOURCE R(X)=sum_j w_j^2 E_j(X),                               (14)
\begin{gather*}
\begin{aligned}
& E(X)=\sum _{j} \vert \nabla _{X} w_{j}\vert ^{2}, \\
& R(X)=\sum _{j} w_{j}^{2} E_{j}(X),
\end{aligned}\tag{14}
\end{gather*}

where \(E_{j}\) is the sum of the ambient
\(\vert \nabla  \theta _{I}\vert ^{2}\) for the single-block select
factors and proper-union reject factors of that sheet. All expressions
are taken before any slow/fast projection.

At a commuting \(X\) on an admitted sheet, \(Y=0\). A pointwise common
diagonalization verifies that the first derivative of every scalar in
(3) in an off-diagonal internal direction is zero. This is a calculation
of an invariant Frechet derivative, not a choice of smoothly varying
eigenvectors. It remains true at repeated joint eigenvalues: trace
derivatives of inverse matrices against off-diagonal perturbations
vanish, and the radius-zero select factor is constant on a neighborhood.
Further, \(K=O_{C}=O_{A}=0\) at \(Y=0\) in the exact kinetic identity.
Therefore the ambient gradients of all raw weights equal their
commuting-base gradients, including at singular strata. The same
statement holds for the angles wherever their derivatives matter.

The whole symmetric commuting identity (2) therefore implies

% DISPLAY j.26
% SOURCE Q=1, dQ=0, E=R                                      (15)
\begin{gather*}
\begin{aligned}
& Q=1, dQ=0, E=R
\end{aligned}\tag{15}
\end{gather*}

at every nonzero commuting tuple. This is stronger than just equality of
the weight values.

On the unit sphere, all supported branch sums and their first two
derivatives are bounded: the incidence Hessian is uniformly invertible
after (7), the flag multiplicity is bounded by (10), all nonconstant
cutoff transition scales \(a_{I}\) obey fixed positive lower bounds, and
the support stays strictly inside the padded branch domain. Accordingly
\(Q\) and E-R are smooth scalar functions in a neighborhood of the
compact commuting unit locus. There is a finite explicitly defined
derivative bound \(K_{N,\kappa ,L,\sigma ,\Theta}\), obtained as the
supremum of their first derivatives on a closed smaller padded
neighborhood, such that

% DISPLAY j.27
% SOURCE |Q(X)-1|+|E(X)-R(X)| <= K dist(X,K_comm)                (16)
\begin{gather*}
\begin{aligned}
& \vert Q(X)-1\vert +\vert E(X)-R(X)\vert  \le  K \operatorname{dist} (X,K_{\mathrm{comm}})
\end{aligned}\tag{16}
\end{gather*}

for unit \(X\) close enough to that locus. This constant can grow with
\(\Delta _{N}^{-1}\); it is not asserted uniform in the shape gap. Its
role is to determine a later near-cone thickness, never to appear as an
unpaid inverse-square coefficient.

For every \(\epsilon >0\) choose a neighborhood thickness \(\eta >0\)
small enough that \(Q\ge 1/2\) and \(K \eta \le \epsilon\). Homogeneity
then proves

% DISPLAY j.28
% SOURCE E(X) <= R(X)+epsilon rho_X^(-2)                        (17)
\begin{gather*}
\begin{aligned}
& E(X) \le  R(X)+\epsilon  \rho _{X}^{-2}
\end{aligned}\tag{17}
\end{gather*}

on that conical neighborhood. If desired, the elementary estimate of
Section \(6\) makes the choice in terms of \(t\): take
\(t_{1}\le (\eta /(2 C_{N}))^{4}\), together with (13) and the
branch-continuation radius. This is a concrete order of choices, not a
simultaneous assumption \(\kappa =\kappa (\Delta _{N})\).

\subsubsection{Why direct leafwise normalization was
insufficient}\label{j-why-direct-leafwise-normalization-was-insufficient}

For a single select factor,
\(\vert d \cos  \theta \vert ^{2}=\sin ^{2} \theta  \vert d \theta \vert ^{2}\)
cannot be bounded by \(\cos ^{2} \theta\) times the same cost near its
zero. Hence one cannot bound \(\vert dw_{j}\vert ^{2}\) by
\(w_{j}^{2} E_{j}\) separately. Nor may one charge every nested
transition to every final leaf: a tiny subcluster transition inside a
larger selected block is not controlled with a shape-independent
coefficient by the latter\textquotesingle s total-radius Hardy estimate.
Only the summed defect argument (15)-(17) repairs this issue.

\subsection{6. An elementary almost-commuting distance
certificate}\label{j-an-elementary-almost-commuting-distance-certificate}

There is no need to assume a smooth retraction onto the singular
commuting cone. For a traceless Hermitian nine-tuple \(X\) of norm one,
put
\(\epsilon _{0}=\max _{i<j}\Vert [X_{i},X_{j}]\Vert _{F}\le \sqrt{t}\).
An elementary spectral-block construction produces a commuting traceless
tuple \(Z\) with

% DISPLAY j.29
% SOURCE ||X-Z|| <= C_N sqrt(epsilon_0) <= C_N t^(1/4),
% SOURCE C_N=729 sqrt(N^3+9)                                   (18)
\begin{gather*}
\begin{aligned}
& \Vert X-Z\Vert  \le  C_{N} \sqrt{\epsilon _{0}} \le  C_{N} t^{1/4}, \\
& C_{N}=729 \sqrt{N^{3}+9}
\end{aligned}\tag{18}
\end{gather*}

as a conservative choice.

To see this, process one component at a time inside the already
constructed invariant blocks. In each current block, group consecutive
eigenvalues across gaps at most \(\eta =\sqrt{\epsilon _{j}}\). Distinct
groups have spectral distance greater than \(\eta\). Replace the
processed component on each new spectral block by its trace-average; its
tuple error is at most \(N^{3/2} \eta\). Delete interblock entries of
all remaining components. The Frobenius norm of each deleted part is at
most \(\epsilon _{j}/\eta\), by its commutator with the processed
component. The new block-diagonal commutator differs from the compressed
old commutator only by two products of deleted off-block parts. Hence
\(\epsilon _{j+1}\le \epsilon _{j}+2(\epsilon _{j}/\eta )^{2}\le 3\epsilon _{j}\).
At most nine steps give (18); all changes preserve traces. Previously
processed components are scalar on their blocks and stay commuting with
subsequent steps.

This proof uses spectral subspaces only to establish a distance
inequality. The actual localization never uses this approximate tuple,
its ordering, or any internal eigenbasis. If the right side of (18) is
at most \(1/2\), normalize \(Z\) to norm one; its distance from \(X\) is
at most \(2C_{N} t^{1/4}\). Thus (16) can be converted into an explicit
\(t^{1/4}\) modulus with the finite derivative constant \(K\).

\subsection{7. Payment and physical
measure}\label{j-payment-and-physical-measure}

The first term \(R\) in (17) has the original leafwise payments:

\begin{enumerate}
\def\labelenumi{\arabic{enumi}.}
\tightlist
\item
  A selected-block transition costs at most
  \(C_{N} L^{-2} \alpha _{a}^{-2}\), with
  \(\alpha _{a}\ge c_{N} \delta  \Delta _{N} r_{\pi}\) by (8). On the
  physical low coefficient the lower-rank soft Hardy estimate pays this
  from an arbitrarily small internal-form fraction, with
  faster-than-inverse-square exterior defect.
\item
  A proper-union transition has
  \(a_{I}^{2}=\sum  \alpha _{a}^{2}+s_{I}^{2}\), where \(s_{I}\) is the
  canonical relative-center radius of that union. Therefore
  \(a_{I}^{-2}\le s_{I}^{-2}\), paid by ordinary Hilbert-valued
  relative-center Hardy in dimension 9(l-1). There are at most \(2^{N}\)
  such terms per sheet. Their coefficient \(C_{N}/L^{2}\) can be made
  arbitrarily small before \(\delta\) and \(\Delta _{N}\) are fixed.
\item
  On the colored complementary coefficient all scalar costs are bounded
  by \(C_{N,\kappa ,L} r_{\pi}^{-2}\), using (6)-(8), and are paid by a
  retained fast gap at sufficiently large radius. The physical
  lower-rank Hardy theorem is not applied to colored vectors.
\item
  The additional \(\epsilon  \rho _{X}^{-2}\) in (17) is at most
  \(\epsilon  r_{\pi}^{-2}\). On the low coefficient it uses an
  arbitrarily small fraction of the center Hardy reserve; on the
  complement it is absorbed by the fast reserve at large radius. If the
  near-cone region is set by \(t\), the existing outer cutoff has its
  separately paid commutator-potential IMS cost.
\end{enumerate}

The only form inputs in this payment are the previously identified
internal soft-Hardy estimates and the separately supplied all-vector
multiblock comparison with sufficient retained reserves. This
construction does not manufacture those reserves.

On each flag sheet the physical density is the exact \(J=\det  M\) from
the intrinsic geometry note. Same-sheet frame changes are unitary
residual-block/normal-frame transitions and introduce no scalar chart
partition. Distinct sheets remain distinct summands. Pull back the
physical section to the incidence cover and multiply by \(w_{j}\). Then

% DISPLAY j.30
% SOURCE ||u||^2=sum_j ||w_j u||^2,
% SOURCE h[u]=sum_j h[w_j u]-integral E(X)|u|^2                  (19)
\begin{gather*}
\begin{aligned}
& \Vert u\Vert ^{2}=\sum _{j} \Vert w_{j} u\Vert ^{2}, \\
& h[u]=\sum _{j} h[w_{j} u]-\int  E(X)\vert u\vert ^{2}
\end{aligned}\tag{19}
\end{gather*}

with sheet integration against the pulled-back physical measure,
equivalently \(J\) dZ dA dY after gauge integration. The \(k\)! factor
in (9) and normalization, rather than discarding branches, enforce
single counting. Local form estimates can therefore be integrated
chart-free on every sheet. No globally smooth projector, angular
gauge-fixing multiplier, or internal bound-state projector is used.

\subsection{8. Noncircular parameter
order}\label{j-noncircular-parameter-order}

For each fixed \(N\) and target loss \(\epsilon\):

\begin{enumerate}
\def\labelenumi{\arabic{enumi}.}
\tightlist
\item
  Fix the necessary fractions of center/internal/fast reserves in the
  conditional local form theorem.
\item
  Choose the geometric endpoint and fast constants; choose tree
  \(\kappa\) so \(\beta\) in (5) lies strictly inside the endpoint
  bound. Any analytic requirement on these constants must itself be
  shape-independent, as required by the local comparison.
\item
  Choose \(L\) so the sum of at most \(2^{N}\) logarithmic costs is paid
  by the allocated slow reserves.
\item
  Set \(\delta =\kappa  e^{-L}, \Delta _{N}\) from (6), padded incidence
  domains, and the finite derivative/\allowbreak{}continuation constants \(K\).
\item
  Choose the near-commuting threshold \(t_{1}\) using (13), (16)-(18),
  so the added IMS defect is at most \(\epsilon  \rho ^{-2}\).
\item
  Choose the outer cutoff parameters and finally the exterior radius
  large enough to pay all fixed-shape constants, complementary
  inverse-square costs, and soft-Hardy remainders.
\end{enumerate}

No rank-uniform statement or effective small radius is asserted. The
point is that every large delta-dependent constant is either multiplied
by a later freely chosen cone thickness or absorbed at a later exterior
radius.

\clearpage
\section{Corrected rotation algebra and support estimates}\label{app:rotation}
This appendix records the algebra, normalization and domain extension used in Section~\ref{sec:rotation}. The rotation anticommutator and kernel-invariance theorem are due to the cited prior work; the quantitative use of a bounded internal radius is the estimate needed here.

\subsection{Two charge conventions and two summation scopes}
Retain the kinetic-one Hamiltonian $h$. Let $\widetilde q$ denote the convention
\begin{equation}
 \{\widetilde q_\alpha,\widetilde q_\beta\}=\delta_{\alpha\beta}h,
 \qquad\sum_{\alpha=1}^{16}\norm{\widetilde q_\alpha g}^2=8\hh[g]
\end{equation}
on the physical core. The manuscript's baseline charges are $q=\sqrt2\widetilde q$ and have stack coefficient $16$. For a fixed $j=(u,v)$ the primitive of \cite[Lemma 2(c)]{HHI} is
\begin{equation}
 \widetilde a_{j,\alpha}=\frac1{16}(X_u\Gamma_v-X_v\Gamma_u)_{\alpha\epsilon}\theta_\epsilon
 +\frac1{56}\sum_{w\notin\{u,v\}}X_w(\Gamma_w\Gamma_{uv})_{\alpha\epsilon}\theta_\epsilon.
 \label{eq:expandedprimitive}
\end{equation}
The color index is contracted in each summand. This follows from $(5b+2c)/7$ because $\Gamma_u\Gamma_{uv}=\Gamma_v$ and $\Gamma_v\Gamma_{uv}=-\Gamma_u$.

All coefficient matrices are real. The first two are symmetric grade-one matrices, the others antisymmetric grade-three matrices. They are mutually orthogonal in the Frobenius inner product and each has squared Frobenius norm $16$. Each $\widetilde a_{j,\alpha}$ is a real linear combination of Majoranas, so its square is scalar by the CAR. Consequently
\begin{align}
 \sum_\alpha\widetilde a_{j,\alpha}^2
 &=\frac12\left[\frac{16}{16^2}(|X_u|^2+|X_v|^2)
          +\frac{16}{56^2}\sum_{w\notin\{u,v\}}|X_w|^2\right]I\\
 &=\left[\frac{|X_u|^2+|X_v|^2}{32}
          +\frac{\sum_{w\notin\{u,v\}}|X_w|^2}{392}\right]I.
\end{align}
For several blocks, concatenate color coordinates; no color-dimension factor occurs. A fixed coordinate belongs to eight of the $36$ pairs and is outside the other $28$, giving
\begin{equation}
 \sum_{u<v}\sum_\alpha\widetilde a_{(u,v),\alpha}^2
 =\left(\frac8{32}+\frac{28}{392}\right)\alpha^2I
 =\frac9{28}\alpha^2I.\label{eq:allsum}
\end{equation}
Thus $1/32$ is the bound for one generator, while $9/28$ is the exact coefficient after summing all generators. They are not alternative estimates for the same sum.

Rescale reciprocally:
\begin{equation}
 q_\alpha=\sqrt2\widetilde q_\alpha,\qquad
 a_{j,\alpha}=\widetilde a_{j,\alpha}/\sqrt2,\qquad
 J_j=\sum_\alpha\{q_\alpha,a_{j,\alpha}\}.
\end{equation}
The fixed-generator bound becomes $\alpha^2/64$, and \eqref{eq:allsum} becomes $9\alpha^2/56$. The products $8(1/32)$ and $16(1/64)$ agree. Rescaling the charge alone would change the generator identity and is not allowed.

\subsection{Direct Clifford verification of the identity}
Write $\widetilde a_\alpha=X_w(M_w)_{\alpha\epsilon}\theta_\epsilon$ for a fixed pair $(u,v)$. In the real gamma convention the coefficient matrices obey
\begin{align}
 \Tr(\Gamma_t^TM_w)&=\delta_{wu}\delta_{tv}-\delta_{wv}\delta_{tu},\\
 \sum_w\Gamma_w^TM_w&=\tfrac14\Gamma_{uv},\\
 \Tr(M_w^T\Gamma_{st})&=0.
\end{align}
The first identity yields the orbital rotation, the second the spin contribution $-i\theta\Gamma_{uv}\theta/4$, and the last removes the interaction anticommutator. These are exact finite Clifford identities. The included integer-matrix program checks them for all $36$ pairs and all indicated indices. It uses a self-contained real symmetric Spin(9) representation constructed from the octonion multiplication table. This test verifies the finite identities only; integration by parts and closed-domain statements require the arguments in the text.

\subsection{A stronger Casimir variant}
The weaker coefficient $1/4$ in Lemma~\ref{lem:ball} is sufficient for the manuscript. For comparison, the all-generator identity proves
\begin{equation}
 H_I[g]\ge\frac7{9R^2}\norm{(1-P_I)g}^2
 \qquad\text{when }\supp g\subset\{\alpha\le R\}.\label{eq:casimirball}
\end{equation}
Use the graded sum of block supercharges, or equivalently a blockwise charge index, with stack coefficient $8$. Its squares add to the complete $H_I$. Let $C_I=\sum_j(J_j^I)^2$ and initially take a physical smooth vector in a Casimir eigenspace $C_Ig=\lambda g$, $\lambda>0$. Put $E=H_I[g]$ and $K=9/28$. Rotational invariance of the complete form gives
\begin{equation}
 \sum_j\norm{J_j^Ig}^2=\lambda\norm g^2,
 \qquad\sum_jH_I[J_j^Ig]=\lambda E.
\end{equation}
Individual supercharges need not commute with rotations; only their averaged square is used. Pair the rotation identity with $J_j^Ig$, sum over $j$, and integrate by parts on the core. The first resulting term is bounded by
$\sqrt{8\lambda E}\sqrt{KR^2\norm g^2}$. For the other, pointwise Cauchy--Schwarz and the fact that each primitive square is scalar give
\begin{equation}
 \sum_\alpha\left\lVert\sum_j\widetilde a_{j,\alpha}(X)(J_j^Ig)(X)\right\rVert^2
 \le K\alpha(X)^2\sum_j\lVert(J_j^Ig)(X)\rVert^2.
\end{equation}
That term has the same bound. Therefore
\begin{equation}
 \lambda\norm g^2\le2\sqrt{8K}\,R\sqrt{\lambda E}\norm g,
 \qquad E\ge\frac{7\lambda}{72R^2}\norm g^2.
\end{equation}
In the generator convention used here, a dominant $B_4$ weight has Casimir
\begin{equation}
 \lambda_C=\sum_{i=1}^4\lambda_i(\lambda_i+9-2i),\qquad
 \lambda_1\ge\lambda_2\ge\lambda_3\ge\lambda_4\ge0,
\end{equation}
with all coordinates integral or all half-integral. The smallest positive integer case $(1,0,0,0)$ has Casimir $8$, and the smallest half-integral case $(1/2,1/2,1/2,1/2)$ has Casimir $9$. Summing positive Casimir sectors proves \eqref{eq:casimirball}.

The rotation-invariant ball preserves all compact-group projections. At finite Casimir cutoff the generators are bounded angular operators, and the projected smooth physical core remains radially compact. Invariant nonnegative forms split orthogonally over the Casimir decomposition. Finite cutoff, approximation in the form norm, and a limit therefore extend the result to the closed form domain. For support in a closed ball use a slightly larger open radius and decrease it to $R$. Spectator Hilbert spaces pass through the same argument by finite sums and closure. Neither $\widetilde a g$ nor $Jg$ is asserted to lie in the Hamiltonian operator domain for a general form vector.

\subsection{Kernel invariance without an assumed positive moment}
For completeness, the domain mechanism behind \cite[Theorem 1(a)]{HHI} is as follows. The physical kernel is invariant under the compact rotation group. Decompose it into isotypes so that $J_j$ is bounded on each finite-dimensional representation factor. Choose a smooth radial cutoff $\chi_R$ that is one in a ball and whose derivative is supported where $R<\rho<3R$. In the generator identity, test with two kernel vectors of finite isotypic type and move the charge through $\chi_R$. The surviving shell terms have bounds of the form
\begin{equation}
 C\norm{\phi}_{\{R<\rho<3R\}}
  \norm{\Psi}_{\{R<\rho<3R\}},
\end{equation}
since $[\widetilde q,\chi_R]=O(R^{-1})$ and $\widetilde a=O(R)$ on that shell. The compatible maximal/self-adjoint realization and physical-core approximation are the ones in the cited theorem. The bound tends to zero by $L^2$ integrability alone. Taking $\phi=J_j\Psi$ within each isotype gives $J_j\Psi=0$. Completing the compact-group decomposition proves invariance of the full kernel. This is an explanation of the cited prior theorem, not a new existence or domain theorem.

\clearpage
\section{Equivariant vacuum and reference form domains}\label{app:model}
This appendix expands the geometric identifications used in Section~\ref{sec:singletmodel}. It addresses the added rotation argument on the local structures supplied by Appendices A--J; it does not replace their collision-uniform estimates.

\subsection{Spatial rotations preserve actual stationary flags}
For an ordered flag $P=(P_1,\ldots,P_k)$, the functional
\begin{equation}
 \Phi_X(P)=\sum_{a,i}n_a^{-1}(\Tr(P_aX_i))^2
\end{equation}
is spatially invariant at the same flag. Thus $\Phi_{RX}$ and $\Phi_X$ are identical functions on the flag manifold, including all their flag derivatives. The map $(X,P)\mapsto(RX,P)$ preserves stationarity and every Hessian admission test. It fixes the flag along each spatial orbit; apparent permutations in a local numbering of sheets are only chart changes.

Center gaps, internal radii, fast norms and commutator norms are spatial scalars. The inverse-trace regularized cluster scales use spatial contractions and are likewise invariant. Raw select/reject weights and their square-sum normalization therefore preserve singlets. No angular cutoff is inserted. Color frames may be chosen as functions of the flag alone, so their transition maps are gauge changes independent of the spatial rotation. Actual normal-bundle frame changes are accounted for by the fast unitaries of Appendix A. Density flattening is scalar and cannot supply an extra angular character.

\subsection{Global two-block anchor including the central element}
At $A=0$, fix a reference center direction $n_0\in S^8$. The fast ground line is a one-dimensional genuine unitary representation of its stabilizer $\operatorname{Spin}(8)$. Every continuous character of a connected compact semisimple group is trivial. The statement includes the central element $-1$, rather than only the infinitesimal Lie algebra.

The fast action is obtained from the physical spatial and Clifford actions, restricted to the fast factor. Its occupied determinant contains $8n_1n_2$ modes and has even parity. There is therefore no hidden projective character or extra fermionic sign on the line. Choose a unit vector $v_*$ in the line over $n_0$ and define
\begin{equation}
 v_0(gn_0)=U_{\rm fast}(g)v_*,\qquad g\in\operatorname{Spin}(9).
\end{equation}
Stabilizer invariance makes this independent of the representative $g$, smooth on the sphere, and exactly equivariant. It is a global normalized section, not only a local choice of phase.

Its Berry one-form is group invariant. The tangent representation of $\operatorname{Spin}(8)$ on $S^8$ has no invariant covector, so the angular one-form is zero. The occupied determinant is independent of the radial center coordinate. The positive normalized Gaussian has zero scalar Berry one-form, because its norm is one and its entries are real. Hence the unperturbed center connection is trivial and its kinetic form is the ordinary $T_x$.

This argument is restricted to the anchor. At nonzero internal coordinates the vacuum may have nonzero Berry curvature; the finite-$A$ comparison is precisely the paid one in Appendices E--G and I. Gauge and Weyl identifications are retained, rather than being discarded by the trivialization.

\subsection{Covariant Kato transport and retained fermions}
The path $tA$, $0\le t\le1$, remains in the local endpoint tube. Let $P_t(x,A)$ be the gapped negative fast spectral projector. Covariance gives
\begin{equation}
 P_t(Rx,RA)=U_{\rm fast}(R)P_t(x,A)U_{\rm fast}(R)^*.
\end{equation}
Differentiating in $t$ shows the same covariance for the anti-Hermitian Kato generator $K_t=[\dot P_t,P_t]$. The solution of $\dot W_t=K_tW_t$, $W_0=I$, is unique, so $W_t$ is covariant. The induced fast Fock transport is even and acts only on fast fermions. Gaussian covariance transport and the scalar normalized radial cutoff are also equivariant. All these maps commute with the retained Clifford action.

Applying these maps to $v_0$ gives the actual injection $U_\sigma$ as an intertwiner. The retained relative module is exactly the $256$-dimensional module of sixteen Majoranas, and the slow generator on the smooth core is
$L_x+S_{\rm rel}+J_I$. This does not make the spatial derivatives of $U_\sigma$ vanish. Their transverse source and scalar Berry contributions stay in the complete comparison already charged in the local ledger; there is no extra Schur dressing or a second completion square.

\subsection{The model form and its extension map}
For the two physical internal Hilbert spaces, the reference model is \eqref{eq:modeldomain}. The internal forms are the complete nonnegative supersymmetric forms, with arbitrary energy and rotation type. The baseline bounded projection and padded support give an actual coefficient $f$ in the Dirichlet form closure of the padded branch. Extend each approximating compactly supported smooth coefficient by zero; the baseline form bounds yield a Cauchy sequence in $\QQ(T_x)\cap\QQ(H_I)$. Its limit is the zero extension of $f$. This establishes the extension for this domain. It asserts no boundary-free extension for an arbitrary section or preservation of $D(h_N)$.

The extended coefficient is supported in $\mathcal C_\beta$ from \eqref{eq:radialtube}. A detailed leaf need not survive internal-only rotations: those may alter extra direction-dependent tests. This is why every original local comparison and leaf-multiplier payment is completed before extension and averaging.

Let $V_I(g)$ rotate both internal blocks, including their fermions, while fixing $x$ and $\FF_{\rm rel}$. It preserves each $\alpha_a$, commutes with $H_I$ and $T_x$, and preserves the radial tube. Individual $V_I(g)$ need not commute with every diagonal rotation. Their Haar average
\begin{equation}
 P_I=\int_{\operatorname{Spin}(9)}V_I(g)\dd g
\end{equation}
does commute, by conjugation invariance of Haar measure. It is a contraction for the norm and the model form norm. The corresponding invariant projection reduces each of $T_x$ and $H_I$, giving
\begin{equation}
 T_x[f]=T_x[P_If]+T_x[(1-P_I)f],\qquad
 H_I[f]=H_I[P_If]+H_I[(1-P_I)f].
\end{equation}
It commutes with center radial weights, so the $s^{-2}$ weighted norm splits as well. Both summands of a diagonal singlet remain diagonal singlets.

Apply these identities only to the already-paid model form. The enlarged support remains in $\mathcal C_\beta$, so for almost every fixed center $x$ the internal vector has support in the ball of radius $\beta_{\rm pad}\nu|x|$. The direct-integral form description gives internal form-domain membership almost everywhere. Lemma~\ref{lem:ball} then applies fiberwise; Fubini and nonnegativity integrate the inequality. Approximation first on compact center annuli, then in the closed form norm, removes any smoothness restriction. No derivative of a radius-dependent internal spectral projection occurs; $P_I$ is independent of $x$.

\subsection{Representation types and residual constraints}
The relative module has highest weights
\begin{equation}
 \mathbf{44}:(2,0,0,0),\qquad
 \mathbf{84}:(1,1,1,0),\qquad
 \mathbf{128}:(3/2,1/2,1/2,1/2).
\end{equation}
A scalar harmonic on $S^8$ of degree $l$ has highest weight $(l,0,0,0)$. Only $\mathbf{44}$ pairs with such a harmonic into a singlet, and it does so at $l=2$. Internal invariants contribute arbitrary multiplicity spaces. The argument does not require each block to be separately invariant: nontrivial block representations can couple to a simultaneous internal singlet.

The relative $\mathbf{128}$ cannot give a singlet with a scalar harmonic and an internal trivial representation. It can contribute when paired with nontrivial spinorial internal representations; these states belong to $(1-P_I)\mathscr K$ and are controlled by the nonsinglet uncertainty bound. They are not omitted. Equal-block exchange and residual physical conditions only reduce the admissible ordered-cover state space and cannot weaken the lower bound.

For three or more final blocks, radial center Hardy already gives at least $64$ in kinetic-one units, so none of these angular identifications is needed for their coercive constant. Their geometry, cutoff errors, colored complements and soft internal payments remain those of the baseline proof. The only additional two-block restriction is the initial padded width $\beta_{\rm pad}\le1/12$.

\end{document}
